% si.tex : Supplementary Information body (English). Assembled 2026-10-05.
% Sources: MANUSCRIPT_SPEC 8 (list); docs/MANUSCRIPT_DRAFT_METHODS_INTRO_2026-09-30.md (Supplementary Methods 1-6,
%   ported text of 30 September 2026, not re-verified at assembly except where stated);
%   docs/MANUSCRIPT_DRAFT_SUPPORT_2026-09-30.md M3.2, M4.2, M5.4-M5.6 (Supplementary Methods 6-7, Note 1);
%   paper/claims/C6_label_placement/README.md 3.1 (Supplementary Methods 8, translated);
%   docs/EXPERIMENT_PLAN_FINAL_BATCH_2026-10-04.md and paper/claims/C2_bias_diagnosis/XA_result_2026-10-04.md (Supplementary Methods 9);
%   docs/QA_FINAL_REVIEW_2026-10-02.md Q14 (Supplementary Note 2, translated).
% Tables: sections/si_tables.tex, generated by tools/make_si_tables.py from the claims READMEs.
% Draft markers [verify], [UPDATE], [RESULT], [CITE], [VERIFY] and new-experiment placeholders [X?: ...] are blue.

\section*{Supplementary Methods}

\subsection*{Supplementary Methods 1. Label parsing rules and processing by data source}

\textbf{Labels.} ALT labels of the main regions come from three compilations: GTN-P/CALM\cite{Streletskiy2025}, the ALLena thaw-depth compilation for the Lena River Delta\cite{Veremeeva2025} and the ABoVE soil moisture and ALT data set for Alaska and Canada, version 2\cite{Moore2025}. Records were aggregated to cells before modelling. For ABoVE and ALLena, a cell is a unique coordinate rounded to four decimal places, and the label is the mean of all observations at that coordinate. For CALM sites, a cell is one site and the label is the multi-year mean. ABoVE records were retained when 0 < ALT < 300 cm \textcolor{blue}{[verify: range rules of the CALM and ALLena parsers]}. The version-3 cell table contains 17,572 rows, of which the analysis uses the 17,467 cells with direct labels. In version 3, direct labels are mechanical probing and GPR. CALM sites whose event metadata name thaw tubes or ground temperature were moved to a separate class and excluded. Twenty cells that the source event metadata identify as temperature-derived remained in the direct class (15 in Alaska, 4 in Canada, 1 on the Tibetan Plateau). The table was frozen before this was found. These cells are flagged in the version-4 label table and are handled like the probe-only sensitivity variant. Most Alaskan labels are GPR-derived \textcolor{blue}{[old result, recheck: 86 \% of Alaskan cells in the audit of 21 September 2026]}. Label years span 1990 to 2024 and are not matched to the covariate reference period. The resulting maps are static and refer to the 2015–2020 climatology.

\textbf{Label unit.} A label row has a different meaning in different regions. In Russia W and Russia E a row is the multi-year mean of one CALM site. In the Lena River Delta a row is a point observation at a rounded coordinate, usually from a single visit. In Alaska and Canada most rows are point observations and some are CALM site means. In the additional regions (below) a row is the mean over an approximately 1 km cell. The label count n in all figures and tables is the number of rows drawn. Table 1 reports for each region the number of rows, the number of distinct 1 km locations, the number of distinct ERA5-Land climate values and the number of 0.5° blocks. Statements that compare regions at equal n state the unit.

\textbf{Scoring mask.} A cell is scored only if its label, its CCI ALT value and its soil thawing degree-days are valid. The same scoring cells are used for all methods.

\subsection*{Supplementary Methods 2. Assembly and eligibility of additional regions (LGD, XF)}

Additional labels were collected from public archives for regions at least 100 km from existing labelled regions (Supplementary Tables S1 and S11). The label rules were fixed before any model was fitted on these regions. A label enters as a target label only if it is a direct measurement (probe, thaw tube, frost tube, GPR, or soil pit, core or auger) of the end-of-season thaw depth. The end of season is established by the data protocol (CALM, thaw tubes), by the maximum over repeated visits after 1 August, by an observation date in August or September for single visits, or by an explicit statement in the data description. Thaw tubes therefore count as direct measurements in the additional regions, whereas CALM thaw-tube sites were excluded in version 3 \textcolor{blue}{[verify: harmonize the thaw-tube rule before the Results section]}. Values estimated from degree-day ratios or models are not used. Disturbed sites (burned areas, water bodies, thermo-erosional landforms) and censored probe readings are excluded. Only annual values from 1990 onward with 0 < ALT ≤ 600 cm are used. New records are aggregated to approximately 1 km cells, indexed by ky = ⌊lat/0.009⌋ and kx = ⌊lon·cos φ/0.009⌋ with φ = (ky + 0.5)·0.009°. New cells within 0.01° of an existing direct-label cell are removed.

Three target regions were defined geographically: Tibet (Tibetan Plateau and Qilian Mountains, 26–40°N, 73–105°E), NAtlantic (Greenland, Svalbard and mainland Norway, Sweden and Finland) and Russia\_C (Russia between 90°E and 140°E, excluding the Lena River Delta window 71.5–73.6°N, 123.3–130.1°E). The North Atlantic grouping was fixed before any result because none of its parts was expected to meet the block criterion alone. Extended versions of Russia W, Russia E and Canada were built as sensitivity cases. A region is eligible for confidence intervals if at least one split is valid and the union of scoring blocks over the used splits contains at least 8 blocks. The regime is shallow if the mean ALT of the target cells is below 150 cm and deep otherwise. Eligibility and regime were determined from the split structure before any model was fitted on these regions. They were recorded with the hash of the input table (18,088 rows = 17,572 version-3 rows + 516 new rows; sha256 4c387b2b…). Tibet (132 cells, 34 blocks, deep regime), NAtlantic (41 cells, 13 blocks, shallow) and Russia\_C (57 cells, 13 blocks, shallow) are eligible. Some NAtlantic splits have only three scoring blocks, and NAtlantic results carry a few-blocks flag. All Tibetan target cells lie above the source-cell elevation range. Contrasts against P0 in Tibet are therefore labelled as predicted outcomes and are not interpreted. Temperature-derived labels in Tibet (39 cells) are used only to test the effect of the label definition (L39).

Three pools are compared. P4 contains the four main regions (version-3 cells, main-run shards). PE1 adds the eligible shallow new regions, and PE2 adds the deep region. Hypotheses L1e, L4e and L8e repeat L1, L4 and L8 on PE1 and PE2 with the same decision rules. Hypotheses L38–L42 test region-level replication, the label definition, the extended versions of existing regions, label-set variants and the source composition. For PE2, stratified means are reported in cm and in units of each region's P0 RMSE, because the deep region can dominate a mean in cm.

\paragraph*{Licence status.} \textbf{Licence status.} Several sources carry no redistribution licence or ask to be consulted before publication (CALM web files of two sub-sites, the CUSP v1.1 synthesis, NSIDC GGD353 and ten report-derived rows of the Yamal transect; Data availability). A licence-verified version of every affected run table was defined before the pooled aggregation (LG revision 15 (m)). It removes all cells flagged as licence-unverified. Counted over target cells, it retains 19 of the 38 new NAtlantic cells, 37 of the 75 new cells of the Canada extension and 7 of the 8 new cells of the Russia W extension. The eligibility rules are applied unchanged to the licence-verified version. If permission for any source has not been obtained at submission, the primary PE1, PE2 and L38–L42 verdicts are reported from the licence-verified version, and the full version is placed in the Supplementary Information with a flag. This choice was fixed before results. \textcolor{blue}{[UPDATE: provider responses and dates.]}

\paragraph*{XF.} The additional-region analysis XF applies the same label rules and eligibility rules, unchanged and without access to label values, to public data sets found after the main runs (survey in Supplementary Table S18), and repeats the core contrasts in each eligible region. New regions are labelled blind (independent-region confirmation); the North Atlantic licence-verified edition enters as a sensitivity case. \textcolor{blue}{[PENDING: XF eligibility results and region list; registered data deadline 11 October 2026, 14:22 KST; no eligible new macro region found so far]}

\subsection*{Supplementary Methods 3. Extension contrasts (LGX) and details of the combination structures}

The extension refits the eight base methods with the same splits, draws and seeds, so that every extension contrast is closed within its own shards. A reproduction check compares the refitted base methods with the main-run shards (per-key cell-weighted RMSE; tolerance 1e-9 cm for physics methods and 0.02 cm for CatBoost methods). The extension covers physics outputs as input features (F1a, F1k, F1n), log-ratio regression (V2) and multiplicative correction (RM). It covers pseudo-label controls: shuffled, target-level constant, source-level constant and linear-TDD pseudo-labels, pseudo-label ratios r ∈ {1, 3, 30} and shuffled target labels. It covers a validation ladder on the pooled data: random cells, coarse covariate groups, 0.05° cells, 0.5° blocks, leave-one-cluster-out with 0, 100 and 500 km buffers, and region holdout. It also covers within-region block cross-validation, a proximity ladder at fixed scoring cells, cell-level predictions with distance strata, learner equivalence over five splits and a physics-baseline ladder (affine Stefan, a Kudryavtsev-type model \textcolor{blue}{[CITE: Kudryavtsev et al. 1974]}, a soil-property Stefan variant, the CCI ALT product and their ensemble). Finally, it covers an error floor from within-group variance, learner capacity (CatBoost with 600 iterations and depth 6, a random forest with 500 trees, and CatBoost tuned by leave-one-source-region-out validation) and sensitivity variants (κ, exclusion of June–July observations, probe-only labels, alternative anchors, inputs without the CCI product, year-matched TDD). The confirmatory hypotheses are L10 (residual structure versus physics-input structure), L12 (additive versus multiplicative correction), L15 (Stefan pseudo-labels versus four placebo pseudo-labels), L19 (conclusions under random splits versus region holdout), L29 (existence of a physics or product baseline better than P0 at n = 0) and L30 (direct ML at n = 0 with higher-capacity learners). If L29 identifies a better baseline P*, the contrasts of L4, L8 and L10 are repeated against P*. The manuscript does not state that a learner improves on the physics baseline when it improves on P0 but not on P*. Holm-adjusted p values within each item are reported as auxiliary columns.

\subsection*{Supplementary Methods 4. Prediction intervals (LGU)}

The uncertainty experiment evaluates predictive distributions in held-out regions. Every method outputs 99 quantiles per scoring cell. Intervals, the continuous ranked probability score (CRPS, grid approximation) and pinball losses are computed from these quantiles\cite{GneitingRaftery2007}. The interval score of a central $(1-\alpha)$ interval $[L, U]$ is
\begin{equation}
\mathrm{IS}_{\alpha} = (U - L) + \frac{2}{\alpha}(L - y)_{+} + \frac{2}{\alpha}(y - U)_{+} \quad (\mathrm{cm}).
\end{equation}
Experiment A is a ladder of hierarchical predictive distributions for $u = \log y - \log(E_0 s)$, modelled as $u = a_r + b_{r,k} + \varepsilon$ with a region offset $a_r$, a block offset $b \sim N(0, \tau_b^2)$ and a cell residual $\varepsilon$. Hyperparameters are estimated from the source only. Each source region's offset is computed against the coefficient estimated without that region, so that source regions mimic a new target. The confirmatory stage (iii) uses a Student-$t$ prior with 4 degrees of freedom for $a_r$ and a half-normal hyperprior on its scale, integrated on a grid. The posterior of $a_r$ given the $n$ labels yields the predictive distribution in unlabelled blocks. Other stages use a normal prior, heteroscedastic residuals from a CatBoost scale model, a normalizing-flow residual distribution that is run only if it passes a registered gate, and a conformal post-correction. The main baseline B4 is centred on $E_n s$ with a constant log half-width. B4 is calibrated by treating each source region in turn as a pseudo-target with the same draw procedure. Experiment B compares normalizers of a hierarchical conformal interval at $n = 0$: constant width, normalizing flow, CatBoost multi-quantile regression, conditional flow matching and a physics-propagated width. Each normalizer is compared with a placebo that permutes it among cells. Calibration groups are source macro-regions with at least 20 cells, pooled with equal weight\cite{Dunn2023}. Experiment C reports spatial-information diagnostics without hypotheses.

Confidence intervals use a two-stage bootstrap. Blocks are resampled within macro-regions on both the calibration side and the scoring side, and the calibration quantities are recomputed in each replicate. The registration specified 2,000 replicates and a timing rule for reducing them; the rule was applied before the main run and set the count to 1,000. The equivalence margin for interval scores is 5 \% of the baseline score (auxiliary 10 \%). If the half-width of an interval exceeds the margin, equivalence is recorded as not testable. A lower interval score is used as support only if the 90 \% coverage of the method is at least 0.80. The confirmatory hypotheses are LGU-A1 (stage (iii) versus B4 at n ∈ {10, 40}, pooled over the Lena River Delta, Canada and Alaska) and LGU-B2 (flow normalizer versus its placebo at n = 0 in the four main regions). A statement that conditional interval width transfers between regions is made only if LGU-B2 meets its four registered conditions. With two to four calibration groups there is no finite-sample coverage guarantee, and coverage is reported as an empirical quantity.

\subsection*{Supplementary Methods 5. Tabular foundation models and source cross-validation tuning of neural networks (LGT, LGF)}

TabPFN v2\cite{Hollmann2025} (package tabpfn 8.0.7 with the public v2 regression weights) and TabICL v2\cite{Qu2026} (package tabicl 2.0.2) were run on the same 27 targets, splits and draws with default settings and without fine-tuning. Both receive a context of at most 10,000 rows, the pre-training range of TabPFN v2. All n target labels enter the context. Source rows fill the remaining budget of 10,000 − max(T, 1,000) rows, where T is the number of target rows, by a region-proportional nested subsample. The source context is therefore identical for all n ≤ 1,000. E0 is always computed from the full source set. Each foundation model is paired with CatBoost fitted on the identical context in the same process. TabICL is additionally run with the full source set as context. The methods are D0, R0 and R1. The augmentation methods (D1, R2, R3) are omitted because the context limit does not allow the pseudo-label ratio of the main experiment.

The tuning experiment re-selects the hyperparameters of the MLP, the multi-head MLP, the reduced FT-Transformer-type model and RealMLP using source data only. For each target and learner, a random search\cite{BergstraBengio2012} over up to 34 configurations including the default is scored by leave-one-source-region-out cross-validation with equal region weight \textcolor{blue}{[verify: configuration counts per learner after LGF revision 5]}. The best candidates are refitted with two new seeds. A non-default configuration is selected only if it improves the fold-wise score consistently over the default, which limits selection bias\cite{CawleyTalbot2010}. Target labels and scoring cells are never used for selection. Default and tuned versions are both refitted on the local server, so that each comparison is closed on one platform. The confirmatory hypotheses are LGF-F1 (TabICL direct prediction at n = 0 under both context conditions versus P0) and LGF-N1 (tuned direct prediction of each of the four neural learners at n = 0 versus P0). Both hypotheses have the form ``not superior''. Support is split into two classes: inferior or equivalent in all contrasts, and no evidence of superiority when a contrast is indeterminate. Only the first class enters statements that a conclusion holds across learner types. Statements about tabular foundation models as a class are limited to the two models tested.

\subsection*{Supplementary Methods 6. Cross-environment rules and registration times}

Two execution environments were used (Supplementary Table S16). The main experiment and the extension runs LGX-A and LGX-B ran on cloud nodes (Rescale iolite-4: 64 CPU cores and four NVIDIA T4 GPUs). LGT, LGF, LGU, LGD, the local continuation of LGX, the LGX aggregation, the WRAPUP aggregation and the map were run on a shared local server with NVIDIA RTX 3090 GPUs (24 GB). The LG aggregation is taken from the cloud job if it completed without error and is otherwise repeated locally \textcolor{blue}{[UPDATE: which case applied]}. New work was moved to the local server following an instruction on 30 September 2026 to avoid further cloud cost. The extension run LGX-A was stopped at 09:02 KST on 30 September. At that time all CPU shards were complete, and 58 of 572 GPU shards, all of the FT-Transformer-type learner, were missing. Primary verdicts are computed from the cloud shards alone. Learner-equivalence verdicts that require the missing splits are recorded as not determined over five splits, and the three-split verdict of LG is retained. \textcolor{blue}{[UPDATE: completion status of the RealMLP shards of the main run.]}

Four comparison rules were fixed before results. First, both sides of a contrast Δ(A − B) come from shards of the same platform. Second, verdicts that combine split distributions use shards of one platform; if platforms are mixed, the verdict is labelled ``cross-environment'' and the single-platform verdict is reported alongside. Third, pooled means that combine regions run on different platforms use only within-region contrasts and carry the cross-environment label; PE1, with P4 from the cloud and the new regions from the local server, is such a pool. Fourth, cross-platform comparison tables are descriptive and are not used for verdicts. Contrasts involving V1r are never formed across platforms.

Three cross-environment checks were defined. Check (i) refits four CPU units of the cloud pre-check and smoke runs on the local server with the same settings. It compares per-key cell-weighted RMSE with a tolerance of 1e-9 cm for physics methods and 0.02 cm for CatBoost and ridge methods. \textcolor{blue}{[RESULT: outcome of check (i) by method class, recorded in LG revision 14.]} Check (ii) reruns the extension aggregation of the pre-check shards on the local server. It requires point estimates and interval limits identical to within 1e-6 cm and identical verdict strings. \textcolor{blue}{[RESULT: check (ii).]} Check (iii) refits four GPU units of the FT-Transformer-type learner twice on the local server and compares them with the cloud shards. The local continuation of the missing shards proceeds only if the physics keys agree to within 1e-9 cm and one of three registered rules holds for the neural keys. Rule (a) requires that no key differs from the cloud value by more than 0.5 cm. Rule (b) applies when the two local refits differ (median absolute difference above 1e-6 cm) and requires the median absolute difference between local and cloud values to be at most three times the median difference between the two local refits. Rule (c) applies when the local refits are identical to within 1e-6 cm and requires a median absolute difference between local and cloud values of at most 0.1 cm, with at most 5 \% of keys differing by more than 0.5 cm. \textcolor{blue}{[RESULT: check (iii) and the continuation decision.]} Verdicts that use locally continued shards are auxiliary, carry the cross-environment label and are not used in confirmatory statements or in the abstract. GPU training is not bitwise deterministic, and the repeat differences are reported with check (iii). CatBoost did not pass check (i), so pooled rows that combine platforms carry a cross-environment flag (Methods, deviations).

\paragraph*{Registered experiments.} The table below lists the experiments of the label-grid family with their registration commits (committer times, Korea Standard Time; commits were not time-stamped by a third party). The workflow experiments were registered as WF0--WF5 in 4168331 (1 October 2026), WF6--WF8 in 1b42b4d (1 October 2026), WF9 and WF10 in 8180632 (2 October 2026) and the soil-grouping sensitivity in 7da0064 (2 October 2026), and the analyses XA to XJ in 2678100 (4 October 2026, 14:22:39) (Methods).

{\fontsize{7.5}{9}\selectfont\setlength{\tabcolsep}{4pt}
\begin{longtable}{@{}>{\raggedright\arraybackslash}p{\dimexpr 0.100\linewidth-8pt\relax}>{\raggedright\arraybackslash}p{\dimexpr 0.360\linewidth-8pt\relax}>{\raggedright\arraybackslash}p{\dimexpr 0.200\linewidth-8pt\relax}>{\raggedright\arraybackslash}p{\dimexpr 0.340\linewidth-8pt\relax}@{}}
\toprule
\textbf{Experiment} & \textbf{Question} & \textbf{Execution environment} & \textbf{Registration commit (KST)} \\
\midrule\endhead
\bottomrule\endlastfoot
LG & Error of 12 methods as a function of n in 27 targets (hypotheses L1–L8) & Cloud (Rescale, NVIDIA T4) & 40be64c, 2026-09-29 15:45:46; run submitted with 19f3f93, 17:04:45 \\
LGX & Combination forms, pseudo-label controls, validation ladder, learner equivalence, sensitivity (L9–L31; confirmatory L10, L12, L15, L19, L29, L30) & Cloud; aggregation and auxiliary continuation local & 8513cc2, 2026-09-29 18:17:29; operational definitions f29bfc8, 21:28:41 \\
LGD & Additional independent regions (L1e, L4e, L8e, L38–L42) & Local CPU & 6799b74, 2026-09-29 22:43:12; input tables and eligibility 896a15e, 2026-09-30 02:35:11; run and aggregation rules 028900a, 04:19:23 \\
LGT & TabPFN v2 on the label-count grid (L32–L37) & Local, NVIDIA RTX 3090 & 6799b74, 2026-09-29 22:43:12; operational definitions 45ef3b6, 2026-09-30 01:39:42 \\
LGU & Predictive distributions and interval evaluation (LGU-A1–A7, B1–B6, C1–C4; confirmatory A1, B2) & Local & a112cef, 2026-09-29 23:30:02 \\
LGF & TabICL v2 and source-tuned neural networks (LGF-F1–F6, N1–N4; confirmatory F1, N1) & Local, NVIDIA RTX 3090 & 3a4fe84, 2026-09-30 01:44:27; revision 1 f208516, 02:22:24 \\
WRAPUP & Reporting rules, abstract bundle AB1–AB10, scenarios SC1w–SC3w, placement L43, AK1w, map & Local & 316714c, 2026-09-30 03:02:42 \\
\end{longtable}}

\paragraph*{Registration events and later amendments.}

{\fontsize{7.5}{9}\selectfont\setlength{\tabcolsep}{4pt}
\begin{longtable}{@{}>{\raggedright\arraybackslash}p{\dimexpr 0.130\linewidth-8pt\relax}>{\raggedright\arraybackslash}p{\dimexpr 0.160\linewidth-8pt\relax}>{\raggedright\arraybackslash}p{\dimexpr 0.280\linewidth-8pt\relax}>{\raggedright\arraybackslash}p{\dimexpr 0.430\linewidth-8pt\relax}@{}}
\toprule
\textbf{Plan (section)} & \textbf{Hypotheses} & \textbf{Registered} & \textbf{Later amendments (no change of hypothesis wording)} \\
\midrule\endhead
\bottomrule\endlastfoot
LG §4 & L1–L8 & \texttt{40be64c} 2026-09-29 15:45:46 (initial text with revisions 1–2) & revision 3 \texttt{249005c} 16:35:29; revision 4 \texttt{19f3f93} 17:04:45 (ZovWo submitted); revision 6 \texttt{d8355d2} 18:18:36 (D0 description ``no physics anchor'') \\
LG 6A (LGX) & L9–L31 & \texttt{8513cc2} 2026-09-29 18:17:29 (revision 5) & revisions 7–8 \texttt{f29bfc8} 21:28:41; revision 11 (first committed as ``revision 9'' in \texttt{6799b74} 22:43:12, renumbered in \texttt{cec5719} 23:29:45); revision 15 (h) \texttt{ff1be02} 2026-09-30 11:44:30 \\
LG 6C (LGT) & L32–L37 & \texttt{6799b74} 2026-09-29 22:43:12 (revision 9; section written 21:45) & revision 12 \texttt{45ef3b6} 2026-09-30 01:39:42 \\
LG 6B (LGD) & L1e, L4e, L8e, L38–L42 & \texttt{6799b74} 2026-09-29 22:43:12 (revision 10; section written 22:42) & revision 13 \texttt{896a15e} 02:35:11; revision 14 \texttt{028900a} 04:19:28; revision 15 (m) \texttt{ff1be02} 11:44:30 \\
LGU & A1–A7, B1–B6, C1–C4 & \texttt{a112cef} 2026-09-29 23:30:02 (written 22:45) & revisions 1–3 \texttt{582be9a} 2026-09-30 03:36:00; revision 4 \texttt{ff1be02} 11:44:30 \\
LGF & F1–F6, N1, N2, N2s, N3, N4 & \texttt{3a4fe84} 2026-09-30 01:44:27 (written 01:30) & revision 1 \texttt{f208516} 02:22:24; revisions 2–4 \texttt{43df2f8} 04:02:52; revision 5 \texttt{b1a4349} 05:00:40; revision 6 \texttt{ff1be02} 11:44:30 \\
WRAPUP & L43, SC1w–SC3w, AK1w, AB1–AB10 & \texttt{316714c} 2026-09-30 03:02:42 (initial draft with revision 1) & revision 2 \texttt{ff1be02} 11:44:30 \\
\end{longtable}}

\subsection*{Supplementary Methods 7. Label unit and label year sensitivity (S-a, S-b) and the sequential stopping rule (SC3w)}

\paragraph*{S-a, label unit.} S-a compared recalibration with point labels and with 1 km location means in Lena, Canada and Alaska (mode x). We used split seeds 1–200 with the split rules of the label grid (valid splits: Lena 197, Canada 200, Alaska 200), five draws per split and n ∈ {3, 10, 40}. In the row condition, n label rows were drawn from the target half A as in the label grid. In the location-mean condition, the rows of A were first averaged within 1 km locations (mean ALT and mean √TDD), and n locations were drawn. In both conditions, E\_ls was the slope through the origin of the drawn labels and E\_n = (n E\_ls + 10 E0)/(n + 10). Both conditions were scored on the same B-half rows. For each split we averaged Δ(P1 − P0) over draws and computed Δ\_unit = Δ\_location-mean − Δ\_row. In Canada at n = 40, the location-mean condition was skipped in the 100 splits where A had at most 40 locations. Median Δ\_unit [10th, 90th percentile] was +0.31 [−0.16, +1.37] cm (n = 3) and +0.34 [−0.26, +1.99] cm (n = 10) in Lena, −0.47 [−1.76, +0.81] cm and −0.86 [−3.48, +1.05] cm in Canada, and +0.07 [−0.19, +0.41] cm and +0.09 [−0.36, +0.58] cm in Alaska. At n = 40 (descriptive), the medians were +0.56 cm (Lena), −0.96 cm (Canada, 100 splits) and +0.08 cm (Alaska). Following the registered rule, the Canadian median at n = 10 (below −0.5 cm) gives the sentence: ``With point labels, the gain from few-label recalibration was smaller than with 1 km location means (Canada, n = 10, −0.86 cm)''. The deployment guideline therefore states its label unit as ``1 km location mean''. Lena and Alaska stayed within ±0.5 cm at n ∈ {3, 10} with the opposite sign, so the sentences are written per region. No median exceeded +0.5 cm at n ∈ {3, 10}. At n ∈ {3, 10}, the 10–90 \% intervals of Lena and Canada did not overlap the Russia W reference interval in either condition, so the sentence ``part of the between-region difference in recalibration gain is explained by the label unit'' is not written. With B also averaged to 1 km locations (auxiliary), the medians at n = 10 were −0.65 cm (Canada), −0.17 cm (Lena) and −0.02 cm (Alaska).

\paragraph*{S-b, label year.} S-b replaced the multi-year CALM site means of Russia W and Russia E by single-year values. From the CALM raw data (PANGAEA 972777; 3,819 rows), we removed values with inequality signs (24 ``>'' and 10 ``<''), inactive rows (117) and empty values (122), leaving 3,663 rows (the conditions can overlap). Label cells were matched to CALM sites by coordinates within 0.005° in latitude and longitude. Of the 31 Russia W cells, 27 matched one site, 4 matched several sites and none was unmatched. Of the 30 Russia E cells, the counts were 27, 3 and 0. Cells with multiple matches kept their multi-year means. For the 54 single-match cells, the multi-year mean of the valid years reproduced the label value (maximum absolute difference below 1e-13 cm). The median number of valid years was 13 (Russia W) and 20 (Russia E). For split seeds 1–200, five draws and n ∈ {3, 10}, one observation year was drawn for each drawn site, E\_ls and E\_n were computed from the single-year values, and scoring used the multi-year means of the B half. Δ\_year = Δ(P1 − P0)\_single-year − Δ(P1 − P0)\_multi-year was paired within split and draw. The share of drawn labels replaced by a single-year value was 0.88 (Russia W) and 0.90 (Russia E). Median Δ\_year [10th, 90th percentile] was +0.02 [−0.62, +0.63] cm (n = 3) and +0.05 [−0.48, +0.52] cm (n = 10) in Russia W, and +0.01 [−0.19, +0.21] cm and +0.04 [−0.25, +0.28] cm in Russia E. Both medians stayed within ±0.5 cm at both n, so the registered sentence applies: ``The recalibration gain from 3–10 labels was maintained with single-year labels.''

\paragraph*{SC3w, sequential stopping rule.} SC3w tested a sequential rule that decides from the target labels alone whether to stop at 3, 10 or 40 labels. For each valid split (1–5) and each of 20 repetitions, a random permutation of A fixed nested label sets of 3, 10 and 40. At 3 and 10 labels, the jackknife standard error of log E (leave-one-label-out least-squares slopes before shrinkage) was compared with τ = 0.05. The rule stopped when the standard error was at most τ and otherwise continued, with 40 labels as the cap. The stopped stage gave E\_n (κ = 10) and P1. The comparison was P1 from the first 40 labels of the same permutation. Δ\_SC3 = RMSE(P1\_rule) − RMSE(P1\_40) was scored on the B half with 10,000 block bootstrap resamples, and the label ratio was the mean number of labels used divided by 40. Support required (i) label ratio ≤ 0.5, (ii) a three-region mean point estimate ≤ 0.3 cm, (iii) both CI upper limits of the three-region mean < 0.5 cm and (iv) no inferior region. A label ratio ≥ 0.9 was classified as ``rule not operating''. Values τ ∈ {0.025, 0.10, 0.20} were descriptive sensitivities. At τ = 0.05 the rule was classified as not operating. The three-region mean label ratio was 0.909 (condition (i) failed; conditions (ii)–(iv) held). The registered sentence is: ``At the main τ of 0.05 the stopping rule rarely operated (label ratio 0.91)''. By the registered rule, Δ\_SC3 is descriptive and the rule is not part of the deployment guideline. The shares of repetitions stopping at 3, 10 and 40 labels were 5, 0 and 95 \% in Lena, 7, 0 and 93 \% in Canada, and 11, 8 and 81 \% in Alaska. The three-region mean Δ\_SC3 was −0.048 cm [−0.079, −0.002] (block-equal −0.077 cm [−0.105, −0.049]), no region was inferior, and each region was equivalent. For τ = 0.025 the label ratio was 0.988 (not operating). For τ = 0.10 it was 0.542, and only condition (i) failed. For τ = 0.20 it was 0.185 and all four conditions held (Δ\_SC3 −0.677 cm [−0.909, −0.353]); this value is descriptive and does not change the verdict. In descriptive targets, Russia W (cap equal to its A size, 14–16 labels) was inferior at τ = 0.05 (+0.29 cm [+0.16, +0.40]), and among the ten sub-regions in mode x only AL-2 was inferior (+0.04 cm [+0.02, +0.18]).

\subsection*{Supplementary Methods 8. Placement strategies and the learned placement policy (XD)}

Candidates are the cells of the label half A of a target (A pool); label values are not used for selection except by S6. Sources: WF plan section 2 (WF2) and revision items 2 and 3 of 1 October 2026; \texttt{scripts/3\_deep\_learning/h54\_workflow.py} (\texttt{strategy\_sets}, \texttt{kcenter\_order}, \texttt{s4\_order}) and \texttt{scripts/3\_deep\_learning/h40\_label\_grid.py} (\texttt{draw\_blocks}).
\begin{description}
\item[S1, random cells.] \texttt{h40.draw\_cells}; the seed is fixed by target, mode, split, $n$ and draw.
\item[S2, block stratification.] The 0.5\textdegree{} blocks of A are put in random order, and one cell per block, in random order within the block, is drawn in turn until $n$ cells are drawn (\texttt{h40.draw\_blocks}).
\item[S3, covariate-cluster stratification.] $k$-means clustering ($k = 8$, 10 initializations) of the A pool in the standardized x25 covariate space, allocation proportional to cluster size by largest remainders, random cells within clusters.
\item[S4, covariate max-min distance.] The x25 covariates of the A pool are median-imputed and standardized with the A-pool mean and standard deviation (no variable weights). The first cell is chosen at random with the draw seed; then the cell with the largest minimum Euclidean distance to the chosen set is added in turn (greedy $k$-centre; ties to the smaller index).
\item[S5, source-extrapolation first.] Cells in descending order of the dissimilarity index of the area of applicability: the standardized x25 nearest-neighbour distance to the source cells (mode x) divided by the mean pairwise distance within the source. Ties are shuffled with the draw seed; one to five draws were stored per (split, $n$). \textcolor{blue}{[미확인: 칸별 동률 수는 세지 않았다]}
\item[S6, sequential active selection.] Start from 20 random cells (the S1 draw at $n = 20$) and add at most 20 cells at a time without crossing the next budget (20, 40, 50, 70, 90, 100, \dots, 500). At each step, choose the cells with the largest prediction variance of a virtual ensemble (10 members) of the residual CatBoost fitted to the current labels.
\item[S7, $\sqrt{\mathrm{TDD}}$ first.] Cells in descending order of $\sqrt{\mathrm{TDD}}$, where a coefficient error propagates most into ALT. Cells in the same ERA5-Land grid cell tie; ties are shuffled with the draw seed (eight cells of the grid have fewer than five stored draws).
\end{description}
Algorithm P (Methods) is the pre-specified placement procedure tested in the workflow analysis; its setting was chosen by a mechanical rule on Alaskan targets (FINAL\_BATCH appendix XC-0). The candidates were S2, S4 and the S8 settings (S8a, $\gamma$ 0 with farthest-first block order and the cell nearest the block centre first; S8b, farthest-first choice within blocks; S8c, $\gamma$ 0.5; S8p, $\gamma$ 1). Rule 2 excluded all candidates, so Algorithm P is S1 (cell-random sampling); the workflow contrasts XC-5b, XC-5c and S-XC3 were therefore not tested and remain in the XC Holm family with $P = 1$. Design-weighted (S8w) and block random-effect (S8r) coefficients reuse the S8a labels.

\paragraph*{Learned placement policy (XD-learn).} Two policies were developed on Alaskan tasks only (Alaska with its parent region excluded and the sub-regions AL-1 to AL-6, splits 1 to 5, $n \in \{10, 20, 40\}$), with 200 label sets per task, split and $n$ (S1 draws 40\%, S2, S4 and S8 sets with other seeds 30\%, and those sets with 20 to 50\% of cells replaced at random 30\%). The utility of a label set is the RMSE of residual ML ($\lambda$ 0.25) relative to the mean over random draws, averaged over the two weightings. Set features use no label values: coverage of the candidates in standardized covariate space, energy distance, block allocation, $\sqrt{\mathrm{TDD}}$ statistics, the unweighted dissimilarity index and the task context. S9-GBM is a CatBoost regressor; S9-DS is a DeepSets network (element features: 25 standardized covariates, $\sqrt{\mathrm{TDD}}$, block share, dissimilarity index and distance to the nearest chosen cell; ranking loss; three seeds). Both policies add five cells at a time, choosing among the 300 farthest-first and 200 random candidates the cells with the smallest predicted utility. The variant with the smaller development cross-validation utility (one Alaskan sub-region left out together with Alaska) was frozen as S9* = S9-DS (commit f4a89ea, 5 October 2026, 03:17:28 KST; freeze record sha256 2751fecf49b8bcbd) before any test task was run. The test tasks outside Alaska (LE-1, LE-2, CA-2, CA-3 and the Tibetan Plateau, mode x, splits 1 to 5, $n$ 20 and 40, policy seeds 0 to 4) were evaluated once against Algorithm P. Because the five tasks form three families (Lena Delta, Canada, Tibetan Plateau), the smallest possible family-level sign-test $P$ is 0.125, so no $P$ values are reported. The test fits ran on the local server (operational deviation). S9* had lower error than Algorithm P in 2 of 5 test tasks (0 of 3 families) and higher error in 2 of 5 tasks (2 of 3 families), in both weightings and without a test (Supplementary Table S13). The candidate pool consists of already measured cells, so gains are not guaranteed to transfer to field candidates. XD-5 (descriptive, built after the XD-4 tables were opened): against the best of 200 candidate label sets, S9* had lower regret than Algorithm P in 3 of 5 test tasks and higher regret in 2; leaving a family out, the cross-validation utility was negative in 2 of 5 tasks for S9-DS and in 5 of 5 for S9-GBM (post hoc comparison). The S9-DS utility is sensitive to the minibatch-order seed, which includes the fold name; this is not a determinism failure, because the frozen model gave identical selections on two GPUs (Supplementary Table S13).

\subsection*{Supplementary Methods 9. Additional registered analyses XA to XL}

The ten analyses were registered in commit 2678100 (4 October 2026, 14:22:39) before any fit (docs/EXPERIMENT\_PLAN\_FINAL\_BATCH\_2026-10-04.md). Their placement in the main text or the Supplementary Information was fixed before results, and results are reported in Supplementary Table S13.

\paragraph*{XA, gain decomposition (main text, post hoc).} The gain of the selection rule W over the source-coefficient model with all labels, $G_{\mathrm{tot}} = \mathrm{P0} - \mathrm{W}$, was split into the recalibration share $G_{\mathrm{recal}} = \mathrm{P0} - \mathrm{P1}$ and the share added by the rule $G_{\mathrm{W}} = \mathrm{P1} - \mathrm{W}$; the ML share beyond recalibration was $G_{\mathrm{ML2}} = \min(\mathrm{P1}, \mathrm{P2}) - \mathrm{R1}(\lambda\,0.25)$ (auxiliary $G_{\mathrm{ML}} = \mathrm{P1} - \mathrm{R1}$ and the shrinkage remainder $G_{\mathrm{shr}} = \mathrm{P1} - \mathrm{P2}$). Coefficient error was the absolute ten-label bias (practical, main) or the absolute log ratio of the all-label target coefficient to the source coefficient (oracle). Spearman rank correlations used the 28 targets of WF4-c. The main interval resampled target families (clusters) jointly with blocks under the same resample index; the auxiliary interval kept the targets fixed. A result was 'positive' or 'negative' only if the main interval excluded zero in both weightings and the subset with at least 100 cells in half A (23 targets) did as well; otherwise it was 'not established'. When the main and auxiliary categories differed, both were recorded and the weaker one was reported. The analysis refitted nothing: it read the stored WF4 shards (132 units), and two reproduction gates (WF4-c, $\rho$ 0.38 [0.17, 0.55]; WF0) passed before the sealed tables were opened. The registered family list had seven families; the North Atlantic target (point estimates only) was added as an eighth single-target family in the implementation. All three registered hypotheses were 'not established' (Supplementary Table S13).

\paragraph*{XB to XJ.} The designs are summarized in Methods (Additional registered analyses). XB learns convex weights of physics models and ALT products from the target labels by cross-fitting; XC applies Algorithm P, the ten-label diagnostic and the method rule in sequence on new splits (seeds 201--210); XD selects the setting of Algorithm P on Alaskan targets and tests a learned placement policy; XE adds public inputs at 10--500~m to within-region residual ML; XF adds new regions from public data (Supplementary Methods 2); XG scores existing ALT maps on the same blocks after masking blocks near product training sites; XH describes the validation ladder in three regions; XI repeats the climate-extrapolation test on the licence-verified expanded Canadian table; XJ tests temperature-derived labels as low-weight auxiliary rows.

\paragraph*{Execution and deviations (XB to XJ).} XB, XD-alg, XH, XI and XJ ran as one cloud job on 5 October 2026 (02:18 to 02:38 KST); XE stage 1 (xh0) and the XD-4 test fits ran on the local server; XG ran locally after a second product extraction (02:13 to 02:15) and a re-check of the downloaded archives against their recorded md5 sums. Each analysis passed its reproduction gate against the earlier shards before its sealed tables were opened (first openings 03:18 to 03:45 KST; FINAL\_BATCH section 8). Deviations: XG compares products under a registration deviation (WRAPUP 10); the XE gate was first scored at the cloud-to-cloud tolerance and rescored at the registered local-to-cloud tolerance (a correction of the gate level, not of the design); the sealed XH reuse table of region-holdout rows was empty, so those values were read from the earlier source table without new fits; the XJ source-side rows were not tested because their licence basis had not been recorded before submission of the job. Flags: all analyses were designed after earlier results were viewed; XB, XD-1 to XD-3, XE (xh0), the CCI rows of XG and XJ-3 are partly unblinded; the Alaska rows of XH and XI are unblinded replications; XI uses warm blocks of the same period as a spatial proxy and is not a true extrapolation test; XD is exploratory and XH descriptive. XC ran on the local server on 5 October 2026 (03:16 to 05:19 KST; 208 units, no failure) after its reproduction gate passed (128 keys), and its sealed tables were first opened at 05:20:39 KST (FINAL\_BATCH 8.10). Because Algorithm P = S1 (Supplementary Methods 8), the placement step of the workflow was random placement and contributed nothing: at ten labels the workflow and random placement with the fixed recipe use the same labels and model, and at 40 and 160 labels the workflow contrast against random placement with the fixed recipe compares the selection rule with the fixed recipe on the same labels (equal to the auxiliary contrast S-XC1). XC-5b, XC-5c, XC-F3, S-XC0 and S-XC3 were therefore not tested (Holm family $m$ = 8, $P$ = 1). All XC statements refer to new random splits of the same regions (re-test in reused regions, partly unblinded: WF2, WF4, WF8, LGX-N4). \textcolor{blue}{[PENDING: appendix XC-F3, registered after the XF data deadline, 11 October 2026, 14:22 KST]} Later local runs on 5 October 2026 completed XD-5 (03:30 to 05:37 KST; opened 05:38:08), a sensitivity edition of XB with the year-matched CCI v5 map added as an eighth candidate (05:20 to 05:38; opened 05:43:31; the registered XB verdicts remain those of the main run, and XB-3 was not retained in this edition) and the XE-e diagnostic with on-site soil moisture (05:41; opened 05:42:35; changes no verdict). \textcolor{blue}{[PENDING: XE-a, XE-b, XE-a0, satellite inputs (stage 2) and the xt2 table; registered deadline 8 October 2026, 14:22 KST]} \textcolor{blue}{[PENDING: XF, new public regions; registered data deadline 11 October 2026, 14:22 KST]}

\paragraph*{XK and XL (added registration).} Both were registered on 5 October 2026 (commit d721d24) before computation, were designed after earlier results were viewed and are descriptive (no verdict words). XK aggregates within-region out-of-block predictions (0.5° block five-fold cross-validation; coefficient refitted on the training folds; $\lambda$ chosen by inner block cross-validation) and labels to grid cells of 0.05°, 0.1° and 0.25°, dropping grid cells with fewer than three label cells, and compares RMSE of P1, R1 and four ALT products at the same grid size (Supplementary Fig.~S15). In Alaska the RMSE difference between P1 and R1 was 0.50~cm at 1~km and 1.54 to 1.67~cm at 0.05°, 0.1° and 0.25° (cell-weighted). At label cells, CCI v5 exceeded the labels on average by 35.63~cm in Alaska, 58.60~cm in Canada and 0.27~cm in the Lena Delta, whereas CCI v4 differed by $-$5.09, +10.35 and $-$6.04~cm (product minus label, 1~km, cell-weighted); no cause is asserted. XL maps ALT at 1~km in the Alaskan permafrost zone and the Lena Delta with the recalibrated anchor plus residual ML trained on all labels and compares the maps with the recalibrated Stefan model and with CCI v5, Wei 2026, Aalto 2018 and Yi--Kimball (Supplementary Figs~S10 to S14). The 1~km map is a display resolution; its information resolution is bounded by the climate (about 10~km) and soil (about 5~km) inputs. The maps carry no accuracy claim; accuracy is compared only at label cells (XG, XK). Values: Supplementary Table S13, parts m and n.

\section*{Supplementary Notes}

\subsection*{Supplementary Note 1. Physical meaning of the Stefan coefficient $E$}

\paragraph*{What $E$ represents.} \textbf{D1. What E represents.} The Stefan coefficient E is a lumped parameter. In the classical Stefan solution, the thaw depth is Z = (2 k\_t I\_s / L)\textasciicircum{}0.5 \cite{Riseborough2008}. Here k\_t is the thermal conductivity of the thawed layer (W m\textasciicircum{}−1 K\textasciicircum{}−1), I\_s is the ground-surface thawing index (K·s), and L = ρ\_w L\_f θ is the volumetric latent heat of the thawed layer (J m\textasciicircum{}−3). ρ\_w is 1,000 kg m\textasciicircum{}−3, L\_f is 3.34 × 10\textasciicircum{}5 J kg\textasciicircum{}−1, and θ is the volumetric water or ice content. When the thawing index is taken from air temperature, I\_s = n\_t I\_a, where n\_t is the thaw n-factor \textcolor{blue}{[CITE: Klene et al. 2001, Arct. Antarct. Alp. Res.]}. In the form ALT = E·√TDD\_air, the coefficient becomes E = (2 k\_t n\_t · 86,400 s d\textasciicircum{}−1 / (ρ\_w L\_f θ))\textasciicircum{}0.5, expressed in cm (°C·d)\textasciicircum{}−0.5. E thus combines the thermal conductivity of the soil, the latent heat set by moisture and ground-ice content, and the surface energy transfer summarized by n\_t.

\paragraph*{Site factors that act through $E$.} \textbf{D2. Site factors that act through E.} Organic surface layers lower k\_t when dry and raise θ when wet. Both effects reduce E. Vegetation, moss and shading lower n\_t. Late-lying snow delays the start of ground thaw, so the air thawing index overstates the surface thawing index and the fitted E decreases. Soil texture and bulk density change both k\_t and θ. Edaphic factors of this kind were the basis of earlier regional Stefan mapping with land-cover specific coefficients \cite{Nelson1997,ShiklomanovNelson2002}.

\paragraph*{Non-physical terms absorbed by $E$.} \textbf{D3. Non-physical terms absorbed by E.} E also absorbs terms that are not soil physics. The air thawing index is computed from monthly means of a 0.1° reanalysis grid for 2015–2020. E therefore absorbs the grid-to-site temperature bias, the elevation mismatch within the grid cell, the difference between monthly and daily degree-day sums, and the difference between the climatology period and the observation years. The label definition enters as well. Probe, ground-penetrating radar (GPR) and temperature-derived labels differ, and a label measured before the end of the thaw season underestimates the seasonal maximum. Finally, E absorbs the error of the square-root law itself, which neglects sensible heat, layered soils and unfrozen water \textcolor{blue}{[CITE: Romanovsky \& Osterkamp 1997, Permafrost Periglac. Process.]}.

\paragraph*{Magnitude.} \textbf{D4. Magnitude.} For illustration, k\_t = 1.0 W m\textasciicircum{}−1 K\textasciicircum{}−1, n\_t = 1.0 and θ = 0.40 give E ≈ 3.6 cm (°C·d)\textasciicircum{}−0.5. The combination k\_t = 0.5 W m\textasciicircum{}−1 K\textasciicircum{}−1, n\_t = 0.8 and θ = 0.60 gives E ≈ 1.9 cm (°C·d)\textasciicircum{}−0.5. The median label ratios ALT/√TDD in the input table are 1.62 cm (°C·d)\textasciicircum{}−0.5 in Alaska, 1.28 cm (°C·d)\textasciicircum{}−0.5 in Lena and 1.37 cm (°C·d)\textasciicircum{}−0.5 in Canada (LG plan 6B.7, label statistics). These medians lie below the second illustrative value. This is consistent with low thermal conductivity, high moisture content, n\_t below 1, or a combination of these. The data used in this study do not separate these factors.

\paragraph*{Why $E$ transfers poorly between regions.} \textbf{D5. Why E transfers poorly between regions.} A coefficient estimated in one set of regions carries that set's mixture of soils, ground ice, vegetation, snow regimes and label types. Air temperature explains the thawing index but not these controls. Three observations in the input data show the size of the variation. First, the standard deviation of ALT among labels in the same 1 km cell is 11.3 cm in Alaska, 13.0 cm in Lena and 17.7 cm in Canada (WRAPUP 2.2). Second, across Canadian 0.5° blocks with at least 3 labels, the coefficient of variation of the block-mean ratio ALT/√TDD is 0.22 (unit metadata viewed in WRAPUP revision 1; WRAPUP 0.1). Third, the label ratio in the Tibetan candidate data is 11.2 cm (°C·d)\textasciicircum{}−0.5 for GPR labels and 12.9 cm (°C·d)\textasciicircum{}−0.5 for temperature-derived labels, about 7–9 times the Arctic values (LG plan 6B.7). The label type also differs by region. Alaskan labels are mostly ABoVE point observations, many from GPR, with some CALM site means. Canadian labels are ABoVE point observations with some CALM site means. Lena labels are mostly single-visit points. Russia W and E labels are multi-year CALM site means (WRAPUP 2.1–2.2).

\paragraph*{What the source-coefficient model tests in this design.} \textbf{D6. What P0 tests in this design.} E0 is pooled over source cells without regional weights. For each of the four main targets, 78.0–94.3 \% of the source cells are Alaskan (Table M3-1). P0 therefore tests how well an Alaska-dominated coefficient transfers. It does not test a generic physical coefficient. The registered covariate-dependent coefficient models test whether covariates recover the between-region variation of E: V1 under LG L7 \textcolor{blue}{[RESULT: L7 verdict]} and V2 under LGX L13 \textcolor{blue}{[RESULT: L13 verdict]}. The wrap-up plan records earlier negative results for models that predict E (WRAPUP 11.1, H27 and V1) and adds no new E model.

\paragraph*{Relation of the two baselines.} \textbf{D7. Relation of the two baselines.} P0 gives the error of a coefficient transferred from other regions with no target label. P1 gives the error after the mean level of E in the target is estimated from n labels, with E0 as a prior worth 10 labels. P1 changes one multiplicative factor. It corrects the regional level of E but not the variation of E within the target. The contrast P1 − P0 at n = 10 (AB4) measures the value of a few labels for the coefficient level \textcolor{blue}{[RESULT: AB4 4-way verdict and Holm-corrected wording]}. The contrasts R1 − P1 at n = 10 (AB5) and at full labels (AB9), and the minimum n of LG L4, measure what residual learning adds after this level correction \textcolor{blue}{[RESULT: AB5, AB9, L4]}. Two pre-result sensitivity analyses concern the label unit and the label year (M5.4, M5.5). A sequential stopping rule based on the jackknife precision of the target coefficient rarely stopped early at the registered threshold (M5.6).

\paragraph*{Limitation.} \textbf{D8. Limitation.} The construction of E0 is a design choice with two consequences. First, P0 and every method anchored on E0 (R0, D1 pseudo-labels at n = 0) inherit the Alaskan dominance of the source pool. Second, their performance relative to direct ML may differ under another pooling rule. We did not evaluate a region-equal-weight E0, so the direction of such a change is unknown. The dependence on E0 decreases as n grows, because E0 carries a fixed weight of 10 labels in E\_n.

\subsection*{Supplementary Note 2. Differences between the earlier contest test and the tests reported here}

An earlier within-Alaska analysis (a contest entry of this project) and the tests reported here answer different questions, and the word `Alaska result' referred to three different tests. In the contest test (six-fold 0.5\textdegree{} block cross-validation within Alaska; five of six block folds for training), the physics residual model reached 13.33~cm against 14.46~cm for the Stefan model (7.8\% lower RMSE), the best of 56 settings. In the within-region retest of this study (WF6; half of the blocks for training, the setting chosen by cross-validation within the labels, 25 splits, comparator the Stefan model recalibrated on the same labels), residual ML reached 13.87~cm against 14.41~cm (3.75\%). In the transfer test (Alaska treated as a new region, trained on the other regions, mostly the Lena Delta, before Alaskan labels were added), the values with all labels were 14.26~cm and 14.54~cm and no difference was established. The two within-region results agree in direction; the smaller gain of the retest follows from the smaller training share (one half instead of five sixths), the absence of selection among several settings and the recalibrated comparator. Relative to the error floor of 11.31~cm (error not reducible with the present covariates), the contest setting reduced the reducible squared error by 38.6\% and the stricter setting by 19.2\% (79.71 to 64.42~cm$^2$; Supplementary Table S7). Under four scoring schemes in Alaska (means of three seeds), the median distance from test cells to the nearest training cell was 0.004~km under random splits and 53.5~km under $k$-fold nearest-neighbour distance matching (kNNDM) cross-validation, which mimics the 81.6~km from map cells to the nearest label; the RMSE of direct ML rose from 11.6 to 17.9~cm, that of the Stefan model stayed at 14.2 to 14.5~cm, and the physics residual model was always below the Stefan model (12.77--13.70~cm; Supplementary Fig.~S1). Source: docs/QA\_FINAL\_REVIEW\_2026-10-02.md, Q14 and Q15, and paper/claims/C4\_sufficient\_labels/README.md 2.6 (values from data/processed/m1/cv\_scheme\_comparison.csv and the WF6 tests); the distance and RMSE values of the scoring schemes were moved here from the Results at the length revision of 5 October 2026.

\clearpage
% si_tables.tex : generated by tools/make_si_tables.py (do not edit by hand; edit the script or the claims READMEs).

% Supplementary Tables S1-S18 (MANUSCRIPT_SPEC 8).

\clearpage
\phantomsection\label{tab:S1}
\noindent\textbf{Supplementary Table S1.} Sub-regions, targets with point estimates only and the North Atlantic group.\par\smallskip
\noindent{\footnotesize Sub-regions of Alaska (AL-1 to AL-6), Canada (CA-2, CA-3) and the Lena Delta (LE-1, LE-2) were defined by $k$-means clustering of 0.5\textdegree{} block centroids without labels (133 blocks, assignment fixed before the main run). Mode x excludes the parent region from the source; mode i keeps the rest of the parent (simulated new area inside a labelled region). Ranges are over valid splits. $E_0$ in cm~(\textdegree{}C~d)$^{-1/2}$. The North Atlantic group is reported only in the full-data edition (licence-unverified rows; D README E19); n.a., not available. Source: paper/claims/D\_data\_and\_design/tables/Table1\_data.csv (copy of the v2 Table 1 data) and README E19.}\par
{\fontsize{7.5}{9}\selectfont\setlength{\tabcolsep}{4pt}\setlength{\LTpre}{2pt}\setlength{\LTpost}{4pt}
\begin{longtable}{@{}>{\raggedright\arraybackslash}p{\dimexpr 0.086\linewidth-8pt\relax}>{\raggedright\arraybackslash}p{\dimexpr 0.099\linewidth-8pt\relax}>{\raggedright\arraybackslash}p{\dimexpr 0.055\linewidth-8pt\relax}>{\raggedright\arraybackslash}p{\dimexpr 0.055\linewidth-8pt\relax}>{\raggedright\arraybackslash}p{\dimexpr 0.054\linewidth-8pt\relax}>{\raggedright\arraybackslash}p{\dimexpr 0.063\linewidth-8pt\relax}>{\raggedright\arraybackslash}p{\dimexpr 0.101\linewidth-8pt\relax}>{\raggedright\arraybackslash}p{\dimexpr 0.101\linewidth-8pt\relax}>{\raggedright\arraybackslash}p{\dimexpr 0.108\linewidth-8pt\relax}>{\raggedright\arraybackslash}p{\dimexpr 0.072\linewidth-8pt\relax}>{\raggedright\arraybackslash}p{\dimexpr 0.207\linewidth-8pt\relax}@{}}
\toprule
\textbf{Target} & \textbf{Group} & \textbf{Modes} & \textbf{Rows} & \textbf{1 km loc.} & \textbf{Blocks} & \textbf{Splits} & \textbf{A cells} & \textbf{Scored cells (blocks)} & \textbf{$E_0$, x / i} & \textbf{Label type} \\
\midrule\endfirsthead
\toprule
\textbf{Target} & \textbf{Group} & \textbf{Modes} & \textbf{Rows} & \textbf{1 km loc.} & \textbf{Blocks} & \textbf{Splits} & \textbf{A cells} & \textbf{Scored cells (blocks)} & \textbf{$E_0$, x / i} & \textbf{Label type} \\
\midrule\endhead
\bottomrule\endlastfoot
Russia C & Point estimate only & x & 7 & 7 & 6 & 5 & 3–4 & 3–4 (2–3) & 1.60 & CALM site mean \\
Greenland & Point estimate only & x & 3 & 2 & 2 & 0 & 1–2 & 0–1 (0–1) & 1.59 & CALM site mean \\
AL-1 & Sub-region & i, x & 5,154 & 51 & 10 & 4 & 288–4,834 & 320–4,866 (2–9) & 1.50 /\allowbreak{} 1.54 & Point, ABoVE (mostly GPR) \\
AL-2 & Sub-region & i, x & 3,491 & 128 & 20 & 5 & 1,333–1,924 & 1,567–2,158 (5–17) & 1.50 /\allowbreak{} 1.54 & Point, ABoVE (mostly GPR) \\
AL-3 & Sub-region & i, x & 651 & 69 & 16 & 5 & 269–361 & 290–382 (3–10) & 1.50 /\allowbreak{} 1.61 & Point, ABoVE (mostly GPR) \\
AL-4 & Sub-region & i, x & 471 & 44 & 4 & 3 & 139–287 & 184–332 (2–3) & 1.50 /\allowbreak{} 1.60 & Point, ABoVE (mostly GPR) \\
AL-5 & Sub-region & i, x & 619 & 27 & 17 & 5 & 270–454 & 165–349 (4–14) & 1.50 /\allowbreak{} 1.59 & Point, ABoVE (mostly GPR) \\
AL-6 & Sub-region & i, x & 3,220 & 24 & 7 & 4 & 1,516–1,705 & 1,515–1,704 (2–5) & 1.50 /\allowbreak{} 1.60 & Point, ABoVE (mostly GPR) \\
CA-2 & Sub-region & i, x & 395 & 27 & 14 & 5 & 161–237 & 158–234 (2–8) & 1.59 /\allowbreak{} 1.60 & Point, ABoVE \\
CA-3 & Sub-region & i, x & 349 & 53 & 19 & 5 & 113–224 & 125–235 (3–13) & 1.59 /\allowbreak{} 1.58 & Point, ABoVE \\
LE-1 & Sub-region & i, x & 1,786 & 135 & 14 & 4 & 448–1,184 & 602–1,335 (3–12) & 1.62 /\allowbreak{} 1.62 & Point, ALLena (single visits) \\
LE-2 & Sub-region & i, x & 1,251 & 66 & 6 & 3 & 183–1,062 & 167–1,014 (4–5) & 1.62 /\allowbreak{} 1.62 & Point, ALLena (single visits) \\
North Atlantic & New region, full edition only & x & 41 & n.a. & 13 & \textcolor{blue}{[S1: from lg\_targets, full edition]} &  &  &  & 1 km cell means \\
\end{longtable}}

\clearpage
\phantomsection\label{tab:S2}
\noindent\textbf{Supplementary Table S2.} Register of registered hypotheses.\par\smallskip
\noindent{\footnotesize Role: C confirmatory; P primary; A auxiliary; D descriptive or diagnostic; N direction-neutral report; abstract, contrast of the abstract bundle. Blinding: B blind; PU partially unblinded; RU replication (unblinded); RE result existed on disk at registration, not viewed; Sa S-a viewing label; NS not listed. Registered: first commit containing the hypothesis text (hash, date, Korea Standard Time). The verdict column gives the first clause of the registered verdict string of each plan (scope = verdict rows of the test tables cited in the claims READMEs); contrasts, $\Delta$, both 95\% CIs, uncorrected and Holm-adjusted $P$ values are in Supplementary Tables S3 to S13. Rows LG to WRAPUP: frame of docs/MANUSCRIPT\_DRAFT\_SUPPORT\_2026-09-30.md M4.3 (88 rows); WF and X rows added at assembly. \textcolor{blue}{[S2: per-hypothesis $\Delta$, CI, $P$ and Holm $P$ columns to be generated from the test tables, MANUSCRIPT\_SPEC 8]}}\par
{\fontsize{7.5}{9}\selectfont\setlength{\tabcolsep}{4pt}\setlength{\LTpre}{2pt}\setlength{\LTpost}{4pt}
\begin{longtable}{@{}>{\raggedright\arraybackslash}p{\dimexpr 0.069\linewidth-8pt\relax}>{\raggedright\arraybackslash}p{\dimexpr 0.100\linewidth-8pt\relax}>{\raggedright\arraybackslash}p{\dimexpr 0.100\linewidth-8pt\relax}>{\raggedright\arraybackslash}p{\dimexpr 0.319\linewidth-8pt\relax}>{\raggedright\arraybackslash}p{\dimexpr 0.115\linewidth-8pt\relax}>{\raggedright\arraybackslash}p{\dimexpr 0.137\linewidth-8pt\relax}>{\raggedright\arraybackslash}p{\dimexpr 0.160\linewidth-8pt\relax}@{}}
\toprule
\textbf{ID} & \textbf{Plan} & \textbf{Role} & \textbf{Statement (abridged)} & \textbf{Registered} & \textbf{Blinding} & \textbf{Registered verdict} \\
\midrule\endfirsthead
\toprule
\textbf{ID} & \textbf{Plan} & \textbf{Role} & \textbf{Statement (abridged)} & \textbf{Registered} & \textbf{Blinding} & \textbf{Registered verdict} \\
\midrule\endhead
\bottomrule\endlastfoot
L1 & LG §4 & P & Direct ML (D0) does not outperform the source-coefficient Stefan model (P0) at n ≤ 40 (at most 1 of 4 main regions with both CI upper bounds < 0) & `40be64c` 09-29 15:45:46 & PU (M1 n = 0 direct ML; P2/\allowbreak{}W3) & supported (regions with both CI upper bounds below 0: 0/4) \\
L2 & LG §4 & P & Augmented combinations R2 or R3 have lower error than R1 at n ∈ \{10, 40, 160\}; otherwise "augmentation is for n = 0 only" & `40be64c` 09-29 15:45:46 & B & rejected (augmentation is for zero labels only); R3 partial \\
L3 & LG §4 & P & Target weighting α ∈ \{10, 100, continued training\}, chosen by within-A cross-validation, gives n* ≤ 40 in ≥ 2 of 4 structure targets & `40be64c` 09-29 15:45:46 & NS & supported \\
L4 & LG §4 & P & Residual ML (R1) outperforms the recalibrated Stefan model (P1); report the minimum n & `40be64c` 09-29 15:45:46 & PU (a2, h25b) & net value of ML with sparse labels (minimum $n$ 10, four-region mean) \\
L5 & LG §4 & P & Neural and generative learners are on par with CatBoost in R1 (manuscript uses L26, WRAPUP 1.5) & `40be64c` 09-29 15:45:46 & NS & learner differences (partial: region count short at $n$ 40, 160) \\
L6 & LG §4 & P & Combinations that pass in Alaskan sub-regions (mode i) also pass in the main 4 regions at the same n (share ≥ 0.5) & `40be64c` 09-29 15:45:46 & NS & not decidable (no passing combination) \\
L7 & LG §4 & P & The covariate-dependent coefficient V1 (or V1r) removes the degradation in Canada and CA-3 (n ≥ 40, point estimate) & `40be64c` 09-29 15:45:46 & NS & rejected \\
L8 & LG §4 & P & With all labels, R1 outperforms P0 (main-4 stratified mean, both weightings) & `40be64c` 09-29 15:45:46 & PU (M1 H13) & supported ($-2.64$ [$-3.47$, $-1.90$]; regions 4/4) \\
L9 & LGX X1 & A & F1a (D0 plus E0·s input) and D0 are indistinguishable & `8513cc2` 09-29 18:17:29 & B & see the Supplementary Table of its family \\
L10 & LGX X1 & C & Residual structure has lower error than physics-input structure (four contrasts at n = 0 and n = 10) & `8513cc2` 09-29 18:17:29 & B & supported \\
L11 & LGX X1 & A & Direct ML with process-model inputs (F1k) does not outperform P0 at n = 0 & `8513cc2` 09-29 18:17:29 & RU (M1 n = 0 direct ML) & see the Supplementary Table of its family \\
L12 & LGX X1 & C & Multiplicative correction RM and additive residual R1 are equivalent (n ∈ \{10, all\}, λ 0.25 and 1.0) & `8513cc2` 09-29 18:17:29 & B & rejected \\
L13 & LGX X1 & A & V2 (log-ratio model) is better than P1 and not different from R1 at n = 10 & `8513cc2` 09-29 18:17:29 & RU for Canada and CA-3 & see the Supplementary Table of its family \\
L14 & LGX X1 & N & Effect of removing physics-related inputs from D0 (D0t, D0c, D0m) & `8513cc2` 09-29 18:17:29 & B & see the Supplementary Table of its family \\
L15 & LGX X2 & C & Stefan pseudo-labels (D1) have lower error than four placebos at n ∈ \{0, 10\} & `8513cc2` 09-29 18:17:29 & RU at n = 0 (M1 H3) & mixed (described by $n$) \\
L16 & LGX X2 & A & Verdicts are insensitive to the pseudo-label ratio r; D1 approaches the physics model as r grows & `8513cc2` 09-29 18:17:29 & B & see the Supplementary Table of its family \\
L17 & LGX X2 & A & R1 and the label-shuffle control R1s are equivalent (n ∈ \{10, all\}) & `8513cc2` 09-29 18:17:29 & B & see the Supplementary Table of its family \\
L18 & LGX X2 & A & With target weighting, the augmented combination is not better than R1 & `8513cc2` 09-29 18:17:29 & B & see the Supplementary Table of its family \\
L19 & LGX X3a & C & Direct ML beats the physics model under random splits (V-R) but not under region hold-out (V-G) & `8513cc2` 09-29 18:17:29 & RU (2026-06 split leakage table, M1) & not supported (no reversal wording; degradation size reported) \\
L20 & LGX X3a & A & D0 RMSE increases in the order V-R < V-B < V-C100 < V-C500 & `8513cc2` 09-29 18:17:29 & RU & supported \\
L21 & LGX X3a & A & Degradation from V-R to V-G is smaller for anchor plus residual than for D0 & `8513cc2` 09-29 18:17:29 & B & supported ($\lambda$ 0.25) \\
L22 & LGX X3b & A & Within a region, R1w outperforms P1w (count of regions out of 3) & `8513cc2` 09-29 18:17:29 & RU for Alaska (M1 X-ak) & lower error in 2 of 3 regions (Lena Delta, Canada) \\
L23 & LGX X3c & A & Labels near scored cells reduce error more than the same number of distant labels & `8513cc2` 09-29 18:17:29 & B & see the Supplementary Table of its family \\
L24 & LGX X5 & A & Residual-learning net value concentrates near labels; recalibration gain does not depend on distance & `8513cc2` 09-29 18:17:29 & B & see the Supplementary Table of its family \\
L25 & LGX X5 & N & R1 90 \% interval coverage and width versus the R0 n = 0 interval (6 cells) & `8513cc2` 09-29 18:17:29 & B & see the Supplementary Table of its family \\
L26 & LGX X6 & A & Each learner minus CatBoost (R1) lies within the 0.5 cm margin at n ∈ \{0, 10, all\}, 5 splits & `8513cc2` 09-29 18:17:29 & B & see the Supplementary Table of its family \\
L27 & LGX X6 & A & The sign of R0 − P0 in Canada does not depend on the learner (≥ 6 of 8 learners) & `8513cc2` 09-29 18:17:29 & B & see the Supplementary Table of its family \\
L28 & LGX X9 & A & 4-way verdicts of core contrasts are unchanged under κ, label-definition, anchor, input-set and matched-year variants & `8513cc2` 09-29 18:17:29 & RU for κ, λ and anchors & see the Supplementary Table of its family \\
L29 & LGX N1 & C & No physics or product baseline is better than P0 at n = 0 & `8513cc2` 09-29 18:17:29 & PU for B:ens − P0 (LG revision 15 (h)) & a baseline better than P0 exists (P* = year-matched Stefan) \\
L30 & LGX N3 & C & The n = 0 direct ML result does not depend on learner capacity (catboost, rf, catboost\_\allowbreak{}tuned) & `8513cc2` 09-29 18:17:29 & RU (M1 n = 0 direct ML) & supported \\
L31 & LGX N4 & A & The full-label R1 − P0 improvement is robust to split composition (50 splits) & `8513cc2` 09-29 18:17:29 & RU for Canada, CA-3, AL-5; Sa & see the Supplementary Table of its family \\
L32 & LGT & A & R1[TabPFN] − R1[CatBoost, same context] lies within the 0.5 cm margin at n ∈ \{0, 10, 40, 160, all\} & `6799b74` 09-29 22:43:12 & RU at n ∈ \{10, 40\} (b4); B at n = 0, 160, all & see the Supplementary Table of its family \\
L33 & LGT & A & TabPFN residual (R1) outperforms P1; report the minimum n & `6799b74` 09-29 22:43:12 & RU at n ∈ \{3, 10, 40\} (b4); B otherwise & see the Supplementary Table of its family \\
L34 & LGT & A & (a) TabPFN direct prediction does not outperform P0 at n = 0; (b) R0[TabPFN] − P0 reported & `6799b74` 09-29 22:43:12 & B; CatBoost rows RU & see the Supplementary Table of its family \\
L35 & LGT & A & With all labels, TabPFN R1 outperforms P0 & `6799b74` 09-29 22:43:12 & B; CatBoost rows RU & see the Supplementary Table of its family \\
L36 & LGT & N & D0[TabPFN] − R1[TabPFN] at five n & `6799b74` 09-29 22:43:12 & B & see the Supplementary Table of its family \\
L37 & LGT & A & Core-contrast verdicts are stable under context variants (d3, d10, rid, tgt) & `6799b74` 09-29 22:43:12 & RU for d3, rid & see the Supplementary Table of its family \\
L1e & LGD & A & L1 recomputed on the extended pools PE1 and PE2 & `6799b74` 09-29 22:43:12 & as L1 for P4; RE under revision 15 (m) & supported (P4, PE1, PE2) \\
L4e & LGD & A & L4 recomputed on PE1 and PE2 (minimum n with region count k/\allowbreak{}N) & `6799b74` 09-29 22:43:12 & as L4 for P4; RE under revision 15 (m) & P4 and PE2: net value with sparse labels; PE1 partial (2/5) \\
L8e & LGD & A & L8 recomputed on PE1 and PE2 & `6799b74` 09-29 22:43:12 & Tibet unblinded (label statistics); RE under revision 15 (m) & supported (P4, PE1, PE2) \\
L38 & LGD & A & In each new eligible region, the region-level verdicts follow the hypothesis directions & `6799b74` 09-29 22:43:12 & Tibet unblinded; NAtlantic, Russia C blind & diverging regions: Central Russia, Tibetan Plateau \\
L39 & LGD & A & In Tibet, changing the label definition leaves core verdicts unchanged & `6799b74` 09-29 22:43:12 & unblinded (label statistics) & robust \\
L40 & LGD & A & Adding new cells to Russia W, Russia E and Canada leaves region-level verdicts unchanged & `6799b74` 09-29 22:43:12 & RU & weakened \\
L41 & LGD & A & New-region verdicts are robust to label-set variants (a)–(c) & `6799b74` 09-29 22:43:12 & as L38 & robust or weakened by region \\
L42 & LGD & A & New-region verdicts are unchanged when other new regions join the source & `6799b74` 09-29 22:43:12 & as L38 & robust or weakened by region \\
L43 & WRAPUP 4 & A & Block-dispersed label placement gives lower P1 error than cell-random placement (n ∈ \{10, 40\}) & `316714c` 09-30 03:02:42 & PU (a2\_\allowbreak{}spread, h29b) & placement effect not established \\
SC1w & WRAPUP 3.3 & A & At T3 and T10, recipe R1 is non-inferior to P0 (0.5 cm) in all 5 independent regions (10 cells) & `316714c` 09-30 03:02:42 & PU; SC1w-P: RE and Sa & safety not confirmed (non-inferior cells 4/10) \\
SC2w & WRAPUP 3.4 & N & Number of regions (Lena, Canada, Alaska x) where R1 beats P1 at n = 40 & `316714c` 09-30 03:02:42 & PU & sentence A (2 of 3 regions at $n$ 40) \\
SC3w & WRAPUP 3.5 & A & A sequential stopping rule (jackknife SE of log E ≤ τ = 0.05) halves labels versus always 40 with error increase ≤ 0.3 cm and no inferior region & `316714c` 09-30 03:02:42 & PU; computed before results (M5.6) & rule not operating (label ratio 0.91) \\
AK1w & WRAPUP 5 & A & Combinations that pass in Alaskan sub-regions with a within-Alaska source (mode i) also pass with an outside-only source (mode x); r ≥ 0.5 & `316714c` 09-30 03:02:42 & PU (h25b, h25x) & majority ($r$ 0.57; 68 pairs, 5 targets) \\
AB1 & WRAPUP 1.1 & abstract & D0 − P0, n = 0, λ 1.0, main 4 (basis L1) & `316714c` 09-30 03:02:42 & PU & see Supplementary Table S3 \\
AB2 & WRAPUP 1.1 & abstract & B:ens − P0, n = 0 (basis L29) & `316714c` 09-30 03:02:42 & PU & see Supplementary Table S3 \\
AB3 & WRAPUP 1.1 & abstract & D1 − D1@shuffle, n = 0, λ 1.0 (basis L15) & `316714c` 09-30 03:02:42 & RU & see Supplementary Table S3 \\
AB4 & WRAPUP 1.1 & abstract & P1 − P0, n = 10 (curve core, L28 core contrast) & `316714c` 09-30 03:02:42 & PU; RE (LGT P0, P1); Sa & see Supplementary Table S3 \\
AB5 & WRAPUP 1.1 & abstract & R1 − P1, n = 10, λ 0.25 (basis L4) & `316714c` 09-30 03:02:42 & PU & see Supplementary Table S3 \\
AB6 & WRAPUP 1.1 & abstract & R2 − R1, n = 10, λ 0.25 (basis L2) & `316714c` 09-30 03:02:42 & B & see Supplementary Table S3 \\
AB7 & WRAPUP 1.1 & abstract & R1 − F1k, n = 10 (basis L10) & `316714c` 09-30 03:02:42 & B & see Supplementary Table S3 \\
AB8 & WRAPUP 1.1 & abstract & R1 − P0, all labels, λ 0.25 (basis L8) & `316714c` 09-30 03:02:42 & PU & see Supplementary Table S3 \\
AB9 & WRAPUP 1.1 & abstract & R1 − P1, all labels, λ 0.25 (basis L4) & `316714c` 09-30 03:02:42 & PU & see Supplementary Table S3 \\
AB10 & WRAPUP 1.1 & abstract & Interval score (α 0.1), rung (iii) − B4, n = 10, Lena, Canada, Alaska x (basis LGU-A1) & `316714c` 09-30 03:02:42 & per LGU §6; RE (LGU revision 4) & see Supplementary Table S3 \\
LGU-A1 & LGU 3.6 & C & The hierarchical predictive interval (rung iii) has a lower interval score than the constant-width calibrated interval B4 at n ∈ \{10, 40\} & `a112cef` 09-29 23:30:02 & NS; RE under revision 4 & see the Supplementary Table of its family \\
LGU-A2 & LGU 3.6 & D & Share of infinite intervals at n = 3 is 0 (implementation check) & `a112cef` 09-29 23:30:02 & RU (C2 wconf) & see the Supplementary Table of its family \\
LGU-A3 & LGU 3.6 & N & Median RMSE, rung (iii) − P1, n ∈ \{3, 10\}, δ 0.5 cm & `a112cef` 09-29 23:30:02 & RU (B2 h32 offsets) & see the Supplementary Table of its family \\
LGU-A4 & LGU 3.6 & gate & CRPS (log scale), rung (v) − (iv), per target, leave-one-region-out & `a112cef` 09-29 23:30:02 & NS & see the Supplementary Table of its family \\
LGU-A5 & LGU 3.6 & D & n = 0 coverage and interval score under hyperprior variants (exploratory) & `a112cef` 09-29 23:30:02 & RU & see the Supplementary Table of its family \\
LGU-A6 & LGU 3.6 & D & Russia W at n ∈ \{3, 10\}: interval score, coverage, width & `a112cef` 09-29 23:30:02 & NS & see the Supplementary Table of its family \\
LGU-A7 & LGU 3.6 & A & Ladder contrasts (ii)−(i), (iii)−(ii), (iv)−(iii), (vi)−(iii), (v)−(iv) & `a112cef` 09-29 23:30:02 & NS & see the Supplementary Table of its family \\
LGU-B1 & LGU 4.4 & A & Pooled generative quantile intervals (nflow, cfm) cover less than 0.85 in new regions (limit check) & `a112cef` 09-29 23:30:02 & unblinded (M1 H15, H17) & see the Supplementary Table of its family \\
LGU-B2 & LGU 4.4 & C & nflow-normalized conformal intervals have a lower interval score than their σ-shuffled placebo at n = 0 (4 conditions) & `a112cef` 09-29 23:30:02 & NS; RE under revision 4 & see the Supplementary Table of its family \\
LGU-B3 & LGU 4.4 & A & Interval score nflow − cbq (5 \% margin) & `a112cef` 09-29 23:30:02 & NS & see the Supplementary Table of its family \\
LGU-B4 & LGU 4.4 & A & Each normalizer − constant, raw and width-matched & `a112cef` 09-29 23:30:02 & NS & see the Supplementary Table of its family \\
LGU-B5 & LGU 4.4 & D & Coverage and width by instrument, season, width quintile and source distance & `a112cef` 09-29 23:30:02 & NS & see the Supplementary Table of its family \\
LGU-B6 & LGU 4.4 & A & nflow − const and cbq − const at n ∈ \{10, 40, 160\} & `a112cef` 09-29 23:30:02 & NS & see the Supplementary Table of its family \\
LGU-C1 & LGU 5 & D & Variogram of the centred log ratio (no verdict) & `a112cef` 09-29 23:30:02 & RU & see the Supplementary Table of its family \\
LGU-C2 & LGU 5 & D & Within-block correlation (no verdict) & `a112cef` 09-29 23:30:02 & RU & see the Supplementary Table of its family \\
LGU-C3 & LGU 5 & D & Share of scored cells with a label within d km (no verdict) & `a112cef` 09-29 23:30:02 & RU & see the Supplementary Table of its family \\
LGU-C4 & LGU 5 & D & Alignment of SAR products with probe labels (no verdict) & `a112cef` 09-29 23:30:02 & RU & see the Supplementary Table of its family \\
LGF-F1 & LGF 3.6 & C & TabICL direct prediction does not outperform P0 at n = 0 (context cap and full source) & `3a4fe84` 09-30 01:44:27 & B; CatBoost rows RU & supported (not superior to P0) \\
LGF-F2 & LGF 3.6 & A & R1[TabICL] − R1[CatBoost, same context] lies within 0.5 cm at five n & `3a4fe84` 09-30 01:44:27 & B & see the Supplementary Table of its family \\
LGF-F3 & LGF 3.6 & N & TabICL versus TabPFN (six contrasts) & `3a4fe84` 09-30 01:44:27 & B; TabPFN column RE at revision 1 & see the Supplementary Table of its family \\
LGF-F4 & LGF 3.6 & A & TabICL residual (R1) outperforms P1; full-label R1 − P0 & `3a4fe84` 09-30 01:44:27 & B; TabPFN rows at n ≤ 40 RU & see the Supplementary Table of its family \\
LGF-F5 & LGF 3.6 & N & n = 0 R0[TabICL] − P0 & `3a4fe84` 09-30 01:44:27 & B & see the Supplementary Table of its family \\
LGF-F6 & LGF 3.6 & N & Effect of the context cap (full source minus 10,000 rows) & `3a4fe84` 09-30 01:44:27 & B & see the Supplementary Table of its family \\
LGF-N1 & LGF 4.5 & C & Direct ML from four tuned discriminative networks does not outperform P0 at n = 0 & `3a4fe84` 09-30 01:44:27 & PU (default-version direction, M1) & supported (not superior to P0; four networks) \\
LGF-N2 & LGF 4.5 & A & Tuned and default networks differ by less than the 0.5 cm margin (9 contrasts per learner) & `3a4fe84` 09-30 01:44:27 & B & see the Supplementary Table of its family \\
LGF-N2s & LGF 4.5 & A & Core-contrast verdicts are the same for default and tuned networks & `3a4fe84` 09-30 01:44:27 & B & see the Supplementary Table of its family \\
LGF-N3 & LGF 4.5 & A & R1[tuned network] − R1[tuned CatBoost] lies within 0.5 cm at n ∈ \{0, 10, all\} & `3a4fe84` 09-30 01:44:27 & B; default rows RU & see the Supplementary Table of its family \\
LGF-N4 & LGF 4.5 & D & Selection tables and tuning diagnostics (no hypothesis) & `3a4fe84` 09-30 01:44:27 & NS & see the Supplementary Table of its family \\
WF0 & WF & D & Zero-label risk table and misspecification diagnostics over 30 targets (post hoc description) & 4168331 10-01 & post hoc & descriptive (Supplementary Table S5) \\
WF1-a & WF & P & Within-region R1 and R2 ($\lambda$ by CV) have lower error than the best physics calibration at $n \geq 1000$ (Alaska, five splits) & 4168331 10-01 & designed after viewing & rejected \\
WF1-b & WF & A & Nine SAR covariates add value to within-region R1 & 4168331 10-01 & designed after viewing & rejected \\
WF1-c & WF & A & Recipes with fewer higher-error $n$ than direct ML & 4168331 10-01 & designed after viewing & partially supported (R2, D1) \\
WF2-a & WF & P & Placement strategies S2 to S7 have lower error than random cells (S1) within regions & 4168331 10-01 & designed after viewing & supported (S3 at $n$ 100; size below 0.5 cm) \\
WF2-b & WF & A & The best Alaskan strategy (S3) transfers to the Lena Delta and Canada & 4168331 10-01 & designed after viewing & rejected \\
WF3-a & WF & A & A covariate-dependent coefficient (Pc) has lower error than P1 at $n \geq 500$ & 4168331 10-01 & designed after viewing & rejected \\
WF4-a & WF & P & Selection rule W is non-inferior to the fixed recipe R1($\lambda$ 0.25) (margin 0.5 cm) & 4168331 10-01 & designed after viewing & partial (regions 2/4): non-inferior at 40, 160 and all labels \\
WF4-b & WF & A & W has lower error than P1 in more targets than R1(0.25) & 4168331 10-01 & designed after viewing & rejected \\
WF4-c & WF & P & The ten-label bias is rank-correlated with the all-label gain of W over P0 & 4168331 10-01 & designed after viewing & supported ($\rho$ 0.38 [0.17, 0.55]) \\
WF6-a & WF & P & Within-region recipes ($\lambda$ by CV) have lower error than P1 at $n$ 200 to all (Alaska; Lena Delta; Canada) & 1b42b4d 10-01 & designed after viewing & Alaska supported; Lena Delta rejected; Canada partial, rejected \\
WF6-b & WF & A & Stratified mean of three targets, R2 versus P1 at $n \geq 500$ & 1b42b4d 10-01 & designed after viewing & partial (regions 2/3): rejected \\
WF8-a & WF & P & In transfer, S4 has lower error than S1 for R1 at $n$ 40 & 1b42b4d 10-01 & designed after viewing & supported (Canada) \\
WF8-b & WF & A & In transfer, S2 and S4 are not worse than S1 for R1 at $n$ 10 & 1b42b4d 10-01 & designed after viewing & supported \\
WF9-a & WF & P & Within-grid gain of R1 ($\lambda$ by CV) over P1 within regions & 8180632 10-02 & designed after viewing & partial (regions 2/3): rejected \\
WF9-c & WF & P & Within-grid gain of R1 ($\lambda$ 0.25) over P1 in transfer & 8180632 10-02 & designed after viewing & partial (regions 2/4): partially supported \\
WF9 soil & WF & A & Soil-$\sqrt{\mathrm{TDD}}$ grouping sensitivity of WF9 & 7da0064 10-02 & designed after viewing & see Supplementary Table S10 \\
WF10-a to c & WF & P & Climate extrapolation (warm variant) & 8180632 10-02 & designed after viewing & not decidable (contrasts 0/2) \\
XA & FINAL\_\allowbreak{}BATCH 2.1 & post hoc & Rank correlation of the recalibration share, the selection-rule share and the ML share beyond recalibration with coefficient error & 2678100 10-04 & replication (unblinded) & not established (XA-1, XA-2, XA-3; Supplementary Table S13) \\
XB & FINAL\_\allowbreak{}BATCH 2.2 & see plan & Label-weighted stacking of physics models and products & 2678100 10-04 & PU & XB-1, XB-2 rejected (undetermined); XB-3 partially supported (lower error at $n$ 500 and 1,000; not retained when CCI v5 was added as a candidate, sensitivity edition); XB-4 supported (non-inferior) \\
XC & FINAL\_\allowbreak{}BATCH 2.3 & see plan & End-to-end workflow on new splits of the same regions & 2678100 10-04 & PU (XC-1, XC-2, XC-5); B (XC-F3) & Algorithm P = S1; XC-1b undetermined; XC-1c lower error (Canada; Lena Delta undetermined); XC-2a, XC-2b, XC-2c non-inferior for the pooled mean (XC-2s lower error at $n$ 40 and 160; regions over 0.5 cm worse: E Russia, Alaska, Lena Delta); XC-5b, XC-5c not tested; XC-F3 not tested \textcolor{blue}{[PENDING: after 11 October 2026]} \\
XD & FINAL\_\allowbreak{}BATCH 2.4 & exploratory & Placement algorithm and learned placement policy & 2678100 10-04 & PU (XD-1 to XD-3); B (XD-4) & Algorithm P = S1; XD-1 rejected; XD-2 rejected (higher error in Alaska and Canada); XD-3 non-inferior 8, higher error 2, not established 6; XD-4, XD-5 descriptive \\
XE & FINAL\_\allowbreak{}BATCH 2.5 & see plan & Sub-grid public covariates within regions & 2678100 10-04 & PU (xh0) & XE-c (xh0, auxiliary) equivalent; XE-e diagnostic (no verdict); XE-a, XE-b \textcolor{blue}{[PENDING]} \\
XF & FINAL\_\allowbreak{}BATCH 2.6 & see plan & New independent regions from public data & 2678100 10-04 & B (new regions) & \textcolor{blue}{[PENDING]} \\
XG & FINAL\_\allowbreak{}BATCH 2.7 & see plan & Existing ALT products on the same blocks & 2678100 10-04 (WRAPUP 10 deviation) & B (Wei rows); PU (CCI rows) & XG-1c higher error; XG-2c, XG-3c lower error; XG-1w to XG-3w undetermined \\
XH & FINAL\_\allowbreak{}BATCH 2.8 & descriptive & Validation ladder & 2678100 10-04 & RU (Alaska); B (Lena Delta, Canada) & descriptive (no verdict words) \\
XI & FINAL\_\allowbreak{}BATCH 2.9 & see plan & Climate-extrapolation retest (warm-block spatial proxy) & 2678100 10-04 & B (Canada); RU (Alaska) & XI-a, XI-b rejected (undetermined); XI-c supported (lower error, non-inferior) \\
XJ & FINAL\_\allowbreak{}BATCH 2.10 & see plan & Temperature-derived auxiliary labels & 2678100 10-04 & PU (L39) & XJ-1, XJ-2 not tested (licence basis not recorded); XJ-3 lower error (six rows) \\
XK & addendum XK & descriptive & Error by grid size within regions & d721d24 10-05 & designed after viewing & descriptive (no verdict words) \\
XL & addendum XL & descriptive & 1 km maps compared with ALT products & d721d24 10-05 & designed after viewing & descriptive (no verdict; no accuracy claim) \\
\end{longtable}}

\clearpage
\phantomsection\label{tab:S3}
\noindent\textbf{Supplementary Table S3.} Abstract contrast bundle AB1 to AB10.\par\smallskip
\noindent{\footnotesize Four main regions, stratified means; AB1 to AB9 with 10,000 block resamples in both weightings, AB10 (interval score of the hierarchical predictive interval against the calibrated constant-width interval B4; Lena Delta, Canada, Alaska) with 1,000 resamples from LGU-A1, for which the copied bundle row has no CI. $P$ values are uncorrected one-sided bootstrap values; Holm adjustment over the ten contrasts. Common-CI verdict: the verdict under the common resampling of both arms; 'yes' marks verdicts that depend on the split-independence assumption. The abstract uses verdict words only for these contrasts. Source: paper/claims/C1\_label0\_safety/tables/lgw\_bundle.csv (copy of data/processed/lgw/lgw\_bundle.csv), rows scope = MEAN. Method codes (internal identifiers allowed in the Supplementary Information): P0 source-coefficient Stefan model; P* year-matched Stefan model; P1 recalibrated Stefan model ($\kappa = 10$); P2 local least-squares Stefan model; P1* or P1@ed soil-property recalibrated Stefan model; R0 source anchor plus residual ML; R1 recalibrated anchor plus residual ML; R2 augmentation plus anchor plus residual; D0 direct ML; D1 physics pseudo-label augmentation; F1a, F1k, F1n physics-input ML; RM multiplicative residual; Re Stefan-CCI mean anchor plus residual; W selection rule (cross-validation within the target labels); S1 random cells, S2 block stratification, S3 covariate-cluster stratification, S4 covariate max-min distance, S5 source-extrapolation first, S6 sequential selection by prediction variance, S7 $\sqrt{\mathrm{TDD}}$ first. Target suffixes: $|x$ parent region excluded from the source, $|i$ rest of the parent kept, $|r$ within-region design. $n$ = number of target labels; all = all labels of half A. Verdicts: lower error, higher error, equivalent ($\pm$0.5~cm), undetermined; both 95\% block-bootstrap CIs (cell-weighted and block-equal) must agree. \dag~a note of the source row in Korean was omitted; see the cited README section.}\par
{\fontsize{7.5}{9}\selectfont\setlength{\tabcolsep}{4pt}\setlength{\LTpre}{2pt}\setlength{\LTpost}{4pt}
\begin{longtable}{@{}>{\raggedright\arraybackslash}p{\dimexpr 0.047\linewidth-8pt\relax}>{\raggedright\arraybackslash}p{\dimexpr 0.152\linewidth-8pt\relax}>{\raggedright\arraybackslash}p{\dimexpr 0.134\linewidth-8pt\relax}>{\raggedright\arraybackslash}p{\dimexpr 0.134\linewidth-8pt\relax}>{\raggedright\arraybackslash}p{\dimexpr 0.062\linewidth-8pt\relax}>{\raggedright\arraybackslash}p{\dimexpr 0.062\linewidth-8pt\relax}>{\raggedright\arraybackslash}p{\dimexpr 0.062\linewidth-8pt\relax}>{\raggedright\arraybackslash}p{\dimexpr 0.108\linewidth-8pt\relax}>{\raggedright\arraybackslash}p{\dimexpr 0.108\linewidth-8pt\relax}>{\raggedright\arraybackslash}p{\dimexpr 0.131\linewidth-8pt\relax}@{}}
\toprule
\textbf{AB} & \textbf{Contrast} & \textbf{$\Delta$ cell-weighted [95\% CI] (cm)} & \textbf{$\Delta$ block-equal [95\% CI] (cm)} & \textbf{$P$ cell} & \textbf{$P$ block} & \textbf{Holm $P$} & \textbf{Verdict} & \textbf{Common-CI verdict} & \textbf{Split-dependent} \\
\midrule\endfirsthead
\toprule
\textbf{AB} & \textbf{Contrast} & \textbf{$\Delta$ cell-weighted [95\% CI] (cm)} & \textbf{$\Delta$ block-equal [95\% CI] (cm)} & \textbf{$P$ cell} & \textbf{$P$ block} & \textbf{Holm $P$} & \textbf{Verdict} & \textbf{Common-CI verdict} & \textbf{Split-dependent} \\
\midrule\endhead
\bottomrule\endlastfoot
AB1 & D0-P0|\allowbreak{}n0 & +2.25 [+1.10, +3.42] & +1.62 [+0.52, +2.69] & 0.0004 & 0.005 & 0.023 & higher error & undetermined & yes \\
AB2 & B:ens-P0|\allowbreak{}n0 & −2.73 [−3.33, −0.91] & +0.05 [−0.94, +1.01] & 0.002 & 0.912 & 1.000 & undetermined & undetermined &  \\
AB3 & D1-D1@shuffle|\allowbreak{}n0 & −1.63 [−1.98, −1.04] & −1.05 [−1.51, −0.57] & 0.0001 & 0.0001 & 0.001 & lower error & lower error &  \\
AB4 & P1-P0|\allowbreak{}n10 & −2.45 [−3.04, −1.88] & −2.62 [−3.30, −1.87] & 0.0001 & 0.0001 & 0.001 & lower error & lower error &  \\
AB5 & R1-P1|\allowbreak{}n10 & −0.18 [−0.53, −0.01] & −0.29 [−0.51, −0.08] & 0.034 & 0.008 & 0.136 & lower error & undetermined & yes \\
AB6 & R2-R1|\allowbreak{}n10 & −0.13 [−0.31, +0.22] & −0.06 [−0.31, +0.20] & 0.604 & 0.656 & 1.000 & equivalent & equivalent &  \\
AB7 & R1-F1k|\allowbreak{}n10 & −3.74 [−5.23, −2.75] & −4.07 [−5.14, −2.98] & 0.0001 & 0.0001 & 0.001 & lower error & lower error &  \\
AB8 & R1-P0|\allowbreak{}all & −2.64 [−3.48, −1.85] & −2.64 [−3.51, −1.69] & 0.0001 & 0.0001 & 0.001 & lower error & lower error &  \\
AB9 & R1-P1|\allowbreak{}all & −0.40 [−0.76, −0.22] & −0.51 [−0.74, −0.27] & 0.0001 & 0.0001 & 0.001 & lower error & lower error &  \\
AB10 & interval score ($\alpha$ 0.1), rung (iii) $-$ B4, n10 & +9.51 [n.a.] & n.a. & 0.072 & 0.092 & 0.277 & undetermined & n.a. &  \\
\end{longtable}}

\clearpage
\phantomsection\label{tab:S4}
\noindent\textbf{Supplementary Table S4.} Transfer without target labels: learners, physics information and physics baselines.\par\smallskip
\noindent{\footnotesize Direct ML minus the source-coefficient Stefan model for ten learners (four main regions), regional rows, the three-region auxiliary column, new regions, placebo pseudo-labels, low-weight residuals and the physics baselines. Holm families differ by row (abstract bundle, LGX, LGT, LGF). Learners 1 to 4 ran on the cloud platform, the other learners and companion rows on the local GPU server; differences between platforms are not computed. Method codes (internal identifiers allowed in the Supplementary Information): P0 source-coefficient Stefan model; P* year-matched Stefan model; P1 recalibrated Stefan model ($\kappa = 10$); P2 local least-squares Stefan model; P1* or P1@ed soil-property recalibrated Stefan model; R0 source anchor plus residual ML; R1 recalibrated anchor plus residual ML; R2 augmentation plus anchor plus residual; D0 direct ML; D1 physics pseudo-label augmentation; F1a, F1k, F1n physics-input ML; RM multiplicative residual; Re Stefan-CCI mean anchor plus residual; W selection rule (cross-validation within the target labels); S1 random cells, S2 block stratification, S3 covariate-cluster stratification, S4 covariate max-min distance, S5 source-extrapolation first, S6 sequential selection by prediction variance, S7 $\sqrt{\mathrm{TDD}}$ first. Target suffixes: $|x$ parent region excluded from the source, $|i$ rest of the parent kept, $|r$ within-region design. $n$ = number of target labels; all = all labels of half A. Verdicts: lower error, higher error, equivalent ($\pm$0.5~cm), undetermined; both 95\% block-bootstrap CIs (cell-weighted and block-equal) must agree. \dag~a note of the source row in Korean was omitted; see the cited README section.}\par

\par\medskip\noindent{\small\textit{a. Direct ML minus source-coefficient Stefan, zero labels, four main regions (C1 README 2.1)}}\par\nopagebreak
{\fontsize{7.5}{9}\selectfont\setlength{\tabcolsep}{4pt}\setlength{\LTpre}{2pt}\setlength{\LTpost}{4pt}
\begin{longtable}{@{}>{\raggedright\arraybackslash}p{\dimexpr 0.073\linewidth-8pt\relax}>{\raggedright\arraybackslash}p{\dimexpr 0.216\linewidth-8pt\relax}>{\raggedright\arraybackslash}p{\dimexpr 0.099\linewidth-8pt\relax}>{\raggedright\arraybackslash}p{\dimexpr 0.086\linewidth-8pt\relax}>{\raggedright\arraybackslash}p{\dimexpr 0.112\linewidth-8pt\relax}>{\raggedright\arraybackslash}p{\dimexpr 0.105\linewidth-8pt\relax}>{\raggedright\arraybackslash}p{\dimexpr 0.191\linewidth-8pt\relax}>{\raggedright\arraybackslash}p{\dimexpr 0.118\linewidth-8pt\relax}@{}}
\toprule
\textbf{\#} & \textbf{Learner} & \textbf{$\Delta$ cell-weighted [CI]} & \textbf{$\Delta$ block-equal [CI]} & \textbf{Verdict} & \textbf{Holm $P$} & \textbf{Common CI (cell / block), verdict} & \textbf{Dependence} \\
\midrule\endfirsthead
\toprule
\textbf{\#} & \textbf{Learner} & \textbf{$\Delta$ cell-weighted [CI]} & \textbf{$\Delta$ block-equal [CI]} & \textbf{Verdict} & \textbf{Holm $P$} & \textbf{Common CI (cell / block), verdict} & \textbf{Dependence} \\
\midrule\endhead
\bottomrule\endlastfoot
1 & CatBoost main learner(catboost\_\allowbreak{}lo) & +2.25 [1.10, 3.42] & +1.62 [0.52, 2.69] & higher error & 0.023(abstract bundle m = 10) & [0.31, 4.19] /\allowbreak{} [−0.21, 3.43], undetermined & depends on split independence \\
2 & CatBoost higher capacity(catboost) & +1.94 [0.88, 3.43] & +1.56 [0.30, 2.76] & higher error & 0.0032 & [0.24, 4.32] /\allowbreak{} [−0.45, 3.57], undetermined & depends on split independence \\
3 & CatBoost source-CV tuning(catboost\_\allowbreak{}tuned) & +1.59 [0.65, 2.69] & +1.63 [0.56, 2.64] & higher error & 0.0032 & [0.07, 3.37] /\allowbreak{} [−0.08, 3.33], undetermined & depends on split independence \\
4 & random forest(rf) & +2.63 [1.47, 3.96] & +2.03 [0.99, 3.06] & higher error & 0.0003 & [0.72, 4.75] /\allowbreak{} [0.31, 3.75], higher error &  \\
5 & TabPFN v2 & +1.02 [0.23, 2.35] & +1.22 [0.27, 2.15] & higher error(holm\_\allowbreak{}note = 'significant before correction') & 0.1168 & [−0.27, 3.13] /\allowbreak{} [−0.34, 2.75], undetermined & depends on split independence \\
6a & TabICL v2, context cap 10,000 rows & +0.93 [0.28, 3.12] & +2.43 [1.35, 3.55] & higher error & 0.0232 & [−0.36, 4.20] /\allowbreak{} [0.72, 4.20], undetermined & depends on split independence \\
6b & TabICL v2, full source context & +0.85 [0.07, 3.19] & +2.22 [1.09, 3.45] & higher error & 0.0374 & [−0.60, 4.26] /\allowbreak{} [0.53, 4.10], undetermined & depends on split independence \\
7 & MLP(leave-one-source-region-out tuning) & +2.45 [1.62, 4.11] & +2.05 [0.95, 3.19] & higher error & 0.0004 & [0.95, 4.90] /\allowbreak{} [0.25, 3.94], higher error &  \\
8 & multi-head MLP(tabm*, not the public TabM) & +2.44 [1.58, 4.73] & +3.94 [2.68, 5.29] & higher error & 0.0004 & [0.87, 5.98] /\allowbreak{} [1.92, 6.33], higher error &  \\
9 & reduced FT-Transformer(ftt*) & +4.13 [3.28, 5.63] & +4.58 [3.44, 5.76] & higher error & 0.0004 & [2.62, 6.52] /\allowbreak{} [2.65, 6.59], higher error &  \\
10 & RealMLP(realmlp*) & +5.21 [4.07, 7.17] & +5.91 [4.60, 7.25] & higher error & 0.0004 & [3.18, 8.31] /\allowbreak{} [3.73, 8.19], higher error &  \\
auxiliary & physics-output input direct ML(F1k) & +2.62 [1.38, 4.36] & +2.73 [1.57, 3.85] & higher error & 0.0023 & [0.40, 5.53] /\allowbreak{} [0.82, 4.62], higher error &  \\
companion & same context CatBoost(local, D0[C]) & +2.64 [1.36, 3.72] & +1.79 [0.69, 2.81] & higher error & none & [0.59, 4.39] /\allowbreak{} [−0.005, 3.51], undetermined & depends on split independence \\
\end{longtable}}
\par\medskip\noindent{\small\textit{b. AB1 regional rows (C1 README 2.2)}}\par\nopagebreak
{\fontsize{7.5}{9}\selectfont\setlength{\tabcolsep}{4pt}\setlength{\LTpre}{2pt}\setlength{\LTpost}{4pt}
\begin{longtable}{@{}>{\raggedright\arraybackslash}p{\dimexpr 0.084\linewidth-8pt\relax}>{\raggedright\arraybackslash}p{\dimexpr 0.160\linewidth-8pt\relax}>{\raggedright\arraybackslash}p{\dimexpr 0.160\linewidth-8pt\relax}>{\raggedright\arraybackslash}p{\dimexpr 0.118\linewidth-8pt\relax}>{\raggedright\arraybackslash}p{\dimexpr 0.360\linewidth-8pt\relax}>{\raggedright\arraybackslash}p{\dimexpr 0.118\linewidth-8pt\relax}@{}}
\toprule
\textbf{Region} & \textbf{$\Delta$ cell-weighted [CI]} & \textbf{$\Delta$ block-equal [CI]} & \textbf{Verdict} & \textbf{Common CI (cell / block), verdict} & \textbf{Dependence} \\
\midrule\endfirsthead
\toprule
\textbf{Region} & \textbf{$\Delta$ cell-weighted [CI]} & \textbf{$\Delta$ block-equal [CI]} & \textbf{Verdict} & \textbf{Common CI (cell / block), verdict} & \textbf{Dependence} \\
\midrule\endhead
\bottomrule\endlastfoot
Lena Delta & +4.02 [1.98, 5.62] & +0.20 [−1.66, 1.97] & undetermined & [0.61, 6.58] /\allowbreak{} [−3.00, 3.18], undetermined &  \\
Canada & −1.56 [−2.70, 1.41] & −0.01 [−1.61, 1.62] & undetermined & [−3.27, 3.23] /\allowbreak{} [−2.85, 2.63], undetermined &  \\
W Russia & +2.05 [−0.94, 4.49] & +2.78 [−0.03, 5.28] & undetermined & [−3.28, 6.35] /\allowbreak{} [−2.11, 7.13], undetermined &  \\
E Russia & +4.48 [1.41, 6.51] & +3.51 [1.14, 5.91] & higher error & [−0.32, 7.50] /\allowbreak{} [−0.10, 7.20], undetermined & depends on split independence \\
\end{longtable}}
\par\medskip\noindent{\small\textit{c. New independent regions, licence-verified edition (L1e, PE2, zero labels) (C1 README 2.2)}}\par\nopagebreak
{\fontsize{7.5}{9}\selectfont\setlength{\tabcolsep}{4pt}\setlength{\LTpre}{2pt}\setlength{\LTpost}{4pt}
\begin{longtable}{@{}>{\raggedright\arraybackslash}p{\dimexpr 0.243\linewidth-8pt\relax}>{\raggedright\arraybackslash}p{\dimexpr 0.130\linewidth-8pt\relax}>{\raggedright\arraybackslash}p{\dimexpr 0.130\linewidth-8pt\relax}>{\raggedright\arraybackslash}p{\dimexpr 0.066\linewidth-8pt\relax}>{\raggedright\arraybackslash}p{\dimexpr 0.073\linewidth-8pt\relax}>{\raggedright\arraybackslash}p{\dimexpr 0.152\linewidth-8pt\relax}>{\raggedright\arraybackslash}p{\dimexpr 0.204\linewidth-8pt\relax}@{}}
\toprule
\textbf{Region} & \textbf{$\Delta$ cell-weighted [CI]} & \textbf{$\Delta$ block-equal [CI]} & \textbf{sig} & \textbf{sig (common)} & \textbf{Flags} & \textbf{Cross-environment} \\
\midrule\endfirsthead
\toprule
\textbf{Region} & \textbf{$\Delta$ cell-weighted [CI]} & \textbf{$\Delta$ block-equal [CI]} & \textbf{sig} & \textbf{sig (common)} & \textbf{Flags} & \textbf{Cross-environment} \\
\midrule\endhead
\bottomrule\endlastfoot
Central Russia(Russia\_\allowbreak{}C|\allowbreak{}x) & +2.75 [2.11, 4.09] & +4.83 [3.17, 6.60] & worse & worse & none & cross-environment; (i) CatBoost not passed \\
Tibetan Plateau(Tibet\_\allowbreak{}LGD|\allowbreak{}x) & −59.98 [−62.90, −57.15] & −62.91 [−65.38, −60.32] & improve & improve & predicted outcome(unblinded); unblinded(label statistics viewed) & cross-environment; (i) CatBoost not passed \\
\end{longtable}}
\par\medskip\noindent{\small\textit{d. Three-region auxiliary column (Lena Delta, Canada, Alaska) (C1 README 2.2)}}\par\nopagebreak
{\fontsize{7.5}{9}\selectfont\setlength{\tabcolsep}{4pt}\setlength{\LTpre}{2pt}\setlength{\LTpost}{4pt}
\begin{longtable}{@{}>{\raggedright\arraybackslash}p{\dimexpr 0.167\linewidth-8pt\relax}>{\raggedright\arraybackslash}p{\dimexpr 0.125\linewidth-8pt\relax}>{\raggedright\arraybackslash}p{\dimexpr 0.120\linewidth-8pt\relax}>{\raggedright\arraybackslash}p{\dimexpr 0.105\linewidth-8pt\relax}>{\raggedright\arraybackslash}p{\dimexpr 0.281\linewidth-8pt\relax}>{\raggedright\arraybackslash}p{\dimexpr 0.105\linewidth-8pt\relax}>{\raggedright\arraybackslash}p{\dimexpr 0.098\linewidth-8pt\relax}@{}}
\toprule
\textbf{Learner} & \textbf{$\Delta$ cell-weighted [CI]} & \textbf{$\Delta$ block-equal [CI]} & \textbf{Verdict} & \textbf{Common CI, verdict} & \textbf{Dependence} & \textbf{Four-region $\Delta$} \\
\midrule\endfirsthead
\toprule
\textbf{Learner} & \textbf{$\Delta$ cell-weighted [CI]} & \textbf{$\Delta$ block-equal [CI]} & \textbf{Verdict} & \textbf{Common CI, verdict} & \textbf{Dependence} & \textbf{Four-region $\Delta$} \\
\midrule\endhead
\bottomrule\endlastfoot
catboost\_\allowbreak{}lo & +7.75 [5.49, 10.11] & +6.23 [5.04, 7.40] & higher error & [4.11, 11.45] /\allowbreak{} [4.23, 8.17], higher error &  & +2.25 \\
catboost & +8.09 [5.89, 9.44] & +5.69 [4.63, 6.74] & higher error & [4.95, 10.28] /\allowbreak{} [3.87, 7.46], higher error &  & +1.94 \\
rf & +6.91 [5.37, 8.07] & +5.17 [4.16, 6.16] & higher error & [4.51, 8.85] /\allowbreak{} [3.39, 6.88], higher error &  & +2.63 \\
catboost\_\allowbreak{}tuned & +4.06 [2.62, 4.85] & +2.61 [1.79, 3.41] & higher error & [2.09, 5.30] /\allowbreak{} [1.22, 3.94], higher error &  & +1.59 \\
TabPFN v2 & +7.83 [5.40, 8.66] & +4.96 [3.94, 5.98] & higher error & [4.67, 9.05] /\allowbreak{} [3.23, 6.62], higher error &  & +1.02 \\
TabICL v2, context cap 10,000 rows & +12.25 [7.55, 14.21] & +7.91 [6.55, 9.26] & higher error & [5.53, 15.38] /\allowbreak{} [5.72, 10.15], higher error &  & +0.93 \\
TabICL v2, full source context & +0.31 [−0.68, 2.37] & +0.70 [−0.35, 1.78] & undetermined & [−1.41, 3.45] /\allowbreak{} [−1.01, 2.58], undetermined &  & +0.85 \\
mlp* & +9.28 [6.30, 10.46] & +6.82 [5.75, 7.90] & higher error & [5.45, 11.12] /\allowbreak{} [5.01, 8.63], higher error &  & +2.45 \\
tabm* & +9.26 [6.64, 10.77] & +7.91 [6.61, 9.18] & higher error & [5.29, 11.77] /\allowbreak{} [5.78, 10.06], higher error &  & +2.44 \\
ftt* & +1.17 [0.48, 2.14] & +0.94 [0.32, 1.55] & higher error & [0.03, 2.76] /\allowbreak{} [−0.14, 2.04], undetermined & depends on split independence & +4.13 \\
realmlp* & +6.54 [4.67, 7.37] & +6.10 [5.07, 7.14] & higher error & [3.78, 8.01] /\allowbreak{} [4.36, 7.82], higher error &  & +5.21 \\
auxiliary F1k & +7.69 [5.48, 10.33] & +7.05 [5.78, 8.33] & higher error & [3.96, 11.97] /\allowbreak{} [4.92, 9.12], higher error &  & +2.62 \\
companion D0[C] & +8.00 [5.55, 10.36] & +6.01 [4.88, 7.15] & higher error & [4.28, 11.50] /\allowbreak{} [4.12, 7.82], higher error &  & +2.64 \\
\end{longtable}}
\par\medskip\noindent{\small\textit{e. Limit row: TabICL v2 with ten labels (C1 README 2.2)}}\par\nopagebreak
{\fontsize{7.5}{9}\selectfont\setlength{\tabcolsep}{4pt}\setlength{\LTpre}{2pt}\setlength{\LTpost}{4pt}
\begin{longtable}{@{}>{\raggedright\arraybackslash}p{\dimexpr 0.250\linewidth-8pt\relax}>{\raggedright\arraybackslash}p{\dimexpr 0.134\linewidth-8pt\relax}>{\raggedright\arraybackslash}p{\dimexpr 0.134\linewidth-8pt\relax}>{\raggedright\arraybackslash}p{\dimexpr 0.085\linewidth-8pt\relax}>{\raggedright\arraybackslash}p{\dimexpr 0.281\linewidth-8pt\relax}>{\raggedright\arraybackslash}p{\dimexpr 0.116\linewidth-8pt\relax}@{}}
\toprule
\textbf{Filter} & \textbf{$\Delta$ cell-weighted [CI]} & \textbf{$\Delta$ block-equal [CI]} & \textbf{Verdict} & \textbf{Common CI, verdict} & \textbf{Dependence} \\
\midrule\endfirsthead
\toprule
\textbf{Filter} & \textbf{$\Delta$ cell-weighted [CI]} & \textbf{$\Delta$ block-equal [CI]} & \textbf{Verdict} & \textbf{Common CI, verdict} & \textbf{Dependence} \\
\midrule\endhead
\bottomrule\endlastfoot
contrast=D0[I]-P0|\allowbreak{}n10, scope=MEAN & −2.91 [−3.58, −1.11] & −1.49 [−2.65, −0.21] & lower error & [−4.07, −0.31] /\allowbreak{} [−3.15, 0.39], undetermined & depends on split independence \\
contrast=D0[I]-P0|\allowbreak{}n10, scope=MEAN3 & +5.98 [2.26, 7.22] & +3.46 [2.52, 4.43] & higher error & [1.85, 7.53] /\allowbreak{} [2.02, 4.99], higher error &  \\
\end{longtable}}
\par\medskip\noindent{\small\textit{f. Source of the physics information: placebo pseudo-labels (AB3, L15) (C1 README 2.3)}}\par\nopagebreak
{\fontsize{7.5}{9}\selectfont\setlength{\tabcolsep}{4pt}\setlength{\LTpre}{2pt}\setlength{\LTpost}{4pt}
\begin{longtable}{@{}>{\raggedright\arraybackslash}p{\dimexpr 0.308\linewidth-8pt\relax}>{\raggedright\arraybackslash}p{\dimexpr 0.212\linewidth-8pt\relax}>{\raggedright\arraybackslash}p{\dimexpr 0.212\linewidth-8pt\relax}>{\raggedright\arraybackslash}p{\dimexpr 0.135\linewidth-8pt\relax}>{\raggedright\arraybackslash}p{\dimexpr 0.134\linewidth-8pt\relax}@{}}
\toprule
\textbf{Contrast} & \textbf{$\Delta$ cell-weighted [CI]} & \textbf{$\Delta$ block-equal [CI]} & \textbf{Verdict} & \textbf{Holm $P$} \\
\midrule\endfirsthead
\toprule
\textbf{Contrast} & \textbf{$\Delta$ cell-weighted [CI]} & \textbf{$\Delta$ block-equal [CI]} & \textbf{Verdict} & \textbf{Holm $P$} \\
\midrule\endhead
\bottomrule\endlastfoot
AB3 D1 − shuffled pseudo-labels, n 0 & −1.63 [−1.98, −1.04] & −1.05 [−1.51, −0.57] & lower error & 0.001(abstract bundle m = 10) \\
L15 shuffle, n 0 & −1.63 [−1.99, −1.04] & −1.05 [−1.52, −0.58] & lower error & 0.003 \\
L15 shuffle, n 10 & −1.15 [−1.43, −0.73] & −0.68 [−1.06, −0.29] & lower error & 0.003 \\
L15 const\_\allowbreak{}t, n 0 & −1.68 [−2.06, −1.12] & −1.06 [−1.55, −0.56] & lower error & 0.003 \\
L15 const\_\allowbreak{}t, n 10 & −1.16 [−1.46, −0.76] & −0.74 [−1.11, −0.37] & lower error & 0.003 \\
L15 const\_\allowbreak{}src, n 0 & −3.63 [−4.33, −2.44] & −2.81 [−3.49, −2.11] & lower error & 0.003 \\
L15 const\_\allowbreak{}src, n 10 & −2.21 [−2.89, −1.27] & −1.54 [−2.23, −0.86] & lower error & 0.003 \\
L15 tddlin, n 0 & −0.48 [−0.62, −0.27] & −0.41 [−0.61, −0.21] & lower error(size below 0.5 cm) & 0.0034 \\
L15 tddlin, n 10 & −0.08 [−0.18, 0.05] & −0.03 [−0.15, 0.09] & equivalent & 1 \\
L15 shuffle, n 0, 3 regions & −3.78 [−4.31, −1.97] & −1.19 [−1.78, −0.61] & lower error & none \\
\end{longtable}}
\par\medskip\noindent{\small\textit{g. Low-weight residuals and combination structure at zero labels (C1 README 2.4)}}\par\nopagebreak
{\fontsize{7.5}{9}\selectfont\setlength{\tabcolsep}{4pt}\setlength{\LTpre}{2pt}\setlength{\LTpost}{4pt}
\begin{longtable}{@{}>{\raggedright\arraybackslash}p{\dimexpr 0.177\linewidth-8pt\relax}>{\raggedright\arraybackslash}p{\dimexpr 0.111\linewidth-8pt\relax}>{\raggedright\arraybackslash}p{\dimexpr 0.097\linewidth-8pt\relax}>{\raggedright\arraybackslash}p{\dimexpr 0.104\linewidth-8pt\relax}>{\raggedright\arraybackslash}p{\dimexpr 0.186\linewidth-8pt\relax}>{\raggedright\arraybackslash}p{\dimexpr 0.133\linewidth-8pt\relax}>{\raggedright\arraybackslash}p{\dimexpr 0.192\linewidth-8pt\relax}@{}}
\toprule
\textbf{Contrast} & \textbf{$\Delta$ cell-weighted [CI]} & \textbf{$\Delta$ block-equal [CI]} & \textbf{Verdict} & \textbf{Common CI, verdict} & \textbf{Dependence} & \textbf{Note} \\
\midrule\endfirsthead
\toprule
\textbf{Contrast} & \textbf{$\Delta$ cell-weighted [CI]} & \textbf{$\Delta$ block-equal [CI]} & \textbf{Verdict} & \textbf{Common CI, verdict} & \textbf{Dependence} & \textbf{Note} \\
\midrule\endhead
\bottomrule\endlastfoot
R0[T] − P0(TabPFN residual, λ 0.25) & −0.31 [−0.48, 0.05] & −0.21 [−0.43, 0.02] & equivalent & [−0.59, 0.23] /\allowbreak{} [−0.58, 0.15], undetermined & depends on split independence & p\_\allowbreak{}eq 0.0163, holm\_\allowbreak{}p\_\allowbreak{}eq none(before correction) \\
same contrast, 3 regions & −0.13 [−0.31, 0.27] & +0.38 [0.10, 0.66] & undetermined & [−0.44, 0.47] /\allowbreak{} [−0.07, 0.84], undetermined &  &  \\
R0[C] − P0(same context CatBoost residual) & −0.05 [−0.42, 0.16] & −0.21 [−0.48, 0.04] & equivalent & [−0.62, 0.31] /\allowbreak{} [−0.64, 0.20], undetermined & depends on split independence & p\_\allowbreak{}eq 0.0145 \\
same contrast, 3 regions & +0.02 [−0.29, 0.53] & +0.39 [0.04, 0.73] & undetermined & [−0.46, 0.76] /\allowbreak{} [−0.22, 0.98], undetermined &  &  \\
R0[I] − P0(TabICL residual) & −0.27 [−0.48, 0.25] & +0.01 [−0.28, 0.30] & equivalent & [−0.67, 0.54] /\allowbreak{} [−0.48, 0.49], undetermined & depends on split independence & eq\_\allowbreak{}p 0.0191, holm\_\allowbreak{}p 1 \\
same contrast, 3 regions & +0.64 [0.17, 1.12] & +0.71 [0.25, 1.18] & higher error & [−0.11, 1.43] /\allowbreak{} [−0.08, 1.49], undetermined & depends on split independence &  \\
R0 − P0(Rescale catboost\_\allowbreak{}lo, λ 0.25), stage T0 & −0.10 [−0.44, 0.14] & CI [−0.48, 0.03] & equivalent & summary table common CI column none. same contrast fixed-composition pool curve rows(below) common-CI verdict undetermined & unknown(summary table column none) & main4\_\allowbreak{}noninf True. m3\_\allowbreak{}delta −0.08(m3\_\allowbreak{}verdict4 undetermined). independent 5 regions: lower error 1, equivalent 1, undetermined 3, higher error 0, non-inferior 2. worst Russia\_\allowbreak{}E|\allowbreak{}x +0.51 [−0.587, 1.098]. 'equivalent 1'(Lena Delta) tables/\allowbreak{}lgw\_\allowbreak{}scenarios.csv common-CI verdict undetermined(depends on split independence), 'lower error 1'(Canada) common-CI verdict lower error \\
same contrast, Alaska cell & +0.42 [−0.30, 1.92] & +2.23 [1.29, 3.13] & undetermined & [−0.72, 2.67] /\allowbreak{} [0.56, 3.82], undetermined &  & noninf False \\
R0 − P0, λ 0.25(fixed-composition pool curve) & −0.10 [−0.45, 0.13] & −0.22 [−0.48, 0.03] & equivalent & [−0.65, 0.29] /\allowbreak{} [−0.66, 0.20], undetermined & depends on split independence(two verdict columns compared) & role = 'descriptive(verdict not used)' \\
R2 − P0, λ 0.25(fixed-composition pool curve) & −0.08 [−0.26, 0.01] & −0.29 [−0.43, −0.14] & equivalent & [−0.37, 0.10] /\allowbreak{} [−0.53, −0.06], undetermined & depends on split independence(two verdict columns compared) & descriptive \\
R0 − F1k(residual structure vs physics-input structure) & −2.72 [−4.34, −1.73] & −2.95 [−3.93, −1.97] & lower error & [−5.37, −0.92] /\allowbreak{} [−4.59, −1.34], lower error &  & Holm 0.0023 \\
\end{longtable}}
\par\medskip\noindent{\small\textit{h. Target curves at zero labels (LG main run; 1,000 resamples) (C1 README 2.6)}}\par\nopagebreak
{\fontsize{7.5}{9}\selectfont\setlength{\tabcolsep}{4pt}\setlength{\LTpre}{2pt}\setlength{\LTpost}{4pt}
\begin{longtable}{@{}>{\raggedright\arraybackslash}p{\dimexpr 0.226\linewidth-8pt\relax}>{\raggedright\arraybackslash}p{\dimexpr 0.115\linewidth-8pt\relax}>{\raggedright\arraybackslash}p{\dimexpr 0.074\linewidth-8pt\relax}>{\raggedright\arraybackslash}p{\dimexpr 0.074\linewidth-8pt\relax}>{\raggedright\arraybackslash}p{\dimexpr 0.218\linewidth-8pt\relax}>{\raggedright\arraybackslash}p{\dimexpr 0.218\linewidth-8pt\relax}>{\raggedright\arraybackslash}p{\dimexpr 0.074\linewidth-8pt\relax}@{}}
\toprule
\textbf{Target} & \textbf{Method} & \textbf{RMSE} & \textbf{RMSE P0} & \textbf{$\Delta$ cell-weighted [CI]} & \textbf{$\Delta$ block-equal [CI]} & \textbf{sig} \\
\midrule\endfirsthead
\toprule
\textbf{Target} & \textbf{Method} & \textbf{RMSE} & \textbf{RMSE P0} & \textbf{$\Delta$ cell-weighted [CI]} & \textbf{$\Delta$ block-equal [CI]} & \textbf{sig} \\
\midrule\endhead
\bottomrule\endlastfoot
Alaska & D0 & 35.33 & 14.54 & +20.79 [14.14, 26.92] & +18.50 [16.09, 20.91] & worse \\
Alaska & D1 & 16.33 & 14.54 & +1.79 [0.37, 2.57] & +1.04 [0.34, 1.85] & worse \\
Alaska & R0 λ 0.25 & 14.96 & 14.54 & +0.42 [−0.31, 1.90] & +2.23 [1.30, 3.12] & ns \\
Alaska & R2 λ 0.25 & 14.13 & 14.54 & −0.41 [−0.60, 0.10] & +0.19 [0.01, 0.40] & ns \\
Lena & D1 & 21.88 & 21.69 & +0.18 [−0.21, 0.42] & −0.45 [−0.99, 0.01] & ns \\
Canada & D1 & 27.94 & 28.16 & −0.22 [−1.13, 0.40] & −1.53 [−2.52, −0.58] & ns \\
Russia\_\allowbreak{}W & D1 & 41.16 & 42.98 & −1.82 [−3.75, −0.27] & −1.45 [−3.09, 0.09] & ns \\
Russia\_\allowbreak{}E & D1 & 30.43 & 29.72 & +0.70 [0.17, 1.36] & +0.69 [−0.02, 1.43] & ns \\
\end{longtable}}
\par\medskip\noindent{\small\textit{i. Physics baselines at zero labels (C1 README 2.7)}}\par\nopagebreak
{\fontsize{7.5}{9}\selectfont\setlength{\tabcolsep}{4pt}\setlength{\LTpre}{2pt}\setlength{\LTpost}{4pt}
\begin{longtable}{@{}>{\raggedright\arraybackslash}p{\dimexpr 0.206\linewidth-8pt\relax}>{\raggedright\arraybackslash}p{\dimexpr 0.117\linewidth-8pt\relax}>{\raggedright\arraybackslash}p{\dimexpr 0.110\linewidth-8pt\relax}>{\raggedright\arraybackslash}p{\dimexpr 0.109\linewidth-8pt\relax}>{\raggedright\arraybackslash}p{\dimexpr 0.109\linewidth-8pt\relax}>{\raggedright\arraybackslash}p{\dimexpr 0.239\linewidth-8pt\relax}>{\raggedright\arraybackslash}p{\dimexpr 0.109\linewidth-8pt\relax}@{}}
\toprule
\textbf{Contrast} & \textbf{$\Delta$ cell-weighted [CI]} & \textbf{$\Delta$ block-equal [CI]} & \textbf{Verdict} & \textbf{Holm $P$} & \textbf{Common CI, verdict} & \textbf{Dependence} \\
\midrule\endfirsthead
\toprule
\textbf{Contrast} & \textbf{$\Delta$ cell-weighted [CI]} & \textbf{$\Delta$ block-equal [CI]} & \textbf{Verdict} & \textbf{Holm $P$} & \textbf{Common CI, verdict} & \textbf{Dependence} \\
\midrule\endhead
\bottomrule\endlastfoot
P* − P0(year-matched degree-days Stefan) & −1.00 [−1.35, −0.32] & −0.41 [−0.74, −0.08] & lower error & 0.0056 & [−1.61, 0.06] /\allowbreak{} [−0.98, 0.19], undetermined & depends on split independence \\
B:ens − P0(physics-model and product ensemble) & −2.73 [−3.33, −0.91] & +0.05 [−0.94, 1.01] & undetermined & 1.0(abstract bundle) & [−4.07, 0.15] /\allowbreak{} [−1.64, 1.63], undetermined &  \\
D0 − P*(direct ML catboost\_\allowbreak{}lo vs year-matched degree-days Stefan) & +3.24 [1.86, 4.29] & +2.03 [0.90, 3.12] & higher error & 0.0050(X9 bundle) & [0.99, 4.99] /\allowbreak{} [0.17, 3.88], higher error &  \\
\end{longtable}}

\clearpage
\phantomsection\label{tab:S5}
\noindent\textbf{Supplementary Table S5.} Zero-label risk table (WF0, post hoc description, no verdict words).\par\smallskip
\noindent{\footnotesize Number of the 30 targets (25 LG targets and five licence-verified LGD targets; seven independent regions and 20 sub-regions) in which the method minus P0 exceeded +2~cm, with median and maximum (cm). Cell-random draws, $\alpha$ = 1, CatBoost. D0\_1.0 direct ML; R1\_1.0 residual weight 1.0; R1\_0.25 low-weight residual (equal to R0 at zero labels); D1\_1.0 physics augmentation; R2\_1.0 augmentation plus anchor plus residual with $\lambda$ 1.0. The count uses the point estimate only. Counting only the seven independent regions, targets above +2~cm at zero labels were D0 5/7, R1\_1.0 2/7, D1 1/7, R1\_0.25 0/7 and R2\_1.0 0/7. Counting targets whose two CIs both exceeded zero (higher error) instead of the 2~cm threshold gives D0\_1.0 14, R1\_1.0 10, D1\_1.0 8, R1\_0.25 6 and R2\_1.0 5 of 30 (post hoc calculation, C1 README 2.8). Source: paper/claims/C1\_label0\_safety/tables/wf0\_risk.csv.}\par
\par\medskip\noindent{\small\textit{Risk table (C1 README 2.5)}}\par\nopagebreak
{\fontsize{7.5}{9}\selectfont\setlength{\tabcolsep}{4pt}\setlength{\LTpre}{2pt}\setlength{\LTpost}{4pt}
\begin{longtable}{@{}>{\raggedright\arraybackslash}p{\dimexpr 0.088\linewidth-8pt\relax}>{\raggedright\arraybackslash}p{\dimexpr 0.513\linewidth-8pt\relax}>{\raggedright\arraybackslash}p{\dimexpr 0.088\linewidth-8pt\relax}>{\raggedright\arraybackslash}p{\dimexpr 0.088\linewidth-8pt\relax}>{\raggedright\arraybackslash}p{\dimexpr 0.104\linewidth-8pt\relax}>{\raggedright\arraybackslash}p{\dimexpr 0.119\linewidth-8pt\relax}@{}}
\toprule
\textbf{$n$} & \textbf{Method} & \textbf{Targets} & \textbf{Targets above +2 cm} & \textbf{Median (cm)} & \textbf{Maximum (cm)} \\
\midrule\endfirsthead
\toprule
\textbf{$n$} & \textbf{Method} & \textbf{Targets} & \textbf{Targets above +2 cm} & \textbf{Median (cm)} & \textbf{Maximum (cm)} \\
\midrule\endhead
\bottomrule\endlastfoot
0 & P1(= P0) & 30 & 0 & +0.00 & +0.00 \\
0 & D0\_\allowbreak{}1.0(direct ML) & 30 & 18 & +2.69 & +38.46 \\
0 & R1\_\allowbreak{}1.0(residual weight 1.0) & 30 & 13 & +1.55 & +38.54 \\
0 & R1\_\allowbreak{}0.25(low-weight residual) & 30 & 5 & +0.15 & +6.29 \\
0 & D1\_\allowbreak{}1.0(physics augmentation) & 30 & 2 & +0.13 & +3.89 \\
0 & R2\_\allowbreak{}1.0(augmentation + anchor + residual) & 30 & 1 & −0.02 & +6.03 \\
10 & P1 & 30 & 2 & +0.06 & +2.27 \\
10 & D0\_\allowbreak{}1.0 & 30 & 12 & +1.49 & +12.89 \\
10 & R1\_\allowbreak{}1.0 & 30 & 11 & +1.34 & +15.19 \\
10 & R1\_\allowbreak{}0.25 & 30 & 2 & +0.49 & +2.71 \\
10 & D1\_\allowbreak{}1.0 & 30 & 5 & +0.43 & +5.00 \\
10 & R2\_\allowbreak{}1.0 & 30 & 2 & +0.31 & +7.50 \\
all & P1 & 30 & 4 & +0.13 & +4.64 \\
all & D0\_\allowbreak{}1.0 & 30 & 7 & +0.27 & +7.35 \\
all & R1\_\allowbreak{}1.0 & 30 & 5 & +0.16 & +9.14 \\
all & R1\_\allowbreak{}0.25 & 30 & 4 & −0.13 & +3.38 \\
all & D1\_\allowbreak{}1.0 & 30 & 5 & −0.17 & +12.72 \\
all & R2\_\allowbreak{}1.0 & 30 & 7 & −0.36 & +10.70 \\
\end{longtable}}

\clearpage
\phantomsection\label{tab:S6}
\noindent\textbf{Supplementary Table S6.} Combination structures and augmentation (L10, L2, L12, L17, L4, recalibration method).\par\smallskip
\noindent{\footnotesize Values: cell-weighted $\Delta$ [95\% CI] / block-equal [95\% CI] in cm. Pools at 40 and 160 labels contain the Lena Delta and Canada only. Method codes (internal identifiers allowed in the Supplementary Information): P0 source-coefficient Stefan model; P* year-matched Stefan model; P1 recalibrated Stefan model ($\kappa = 10$); P2 local least-squares Stefan model; P1* or P1@ed soil-property recalibrated Stefan model; R0 source anchor plus residual ML; R1 recalibrated anchor plus residual ML; R2 augmentation plus anchor plus residual; D0 direct ML; D1 physics pseudo-label augmentation; F1a, F1k, F1n physics-input ML; RM multiplicative residual; Re Stefan-CCI mean anchor plus residual; W selection rule (cross-validation within the target labels); S1 random cells, S2 block stratification, S3 covariate-cluster stratification, S4 covariate max-min distance, S5 source-extrapolation first, S6 sequential selection by prediction variance, S7 $\sqrt{\mathrm{TDD}}$ first. Target suffixes: $|x$ parent region excluded from the source, $|i$ rest of the parent kept, $|r$ within-region design. $n$ = number of target labels; all = all labels of half A. Verdicts: lower error, higher error, equivalent ($\pm$0.5~cm), undetermined; both 95\% block-bootstrap CIs (cell-weighted and block-equal) must agree. \dag~a note of the source row in Korean was omitted; see the cited README section.}\par

\par\medskip\noindent{\small\textit{a. Residual structure versus physics-input structure, 0 to 10 labels (C3 README 2.1)}}\par\nopagebreak
{\fontsize{7.5}{9}\selectfont\setlength{\tabcolsep}{4pt}\setlength{\LTpre}{2pt}\setlength{\LTpost}{4pt}
\begin{longtable}{@{}>{\raggedright\arraybackslash}p{\dimexpr 0.053\linewidth-8pt\relax}>{\raggedright\arraybackslash}p{\dimexpr 0.259\linewidth-8pt\relax}>{\raggedright\arraybackslash}p{\dimexpr 0.344\linewidth-8pt\relax}>{\raggedright\arraybackslash}p{\dimexpr 0.344\linewidth-8pt\relax}@{}}
\toprule
\textbf{ID} & \textbf{Contrast or quantity} & \textbf{Value (cm unless stated): cell-weighted [CI] / block-equal [CI]} & \textbf{Verdict} \\
\midrule\endfirsthead
\toprule
\textbf{ID} & \textbf{Contrast or quantity} & \textbf{Value (cm unless stated): cell-weighted [CI] / block-equal [CI]} & \textbf{Verdict} \\
\midrule\endhead
\bottomrule\endlastfoot
E01 & R0 − F1k, n 0 & −2.72 [−4.34, −1.73] /\allowbreak{} [−3.93, −1.97] & lower error(pool 4/\allowbreak{}4. region rows lower error 2/\allowbreak{}4, E08) \\
E02 & R0 − F1a, n 0 & −2.49 [−3.35, −1.56] /\allowbreak{} [−2.97, −1.11] & lower error(pool 4/\allowbreak{}4. region rows lower error 2/\allowbreak{}4, E08b) \\
E03 & R1 − F1k, n 10 & −3.74 [−5.18, −2.70] /\allowbreak{} [−5.12, −2.96] & lower error(pool 4/\allowbreak{}4. region rows lower error 2/\allowbreak{}4, E07) \\
E04 & R1 − F1n, n 10 & −2.52 [−3.26, −1.56] /\allowbreak{} [−2.90, −1.16] & lower error(pool 4/\allowbreak{}4. region rows lower error 2/\allowbreak{}4, higher error 1/\allowbreak{}4 Canada, E08a) \\
E05 & L10 hypothesis verdict & holm\_\allowbreak{}p 0.0023(within item auxiliary column)\textsuperscript{\dag} & supported(blind) \\
E06 & AB7 = R1 − F1k, n 10 & −3.74 [−5.23, −2.75] /\allowbreak{} [−5.14, −2.98]; holm\_\allowbreak{}p 0.001; '(a) direction stated'. note: F1k inputs have recalibration coefficient E\_\allowbreak{}n none(LG 6A.3). same n P1 − P0 −2.45(E36) & lower error \\
E07 & AB7 regions rows & Lena Delta −1.09 [−2.37, 0.19] /\allowbreak{} [−0.24, 2.27]; Canada +3.44 [−1.57, 5.38] /\allowbreak{} [−1.67, 1.90]; W Russia −13.98 [−16.60, −11.22] /\allowbreak{} [−18.68, −12.13]; E Russia −3.33 [−4.78, −0.76] /\allowbreak{} [−3.92, −0.12] ‡ & undetermined, undetermined, lower error, lower error \\
E08 & R0 − F1k, n 0 regions rows & Lena Delta −3.00 [−4.28, −1.60] /\allowbreak{} [−1.63, 1.22]; Canada −0.31 [−6.00, 2.11] /\allowbreak{} [−6.30, −1.90]; W Russia −3.96 [−5.81, −1.73] /\allowbreak{} [−6.54, −2.57] ‡; E Russia −3.61 [−5.53, −1.20] /\allowbreak{} [−5.12, −0.75] ‡ & undetermined, undetermined, lower error, lower error \\
E08a & R1 − F1n, n 10 regions rows & Lena Delta −1.86 [−3.26, −0.22] /\allowbreak{} [−0.60, 2.11]; Canada +3.12 [1.42, 3.97] /\allowbreak{} [0.07, 2.74] ‡; W Russia −7.50 [−9.08, −5.83] /\allowbreak{} [−10.45, −6.14]; E Russia −3.83 [−5.27, −0.72] /\allowbreak{} [−3.97, −0.03] ‡ & undetermined, higher error, lower error, lower error \\
E08b & R0 − F1a, n 0 regions rows & Lena Delta −2.83 [−4.12, −1.61] /\allowbreak{} [−1.59, 1.18]; Canada +1.53 [−0.04, 2.29] /\allowbreak{} [−1.09, 1.58]; W Russia −4.12 [−6.16, −1.57] /\allowbreak{} [−6.84, −2.24] ‡; E Russia −4.55 [−6.27, −1.90] /\allowbreak{} [−5.86, −1.40] & undetermined, undetermined, lower error, lower error \\
E09 & R1 − F1k, R1 − F1n, n 3(auxiliary) & −3.02 [−4.65, −1.93] /\allowbreak{} [−4.14, −2.16]; −3.17 [−4.00, −2.10] /\allowbreak{} [−3.39, −1.51] & lower error, lower error(both pool 4/\allowbreak{}4) \\
E09a & same contrast n 3 regions rows(auxiliary) & R1 − F1k: Lena Delta −2.05 [−3.25, −0.70] /\allowbreak{} [−0.88, 1.73]; Canada +1.40 [−4.62, 3.83] /\allowbreak{} [−4.73, −0.52]; W Russia −7.96 [−10.07, −5.50] /\allowbreak{} [−10.88, −6.18]; E Russia −3.44 [−5.01, −0.78] /\allowbreak{} [−3.91, 0.11]. R1 − F1n: Lena Delta −2.61 [−3.92, −1.07] /\allowbreak{} [−1.11, 1.68]; Canada +2.02 [0.55, 2.81] /\allowbreak{} [−0.57, 1.98]; W Russia −7.95 [−10.14, −5.35] /\allowbreak{} [−10.96, −5.98]; E Russia −4.13 [−5.79, −1.05] /\allowbreak{} [−4.41, −0.20] ‡ & F1k: undetermined, undetermined, lower error, undetermined. F1n: undetermined, undetermined, lower error, lower error \\
E10 & R1 − F1k, R1 − F1n, all(auxiliary) & −2.93 [−3.90, −1.73] /\allowbreak{} [−3.88, −1.80]; −0.77 [−1.60, 0.37] /\allowbreak{} [−1.21, 0.48] & lower error, undetermined \\
E11 & missing substitute sensitivity R1 − F1k@tr, n 10 & −3.75 [−5.21, −2.71] /\allowbreak{} [−5.16, −2.99]; 'missing substitute statistics training side median changed four-way verdict same' & lower error \\
E11a & L29 P* companion: R0@tddm − F1k(n 0), R1@tddm − F1k(n 10) & −3.45 [−4.73, −2.37] /\allowbreak{} [−3.98, −2.09]; −4.37 [−5.63, −3.29] /\allowbreak{} [−5.16, −3.01]. region rows: R0@tddm − F1k Lena Delta −3.92, Canada −1.28, W Russia −4.61, E Russia −3.98; R1@tddm − F1k Lena Delta −2.29, Canada +1.88, W Russia −13.44, E Russia −3.61. verdict:\textsuperscript{\dag} & lower error, lower error(auxiliary rows, pool 4/\allowbreak{}4. region rows lower error 2/\allowbreak{}4, W Russia·E) \\
\end{longtable}}
\par\medskip\noindent{\small\textit{b. 40 to 160 labels: reversal and regional rows (C3 README 2.2)}}\par\nopagebreak
{\fontsize{7.5}{9}\selectfont\setlength{\tabcolsep}{4pt}\setlength{\LTpre}{2pt}\setlength{\LTpost}{4pt}
\begin{longtable}{@{}>{\raggedright\arraybackslash}p{\dimexpr 0.050\linewidth-8pt\relax}>{\raggedright\arraybackslash}p{\dimexpr 0.242\linewidth-8pt\relax}>{\raggedright\arraybackslash}p{\dimexpr 0.354\linewidth-8pt\relax}>{\raggedright\arraybackslash}p{\dimexpr 0.354\linewidth-8pt\relax}@{}}
\toprule
\textbf{ID} & \textbf{Contrast or quantity} & \textbf{Value (cm unless stated): cell-weighted [CI] / block-equal [CI]} & \textbf{Verdict} \\
\midrule\endfirsthead
\toprule
\textbf{ID} & \textbf{Contrast or quantity} & \textbf{Value (cm unless stated): cell-weighted [CI] / block-equal [CI]} & \textbf{Verdict} \\
\midrule\endhead
\bottomrule\endlastfoot
E12 & R1 − F1k, n 40 & +2.38 [0.52, 3.25] /\allowbreak{} [1.17, 2.95]; 'depends on split independence' & higher error(pool 2/\allowbreak{}4: Lena Delta and Canada. region rows higher error 1/\allowbreak{}2, Canada) \\
E13 & R1 − F1n, n 40 & +1.85 [0.73, 2.61] /\allowbreak{} [1.34, 3.09] & higher error(pool 2/\allowbreak{}4. region rows higher error 1/\allowbreak{}2, Canada) \\
E14 & R1 − F1k, n 160 & +2.23 [0.75, 3.04] /\allowbreak{} [1.30, 3.09] & higher error(pool 2/\allowbreak{}4. region rows higher error 1/\allowbreak{}2, Canada) \\
E15 & R1 − F1n, n 160 & +1.98 [0.80, 2.73] /\allowbreak{} [1.35, 3.12] & higher error(pool 2/\allowbreak{}4. region rows higher error 1/\allowbreak{}2, Canada) \\
E16 & of the four above region rows, n 40 & F1k: Lena Delta −0.38 [−1.38, 0.73] /\allowbreak{} [0.64, 2.71], Canada +5.14 [1.58, 6.60] /\allowbreak{} [0.97, 3.91] ‡, Alaska(x) −7.46 [−9.21, −3.03] /\allowbreak{} [−3.66, −1.83]. F1n: Lena Delta −0.76, Canada +4.46 [2.54, 5.41] /\allowbreak{} [1.50, 4.09], Alaska(x) −7.12 [−8.73, −3.31] /\allowbreak{} [−4.38, −2.62] & Lena Delta undetermined, Canada higher error, Alaska(x) lower error \\
E17 & same region rows, n 160 & F1k: Lena Delta −0.39, Canada +4.86 [2.04, 6.10] /\allowbreak{} [1.28, 4.17], Alaska(x) −7.20 [−9.12, −2.07] /\allowbreak{} [−2.37, −0.54]. F1n: Lena Delta −0.71, Canada +4.67 [2.52, 5.72] /\allowbreak{} [1.52, 4.16], Alaska(x) −6.06 [−7.64, −1.95] /\allowbreak{} [−2.59, −0.91] & Lena Delta undetermined, Canada higher error, Alaska(x) lower error \\
E18 & 3 regions mean(Alaska included), n 40·160 & F1k n 40 −0.90 [−2.19, 0.63] /\allowbreak{} [−0.22, 1.12]; F1n n 40 −1.14; F1k n 160 −0.91; F1n n 160 −0.70 & all undetermined \\
E19 & P1* contrast, n 40·160 & P1@ed − P1: −1.87 [−2.64, −0.34] /\allowbreak{} [−1.32, −0.13] ‡, −1.93 [−2.63, −0.56] /\allowbreak{} [−1.20, −0.07] ‡. R1 − P1@ed: +1.31 [−0.13, 1.96] /\allowbreak{} [−0.54, 0.57], +1.30 [0.01, 1.86] /\allowbreak{} [−0.68, 0.41] & P1@ed lower error(pool 2/\allowbreak{}4); R1 − P1@ed undetermined(pool 2/\allowbreak{}4) \\
E20 & TabPFN direct − TabPFN residual(L36) & n 0 +1.33 [0.63, 2.39] /\allowbreak{} [0.63, 2.23]; n 10 +0.97; n 40 −1.79 [−2.49, −0.86] /\allowbreak{} [−2.98, −1.19]; n 160 −2.45 [−3.48, −0.97] /\allowbreak{} [−3.41, −1.42]; all −0.60. regions n 40: Lena Delta −0.10, Canada −3.47, Alaska(x) +8.57; n 160: Lena Delta −0.57, Canada −4.33, Alaska(x) +6.52 & higher error, higher error, lower error(pool 2/\allowbreak{}4), lower error(pool 2/\allowbreak{}4), undetermined. regions: undetermined, lower error, higher error \\
\end{longtable}}
\par\medskip\noindent{\small\textit{c. Augmented combination and stacking (C3 README 2.3)}}\par\nopagebreak
{\fontsize{7.5}{9}\selectfont\setlength{\tabcolsep}{4pt}\setlength{\LTpre}{2pt}\setlength{\LTpost}{4pt}
\begin{longtable}{@{}>{\raggedright\arraybackslash}p{\dimexpr 0.050\linewidth-8pt\relax}>{\raggedright\arraybackslash}p{\dimexpr 0.276\linewidth-8pt\relax}>{\raggedright\arraybackslash}p{\dimexpr 0.337\linewidth-8pt\relax}>{\raggedright\arraybackslash}p{\dimexpr 0.337\linewidth-8pt\relax}@{}}
\toprule
\textbf{ID} & \textbf{Contrast or quantity} & \textbf{Value (cm unless stated): cell-weighted [CI] / block-equal [CI]} & \textbf{Verdict} \\
\midrule\endfirsthead
\toprule
\textbf{ID} & \textbf{Contrast or quantity} & \textbf{Value (cm unless stated): cell-weighted [CI] / block-equal [CI]} & \textbf{Verdict} \\
\midrule\endhead
\bottomrule\endlastfoot
E21 & R2 − R1, n 10·40·160 & n 10 −0.13 [−0.31, 0.24] /\allowbreak{} [−0.32, 0.20]; n 40 +0.34 [0.25, 0.49] /\allowbreak{} [0.21, 0.50]; n 160 +0.08 [−0.09, 0.42] /\allowbreak{} [0.02, 0.31] & ns, worse, ns. hypothesis verdict rows(scope=verdict): 'rejected(augmentation zero labels only)' \\
E22 & four-way auxiliary column of the same contrast & P4: n 10 −0.13 equivalent, n 40 +0.34 higher error('size below 0.5 cm'), n 160 +0.08 equivalent ‡. MEAN3: +0.01, +0.04, −0.09 all equivalent & as above same \\
E23 & AB6 = R2 − R1, n 10 & −0.13 [−0.31, 0.22] /\allowbreak{} [−0.31, 0.20]; holm\_\allowbreak{}p\_\allowbreak{}eq 0.002; '(c) limit stated('0.5 cm same within')' & equivalent(blind) \\
E24 & AB6 regions rows & Lena Delta +0.02 [−0.08, 0.19] /\allowbreak{} [0.04, 0.35]; Canada +0.76 [0.56, 1.00] /\allowbreak{} [0.48, 0.98]; W Russia −0.25; E Russia −1.05 [−1.53, 0.15] /\allowbreak{} [−1.31, −0.06]; Alaska(x) −0.76 [−1.30, −0.34] /\allowbreak{} [−1.68, −0.76] & equivalent, higher error, undetermined, undetermined, lower error \\
E25 & R3 − R1, n 10·40·160 & n 10 +0.49 [−0.17, 1.41] /\allowbreak{} [−0.39, 1.38]; n 40 −1.23 [−1.63, −0.50] /\allowbreak{} [−1.90, −0.79]; n 160 −1.83 [−2.59, −0.40] /\allowbreak{} [−2.10, −0.61] & ns, improve(2 regions), improve(2 regions). hypothesis verdict rows: 'partial' \\
E26 & four-way auxiliary column of the same contrast & n 10 +0.49 /\allowbreak{} n 40 −1.23 /\allowbreak{} n 160 −1.83 & undetermined, lower error, lower error \\
E27 & R3 − R1, 3 regions mean & n 40 +0.30 [−0.67, 0.84] /\allowbreak{} [−1.23, −0.37]; n 160 −0.02 [−1.17, 0.92] /\allowbreak{} [−1.29, −0.24] & undetermined, undetermined \\
E28 & R3 − R1 regions rows & n 40: Lena Delta −0.00, Canada −2.46 [−3.19, −1.03] /\allowbreak{} [−2.82, −1.00], Alaska(x) +3.35 [0.56, 4.25] /\allowbreak{} [−0.42, 0.91]. n 160: Lena Delta +0.20, Canada −3.86 [−5.22, −1.02] /\allowbreak{} [−3.03, −0.66], Alaska(x) +3.59 [0.53, 4.57] /\allowbreak{} [−0.29, 1.02] & Lena Delta undetermined, Canada lower error, Alaska(x) undetermined \\
\end{longtable}}
\par\medskip\noindent{\small\textit{d. Multiplicative residual and label shuffle (C3 README 2.4)}}\par\nopagebreak
{\fontsize{7.5}{9}\selectfont\setlength{\tabcolsep}{4pt}\setlength{\LTpre}{2pt}\setlength{\LTpost}{4pt}
\begin{longtable}{@{}>{\raggedright\arraybackslash}p{\dimexpr 0.068\linewidth-8pt\relax}>{\raggedright\arraybackslash}p{\dimexpr 0.254\linewidth-8pt\relax}>{\raggedright\arraybackslash}p{\dimexpr 0.522\linewidth-8pt\relax}>{\raggedright\arraybackslash}p{\dimexpr 0.156\linewidth-8pt\relax}@{}}
\toprule
\textbf{ID} & \textbf{Contrast or quantity} & \textbf{Value (cm unless stated): cell-weighted [CI] / block-equal [CI]} & \textbf{Verdict} \\
\midrule\endfirsthead
\toprule
\textbf{ID} & \textbf{Contrast or quantity} & \textbf{Value (cm unless stated): cell-weighted [CI] / block-equal [CI]} & \textbf{Verdict} \\
\midrule\endhead
\bottomrule\endlastfoot
E29 & RM − R1, n 10, λ 0.25 & +0.35 [0.08, 0.46] /\allowbreak{} [0.12, 0.44]; small\_\allowbreak{}effect True; 'depends on split independence'; holm\_\allowbreak{}p 0.0768 'significant before correction' & higher error \\
E30 & RM − R1, n 10, λ 1.0 & +1.09 [0.19, 1.38] /\allowbreak{} [0.43, 1.62]; 'depends on split independence'; holm\_\allowbreak{}p 0.1904 'significant before correction' & higher error \\
E31 & RM − R1, all, λ 0.25 & +0.25 [0.07, 0.37] /\allowbreak{} [0.18, 0.55]; small\_\allowbreak{}effect True; 'depends on split independence'; holm\_\allowbreak{}p 0.0900 'significant before correction' & higher error \\
E32 & RM − R1, all, λ 1.0 & +0.70 [−0.12, 1.13] /\allowbreak{} [0.45, 1.78] & undetermined \\
E33 & RM − R1, 3 regions mean & n 10 λ 0.25 +0.04 equivalent; n 10 λ 1.0 −0.94 [−2.19, −0.07] /\allowbreak{} [−1.06, −0.18] lower error ‡; all λ 0.25 +0.13 higher error ‡('size below 0.5 cm'); all λ 1.0 +0.16 undetermined & rowsby \\
E34 & R1 − R1s(label shuffle) & n 10: −0.12, −0.26; all: −0.60 [−0.80, −0.36] /\allowbreak{} [−0.64, −0.32], −1.69 [−2.48, −0.84] /\allowbreak{} [−2.01, −0.94]\textsuperscript{\dag} & equivalent, equivalent, lower error, lower error \\
\end{longtable}}
\par\medskip\noindent{\small\textit{e. Net value of the residual over the recalibrated model (C3 README 2.5)}}\par\nopagebreak
{\fontsize{7.5}{9}\selectfont\setlength{\tabcolsep}{4pt}\setlength{\LTpre}{2pt}\setlength{\LTpost}{4pt}
\begin{longtable}{@{}>{\raggedright\arraybackslash}p{\dimexpr 0.053\linewidth-8pt\relax}>{\raggedright\arraybackslash}p{\dimexpr 0.204\linewidth-8pt\relax}>{\raggedright\arraybackslash}p{\dimexpr 0.372\linewidth-8pt\relax}>{\raggedright\arraybackslash}p{\dimexpr 0.372\linewidth-8pt\relax}@{}}
\toprule
\textbf{ID} & \textbf{Contrast or quantity} & \textbf{Value (cm unless stated): cell-weighted [CI] / block-equal [CI]} & \textbf{Verdict} \\
\midrule\endfirsthead
\toprule
\textbf{ID} & \textbf{Contrast or quantity} & \textbf{Value (cm unless stated): cell-weighted [CI] / block-equal [CI]} & \textbf{Verdict} \\
\midrule\endhead
\bottomrule\endlastfoot
E35 & R1 − P1(CatBoost, Rescale) & n 3 −0.08 [−0.43, 0.11] /\allowbreak{} [−0.47, −0.00]; n 10 −0.18 [−0.52, −0.02] /\allowbreak{} [−0.52, −0.09]; n 40 −0.56; n 160 −0.64; n 320 −0.83; all −0.40 [−0.76, −0.23] /\allowbreak{} [−0.76, −0.30]. 'net value of ML with sparse labels' & ns, improve(pool 4/\allowbreak{}4), improve(pool 2/\allowbreak{}4) ×3, improve(pool 4/\allowbreak{}4). P1* flag: n 3·320 'P1* confirmednot n', n 40·160 E19 R1 − P1@ed(undetermined) companion \\
E36 & AB4 = P1 − P0, AB5 = R1 − P1(n 10) & AB4 −2.45 [−3.04, −1.88] /\allowbreak{} [−3.30, −1.87], holm\_\allowbreak{}p 0.001, 'S-a viewing after cited'. AB5 −0.18 [−0.53, −0.01] /\allowbreak{} [−0.51, −0.08] ‡, 'distinguishable, size below 0.5 cm', holm\_\allowbreak{}p 0.136, '(b) abstract 'no difference established', main text 'significant before correction'' & AB4 lower error(pool 4/\allowbreak{}4. region rows lower error 1/\allowbreak{}4, E36a), AB5 lower error \\
E36a & AB4 regions rows & Lena Delta +0.22 [−0.32, 0.92] /\allowbreak{} [0.25, 1.14]; Canada +2.26 [0.71, 2.94] /\allowbreak{} [0.67, 2.15]; W Russia −11.96 [−13.62, −10.53] /\allowbreak{} [−15.71, −10.54]; E Russia −0.34 [−0.90, 1.30] /\allowbreak{} [−0.08, 1.48] & undetermined, higher error, lower error, undetermined(lower error regions 1/\allowbreak{}4) \\
E37 & extended pool R1 − P1, n 10 & P4 −0.18 [−0.52, −0.02] /\allowbreak{} [−0.51, −0.08] ‡; PE1 −0.13 [−0.42, 0.01] /\allowbreak{} [−0.41, 0.01] ‡, l28 'weakened', cross-environment flag; Central Russia(PE1-new) +0.06 [−0.40, 0.51] /\allowbreak{} [−0.39, 0.81], cross-environment flag & lower error, equivalent, undetermined \\
E38 & TabPFN residual R1[T] − P1 & n 3 −0.32 [−0.46, −0.02] /\allowbreak{} [−0.38, −0.00] ‡, holm\_\allowbreak{}p 0.2366 'significant before correction'; n 10 −0.42 [−0.54, −0.18] /\allowbreak{} [−0.47, −0.14], holm\_\allowbreak{}p 0.0036; n 40 −0.62; n 160 −0.93; n 320 −1.10; all −0.80. MEAN3: n 3 −0.18, n 10 −0.18. companion R1[C] − P1 n 10 −0.16 ‡ & lower error(n 3·10 'size below 0.5 cm', blind False replication (unblinded); n 3 within item Holm auxiliary column 'significant before correction'). P1* flag: n 3·320 'P1* confirmednot n', n 40·160 . MEAN3 equivalent·equivalent\textsuperscript{\dag} \\
E39 & TabICL residual R1[I] − P1 & n 3 −0.42 [−0.59, 0.01] /\allowbreak{} [−0.40, 0.11]; n 10 −0.55 [−0.71, −0.21] /\allowbreak{} [−0.63, −0.10] ‡, holm\_\allowbreak{}p 0.0180(LGF auxiliary bundle); n 40 −0.79; n 160 −1.18; n 320 −1.43; all −1.00 [−1.23, −0.63] /\allowbreak{} [−1.15, −0.60]. MEAN3: n 10 −0.02, n 40 −0.32 ‡. companion R1[C] − P1 n 10 −0.16 ‡. platform local-3090-torch2.6.0-numpy1.26.4 & n 3 undetermined, n ≥ 10 lower error. P1* flag: n 3·320 'P1* confirmednot n', n 40·160 . MEAN3 n 10 equivalent, n 40 lower error('size below 0.5 cm')\textsuperscript{\dag} \\
\end{longtable}}
\par\medskip\noindent{\small\textit{f. Shrinkage recalibration versus local least squares (C3 README 2.6)}}\par\nopagebreak
{\fontsize{7.5}{9}\selectfont\setlength{\tabcolsep}{4pt}\setlength{\LTpre}{2pt}\setlength{\LTpost}{4pt}
\begin{longtable}{@{}>{\raggedright\arraybackslash}p{\dimexpr 0.051\linewidth-8pt\relax}>{\raggedright\arraybackslash}p{\dimexpr 0.207\linewidth-8pt\relax}>{\raggedright\arraybackslash}p{\dimexpr 0.388\linewidth-8pt\relax}>{\raggedright\arraybackslash}p{\dimexpr 0.354\linewidth-8pt\relax}@{}}
\toprule
\textbf{ID} & \textbf{Contrast or quantity} & \textbf{Value (cm unless stated): cell-weighted [CI] / block-equal [CI]} & \textbf{Verdict} \\
\midrule\endfirsthead
\toprule
\textbf{ID} & \textbf{Contrast or quantity} & \textbf{Value (cm unless stated): cell-weighted [CI] / block-equal [CI]} & \textbf{Verdict} \\
\midrule\endhead
\bottomrule\endlastfoot
E40 & P2 − P1, n 3·10, 4 regions mean & n 3 +1.76 [0.90, 3.05] /\allowbreak{} [1.62, 3.55]; n 10 +0.84 [0.03, 1.91] /\allowbreak{} [0.42, 2.33]; role 'descriptive(verdict not used)' & higher error, higher error \\
E41 & P2 − P1, n 3, region rows & Lena Delta +2.91 [2.17, 3.87] /\allowbreak{} [3.03, 4.47], RMSE 24.76 vs 21.84; Canada +5.40 [3.91, 6.23] /\allowbreak{} [4.04, 5.76], 34.54 vs 29.14; W Russia −4.97 [−7.99, −1.62] /\allowbreak{} [−8.82, −1.78], 33.36 vs 38.33; E Russia +3.69 [2.58, 7.18] /\allowbreak{} [5.83, 7.97], 33.32 vs 29.64 & higher error, higher error, lower error, higher error \\
E42 & P2 − P1, n 10, region rows & Lena Delta +1.01; Canada +3.25; W Russia −1.44 [−4.62, 2.35] /\allowbreak{} [−4.41, 2.93]; E Russia +0.55 & higher error, higher error, undetermined, undetermined \\
E43 & Tibetan Plateau P2 − P1 & n 3 −81.77 [−96.22, −64.75] /\allowbreak{} [−81.01, −42.14]; n 10 −40.02 [−53.74, −24.18] /\allowbreak{} [−35.99, −0.51] ‡. two rows all cross-environment flag. RMSE n 3: P2 119.28, P1 201.05, P0 244.51. n 10: P2 108.12, P1 148.14 & improve, improve \\
E44 & Central Russia P2 − P1 & n 3 +2.32 [−0.43, 4.72] /\allowbreak{} [0.82, 5.95]; n 10 −0.57 [−2.14, 1.00] /\allowbreak{} [−1.14, 1.95]. two rows all cross-environment flag & ns, ns \\
\end{longtable}}

\clearpage
\phantomsection\label{tab:S7}
\noindent\textbf{Supplementary Table S7.} Within-region label-count tests (WF1, WF6, SAR).\par\smallskip
\noindent{\footnotesize Designed after the transfer results had been viewed and registered before running. WF6-a: 25 splits (24 valid in the Lena Delta), comparator P1. WF1: first test, five splits, comparator the best cross-validated physics calibration (Pbest); WF1-b adds nine SAR covariates (x34) in Alaska. Method codes (internal identifiers allowed in the Supplementary Information): P0 source-coefficient Stefan model; P* year-matched Stefan model; P1 recalibrated Stefan model ($\kappa = 10$); P2 local least-squares Stefan model; P1* or P1@ed soil-property recalibrated Stefan model; R0 source anchor plus residual ML; R1 recalibrated anchor plus residual ML; R2 augmentation plus anchor plus residual; D0 direct ML; D1 physics pseudo-label augmentation; F1a, F1k, F1n physics-input ML; RM multiplicative residual; Re Stefan-CCI mean anchor plus residual; W selection rule (cross-validation within the target labels); S1 random cells, S2 block stratification, S3 covariate-cluster stratification, S4 covariate max-min distance, S5 source-extrapolation first, S6 sequential selection by prediction variance, S7 $\sqrt{\mathrm{TDD}}$ first. Target suffixes: $|x$ parent region excluded from the source, $|i$ rest of the parent kept, $|r$ within-region design. $n$ = number of target labels; all = all labels of half A. Verdicts: lower error, higher error, equivalent ($\pm$0.5~cm), undetermined; both 95\% block-bootstrap CIs (cell-weighted and block-equal) must agree. \dag~a note of the source row in Korean was omitted; see the cited README section.}\par

\par\medskip\noindent{\small\textit{a. WF6-a, Alaska (C4 README 2.1)}}\par\nopagebreak
{\fontsize{7.5}{9}\selectfont\setlength{\tabcolsep}{4pt}\setlength{\LTpre}{2pt}\setlength{\LTpost}{4pt}
\begin{longtable}{@{}>{\raggedright\arraybackslash}p{\dimexpr 0.107\linewidth-8pt\relax}>{\raggedright\arraybackslash}p{\dimexpr 0.193\linewidth-8pt\relax}>{\raggedright\arraybackslash}p{\dimexpr 0.133\linewidth-8pt\relax}>{\raggedright\arraybackslash}p{\dimexpr 0.133\linewidth-8pt\relax}>{\raggedright\arraybackslash}p{\dimexpr 0.157\linewidth-8pt\relax}>{\raggedright\arraybackslash}p{\dimexpr 0.115\linewidth-8pt\relax}>{\raggedright\arraybackslash}p{\dimexpr 0.162\linewidth-8pt\relax}@{}}
\toprule
\textbf{ID} & \textbf{Contrast} & \textbf{$\Delta$ cell [CI]} & \textbf{$\Delta$ block-equal [CI]} & \textbf{Verdict} & \textbf{Common-CI verdict} & \textbf{RMSE A / B} \\
\midrule\endfirsthead
\toprule
\textbf{ID} & \textbf{Contrast} & \textbf{$\Delta$ cell [CI]} & \textbf{$\Delta$ block-equal [CI]} & \textbf{Verdict} & \textbf{Common-CI verdict} & \textbf{RMSE A / B} \\
\midrule\endhead
\bottomrule\endlastfoot
E1.R1.n200 & R1(λ cv)-P1|\allowbreak{}n200 & −0.18 [−0.34, 0.01] & −0.25 [−0.42, −0.07] & equivalent & undetermined & 14.22 /\allowbreak{} 14.40 \\
E1.R1.n500 & R1(λ cv)-P1|\allowbreak{}n500 & −0.39 [−0.54, −0.21] & −0.39 [−0.57, −0.21] & lower error(flag: size below 0.5 cm) & undetermined & 13.99 /\allowbreak{} 14.38 \\
E1.R1.n1000 & R1(λ cv)-P1|\allowbreak{}n1000 & −0.54 [−0.69, −0.32] & −0.60 [−0.80, −0.40] & lower error & lower error & 13.87 /\allowbreak{} 14.41 \\
E1.R1.nall & R1(λ cv)-P1|\allowbreak{}nall & −0.52 [−0.70, −0.31] & −0.68 [−0.91, −0.44] & lower error & undetermined & 13.88 /\allowbreak{} 14.40 \\
E1.R2.n200 & R2(λ cv)-P1|\allowbreak{}n200 & −0.20 [−0.24, −0.11] & −0.07 [−0.14, −0.01] & lower error(below 0.5 cm) & equivalent & 14.20 /\allowbreak{} 14.40 \\
E1.R2.n500 & R2(λ cv)-P1|\allowbreak{}n500 & −0.31 [−0.38, −0.18] & −0.17 [−0.26, −0.07] & lower error(below 0.5 cm) & equivalent & 14.07 /\allowbreak{} 14.38 \\
E1.R2.n1000 & R2(λ cv)-P1|\allowbreak{}n1000 & −0.42 [−0.51, −0.26] & −0.33 [−0.44, −0.21] & lower error(below 0.5 cm) & lower error & 13.99 /\allowbreak{} 14.41 \\
E1.R2.nall & R2(λ cv)-P1|\allowbreak{}nall & 0.70 [0.06, 0.90] & −0.48 [−0.88, −0.08] & undetermined & undetermined & 15.10 /\allowbreak{} 14.40 \\
E1.Re.n200 & Re(λ cv)-P1|\allowbreak{}n200 & 1.11 [0.78, 1.44] & 1.30 [0.88, 1.75] & higher error & higher error & 15.51 /\allowbreak{} 14.40 \\
E1.Re.n500 & Re(λ cv)-P1|\allowbreak{}n500 & 0.80 [0.53, 1.17] & 1.05 [0.63, 1.48] & higher error & undetermined & 15.19 /\allowbreak{} 14.38 \\
E1.Re.n1000 & Re(λ cv)-P1|\allowbreak{}n1000 & 0.48 [0.20, 0.87] & 0.52 [0.12, 0.94] & higher error(below 0.5 cm) & undetermined & 14.89 /\allowbreak{} 14.41 \\
E1.Re.nall & Re(λ cv)-P1|\allowbreak{}nall & 0.50 [0.22, 0.85] & 0.35 [−0.07, 0.77] & undetermined & undetermined & 14.90 /\allowbreak{} 14.40 \\
E1.D0.n200 & D0(catboost)-P1|\allowbreak{}n200 & 3.51 [1.71, 3.56] & −0.21 [−0.71, 0.30] & undetermined & undetermined & 17.91 /\allowbreak{} 14.40 \\
E1.D0.n500 & D0(catboost)-P1|\allowbreak{}n500 & 3.16 [1.46, 3.20] & −0.56 [−1.08, −0.04] & undetermined & undetermined & 17.55 /\allowbreak{} 14.38 \\
E1.D0.n1000 & D0(catboost)-P1|\allowbreak{}n1000 & 2.76 [1.16, 2.78] & −0.65 [−1.15, −0.14] & undetermined & undetermined & 17.17 /\allowbreak{} 14.41 \\
E1.D0.nall & D0(catboost)-P1|\allowbreak{}nall & 1.27 [0.30, 1.37] & −1.06 [−1.56, −0.58] & undetermined & undetermined & 15.66 /\allowbreak{} 14.40 \\
\end{longtable}}
\par\medskip\noindent{\small\textit{b. WF6-a, Lena Delta and Canada (C4 README 2.2)}}\par\nopagebreak
{\fontsize{7.5}{9}\selectfont\setlength{\tabcolsep}{4pt}\setlength{\LTpre}{2pt}\setlength{\LTpost}{4pt}
\begin{longtable}{@{}>{\raggedright\arraybackslash}p{\dimexpr 0.154\linewidth-8pt\relax}>{\raggedright\arraybackslash}p{\dimexpr 0.154\linewidth-8pt\relax}>{\raggedright\arraybackslash}p{\dimexpr 0.216\linewidth-8pt\relax}>{\raggedright\arraybackslash}p{\dimexpr 0.148\linewidth-8pt\relax}>{\raggedright\arraybackslash}p{\dimexpr 0.148\linewidth-8pt\relax}>{\raggedright\arraybackslash}p{\dimexpr 0.114\linewidth-8pt\relax}>{\raggedright\arraybackslash}p{\dimexpr 0.066\linewidth-8pt\relax}@{}}
\toprule
\textbf{ID} & \textbf{Target} & \textbf{Contrast} & \textbf{$\Delta$ cell [CI]} & \textbf{$\Delta$ block-equal [CI]} & \textbf{Verdict} & \textbf{Splits} \\
\midrule\endfirsthead
\toprule
\textbf{ID} & \textbf{Target} & \textbf{Contrast} & \textbf{$\Delta$ cell [CI]} & \textbf{$\Delta$ block-equal [CI]} & \textbf{Verdict} & \textbf{Splits} \\
\midrule\endhead
\bottomrule\endlastfoot
E2.Lena.R1.n200 & Lena|\allowbreak{}r & R1(λ cv)-P1|\allowbreak{}n200 & 0.22 [−0.79, 0.64] & −0.20 [−0.75, 0.34] & undetermined & 24 \\
E2.Lena.R1.n500 & Lena|\allowbreak{}r & R1(λ cv)-P1|\allowbreak{}n500 & 0.34 [−0.76, 0.78] & −0.36 [−1.16, 0.38] & undetermined & 24 \\
E2.Lena.R1.n1000 & Lena|\allowbreak{}r & R1(λ cv)-P1|\allowbreak{}n1000 & 0.61 [−0.57, 1.38] & −0.08 [−1.04, 0.81] & undetermined & 23 \\
E2.Lena.R1.nall & Lena|\allowbreak{}r & R1(λ cv)-P1|\allowbreak{}nall & 0.47 [−0.56, 1.28] & −0.14 [−1.07, 0.72] & undetermined & 24 \\
E2.Lena.Re.n200 & Lena|\allowbreak{}r & Re(λ cv)-P1|\allowbreak{}n200 & −1.03 [−1.80, −0.43] & 0.37 [−0.12, 0.86] & undetermined & 24 \\
E2.Lena.D0.n1000 & Lena|\allowbreak{}r & D0(catboost)-P1|\allowbreak{}n1000 & −0.78 [−1.94, 0.17] & −1.08 [−2.00, −0.21] & undetermined & 23 \\
E2.Canada.R1.n200 & Canada|\allowbreak{}r & R1(λ cv)-P1|\allowbreak{}n200 & −0.40 [−1.05, 0.19] & −0.84 [−1.36, −0.28] & undetermined & 25 \\
E2.Canada.R1.nall & Canada|\allowbreak{}r & R1(λ cv)-P1|\allowbreak{}nall & 0.45 [−0.73, 1.24] & −0.23 [−1.09, 0.65] & undetermined & 25 \\
E2.Canada.Re.n200 & Canada|\allowbreak{}r & Re(λ cv)-P1|\allowbreak{}n200 & −2.13 [−2.72, −0.23] & 2.17 [1.29, 3.06] & undetermined & 25 \\
E2.Canada.Re.nall & Canada|\allowbreak{}r & Re(λ cv)-P1|\allowbreak{}nall & −1.73 [−2.45, 0.09] & 1.86 [0.92, 2.83] & undetermined & 25 \\
E2.Canada.D0.n200 & Canada|\allowbreak{}r & D0(catboost)-P1|\allowbreak{}n200 & −1.50 [−1.99, −0.24] & −0.03 [−1.03, 1.01] & undetermined & 25 \\
E2.Canada.D0.nall & Canada|\allowbreak{}r & D0(catboost)-P1|\allowbreak{}nall & −1.50 [−2.19, −0.12] & −0.30 [−1.44, 0.86] & undetermined & 25 \\
\end{longtable}}
\par\medskip\noindent{\small\textit{c. WF6-b, stratified mean of three targets (C4 README 2.3)}}\par\nopagebreak
{\fontsize{7.5}{9}\selectfont\setlength{\tabcolsep}{4pt}\setlength{\LTpre}{2pt}\setlength{\LTpost}{4pt}
\begin{longtable}{@{}>{\raggedright\arraybackslash}p{\dimexpr 0.202\linewidth-8pt\relax}>{\raggedright\arraybackslash}p{\dimexpr 0.083\linewidth-8pt\relax}>{\raggedright\arraybackslash}p{\dimexpr 0.257\linewidth-8pt\relax}>{\raggedright\arraybackslash}p{\dimexpr 0.268\linewidth-8pt\relax}>{\raggedright\arraybackslash}p{\dimexpr 0.190\linewidth-8pt\relax}@{}}
\toprule
\textbf{ID} & \textbf{$n$} & \textbf{$\Delta$ cell [CI]} & \textbf{$\Delta$ block-equal [CI]} & \textbf{Verdict} \\
\midrule\endfirsthead
\toprule
\textbf{ID} & \textbf{$n$} & \textbf{$\Delta$ cell [CI]} & \textbf{$\Delta$ block-equal [CI]} & \textbf{Verdict} \\
\midrule\endhead
\bottomrule\endlastfoot
E3.WF6-b.n500 & 500 & 0.02 [−0.30, 0.23] & −0.25 [−0.55, 0.03] & undetermined \\
E3.WF6-b.n1000 & 1,000 & −0.02 [−0.53, 0.54] & −0.03 [−0.53, 0.43] & undetermined \\
E3.WF6-b.nall & all & 0.31 [−0.28, 1.16] & 0.05 [−0.57, 0.65] & undetermined \\
\end{longtable}}
\par\medskip\noindent{\small\textit{d. WF1-a, first test (C4 README 2.4)}}\par\nopagebreak
{\fontsize{7.5}{9}\selectfont\setlength{\tabcolsep}{4pt}\setlength{\LTpre}{2pt}\setlength{\LTpost}{4pt}
\begin{longtable}{@{}>{\raggedright\arraybackslash}p{\dimexpr 0.160\linewidth-8pt\relax}>{\raggedright\arraybackslash}p{\dimexpr 0.132\linewidth-8pt\relax}>{\raggedright\arraybackslash}p{\dimexpr 0.185\linewidth-8pt\relax}>{\raggedright\arraybackslash}p{\dimexpr 0.127\linewidth-8pt\relax}>{\raggedright\arraybackslash}p{\dimexpr 0.127\linewidth-8pt\relax}>{\raggedright\arraybackslash}p{\dimexpr 0.113\linewidth-8pt\relax}>{\raggedright\arraybackslash}p{\dimexpr 0.155\linewidth-8pt\relax}@{}}
\toprule
\textbf{ID} & \textbf{Target} & \textbf{Contrast} & \textbf{$\Delta$ cell [CI]} & \textbf{$\Delta$ block-equal [CI]} & \textbf{Verdict} & \textbf{RMSE A / B} \\
\midrule\endfirsthead
\toprule
\textbf{ID} & \textbf{Target} & \textbf{Contrast} & \textbf{$\Delta$ cell [CI]} & \textbf{$\Delta$ block-equal [CI]} & \textbf{Verdict} & \textbf{RMSE A / B} \\
\midrule\endhead
\bottomrule\endlastfoot
E6.Alaska.R1.n1000 & Alaska|\allowbreak{}r & R1(λ cv)-Pbest|\allowbreak{}n1000 & −0.49 [−0.82, −0.01] & 0.24 [−0.23, 0.72] & undetermined & 14.26 /\allowbreak{} 14.75 \\
E6.Alaska.R1.n2000 & Alaska|\allowbreak{}r & R1(λ cv)-Pbest|\allowbreak{}n2000 & −0.58 [−0.89, −0.22] & −0.10 [−0.57, 0.41] & undetermined & 14.13 /\allowbreak{} 14.71 \\
E6.Alaska.R1.n5000 & Alaska|\allowbreak{}r & R1(λ cv)-Pbest|\allowbreak{}n5000 & −0.37 [−0.66, 0.03] & 0.26 [−0.22, 0.74] & undetermined & 14.14 /\allowbreak{} 14.51 \\
E6.Alaska.R1.nall & Alaska|\allowbreak{}r & R1(λ cv)-Pbest|\allowbreak{}nall & −0.77 [−2.18, 0.08] & 0.28 [−0.52, 1.12] & undetermined & 14.17 /\allowbreak{} 14.93 \\
E6.Alaska.R2.nall & Alaska|\allowbreak{}r & R2(λ cv)-Pbest|\allowbreak{}nall & 1.05 [−0.66, 2.05] & 1.00 [0.18, 1.89] & undetermined & 15.98 /\allowbreak{} 14.93 \\
E6.AL-2.R1.n1000 & AL-2|\allowbreak{}r & R1(λ cv)-Pbest|\allowbreak{}n1000 & −1.35 [−1.64, −0.47] & −0.85 [−1.53, −0.19] & lower error & 12.38 /\allowbreak{} 13.73 \\
E6.AL-2.R2.n1000 & AL-2|\allowbreak{}r & R2(λ cv)-Pbest|\allowbreak{}n1000 & −1.42 [−1.70, −0.54] & −0.99 [−1.73, −0.29] & lower error & 12.31 /\allowbreak{} 13.73 \\
E6.AL-2.R1.nall & AL-2|\allowbreak{}r & R1(λ cv)-Pbest|\allowbreak{}nall & −0.83 [−1.31, −0.18] & −0.99 [−1.70, −0.29] & lower error & 12.32 /\allowbreak{} 13.15 \\
E6.AL-2.R2.nall & AL-2|\allowbreak{}r & R2(λ cv)-Pbest|\allowbreak{}nall & −0.74 [−1.20, 0.04] & −0.75 [−1.47, −0.04] & undetermined & 12.41 /\allowbreak{} 13.15 \\
\end{longtable}}
\par\medskip\noindent{\small\textit{e. WF1-b (SAR) and WF1-c (risk counts) (C4 README 2.5)}}\par\nopagebreak
{\fontsize{7.5}{9}\selectfont\setlength{\tabcolsep}{4pt}\setlength{\LTpre}{2pt}\setlength{\LTpost}{4pt}
\begin{longtable}{@{}>{\raggedright\arraybackslash}p{\dimexpr 0.136\linewidth-8pt\relax}>{\raggedright\arraybackslash}p{\dimexpr 0.311\linewidth-8pt\relax}>{\raggedright\arraybackslash}p{\dimexpr 0.120\linewidth-8pt\relax}>{\raggedright\arraybackslash}p{\dimexpr 0.160\linewidth-8pt\relax}>{\raggedright\arraybackslash}p{\dimexpr 0.154\linewidth-8pt\relax}>{\raggedright\arraybackslash}p{\dimexpr 0.119\linewidth-8pt\relax}@{}}
\toprule
\textbf{ID} & \textbf{Contrast} & \textbf{Role} & \textbf{$\Delta$ cell [CI]} & \textbf{$\Delta$ block-equal [CI]} & \textbf{Verdict} \\
\midrule\endfirsthead
\toprule
\textbf{ID} & \textbf{Contrast} & \textbf{Role} & \textbf{$\Delta$ cell [CI]} & \textbf{$\Delta$ block-equal [CI]} & \textbf{Verdict} \\
\midrule\endhead
\bottomrule\endlastfoot
E7.R1x34.n500 & R1@x34-R1|\allowbreak{}n500|\allowbreak{}λ cv & decision rule & 0.07 [−0.04, 0.18] & 0.20 [0.04, 0.36] & equivalent \\
E7.R1x34.n1000 & R1@x34-R1|\allowbreak{}n1000|\allowbreak{}λ cv & decision rule & −0.11 [−0.21, 0.05] & 0.06 [−0.15, 0.28] & equivalent \\
E7.R1x34.n2000 & R1@x34-R1|\allowbreak{}n2000|\allowbreak{}λ cv & decision rule & −0.09 [−0.21, 0.15] & 0.08 [−0.23, 0.37] & equivalent \\
E7.R1x34.n5000 & R1@x34-R1|\allowbreak{}n5000|\allowbreak{}λ cv & decision rule & −0.08 [−0.18, 0.08] & 0.15 [−0.09, 0.41] & equivalent \\
E7.R1x34.nall & R1@x34-R1|\allowbreak{}nall|\allowbreak{}λ cv & decision rule & −0.04 [−0.20, 0.21] & 0.23 [−0.08, 0.56] & undetermined \\
E7.D1x34.n500 & D1@x34-D1|\allowbreak{}n500 & auxiliary & 0.24 [0.01, 0.35] & 0.11 [0.01, 0.22] & higher error \\
E7.D1x34.n1000 & D1@x34-D1|\allowbreak{}n1000 & auxiliary & 0.38 [0.03, 0.53] & 0.16 [0.01, 0.32] & higher error \\
E7.D1x34.n2000 & D1@x34-D1|\allowbreak{}n2000 & auxiliary & 0.41 [0.10, 0.55] & 0.34 [0.07, 0.62] & higher error \\
E7.D1x34.n5000 & D1@x34-D1|\allowbreak{}n5000 & auxiliary & 0.81 [0.05, 1.08] & 0.61 [0.04, 1.25] & higher error \\
E7.D1x34.nall & D1@x34-D1|\allowbreak{}nall & auxiliary & 0.53 [−0.21, 0.82] & 0.57 [−0.12, 1.28] & undetermined \\
\end{longtable}}
\par\medskip\noindent{\small\textit{f. Error floors (covariate-conditional, scoring cells) and contest values (cm) (C4 README 2.6)}}\par\nopagebreak
{\fontsize{7.5}{9}\selectfont\setlength{\tabcolsep}{4pt}\setlength{\LTpre}{2pt}\setlength{\LTpost}{4pt}
\begin{longtable}{@{}>{\raggedright\arraybackslash}p{\dimexpr 0.312\linewidth-8pt\relax}>{\raggedright\arraybackslash}p{\dimexpr 0.448\linewidth-8pt\relax}>{\raggedright\arraybackslash}p{\dimexpr 0.240\linewidth-8pt\relax}@{}}
\toprule
\textbf{ID} & \textbf{Column} & \textbf{Value} \\
\midrule\endfirsthead
\toprule
\textbf{ID} & \textbf{Column} & \textbf{Value} \\
\midrule\endhead
\bottomrule\endlastfoot
E8.floor & floor\_\allowbreak{}rmse\_\allowbreak{}cm, n\_\allowbreak{}cells\_\allowbreak{}used & 11.31, 13,333 \\
E8.floor.Lena /\allowbreak{} Canada & floor\_\allowbreak{}rmse\_\allowbreak{}cm & 14.83 /\allowbreak{} 17.90 \\
E8.contest.stefan & rmse\_\allowbreak{}cm & 14.46(file 14.457) \\
E8.contest.resid & rmse\_\allowbreak{}cm & 13.33 \\
E8.contest.min & rmse\_\allowbreak{}cm & 13.33 \\
E8.contest.nrows\_\allowbreak{}lam\_\allowbreak{}pos & row count & 56 \\
E8.cv.block\_\allowbreak{}canonical.* & rmse\_\allowbreak{}cm, nnd\_\allowbreak{}median\_\allowbreak{}km & stefan 14.46, stefan\_\allowbreak{}ridge\_\allowbreak{}l075 13.33, ridge 13.62, catboost\_\allowbreak{}lo 16.12, median nearest distance 28.37 km \\
\end{longtable}}
\par\medskip\noindent{\small\textit{g. Derived values: relative RMSE change and reducible squared error (F = error floor) (C4 README 2.6)}}\par\nopagebreak
{\fontsize{7.5}{9}\selectfont\setlength{\tabcolsep}{4pt}\setlength{\LTpre}{2pt}\setlength{\LTpost}{4pt}
\begin{longtable}{@{}>{\raggedright\arraybackslash}p{\dimexpr 0.385\linewidth-8pt\relax}>{\raggedright\arraybackslash}p{\dimexpr 0.385\linewidth-8pt\relax}>{\raggedright\arraybackslash}p{\dimexpr 0.229\linewidth-8pt\relax}@{}}
\toprule
\textbf{ID} & \textbf{Calculation} & \textbf{Value} \\
\midrule\endfirsthead
\toprule
\textbf{ID} & \textbf{Calculation} & \textbf{Value} \\
\midrule\endhead
\bottomrule\endlastfoot
D1.wf6.rmse\_\allowbreak{}rel\_\allowbreak{}P1\_\allowbreak{}pct & (rmse\_\allowbreak{}B − rmse\_\allowbreak{}A)/\allowbreak{}rmse\_\allowbreak{}B × 100, E1.R1.n1000 rows & 3.75 \% \\
D1.wf6.reducible\_\allowbreak{}P1 /\allowbreak{} R1 & rmse\_\allowbreak{}B² − F², rmse\_\allowbreak{}A² − F² & 79.71 /\allowbreak{} 64.42 cm² \\
D1.wf6.reducible\_\allowbreak{}cut\_\allowbreak{}pct & (79.71 − 64.42)/\allowbreak{}79.71 × 100 & 19.17 \% \\
D2.contest.rmse\_\allowbreak{}rel\_\allowbreak{}pct & (14.457 − 13.330)/\allowbreak{}14.457 × 100 & 7.80 \% \\
D2.contest.reducible\_\allowbreak{}stefan /\allowbreak{} resid & 14.457² − F², 13.330² − F² & 81.13 /\allowbreak{} 49.82 cm² \\
D2.contest.reducible\_\allowbreak{}cut\_\allowbreak{}pct & (81.13 − 49.82)/\allowbreak{}81.13 × 100 & 38.60 \% \\
D3.c2.rho\_\allowbreak{}recal\_\allowbreak{}vs\_\allowbreak{}mlvsp1 & Spearman(recal\_\allowbreak{}gain, ml\_\allowbreak{}vs\_\allowbreak{}p1), tables/\allowbreak{}wf0\_\allowbreak{}misspec.csv 30 rows. signs as in the source column & ρ 0.08, p 0.67 \\
D3.c2.rho\_\allowbreak{}recal\_\allowbreak{}vs\_\allowbreak{}mlbest & Spearman(recal\_\allowbreak{}gain, ml\_\allowbreak{}best\_\allowbreak{}delta). signs as in the source column & ρ −0.66 \\
(revision check) gain basis & ml\_\allowbreak{}vs\_\allowbreak{}p1 = ml\_\allowbreak{}best\_\allowbreak{}delta − P1(best ML − P1, maximum difference 3.6e-15)so Spearman(recal\_\allowbreak{}gain, P1 − best ML) = −0.08, Spearman(recal\_\allowbreak{}gain, P0 − best ML) = 0.66 & ρ −0.08(p 0.67), ρ 0.66 \\
\end{longtable}}

\clearpage
\phantomsection\label{tab:S8}
\noindent\textbf{Supplementary Table S8.} Selection rule W and bias diagnosis (WF4-a, WF4-b, WF4-c, selection counts).\par\smallskip
\noindent{\footnotesize W chooses among seven candidates by block cross-validation within the target labels (Methods); fixed recipe R1 with $\lambda$ 0.25. Non-inferiority margin 0.5~cm. Designed after the label-grid results had been viewed and registered before running. Method codes (internal identifiers allowed in the Supplementary Information): P0 source-coefficient Stefan model; P* year-matched Stefan model; P1 recalibrated Stefan model ($\kappa = 10$); P2 local least-squares Stefan model; P1* or P1@ed soil-property recalibrated Stefan model; R0 source anchor plus residual ML; R1 recalibrated anchor plus residual ML; R2 augmentation plus anchor plus residual; D0 direct ML; D1 physics pseudo-label augmentation; F1a, F1k, F1n physics-input ML; RM multiplicative residual; Re Stefan-CCI mean anchor plus residual; W selection rule (cross-validation within the target labels); S1 random cells, S2 block stratification, S3 covariate-cluster stratification, S4 covariate max-min distance, S5 source-extrapolation first, S6 sequential selection by prediction variance, S7 $\sqrt{\mathrm{TDD}}$ first. Target suffixes: $|x$ parent region excluded from the source, $|i$ rest of the parent kept, $|r$ within-region design. $n$ = number of target labels; all = all labels of half A. Verdicts: lower error, higher error, equivalent ($\pm$0.5~cm), undetermined; both 95\% block-bootstrap CIs (cell-weighted and block-equal) must agree. \dag~a note of the source row in Korean was omitted; see the cited README section.}\par

\par\medskip\noindent{\small\textit{a. WF4-a pooled contrast W $-$ R1(0.25) (C5 README 2.1)}}\par\nopagebreak
{\fontsize{7.5}{9}\selectfont\setlength{\tabcolsep}{4pt}\setlength{\LTpre}{2pt}\setlength{\LTpost}{4pt}
\begin{longtable}{@{}>{\raggedright\arraybackslash}p{\dimexpr 0.041\linewidth-8pt\relax}>{\raggedright\arraybackslash}p{\dimexpr 0.176\linewidth-8pt\relax}>{\raggedright\arraybackslash}p{\dimexpr 0.052\linewidth-8pt\relax}>{\raggedright\arraybackslash}p{\dimexpr 0.157\linewidth-8pt\relax}>{\raggedright\arraybackslash}p{\dimexpr 0.103\linewidth-8pt\relax}>{\raggedright\arraybackslash}p{\dimexpr 0.205\linewidth-8pt\relax}>{\raggedright\arraybackslash}p{\dimexpr 0.103\linewidth-8pt\relax}>{\raggedright\arraybackslash}p{\dimexpr 0.110\linewidth-8pt\relax}>{\raggedright\arraybackslash}p{\dimexpr 0.052\linewidth-8pt\relax}@{}}
\toprule
\textbf{ID} & \textbf{$n$ (pool)} & \textbf{$\Delta$} & \textbf{CI cell / block-equal} & \textbf{Verdict} & \textbf{Common CI, verdict} & \textbf{Holm $P$} & \textbf{$P_{\mathrm{ni}}$, non-inferior} & \textbf{$\Delta$ (\% of P0)} \\
\midrule\endfirsthead
\toprule
\textbf{ID} & \textbf{$n$ (pool)} & \textbf{$\Delta$} & \textbf{CI cell / block-equal} & \textbf{Verdict} & \textbf{Common CI, verdict} & \textbf{Holm $P$} & \textbf{$P_{\mathrm{ni}}$, non-inferior} & \textbf{$\Delta$ (\% of P0)} \\
\midrule\endhead
\bottomrule\endlastfoot
E1 & 10(regions 4/\allowbreak{}4) & -0.25 & [-1.11, 0.73] /\allowbreak{} [-1.09, 0.69] & undetermined & [-1.56, 1.30] /\allowbreak{} [-1.63, 1.20], undetermined & 0.66 & 0.07, False & -0.82 \\
E2 & 40(regions 2/\allowbreak{}4, Lena Delta and Canada) & -0.87 & [-1.20, -0.27] /\allowbreak{} [-1.28, -0.46] & lower error & [-1.31, 0.07] /\allowbreak{} [-1.47, -0.29], undetermined(ci\_\allowbreak{}dependence = depends on split independence) & 0.01(0.0054) & 0.00, True & -2.85 \\
E3 & 160(regions 2/\allowbreak{}4, Lena Delta and Canada) & -1.32 & [-1.70, -0.68] /\allowbreak{} [-1.77, -0.83] & lower error & [-1.86, -0.34] /\allowbreak{} [-2.02, -0.63], lower error & 0.00(0.0004) & 0.00, True & -4.31 \\
E4 & all(regions 4/\allowbreak{}4) & -0.74 & [-1.65, 0.19] /\allowbreak{} [-1.75, -0.07] & undetermined & [-2.02, 0.62] /\allowbreak{} [-2.18, 0.33], undetermined\textsuperscript{\dag} & 0.25 & 0.00(0.0047), True & -2.41 \\
\end{longtable}}
\par\medskip\noindent{\small\textit{b. WF4-a regional rows (C5 README 2.2)}}\par\nopagebreak
{\fontsize{7.5}{9}\selectfont\setlength{\tabcolsep}{4pt}\setlength{\LTpre}{2pt}\setlength{\LTpost}{4pt}
\begin{longtable}{@{}>{\raggedright\arraybackslash}p{\dimexpr 0.050\linewidth-8pt\relax}>{\raggedright\arraybackslash}p{\dimexpr 0.088\linewidth-8pt\relax}>{\raggedright\arraybackslash}p{\dimexpr 0.068\linewidth-8pt\relax}>{\raggedright\arraybackslash}p{\dimexpr 0.095\linewidth-8pt\relax}>{\raggedright\arraybackslash}p{\dimexpr 0.338\linewidth-8pt\relax}>{\raggedright\arraybackslash}p{\dimexpr 0.181\linewidth-8pt\relax}>{\raggedright\arraybackslash}p{\dimexpr 0.181\linewidth-8pt\relax}@{}}
\toprule
\textbf{ID} & \textbf{Target} & \textbf{$n$} & \textbf{$\Delta$} & \textbf{CI cell / block-equal} & \textbf{Verdict} & \textbf{Common-CI verdict} \\
\midrule\endfirsthead
\toprule
\textbf{ID} & \textbf{Target} & \textbf{$n$} & \textbf{$\Delta$} & \textbf{CI cell / block-equal} & \textbf{Verdict} & \textbf{Common-CI verdict} \\
\midrule\endhead
\bottomrule\endlastfoot
E6 & Canada & 40 & -2.10 & [-2.71, -0.95] /\allowbreak{} [-2.56, -0.98] & lower error & lower error \\
E7 & Canada & 160 & -2.72 & [-3.41, -1.51] /\allowbreak{} [-3.21, -1.40] & lower error & lower error \\
E8 & Canada & 10 & -0.86 & [-1.35, 0.11] /\allowbreak{} [-1.53, -0.15] & undetermined & undetermined \\
E9 & Canada & all & -3.06 & [-3.81, -1.74] /\allowbreak{} [-3.50, -1.49] & lower error & lower error \\
E10 & Lena Delta & 40 & 0.35 & [0.03, 0.67] /\allowbreak{} [-0.21, 0.26] & undetermined & undetermined \\
E11 & Lena Delta & 160 & 0.07 & [-0.36, 0.48] /\allowbreak{} [-0.61, -0.02] & undetermined & undetermined \\
E12 & Lena Delta & 10 & 0.70 & [0.56, 1.03] /\allowbreak{} [0.80, 1.31] & higher error & higher error \\
E13 & Lena Delta & all & 0.46 & [-1.15, 1.81] /\allowbreak{} [-0.69, 1.02] & undetermined & undetermined \\
E14 & W Russia & 10, all & -0.42, 0.49 & [-3.47, 3.20] /\allowbreak{} [-3.20, 3.47]; [-1.90, 3.43] /\allowbreak{} [-1.71, 3.93] & undetermined, undetermined & undetermined, undetermined \\
E15 & E Russia & 10, all & -0.43, -0.84 & [-2.08, 0.69] /\allowbreak{} [-2.12, -0.19]; [-3.23, 0.46] /\allowbreak{} [-3.60, -1.14] & undetermined, undetermined & undetermined, undetermined \\
\end{longtable}}
\par\medskip\noindent{\small\textit{c. WF4-b, targets with lower error than P1 (C5 README 2.3)}}\par\nopagebreak
{\fontsize{7.5}{9}\selectfont\setlength{\tabcolsep}{4pt}\setlength{\LTpre}{2pt}\setlength{\LTpost}{4pt}
\begin{longtable}{@{}>{\raggedright\arraybackslash}p{\dimexpr 0.050\linewidth-8pt\relax}>{\raggedright\arraybackslash}p{\dimexpr 0.173\linewidth-8pt\relax}>{\raggedright\arraybackslash}p{\dimexpr 0.184\linewidth-8pt\relax}>{\raggedright\arraybackslash}p{\dimexpr 0.247\linewidth-8pt\relax}>{\raggedright\arraybackslash}p{\dimexpr 0.184\linewidth-8pt\relax}>{\raggedright\arraybackslash}p{\dimexpr 0.162\linewidth-8pt\relax}@{}}
\toprule
\textbf{ID} & \textbf{$n$} & \textbf{W lower error} & \textbf{R1(0.25) lower error} & \textbf{W higher error} & \textbf{R1(0.25) higher error} \\
\midrule\endfirsthead
\toprule
\textbf{ID} & \textbf{$n$} & \textbf{W lower error} & \textbf{R1(0.25) lower error} & \textbf{W higher error} & \textbf{R1(0.25) higher error} \\
\midrule\endhead
\bottomrule\endlastfoot
E18 & 160(17/\allowbreak{}30) & 3(AL-5 i, CA-3 i, Canada) & 4(AL-5 i, CA-2 x, Canada, Lena Delta) & 3(AL-1 x, AL-2 i, AL-2 x) & 3(AL-1 i, AL-1 x, AL-2 i) \\
E19 & all(21/\allowbreak{}30) & 2(Canada, Tibetan Plateau) & 8(AL-2 x, CA-2 i, CA-2 x, Canada, LE-1 x, Lena Delta, Central Russia(LGD edition, Russia\_\allowbreak{}C\textasciitilde{}lgd), Tibetan Plateau) & 3(AL-1 i, AL-2 x, AL-3 x) & 1(AL-1 x) \\
\end{longtable}}
\par\medskip\noindent{\small\textit{d. W against the physics baselines (curve table, auxiliary) (C5 README 2.4)}}\par\nopagebreak
{\fontsize{7.5}{9}\selectfont\setlength{\tabcolsep}{4pt}\setlength{\LTpre}{2pt}\setlength{\LTpost}{4pt}
\begin{longtable}{@{}>{\raggedright\arraybackslash}p{\dimexpr 0.050\linewidth-8pt\relax}>{\raggedright\arraybackslash}p{\dimexpr 0.091\linewidth-8pt\relax}>{\raggedright\arraybackslash}p{\dimexpr 0.047\linewidth-8pt\relax}>{\raggedright\arraybackslash}p{\dimexpr 0.297\linewidth-8pt\relax}>{\raggedright\arraybackslash}p{\dimexpr 0.115\linewidth-8pt\relax}>{\raggedright\arraybackslash}p{\dimexpr 0.284\linewidth-8pt\relax}>{\raggedright\arraybackslash}p{\dimexpr 0.115\linewidth-8pt\relax}@{}}
\toprule
\textbf{ID} & \textbf{Target} & \textbf{$n$} & \textbf{W $-$ P0 cell / block-equal} & \textbf{Verdict} & \textbf{W $-$ P1 cell / block-equal} & \textbf{Verdict} \\
\midrule\endfirsthead
\toprule
\textbf{ID} & \textbf{Target} & \textbf{$n$} & \textbf{W $-$ P0 cell / block-equal} & \textbf{Verdict} & \textbf{W $-$ P1 cell / block-equal} & \textbf{Verdict} \\
\midrule\endhead
\bottomrule\endlastfoot
E20 & Lena Delta & 10 & 0.89 [0.33, 1.72] /\allowbreak{} [0.87, 1.94] & higher error & 0.67 [0.47, 0.89] /\allowbreak{} [0.47, 0.94] & higher error \\
E21 & Lena Delta & 40 & 0.20 [-0.37, 1.04] /\allowbreak{} [0.34, 1.39] & undetermined & 0.22 [-0.20, 0.48] /\allowbreak{} [-0.70, -0.08] & undetermined \\
E22 & Lena Delta & 160 & 0.02 [-0.51, 0.81] /\allowbreak{} [0.07, 1.13] & undetermined & -0.07 [-0.64, 0.29] /\allowbreak{} [-1.19, -0.39] & undetermined \\
E23 & Canada & 10 & 0.51 [-0.97, 1.52] /\allowbreak{} [-1.54, 0.62] & undetermined & -1.75 [-2.50, -0.63] /\allowbreak{} [-2.84, -0.95] & lower error \\
E24 & Canada & 40 & 0.90 [-0.53, 1.76] /\allowbreak{} [-1.24, 1.28] & undetermined & -3.08 [-3.93, -1.66] /\allowbreak{} [-3.82, -1.71] & lower error \\
E25 & Canada & 160 & 0.35 [-1.20, 1.11] /\allowbreak{} [-1.84, 0.83] & undetermined & -3.84 [-4.80, -2.29] /\allowbreak{} [-4.55, -2.22] & lower error \\
E26 & Canada & all & 0.32 [-1.34, 1.14] /\allowbreak{} [-1.91, 0.99] & undetermined & -4.32 [-5.35, -2.67] /\allowbreak{} [-4.96, -2.41] & lower error \\
E27a & Alaska(x) & 10 & 0.91 [0.58, 1.42] /\allowbreak{} [1.14, 2.13] & higher error & 0.93 [0.67, 1.19] /\allowbreak{} [0.55, 0.98] & higher error \\
E27b & Alaska(x) & 40 & 0.98 [0.44, 1.57] /\allowbreak{} [0.93, 1.91] & higher error & 0.75 [0.28, 1.09] /\allowbreak{} [-0.17, 0.29] & undetermined \\
E27c & Alaska(x) & 160 & 2.14 [0.28, 3.18] /\allowbreak{} [0.71, 1.94] & higher error & 2.19 [0.23, 3.11] /\allowbreak{} [-0.45, 0.64] & undetermined \\
E27d & Alaska(x) & all & 0.26 [-0.41, 1.12] /\allowbreak{} [-0.04, 1.80] & undetermined & 0.19 [-0.54, 0.82] /\allowbreak{} [-1.30, 0.30] & undetermined \\
E28a & W Russia & 10 & -12.94 [-17.44, -8.02] /\allowbreak{} [-18.63, -7.89] & lower error & -0.97 [-4.09, 2.72] /\allowbreak{} [-3.64, 3.42] & undetermined \\
E28b & W Russia & all & -13.49 [-17.64, -8.93] /\allowbreak{} [-19.24, -8.62] & lower error & 0.01 [-2.48, 2.99] /\allowbreak{} [-2.09, 3.85] & undetermined \\
E28c & E Russia & 10 & 0.00 [0.00, 0.00] /\allowbreak{} [0.00, 0.00](file value absolute value 1e-14 below) & equivalent & 0.34 [-1.30, 0.92] /\allowbreak{} [-1.47, 0.10] & undetermined \\
E28d & E Russia & all & 0.00 [0.00, 0.00] /\allowbreak{} [0.00, 0.00](same) & equivalent & 0.09 [-2.46, 1.09] /\allowbreak{} [-2.92, -0.58] & undetermined \\
\end{longtable}}
\par\medskip\noindent{\small\textit{e. Methods chosen by W (counts over splits and draws; this check) (C5 README 2.5)}}\par\nopagebreak
{\fontsize{7.5}{9}\selectfont\setlength{\tabcolsep}{4pt}\setlength{\LTpre}{2pt}\setlength{\LTpost}{4pt}
\begin{longtable}{@{}>{\raggedright\arraybackslash}p{\dimexpr 0.067\linewidth-8pt\relax}>{\raggedright\arraybackslash}p{\dimexpr 0.113\linewidth-8pt\relax}>{\raggedright\arraybackslash}p{\dimexpr 0.233\linewidth-8pt\relax}>{\raggedright\arraybackslash}p{\dimexpr 0.233\linewidth-8pt\relax}>{\raggedright\arraybackslash}p{\dimexpr 0.181\linewidth-8pt\relax}>{\raggedright\arraybackslash}p{\dimexpr 0.172\linewidth-8pt\relax}@{}}
\toprule
\textbf{ID} & \textbf{Target} & \textbf{$n$ 10} & \textbf{$n$ 40} & \textbf{$n$ 160} & \textbf{All} \\
\midrule\endfirsthead
\toprule
\textbf{ID} & \textbf{Target} & \textbf{$n$ 10} & \textbf{$n$ 40} & \textbf{$n$ 160} & \textbf{All} \\
\midrule\endhead
\bottomrule\endlastfoot
E29 & Canada & R1@1.0 19, P0 4, P2 2 & R1@1.0 22, P0 2, P2 1 & R1@1.0 24, P0 1 & R1@1.0 5 \\
E30 & Lena Delta & P0 8, P2 7, R1@1.0 4, R1@0.25 2, D1 2, R2@0.25 2 & P0 11, R1@1.0 5, P2 4, R1@0.25 3, R2@0.25 2 & P0 12, R1@1.0 8, R1@0.25 3, P2 2 & P0 2, R1@0.25 2, R1@1.0 1 \\
E31 & W Russia & P2 23, R1@1.0 2 & none & none & P2 4, R1@1.0 1 \\
E32 & E Russia & P0 25 & none & none & P0 5 \\
\end{longtable}}
\par\medskip\noindent{\small\textit{f. Point estimates against fixed recipes (this check) (C5 README 2.5)}}\par\nopagebreak
{\fontsize{7.5}{9}\selectfont\setlength{\tabcolsep}{4pt}\setlength{\LTpre}{2pt}\setlength{\LTpost}{4pt}
\begin{longtable}{@{}>{\raggedright\arraybackslash}p{\dimexpr 0.077\linewidth-8pt\relax}>{\raggedright\arraybackslash}p{\dimexpr 0.142\linewidth-8pt\relax}>{\raggedright\arraybackslash}p{\dimexpr 0.218\linewidth-8pt\relax}>{\raggedright\arraybackslash}p{\dimexpr 0.277\linewidth-8pt\relax}>{\raggedright\arraybackslash}p{\dimexpr 0.286\linewidth-8pt\relax}@{}}
\toprule
\textbf{ID} & \textbf{Target} & \textbf{$n$} & \textbf{W $-$ R1(0.25)} & \textbf{W $-$ R1(1.0)} \\
\midrule\endfirsthead
\toprule
\textbf{ID} & \textbf{Target} & \textbf{$n$} & \textbf{W $-$ R1(0.25)} & \textbf{W $-$ R1(1.0)} \\
\midrule\endhead
\bottomrule\endlastfoot
E33 & Canada & 10, 40, 160, all & -0.86, -2.10, -2.72, -3.06 & 0.98, 0.12, 0.00, 0.00 \\
E34 & Lena Delta & 10, 40, 160, all & 0.70, 0.35, 0.07, 0.46 & 0.30, 0.35, -0.12, 1.11 \\
E35 & Alaska(x) & 10, 40, 160, all & 0.36, 0.29, 1.89, 0.54 & -8.78, -5.88, -3.91, -0.15 \\
E36 & E Russia & 10, all & -0.43, -0.84 & -4.63, -5.29 \\
E37 & W Russia & 10, all & -0.42, 0.49 & 0.31, 1.04 \\
\end{longtable}}
\par\medskip\noindent{\small\textit{g. Bias diagnosis records (WF4-c, WF0) (C2 README 2.1)}}\par\nopagebreak
{\fontsize{7.5}{9}\selectfont\setlength{\tabcolsep}{4pt}\setlength{\LTpre}{2pt}\setlength{\LTpost}{4pt}
\begin{longtable}{@{}>{\raggedright\arraybackslash}p{\dimexpr 0.046\linewidth-8pt\relax}>{\raggedright\arraybackslash}p{\dimexpr 0.294\linewidth-8pt\relax}>{\raggedright\arraybackslash}p{\dimexpr 0.294\linewidth-8pt\relax}>{\raggedright\arraybackslash}p{\dimexpr 0.365\linewidth-8pt\relax}@{}}
\toprule
\textbf{ID} & \textbf{Contrast or quantity} & \textbf{Value (cm unless stated): cell-weighted [CI] / block-equal [CI]} & \textbf{Verdict} \\
\midrule\endfirsthead
\toprule
\textbf{ID} & \textbf{Contrast or quantity} & \textbf{Value (cm unless stated): cell-weighted [CI] / block-equal [CI]} & \textbf{Verdict} \\
\midrule\endhead
\bottomrule\endlastfoot
A1 & recalibration gain(P0 − P1) best ML gain(P0 − best ML) rank correlation, all labels & 0.66, p 8.15e-05, 30 targets & none(label: 'post hoc descriptive(verdict word none)'). note: best ML five recipe among post hoc selection optimistic \\
A2 & P0 RMSE best ML gain rank correlation & 0.51, p 0.00362 & none(post hoc descriptive) \\
A3 & labels 10 diagnostic rule W gain(P0 − W, all labels) rank correlation & 0.38 [0.17, 0.55], 28 targets, resamples 10000 times, one-sided p 0.0002 & supported. scope == 'verdict' rows verdict column: 'supported: positive rank correlation(ρ 0.38, CI [0.17, 0.55])designed after viewing results'\textsuperscript{\dag} \\
A4 & signed bias gain rank correlation & −0.01 [−0.17, 0.18] & auxiliary(verdict word none) \\
A5 & diagnostic fixed R1(λ 0.25) gain(P0 − R1) rank correlation & 0.47(CI column none) & auxiliary \\
A6 & P1 − P0, labels 10, main 4 regions mean(AB4) & −2.45 [−3.04, −1.88], block-equal −2.62, Holm p 0.001 & lower error. abstract rule (a). S-a calculation same quantity(draws seed differs) analytic P1 − P0 splits distribution viewed(2026-09-30 03:25) \\
A7 & R1 − P1, labels 10, main 4 regions mean(AB5) & −0.18 [−0.53, −0.01], block-equal −0.29, Holm p 0.136 & lower error(before correction). abstract rule (b) 'no difference established' \\
A8 & R1 − P0, all labels, main 4 regions mean(AB8) & −2.64 [−3.48, −1.85], block-equal −2.64, Holm p 0.001 & lower error. abstract rule (a) \\
A9 & R1 − P1, all labels, main 4 regions mean(AB9) & −0.40 [−0.76, −0.22], block-equal −0.51, Holm p 0.001 & lower error. size below 0.5 cm \\
A10 & A8 − A9 = P1 − P0(all labels, main 4 regions mean) & −2.23(A8 85 \%) & none(arithmetic) \\
A11 & AB8 regions rows(R1 − P0, all labels) & Lena Delta −0.78 [−1.38, 0.08], Canada +3.38 [1.10, 4.40], W Russia −13.98 [−16.05, −12.06], E Russia +0.84 [−0.43, 3.19] & W Russia lower error, Canada higher error, Lena Delta·E Russia undetermined(improvement significant regions 1/\allowbreak{}4) \\
A12 & AB9 regions rows(R1 − P1, all labels) & Lena Delta −0.80 [−1.51, −0.44], Canada −1.26 [−1.57, −0.90], W Russia −0.48 [−1.05, 0.01], E Russia +0.93 [0.00(fourth decimal 0.0002), 1.28] & Lena Delta and Canada lower error, W Russia undetermined, E Russia higher error. same contrast LGX auxiliary column(data/\allowbreak{}processed/\allowbreak{}paper\_\allowbreak{}figs/\allowbreak{}fig3\_\allowbreak{}c.csv, L4 rows) E Russia [−0.02, 1.28] dependence)\textsuperscript{\dag} \\
A13 & R1 − P0, all labels, Lena Delta and Canada·Alaska(x) 3 regions mean & +0.77 [0.02, 1.28], block-equal +1.15 [0.66, 1.63] & higher error. ci\_\allowbreak{}dependence 'depends on split independence' \\
A14 & R1 − P*, all labels, main 4 regions mean(P* = P0@tddm) & −1.64 [−2.95, −0.74], block-equal −2.23 [−3.17, −1.24] & lower error \\
\end{longtable}}
\par\medskip\noindent{\small\textit{h. Source cross-validation tuning of neural networks, tuned minus default (LGF-N2; cm, four main regions) (C5 README 2.6)}}\par\nopagebreak
{\fontsize{7.5}{9}\selectfont\setlength{\tabcolsep}{4pt}\setlength{\LTpre}{2pt}\setlength{\LTpost}{4pt}
\begin{longtable}{@{}>{\raggedright\arraybackslash}p{\dimexpr 0.040\linewidth-8pt\relax}>{\raggedright\arraybackslash}p{\dimexpr 0.085\linewidth-8pt\relax}>{\raggedright\arraybackslash}p{\dimexpr 0.092\linewidth-8pt\relax}>{\raggedright\arraybackslash}p{\dimexpr 0.099\linewidth-8pt\relax}>{\raggedright\arraybackslash}p{\dimexpr 0.092\linewidth-8pt\relax}>{\raggedright\arraybackslash}p{\dimexpr 0.099\linewidth-8pt\relax}>{\raggedright\arraybackslash}p{\dimexpr 0.099\linewidth-8pt\relax}>{\raggedright\arraybackslash}p{\dimexpr 0.099\linewidth-8pt\relax}>{\raggedright\arraybackslash}p{\dimexpr 0.099\linewidth-8pt\relax}>{\raggedright\arraybackslash}p{\dimexpr 0.099\linewidth-8pt\relax}>{\raggedright\arraybackslash}p{\dimexpr 0.099\linewidth-8pt\relax}@{}}
\toprule
\textbf{ID} & \textbf{Learner} & \textbf{R1 $n$ 0 $\lambda$ 0.25} & \textbf{R1 $n$ 10 $\lambda$ 0.25} & \textbf{R1 all $\lambda$ 0.25} & \textbf{D0 $n$ 0} & \textbf{D0 $n$ 10} & \textbf{D0 all} & \textbf{R1 $n$ 0 $\lambda$ 1.0} & \textbf{R1 $n$ 10 $\lambda$ 1.0} & \textbf{R1 all $\lambda$ 1.0} \\
\midrule\endfirsthead
\toprule
\textbf{ID} & \textbf{Learner} & \textbf{R1 $n$ 0 $\lambda$ 0.25} & \textbf{R1 $n$ 10 $\lambda$ 0.25} & \textbf{R1 all $\lambda$ 0.25} & \textbf{D0 $n$ 0} & \textbf{D0 $n$ 10} & \textbf{D0 all} & \textbf{R1 $n$ 0 $\lambda$ 1.0} & \textbf{R1 $n$ 10 $\lambda$ 1.0} & \textbf{R1 all $\lambda$ 1.0} \\
\midrule\endhead
\bottomrule\endlastfoot
E38 & MLP & 0.06 equivalent & 0.43 undetermined & 0.25 equivalent‡ & 1.51 higher error‡ & 2.96 higher error & 1.10 higher error‡ & -0.27 undetermined & 0.70 undetermined & -0.01 undetermined \\
E39 & multi-head MLP & 0.24 higher error & 0.05 equivalent & -0.02 equivalent & 1.33 higher error & 2.56 higher error & 1.67 higher error & 0.75 higher error‡ & -0.22 undetermined & 0.12 undetermined \\
E40 & reduced FT-T & 0.01 equivalent & 0.02 equivalent & -0.19 lower error‡ & -0.82 undetermined & -0.36 undetermined & -0.16 undetermined & 0.20 undetermined & 0.10 undetermined & -0.76 lower error‡ \\
E41 & RealMLP & -0.06 equivalent‡ & 0.19 equivalent‡ & -0.08 equivalent & 0.51 undetermined & -0.23 undetermined & -0.13 undetermined & 0.23 undetermined & 0.74 undetermined & -0.34 undetermined \\
\end{longtable}}
\par\medskip\noindent{\small\textit{i. LGF-N2 verdicts by learner (C5 README 2.6)}}\par\nopagebreak
{\fontsize{7.5}{9}\selectfont\setlength{\tabcolsep}{4pt}\setlength{\LTpre}{2pt}\setlength{\LTpost}{4pt}
\begin{longtable}{@{}>{\raggedright\arraybackslash}p{\dimexpr 0.061\linewidth-8pt\relax}>{\raggedright\arraybackslash}p{\dimexpr 0.140\linewidth-8pt\relax}>{\raggedright\arraybackslash}p{\dimexpr 0.467\linewidth-8pt\relax}>{\raggedright\arraybackslash}p{\dimexpr 0.332\linewidth-8pt\relax}@{}}
\toprule
\textbf{ID} & \textbf{Learner} & \textbf{Verdict wording (summary of the file values)} & \textbf{Default kept (k\_default)} \\
\midrule\endfirsthead
\toprule
\textbf{ID} & \textbf{Learner} & \textbf{Verdict wording (summary of the file values)} & \textbf{Default kept (k\_default)} \\
\midrule\endhead
\bottomrule\endlastfoot
E42 & MLP & source-CV tuning increased the transfer error(D0 n 0·10·all higher error)(contrasts with insufficient precision 6/\allowbreak{}9) & R 0/\allowbreak{}4, D 1/\allowbreak{}4 \\
E43 & multi-head MLP & source-CV tuning increased the transfer error(R1 n 0 λ 0.25·λ 1.0, D0 n 0·10·all higher error)(contrasts with insufficient precision 4/\allowbreak{}9) & R 2/\allowbreak{}4, D 1/\allowbreak{}4 \\
E44 & reduced FT-T & the tuned version had lower error(R1 all λ 0.25·λ 1.0 lower error)(contrasts with insufficient precision 6/\allowbreak{}9) & R 2/\allowbreak{}4, D 2/\allowbreak{}4 \\
E45 & RealMLP & precision insufficient(equivalence not decidable, contrast 6/\allowbreak{}9) & R 2/\allowbreak{}4, D 2/\allowbreak{}4 \\
E46 & all & condition for a learner-general sentence not met(learners meeting the robust-sentence rule: none) &  \\
\end{longtable}}

\clearpage
\phantomsection\label{tab:S9}
\noindent\textbf{Supplementary Table S9.} Label placement strategies (WF2, WF8, L43).\par\smallskip
\noindent{\footnotesize Strategy minus random cells (S1) at the same budget; within-region designs (target suffix $|r$) and transfer ($|x$). Values in cm; RMSE A / B are the RMSE of the strategy and of random cells. WF2-a rows marked descriptive are reported without verdict words in the main text. Method codes (internal identifiers allowed in the Supplementary Information): P0 source-coefficient Stefan model; P* year-matched Stefan model; P1 recalibrated Stefan model ($\kappa = 10$); P2 local least-squares Stefan model; P1* or P1@ed soil-property recalibrated Stefan model; R0 source anchor plus residual ML; R1 recalibrated anchor plus residual ML; R2 augmentation plus anchor plus residual; D0 direct ML; D1 physics pseudo-label augmentation; F1a, F1k, F1n physics-input ML; RM multiplicative residual; Re Stefan-CCI mean anchor plus residual; W selection rule (cross-validation within the target labels); S1 random cells, S2 block stratification, S3 covariate-cluster stratification, S4 covariate max-min distance, S5 source-extrapolation first, S6 sequential selection by prediction variance, S7 $\sqrt{\mathrm{TDD}}$ first. Target suffixes: $|x$ parent region excluded from the source, $|i$ rest of the parent kept, $|r$ within-region design. $n$ = number of target labels; all = all labels of half A. Verdicts: lower error, higher error, equivalent ($\pm$0.5~cm), undetermined; both 95\% block-bootstrap CIs (cell-weighted and block-equal) must agree. \dag~a note of the source row in Korean was omitted; see the cited README section.}\par

\par\medskip\noindent{\small\textit{a. Block stratification (S2) and covariate max-min distance (S4), WF2-a (C6 README 2.1)}}\par\nopagebreak
{\fontsize{7.5}{9}\selectfont\setlength{\tabcolsep}{4pt}\setlength{\LTpre}{2pt}\setlength{\LTpost}{4pt}
\begin{longtable}{@{}>{\raggedright\arraybackslash}p{\dimexpr 0.049\linewidth-8pt\relax}>{\raggedright\arraybackslash}p{\dimexpr 0.345\linewidth-8pt\relax}>{\raggedright\arraybackslash}p{\dimexpr 0.151\linewidth-8pt\relax}>{\raggedright\arraybackslash}p{\dimexpr 0.151\linewidth-8pt\relax}>{\raggedright\arraybackslash}p{\dimexpr 0.184\linewidth-8pt\relax}>{\raggedright\arraybackslash}p{\dimexpr 0.119\linewidth-8pt\relax}@{}}
\toprule
\textbf{ID} & \textbf{Target; contrast (role)} & \textbf{$\Delta$ cell [CI]} & \textbf{$\Delta$ block-equal [CI]} & \textbf{RMSE A / B} & \textbf{Verdict} \\
\midrule\endfirsthead
\toprule
\textbf{ID} & \textbf{Target; contrast (role)} & \textbf{$\Delta$ cell [CI]} & \textbf{$\Delta$ block-equal [CI]} & \textbf{RMSE A / B} & \textbf{Verdict} \\
\midrule\endhead
\bottomrule\endlastfoot
A1 & Canada|\allowbreak{}r; S2-S1|\allowbreak{}R1(0.25)|\allowbreak{}n20 (descriptive) & -3.46 [-4.80, -0.69] & -2.66 [-3.92, -1.40] & 27.76 /\allowbreak{} 31.22 & lower error \\
A2 & Canada|\allowbreak{}r; S2-S1|\allowbreak{}R1(0.25)|\allowbreak{}n50 (descriptive) & -2.74 [-3.52, -0.87] & -1.73 [-2.69, -0.79] & 28.27 /\allowbreak{} 31.01 & lower error \\
A3 & Canada|\allowbreak{}r; S2-S1|\allowbreak{}R1(0.25)|\allowbreak{}n100 (descriptive) & -2.54 [-3.22, -1.12] & -1.98 [-2.69, -1.30] & 28.75 /\allowbreak{} 31.29 & lower error \\
A4 & Canada|\allowbreak{}r; S2-S1|\allowbreak{}R1(0.25)|\allowbreak{}n200 (descriptive) & -1.44 [-2.00, -0.42] & -1.28 [-1.75, -0.80] & 29.16 /\allowbreak{} 30.59 & lower error \\
A5 & Canada|\allowbreak{}r; S4-S1|\allowbreak{}R1(0.25)|\allowbreak{}n20 (descriptive) & -2.24 [-3.41, 0.03] & -2.10 [-3.28, -0.95] & 28.98 /\allowbreak{} 31.22 & undetermined \\
A6 & Canada|\allowbreak{}r; S4-S1|\allowbreak{}R1(0.25)|\allowbreak{}n50 (descriptive) & -2.63 [-3.49, -0.70] & -1.64 [-2.55, -0.76] & 28.38 /\allowbreak{} 31.01 & lower error \\
A7 & Canada|\allowbreak{}r; S4-S1|\allowbreak{}R1(0.25)|\allowbreak{}n100 (descriptive) & -3.21 [-4.27, -1.10] & -2.24 [-3.09, -1.40] & 28.09 /\allowbreak{} 31.29 & lower error \\
A8 & Canada|\allowbreak{}r; S4-S1|\allowbreak{}R1(0.25)|\allowbreak{}n200 (descriptive) & -1.25 [-1.75, -0.37] & -1.20 [-1.63, -0.77] & 29.34 /\allowbreak{} 30.59 & lower error \\
A9 & AL-1|\allowbreak{}r; S4-S1|\allowbreak{}R1(0.25)|\allowbreak{}n50 (descriptive) & -0.77 [-2.64, -0.36] & -1.82 [-2.87, -0.77] & 16.36 /\allowbreak{} 17.13 & lower error \\
A10 & AL-1|\allowbreak{}r; S4-S1|\allowbreak{}R1(0.25)|\allowbreak{}n100 (descriptive) & -0.50 [-2.29, -0.20] & -1.99 [-2.86, -1.12] & 16.13 /\allowbreak{} 16.64 & lower error \\
A11 & AL-1|\allowbreak{}r; S4-S1|\allowbreak{}R1(0.25)|\allowbreak{}n200 (descriptive) & -0.32 [-1.09, -0.16] & -1.04 [-1.46, -0.62] & 16.23 /\allowbreak{} 16.55 & lower error \\
A12 & AL-2|\allowbreak{}r; S2-S1|\allowbreak{}R1(0.25)|\allowbreak{}n50 (descriptive) & -0.70 [-1.29, -0.53] & -0.67 [-1.34, -0.01] & 12.34 /\allowbreak{} 13.04 & lower error \\
A13 & AL-2|\allowbreak{}r; S2-S1|\allowbreak{}R1(0.25)|\allowbreak{}n100 (descriptive) & -0.65 [-1.31, -0.45] & -0.73 [-1.45, -0.03] & 12.34 /\allowbreak{} 12.98 & lower error \\
A14 & AL-2|\allowbreak{}r; S2-S1|\allowbreak{}R1(0.25)|\allowbreak{}n200 (descriptive) & -0.36 [-1.02, -0.21] & -0.55 [-1.14, 0.02] & 12.31 /\allowbreak{} 12.67 & undetermined \\
A15 & AL-2|\allowbreak{}r; S4-S1|\allowbreak{}R1(0.25)|\allowbreak{}n50 (descriptive) & -0.39 [-0.97, -0.25] & -0.51 [-0.94, -0.09] & 12.65 /\allowbreak{} 13.04 & lower error \\
A16 & Lena|\allowbreak{}r; S2-S1|\allowbreak{}R1(0.25)|\allowbreak{}n50 (descriptive) & 1.47 [0.37, 2.29] & -0.05 [-1.14, 0.98] & 23.12 /\allowbreak{} 21.65 & undetermined \\
A17 & Lena|\allowbreak{}r; S2-S1|\allowbreak{}R1(0.25)|\allowbreak{}n100 (descriptive) & 1.65 [0.90, 2.40] & 0.88 [0.15, 1.57] & 22.83 /\allowbreak{} 21.17 & higher error \\
A18 & Lena|\allowbreak{}r; S2-S1|\allowbreak{}R1(0.25)|\allowbreak{}n200 (descriptive) & 1.79 [0.84, 2.46] & 0.73 [-0.12, 1.54] & 22.52 /\allowbreak{} 20.73 & undetermined \\
A19 & Lena|\allowbreak{}r; S2-S1|\allowbreak{}R1(0.25)|\allowbreak{}n500 (descriptive) & 1.48 [0.70, 2.12] & 0.76 [-0.04, 1.52] & 22.23 /\allowbreak{} 20.75 & undetermined \\
A20 & Lena|\allowbreak{}r; S4-S1|\allowbreak{}R1(0.25)|\allowbreak{}n50 (descriptive) & 0.11 [-1.12, 1.01] & -0.62 [-1.35, 0.09] & 21.75 /\allowbreak{} 21.65 & undetermined \\
A21 & Lena|\allowbreak{}r; S4-S1|\allowbreak{}R1(0.25)|\allowbreak{}n100 (descriptive) & 0.13 [-0.57, 0.72] & -0.12 [-0.61, 0.36] & 21.30 /\allowbreak{} 21.17 & undetermined \\
A22 & Lena|\allowbreak{}r; S4-S1|\allowbreak{}R1(0.25)|\allowbreak{}n200 (descriptive) & 0.44 [-0.14, 0.78] & -0.20 [-0.65, 0.24] & 21.17 /\allowbreak{} 20.73 & undetermined \\
A23 & Alaska|\allowbreak{}r; S2-S1|\allowbreak{}R1(0.25)|\allowbreak{}n50 (main) & 0.45 [-0.23, 0.77] & -1.30 [-1.97, -0.62] & 14.73 /\allowbreak{} 14.28 & undetermined \\
A24 & Alaska|\allowbreak{}r; S2-S1|\allowbreak{}R1(0.25)|\allowbreak{}n100 (main) & 0.34 [-0.45, 0.66] & -1.73 [-2.47, -0.97] & 14.68 /\allowbreak{} 14.33 & undetermined \\
A25 & Alaska|\allowbreak{}r; S2-S1|\allowbreak{}R1(0.25)|\allowbreak{}n200 (main) & 0.61 [-0.05, 0.89] & -1.08 [-1.64, -0.51] & 14.66 /\allowbreak{} 14.05 & undetermined \\
A26 & Alaska|\allowbreak{}r; S3-S1|\allowbreak{}R1(0.25)|\allowbreak{}n50 (main) & -0.04 [-0.23, 0.08] & -0.35 [-0.58, -0.14] & 14.23 /\allowbreak{} 14.28 & undetermined \\
A27 & Alaska|\allowbreak{}r; S3-S1|\allowbreak{}R1(0.25)|\allowbreak{}n100 (main) & -0.11 [-0.33, -0.03] & -0.62 [-0.88, -0.37] & 14.23 /\allowbreak{} 14.33 & lower error \\
A28 & Alaska|\allowbreak{}r; S3-S1|\allowbreak{}R1(0.25)|\allowbreak{}n200 (main) & 0.08 [-0.01, 0.17] & 0.16 [0.02, 0.29] & 14.13 /\allowbreak{} 14.05 & equivalent \\
A29 & Alaska|\allowbreak{}r; S4-S1|\allowbreak{}R1(0.25)|\allowbreak{}n50 (main) & 0.28 [-0.34, 0.62] & -1.15 [-1.79, -0.49] & 14.56 /\allowbreak{} 14.28 & undetermined \\
A30 & Alaska|\allowbreak{}r; S4-S1|\allowbreak{}R1(0.25)|\allowbreak{}n100 (main) & 0.44 [-0.40, 0.81] & -1.71 [-2.48, -0.94] & 14.77 /\allowbreak{} 14.33 & undetermined \\
A31 & Alaska|\allowbreak{}r; S4-S1|\allowbreak{}R1(0.25)|\allowbreak{}n200 (main) & 0.74 [0.01, 1.09] & -1.09 [-1.68, -0.49] & 14.79 /\allowbreak{} 14.05 & undetermined \\
\end{longtable}}
\par\medskip\noindent{\small\textit{b. Sequential active selection (S6), WF2-a (C6 README 2.2)}}\par\nopagebreak
{\fontsize{7.5}{9}\selectfont\setlength{\tabcolsep}{4pt}\setlength{\LTpre}{2pt}\setlength{\LTpost}{4pt}
\begin{longtable}{@{}>{\raggedright\arraybackslash}p{\dimexpr 0.049\linewidth-8pt\relax}>{\raggedright\arraybackslash}p{\dimexpr 0.345\linewidth-8pt\relax}>{\raggedright\arraybackslash}p{\dimexpr 0.151\linewidth-8pt\relax}>{\raggedright\arraybackslash}p{\dimexpr 0.151\linewidth-8pt\relax}>{\raggedright\arraybackslash}p{\dimexpr 0.184\linewidth-8pt\relax}>{\raggedright\arraybackslash}p{\dimexpr 0.119\linewidth-8pt\relax}@{}}
\toprule
\textbf{ID} & \textbf{Target; contrast (role)} & \textbf{$\Delta$ cell [CI]} & \textbf{$\Delta$ block-equal [CI]} & \textbf{RMSE A / B} & \textbf{Verdict} \\
\midrule\endfirsthead
\toprule
\textbf{ID} & \textbf{Target; contrast (role)} & \textbf{$\Delta$ cell [CI]} & \textbf{$\Delta$ block-equal [CI]} & \textbf{RMSE A / B} & \textbf{Verdict} \\
\midrule\endhead
\bottomrule\endlastfoot
B1 & Alaska|\allowbreak{}r; S6-S1|\allowbreak{}R1(0.25)|\allowbreak{}n50 (main) & 1.88 [1.58, 2.29] & 2.40 [1.97, 2.81] & 16.15 /\allowbreak{} 14.28 & higher error \\
B2 & Alaska|\allowbreak{}r; S6-S1|\allowbreak{}R1(0.25)|\allowbreak{}n100 (main) & 1.25 [1.02, 1.50] & 1.15 [0.92, 1.39] & 15.58 /\allowbreak{} 14.33 & higher error \\
B3 & Alaska|\allowbreak{}r; S6-S1|\allowbreak{}R1(0.25)|\allowbreak{}n200 (main) & 0.69 [0.50, 0.91] & 0.60 [0.29, 0.89] & 14.74 /\allowbreak{} 14.05 & higher error \\
B4 & Lena|\allowbreak{}r; S6-S1|\allowbreak{}R1(0.25)|\allowbreak{}n50 (descriptive) & 2.12 [1.07, 2.66] & 1.19 [0.38, 1.93] & 23.77 /\allowbreak{} 21.65 & higher error \\
B5 & Lena|\allowbreak{}r; S6-S1|\allowbreak{}R1(0.25)|\allowbreak{}n100 (descriptive) & 2.56 [1.46, 3.18] & 1.81 [0.86, 2.62] & 23.73 /\allowbreak{} 21.17 & higher error \\
B6 & Lena|\allowbreak{}r; S6-S1|\allowbreak{}R1(0.25)|\allowbreak{}n200 (descriptive) & 2.79 [1.48, 3.45] & 1.29 [0.14, 2.37] & 23.52 /\allowbreak{} 20.73 & higher error \\
B7 & AL-1|\allowbreak{}r; S6-S1|\allowbreak{}R1(0.25)|\allowbreak{}n50 (descriptive) & 0.44 [0.39, 1.86] & 1.34 [0.67, 2.01] & 17.57 /\allowbreak{} 17.13 & higher error \\
B8 & AL-1|\allowbreak{}r; S6-S1|\allowbreak{}R1(0.25)|\allowbreak{}n100 (descriptive) & 0.89 [0.85, 1.88] & 1.61 [1.15, 2.07] & 17.52 /\allowbreak{} 16.64 & higher error \\
B9 & AL-1|\allowbreak{}r; S6-S1|\allowbreak{}R1(0.25)|\allowbreak{}n200 (descriptive) & 0.12 [0.07, 0.82] & 0.52 [0.19, 0.86] & 16.63 /\allowbreak{} 16.50 & higher error \\
B10 & AL-2|\allowbreak{}r; S6-S1|\allowbreak{}R1(0.25)|\allowbreak{}n50 (descriptive) & 0.21 [-0.65, 0.38] & -0.46 [-1.13, 0.20] & 13.25 /\allowbreak{} 13.04 & undetermined \\
B11 & AL-2|\allowbreak{}r; S6-S1|\allowbreak{}R1(0.25)|\allowbreak{}n100 (descriptive) & -0.06 [-0.62, 0.04] & -0.42 [-0.89, 0.05] & 12.93 /\allowbreak{} 12.98 & undetermined \\
B12 & AL-2|\allowbreak{}r; S6-S1|\allowbreak{}R1(0.25)|\allowbreak{}n200 (descriptive) & 0.13 [-0.27, 0.19] & -0.15 [-0.44, 0.13] & 12.80 /\allowbreak{} 12.67 & equivalent \\
B13 & Canada|\allowbreak{}r; S6-S1|\allowbreak{}R1(0.25)|\allowbreak{}n50 (descriptive) & -0.98 [-1.43, -0.10] & -0.14 [-0.60, 0.32] & 30.03 /\allowbreak{} 31.01 & undetermined \\
B14 & Canada|\allowbreak{}r; S6-S1|\allowbreak{}R1(0.25)|\allowbreak{}n100 (descriptive) & -1.50 [-2.00, -0.54] & -0.95 [-1.38, -0.52] & 29.79 /\allowbreak{} 31.29 & lower error \\
B15 & Canada|\allowbreak{}r; S6-S1|\allowbreak{}R1(0.25)|\allowbreak{}n200 (descriptive) & -0.68 [-1.03, -0.09] & -0.71 [-1.06, -0.36] & 29.91 /\allowbreak{} 30.59 & lower error \\
\end{longtable}}
\par\medskip\noindent{\small\textit{c. Source-extrapolation first (S5) and $\sqrt{\mathrm{TDD}}$ first (S7), WF2-a (C6 README 2.3)}}\par\nopagebreak
{\fontsize{7.5}{9}\selectfont\setlength{\tabcolsep}{4pt}\setlength{\LTpre}{2pt}\setlength{\LTpost}{4pt}
\begin{longtable}{@{}>{\raggedright\arraybackslash}p{\dimexpr 0.049\linewidth-8pt\relax}>{\raggedright\arraybackslash}p{\dimexpr 0.347\linewidth-8pt\relax}>{\raggedright\arraybackslash}p{\dimexpr 0.147\linewidth-8pt\relax}>{\raggedright\arraybackslash}p{\dimexpr 0.152\linewidth-8pt\relax}>{\raggedright\arraybackslash}p{\dimexpr 0.186\linewidth-8pt\relax}>{\raggedright\arraybackslash}p{\dimexpr 0.119\linewidth-8pt\relax}@{}}
\toprule
\textbf{ID} & \textbf{Target; contrast (role)} & \textbf{$\Delta$ cell [CI]} & \textbf{$\Delta$ block-equal [CI]} & \textbf{RMSE A / B} & \textbf{Verdict} \\
\midrule\endfirsthead
\toprule
\textbf{ID} & \textbf{Target; contrast (role)} & \textbf{$\Delta$ cell [CI]} & \textbf{$\Delta$ block-equal [CI]} & \textbf{RMSE A / B} & \textbf{Verdict} \\
\midrule\endhead
\bottomrule\endlastfoot
F1 & Alaska|\allowbreak{}r; S5-S1|\allowbreak{}R1(0.25)|\allowbreak{}n100 (main) & 0.23 [0.08, 0.61] & 0.61 [0.31, 0.92] & 14.58 /\allowbreak{} 14.35 & higher error \\
F2 & Alaska|\allowbreak{}r; S5-S1|\allowbreak{}R1(0.25)|\allowbreak{}n200 (main) & 1.17 [0.70, 1.81] & 1.57 [0.89, 2.25] & 15.19 /\allowbreak{} 14.02 & higher error \\
F3 & AL-2|\allowbreak{}r; S5-S1|\allowbreak{}R1(0.25)|\allowbreak{}n50 (descriptive) & 1.55 [1.16, 1.73] & 1.05 [0.53, 1.52] & 14.59 /\allowbreak{} 13.04 & higher error \\
F4 & AL-2|\allowbreak{}r; S5-S1|\allowbreak{}R1(0.25)|\allowbreak{}n200 (descriptive) & 1.04 [0.31, 1.21] & 0.50 [0.01, 1.00] & 13.69 /\allowbreak{} 12.65 & higher error \\
F5 & Lena|\allowbreak{}r; S5-S1|\allowbreak{}R1(0.25)|\allowbreak{}n100 (descriptive) & -0.20 [-0.55, 0.45] & -0.09 [-0.63, 0.48] & 20.97 /\allowbreak{} 21.17 & undetermined \\
F6 & Canada|\allowbreak{}r; S5-S1|\allowbreak{}R1(0.25)|\allowbreak{}n100 (descriptive) & -2.63 [-3.70, -0.37] & -0.89 [-1.79, 0.02] & 28.67 /\allowbreak{} 31.30 & undetermined \\
F7 & AL-1|\allowbreak{}r; S5-S1|\allowbreak{}R1(0.25)|\allowbreak{}n100 (descriptive) & -0.29 [-4.10, 0.43] & -3.72 [-5.67, -1.77] & 16.37 /\allowbreak{} 16.66 & undetermined \\
F8 & Lena|\allowbreak{}r; S7-S1|\allowbreak{}R1(0.25)|\allowbreak{}n50 (descriptive) & 3.61 [2.63, 5.20] & 4.54 [3.13, 5.76] & 25.26 /\allowbreak{} 21.65 & higher error \\
F9 & Lena|\allowbreak{}r; S7-S1|\allowbreak{}R1(0.25)|\allowbreak{}n100 (descriptive) & 5.06 [3.92, 7.16] & 6.72 [5.11, 8.20] & 26.24 /\allowbreak{} 21.17 & higher error \\
F10 & Lena|\allowbreak{}r; S7-S1|\allowbreak{}R1(0.25)|\allowbreak{}n200 (descriptive) & 2.91 [1.85, 4.51] & 3.89 [2.75, 4.90] & 23.64 /\allowbreak{} 20.73 & higher error \\
F11 & Alaska|\allowbreak{}r; S7-S1|\allowbreak{}R1(0.25)|\allowbreak{}n100 (main) & 0.47 [-0.03, 0.90] & -0.69 [-1.34, -0.01] & 14.80 /\allowbreak{} 14.33 & undetermined \\
F12 & Canada|\allowbreak{}r; S7-S1|\allowbreak{}R1(0.25)|\allowbreak{}n100 (descriptive) & -2.04 [-3.22, -0.00] & -0.93 [-1.64, -0.22] & 29.25 /\allowbreak{} 31.29 & lower error \\
F13 & AL-1|\allowbreak{}r; S7-S1|\allowbreak{}R1(0.25)|\allowbreak{}n100 (descriptive) & -0.66 [-1.87, 1.23] & 0.75 [-0.23, 1.70] & 15.98 /\allowbreak{} 16.64 & undetermined \\
F14 & AL-2|\allowbreak{}r; S7-S1|\allowbreak{}R1(0.25)|\allowbreak{}n100 (descriptive) & -0.47 [-1.02, 0.00] & -0.44 [-0.99, 0.11] & 12.52 /\allowbreak{} 12.98 & undetermined \\
\end{longtable}}
\par\medskip\noindent{\small\textit{d. Transfer of the best Alaskan strategy, WF2-b (C6 README 2.4)}}\par\nopagebreak
{\fontsize{7.5}{9}\selectfont\setlength{\tabcolsep}{4pt}\setlength{\LTpre}{2pt}\setlength{\LTpost}{4pt}
\begin{longtable}{@{}>{\raggedright\arraybackslash}p{\dimexpr 0.050\linewidth-8pt\relax}>{\raggedright\arraybackslash}p{\dimexpr 0.357\linewidth-8pt\relax}>{\raggedright\arraybackslash}p{\dimexpr 0.145\linewidth-8pt\relax}>{\raggedright\arraybackslash}p{\dimexpr 0.157\linewidth-8pt\relax}>{\raggedright\arraybackslash}p{\dimexpr 0.191\linewidth-8pt\relax}>{\raggedright\arraybackslash}p{\dimexpr 0.100\linewidth-8pt\relax}@{}}
\toprule
\textbf{ID} & \textbf{Target; contrast (role)} & \textbf{$\Delta$ cell [CI]} & \textbf{$\Delta$ block-equal [CI]} & \textbf{RMSE A / B} & \textbf{Verdict} \\
\midrule\endfirsthead
\toprule
\textbf{ID} & \textbf{Target; contrast (role)} & \textbf{$\Delta$ cell [CI]} & \textbf{$\Delta$ block-equal [CI]} & \textbf{RMSE A / B} & \textbf{Verdict} \\
\midrule\endhead
\bottomrule\endlastfoot
C1 & Lena|\allowbreak{}r; S3-S1|\allowbreak{}R1(0.25)|\allowbreak{}n100 (main) & 0.33 [0.04, 0.77] & 0.68 [0.36, 1.00] & 21.50 /\allowbreak{} 21.17 & higher error \\
C2 & Canada|\allowbreak{}r; S3-S1|\allowbreak{}R1(0.25)|\allowbreak{}n100 (main) & -0.48 [-0.65, -0.25] & -0.24 [-0.44, -0.02] & 30.81 /\allowbreak{} 31.29 & lower error \\
\end{longtable}}
\par\medskip\noindent{\small\textit{e. Transfer conditions, WF8 (C6 README 2.5)}}\par\nopagebreak
{\fontsize{7.5}{9}\selectfont\setlength{\tabcolsep}{4pt}\setlength{\LTpre}{2pt}\setlength{\LTpost}{4pt}
\begin{longtable}{@{}>{\raggedright\arraybackslash}p{\dimexpr 0.049\linewidth-8pt\relax}>{\raggedright\arraybackslash}p{\dimexpr 0.345\linewidth-8pt\relax}>{\raggedright\arraybackslash}p{\dimexpr 0.151\linewidth-8pt\relax}>{\raggedright\arraybackslash}p{\dimexpr 0.151\linewidth-8pt\relax}>{\raggedright\arraybackslash}p{\dimexpr 0.184\linewidth-8pt\relax}>{\raggedright\arraybackslash}p{\dimexpr 0.119\linewidth-8pt\relax}@{}}
\toprule
\textbf{ID} & \textbf{Target; contrast (role)} & \textbf{$\Delta$ cell [CI]} & \textbf{$\Delta$ block-equal [CI]} & \textbf{RMSE A / B} & \textbf{Verdict} \\
\midrule\endfirsthead
\toprule
\textbf{ID} & \textbf{Target; contrast (role)} & \textbf{$\Delta$ cell [CI]} & \textbf{$\Delta$ block-equal [CI]} & \textbf{RMSE A / B} & \textbf{Verdict} \\
\midrule\endhead
\bottomrule\endlastfoot
D1 & Canada|\allowbreak{}x; S4-S1|\allowbreak{}R1(0.25)|\allowbreak{}n40 (main) & -2.39 [-3.40, -0.49] & -1.64 [-2.49, -0.76] & 28.77 /\allowbreak{} 31.16 & lower error \\
D2 & Lena|\allowbreak{}x; S4-S1|\allowbreak{}R1(0.25)|\allowbreak{}n40 (main) & 0.55 [-0.44, 1.27] & -0.41 [-1.09, 0.24] & 22.10 /\allowbreak{} 21.55 & undetermined \\
D3 & Lena|\allowbreak{}x; S2-S1|\allowbreak{}R1(0.25)|\allowbreak{}n10 (main) & 1.89 [0.30, 3.08] & 0.86 [-0.27, 1.93] & 23.78 /\allowbreak{} 21.89 & undetermined \\
D4 & Lena|\allowbreak{}x; S4-S1|\allowbreak{}R1(0.25)|\allowbreak{}n10 (main) & 1.90 [-0.02, 3.37] & 1.09 [-0.15, 2.28] & 23.79 /\allowbreak{} 21.89 & undetermined \\
D5 & Canada|\allowbreak{}x; S2-S1|\allowbreak{}R1(0.25)|\allowbreak{}n10 (main) & -1.35 [-3.08, 2.20] & 0.18 [-1.32, 1.67] & 28.18 /\allowbreak{} 29.53 & undetermined \\
D6 & Canada|\allowbreak{}x; S4-S1|\allowbreak{}R1(0.25)|\allowbreak{}n10 (main) & -1.08 [-2.70, 2.02] & 0.06 [-1.39, 1.48] & 28.46 /\allowbreak{} 29.53 & undetermined \\
D7 & Russia\_\allowbreak{}W|\allowbreak{}x; S2-S1|\allowbreak{}R1(0.25)|\allowbreak{}n10 (main) & -0.72 [-1.13, -0.32] & -0.80 [-1.41, -0.18] & 29.73 /\allowbreak{} 30.46 & lower error \\
D8 & Russia\_\allowbreak{}W|\allowbreak{}x; S4-S1|\allowbreak{}R1(0.25)|\allowbreak{}n10 (main) & -0.38 [-0.89, 0.12] & -0.42 [-1.11, 0.29] & 30.08 /\allowbreak{} 30.46 & undetermined \\
D9 & Russia\_\allowbreak{}E|\allowbreak{}x; S2-S1|\allowbreak{}R1(0.25)|\allowbreak{}n10 (main) & 0.36 [-0.94, 0.82] & -0.63 [-1.29, 0.07] & 30.52 /\allowbreak{} 30.15 & undetermined \\
D10 & Russia\_\allowbreak{}E|\allowbreak{}x; S4-S1|\allowbreak{}R1(0.25)|\allowbreak{}n10 (main) & 0.06 [-0.33, 0.54] & 0.36 [0.08, 0.62] & 30.21 /\allowbreak{} 30.15 & undetermined \\
D11 & Alaska|\allowbreak{}x; S2-S1|\allowbreak{}R1(0.25)|\allowbreak{}n10 (main) & -0.79 [-1.33, -0.26] & -1.37 [-1.82, -0.91] & 14.30 /\allowbreak{} 15.09 & lower error \\
D12 & Alaska|\allowbreak{}x; S4-S1|\allowbreak{}R1(0.25)|\allowbreak{}n10 (main) & -0.74 [-1.34, -0.12] & -1.02 [-1.50, -0.54] & 14.35 /\allowbreak{} 15.09 & lower error \\
D13 & Canada|\allowbreak{}x; S2-S1|\allowbreak{}R1(0.25)|\allowbreak{}n40 (descriptive) & -2.32 [-3.43, -0.18] & -1.43 [-2.39, -0.45] & 28.85 /\allowbreak{} 31.16 & lower error \\
D14 & Lena|\allowbreak{}x; S2-S1|\allowbreak{}R1(0.25)|\allowbreak{}n40 (descriptive) & 2.18 [0.84, 2.88] & 0.74 [-0.38, 1.75] & 23.72 /\allowbreak{} 21.55 & undetermined \\
D15 & Alaska|\allowbreak{}x; S2-S1|\allowbreak{}R1(0.25)|\allowbreak{}n40 (descriptive) & -1.13 [-2.06, -0.48] & -2.17 [-2.94, -1.41] & 14.09 /\allowbreak{} 15.22 & lower error \\
D16 & Alaska|\allowbreak{}x; S4-S1|\allowbreak{}R1(0.25)|\allowbreak{}n40 (descriptive) & -1.19 [-2.08, -0.52] & -2.11 [-2.87, -1.36] & 14.04 /\allowbreak{} 15.22 & lower error \\
\end{longtable}}
\par\medskip\noindent{\small\textit{f. Block-dispersed versus cell-random placement, L43 (C6 README 2.6)}}\par\nopagebreak
{\fontsize{7.5}{9}\selectfont\setlength{\tabcolsep}{4pt}\setlength{\LTpre}{2pt}\setlength{\LTpost}{4pt}
\begin{longtable}{@{}>{\raggedright\arraybackslash}p{\dimexpr 0.054\linewidth-8pt\relax}>{\raggedright\arraybackslash}p{\dimexpr 0.410\linewidth-8pt\relax}>{\raggedright\arraybackslash}p{\dimexpr 0.180\linewidth-8pt\relax}>{\raggedright\arraybackslash}p{\dimexpr 0.180\linewidth-8pt\relax}>{\raggedright\arraybackslash}p{\dimexpr 0.054\linewidth-8pt\relax}>{\raggedright\arraybackslash}p{\dimexpr 0.123\linewidth-8pt\relax}@{}}
\toprule
\textbf{ID} & \textbf{Target; contrast (role)} & \textbf{$\Delta$ cell [CI]} & \textbf{$\Delta$ block-equal [CI]} & \textbf{RMSE A / B} & \textbf{Verdict} \\
\midrule\endfirsthead
\toprule
\textbf{ID} & \textbf{Target; contrast (role)} & \textbf{$\Delta$ cell [CI]} & \textbf{$\Delta$ block-equal [CI]} & \textbf{RMSE A / B} & \textbf{Verdict} \\
\midrule\endhead
\bottomrule\endlastfoot
E1 & MEAN[Lena|\allowbreak{}x,Canada|\allowbreak{}x,Russia\_\allowbreak{}W|\allowbreak{}x,Russia\_\allowbreak{}E|\allowbreak{}x]; P1, n 10 (main) & -0.07 [-0.86, 0.89] & -0.21 [-0.81, 0.41] &  & undetermined \\
E2 & MEAN[Lena|\allowbreak{}x,Canada|\allowbreak{}x]; P1, n 40 (main) & -0.24 [-1.19, 0.99] & -0.52 [-1.36, 0.30] &  & undetermined \\
E3 & MEAN[Lena|\allowbreak{}x,Canada|\allowbreak{}x]; R1, n 40 (auxiliary) & -0.07 [-1.02, 1.15] & -0.34 [-1.19, 0.48] &  & undetermined \\
E4 & MEAN[Lena|\allowbreak{}x,Canada|\allowbreak{}x]; R1, n 160 (auxiliary) & 0.26 [-0.32, 0.80] & -0.20 [-0.78, 0.35] &  & undetermined \\
E5 & MEAN[Lena|\allowbreak{}x,Canada|\allowbreak{}x]; P1, n 160 (auxiliary) & 0.32 [-0.33, 0.90] & -0.18 [-0.81, 0.40] &  & undetermined \\
E6 & Lena|\allowbreak{}x; P1, n 10 (main) & 1.55 [-0.18, 2.97] & 0.58 [-0.60, 1.85] &  & undetermined \\
E7 & Canada|\allowbreak{}x; P1, n 10 (main) & -1.62 [-3.57, 2.06] & -0.28 [-1.98, 1.43] &  & undetermined \\
E8 & Russia\_\allowbreak{}W|\allowbreak{}x; P1, n 10 (main) & -0.77 [-1.58, -0.05] & -0.85 [-1.82, 0.00] &  & undetermined \\
E9 & Russia\_\allowbreak{}E|\allowbreak{}x; P1, n 10 (main) & 0.56 [-0.82, 1.13] & -0.27 [-1.10, 0.56] &  & undetermined \\
E10 & Lena|\allowbreak{}x; P1, n 40 (main) & 2.10 [0.73, 2.97] & 0.69 [-0.41, 1.78] &  & undetermined \\
E11 & Canada|\allowbreak{}x; P1, n 40 (main) & -2.59 [-4.01, -0.24] & -1.72 [-2.96, -0.55] &  & lower error \\
E12 & Canada|\allowbreak{}x; R1, n 40 (auxiliary) & -2.32 [-3.74, -0.04] & -1.43 [-2.64, -0.27] &  & lower error \\
E13 & Alaska|\allowbreak{}x; R1, n 40 (auxiliary) & -1.13 [-2.36, -0.36] & -2.17 [-3.15, -1.27] &  & lower error \\
\end{longtable}}
\par\medskip\noindent{\small\textit{g. Verdict counts by strategy (post hoc count) (C6 README 2.9)}}\par\nopagebreak
{\fontsize{7.5}{9}\selectfont\setlength{\tabcolsep}{4pt}\setlength{\LTpre}{2pt}\setlength{\LTpost}{4pt}
\begin{longtable}{@{}>{\raggedright\arraybackslash}p{\dimexpr 0.362\linewidth-8pt\relax}>{\raggedright\arraybackslash}p{\dimexpr 0.080\linewidth-8pt\relax}>{\raggedright\arraybackslash}p{\dimexpr 0.096\linewidth-8pt\relax}>{\raggedright\arraybackslash}p{\dimexpr 0.111\linewidth-8pt\relax}>{\raggedright\arraybackslash}p{\dimexpr 0.080\linewidth-8pt\relax}>{\raggedright\arraybackslash}p{\dimexpr 0.080\linewidth-8pt\relax}>{\raggedright\arraybackslash}p{\dimexpr 0.111\linewidth-8pt\relax}>{\raggedright\arraybackslash}p{\dimexpr 0.080\linewidth-8pt\relax}@{}}
\toprule
\textbf{Strategy} & \textbf{WF2 lower} & \textbf{WF2 equivalent} & \textbf{WF2 undetermined} & \textbf{WF2 higher} & \textbf{WF8 lower} & \textbf{WF8 undetermined} & \textbf{WF8 higher} \\
\midrule\endfirsthead
\toprule
\textbf{Strategy} & \textbf{WF2 lower} & \textbf{WF2 equivalent} & \textbf{WF2 undetermined} & \textbf{WF2 higher} & \textbf{WF8 lower} & \textbf{WF8 undetermined} & \textbf{WF8 higher} \\
\midrule\endhead
\bottomrule\endlastfoot
S2 block stratification & 8 & 1 & 14 & 1 & 4 & 4 & 0 \\
S3 covariate-cluster stratification & 3 & 12 & 8 & 1 &  &  &  \\
S4 covariate max-min distance & 8 & 2 & 14 & 0 & 3 & 5 & 0 \\
S5 source extrapolation first & 2 & 0 & 15 & 7 &  &  &  \\
S6 sequential active selection & 3 & 8 & 4 & 9 &  &  &  \\
S7 √TDD first & 2 & 2 & 15 & 5 &  &  &  \\
\end{longtable}}

\clearpage
\phantomsection\label{tab:S10}
\noindent\textbf{Supplementary Table S10.} Gain decomposition and climate extrapolation (WF9, WF10).\par\smallskip
\noindent{\footnotesize Scoring cells grouped by 0.5\textdegree{} block and air $\sqrt{\mathrm{TDD}}$ (registered: soil $\sqrt{\mathrm{TDD}}$, sensitivity rows b and e). Shares are computed from the squared-error sums (Methods, equation for SSE). Explained within-grid share = 1 $-$ (within-grid RMSE of the method / within-grid RMSE of P1)$^2$. Method codes (internal identifiers allowed in the Supplementary Information): P0 source-coefficient Stefan model; P* year-matched Stefan model; P1 recalibrated Stefan model ($\kappa = 10$); P2 local least-squares Stefan model; P1* or P1@ed soil-property recalibrated Stefan model; R0 source anchor plus residual ML; R1 recalibrated anchor plus residual ML; R2 augmentation plus anchor plus residual; D0 direct ML; D1 physics pseudo-label augmentation; F1a, F1k, F1n physics-input ML; RM multiplicative residual; Re Stefan-CCI mean anchor plus residual; W selection rule (cross-validation within the target labels); S1 random cells, S2 block stratification, S3 covariate-cluster stratification, S4 covariate max-min distance, S5 source-extrapolation first, S6 sequential selection by prediction variance, S7 $\sqrt{\mathrm{TDD}}$ first. Target suffixes: $|x$ parent region excluded from the source, $|i$ rest of the parent kept, $|r$ within-region design. $n$ = number of target labels; all = all labels of half A. Verdicts: lower error, higher error, equivalent ($\pm$0.5~cm), undetermined; both 95\% block-bootstrap CIs (cell-weighted and block-equal) must agree. \dag~a note of the source row in Korean was omitted; see the cited README section.}\par

\par\medskip\noindent{\small\textit{a. Decomposition of the within-region gain in Alaska (all labels, $\lambda$ by CV) (C8 README 2.1)}}\par\nopagebreak
{\fontsize{7.5}{9}\selectfont\setlength{\tabcolsep}{4pt}\setlength{\LTpre}{2pt}\setlength{\LTpost}{4pt}
\begin{longtable}{@{}>{\raggedright\arraybackslash}p{\dimexpr 0.058\linewidth-8pt\relax}>{\raggedright\arraybackslash}p{\dimexpr 0.398\linewidth-8pt\relax}>{\raggedright\arraybackslash}p{\dimexpr 0.420\linewidth-8pt\relax}>{\raggedright\arraybackslash}p{\dimexpr 0.124\linewidth-8pt\relax}@{}}
\toprule
\textbf{ID} & \textbf{Contrast or quantity} & \textbf{Value (cm unless stated): cell-weighted [CI] / block-equal [CI]} & \textbf{Verdict} \\
\midrule\endfirsthead
\toprule
\textbf{ID} & \textbf{Contrast or quantity} & \textbf{Value (cm unless stated): cell-weighted [CI] / block-equal [CI]} & \textbf{Verdict} \\
\midrule\endhead
\bottomrule\endlastfoot
A1 & Alaska within region all, R1(cross-validation λ) − P1, total & -0.52 [-0.70, -0.31] /\allowbreak{} -0.68 [-0.91, -0.44] & lower error \\
A2 & same contrast, between grids store & -1.03 [-1.30, -0.59] /\allowbreak{} -0.89 [-1.16, -0.61] & lower error \\
A3 & same contrast, within grid store(WF9-a targets rows) & -0.01 [-0.03, 0.04] /\allowbreak{} 0.04 [0.01, 0.06] & equivalent \\
A4 & soil definition: total & -0.52 [-0.70, -0.31] /\allowbreak{} -0.68 [-0.91, -0.44] & lower error \\
A5 & soil definition: between grids & -1.03 [-1.30, -0.59] /\allowbreak{} -0.89 [-1.16, -0.61] & lower error \\
A6 & soil definition: within grid & -0.01 [-0.03, 0.04] /\allowbreak{} 0.04 [0.01, 0.06] & equivalent \\
\end{longtable}}
\par\medskip\noindent{\small\textit{b. Shares of the squared error (descriptive, no CI) (C8 README 2.2)}}\par\nopagebreak
{\fontsize{7.5}{9}\selectfont\setlength{\tabcolsep}{4pt}\setlength{\LTpre}{2pt}\setlength{\LTpost}{4pt}
\begin{longtable}{@{}>{\raggedright\arraybackslash}p{\dimexpr 0.072\linewidth-8pt\relax}>{\raggedright\arraybackslash}p{\dimexpr 0.550\linewidth-8pt\relax}>{\raggedright\arraybackslash}p{\dimexpr 0.223\linewidth-8pt\relax}>{\raggedright\arraybackslash}p{\dimexpr 0.154\linewidth-8pt\relax}@{}}
\toprule
\textbf{ID} & \textbf{Contrast or quantity} & \textbf{Value (cm unless stated): cell-weighted [CI] / block-equal [CI]} & \textbf{Verdict} \\
\midrule\endfirsthead
\toprule
\textbf{ID} & \textbf{Contrast or quantity} & \textbf{Value (cm unless stated): cell-weighted [CI] / block-equal [CI]} & \textbf{Verdict} \\
\midrule\endhead
\bottomrule\endlastfoot
B1 & P1 total SSE of within grid share, Alaska within region all & 0.72 (71.9 \%) & descriptive \\
B2 & same, Lena Delta & 0.62 (61.7 \%) & descriptive \\
B3 & same, Canada & 0.37 (37.4 \%) & descriptive \\
B4 & R1(cross-validation λ) total SSE gain of within grid component share, Alaska within region & 0.01 (0.9 \%) & descriptive \\
B5 & R1(λ 0.25), Alaska within region & 0.07 (6.7 \%) & descriptive \\
B6 & R1(λ 0.25), Lena Delta within region & 0.53 (53.5 \%) & descriptive \\
B7 & R1(λ 0.25), Canada within region & -0.00 (-0.1 \%) & descriptive \\
B8 & R1(λ 0.25), Lena Delta transfer all & 0.25 (25.2 \%) & descriptive \\
B9 & R1(λ 0.25), Alaska transfer all & 0.36 (35.6 \%) & descriptive \\
B10 & soil definition: R1(cross-validation λ), Alaska within region & 0.01 (1.0 \%) & descriptive \\
B11 & soil definition: R1(λ 0.25), Lena Delta within region & 0.24 (23.7 \%) & descriptive \\
B12 & P1 decomp RMSE column(total; between grids; within grid), Alaska within region & 14.53; 7.68; 12.37 & descriptive \\
B13 & R1(cross-validation λ) decomp RMSE column, Alaska within region & 14.06; 6.76; 12.36 & descriptive \\
B14 & R1(cross-validation λ) total SSE gain within grid share, Canada within region & -12,133.57; 0.21 & descriptive \\
\end{longtable}}
\par\medskip\noindent{\small\textit{c. Explained within-grid share (descriptive, no CI) (C8 README 2.3)}}\par\nopagebreak
{\fontsize{7.5}{9}\selectfont\setlength{\tabcolsep}{4pt}\setlength{\LTpre}{2pt}\setlength{\LTpost}{4pt}
\begin{longtable}{@{}>{\raggedright\arraybackslash}p{\dimexpr 0.076\linewidth-8pt\relax}>{\raggedright\arraybackslash}p{\dimexpr 0.557\linewidth-8pt\relax}>{\raggedright\arraybackslash}p{\dimexpr 0.205\linewidth-8pt\relax}>{\raggedright\arraybackslash}p{\dimexpr 0.162\linewidth-8pt\relax}@{}}
\toprule
\textbf{ID} & \textbf{Contrast or quantity} & \textbf{Value (cm unless stated): cell-weighted [CI] / block-equal [CI]} & \textbf{Verdict} \\
\midrule\endfirsthead
\toprule
\textbf{ID} & \textbf{Contrast or quantity} & \textbf{Value (cm unless stated): cell-weighted [CI] / block-equal [CI]} & \textbf{Verdict} \\
\midrule\endhead
\bottomrule\endlastfoot
C1 & R1(cross-validation λ), Alaska within region & 0.00 (0.1 \%) & descriptive \\
C2 & R1(λ 0.25), Alaska within region & 0.01 (0.7 \%) & descriptive \\
C3 & D0(catboost), Alaska within region & -0.02 (-2.0 \%) & descriptive \\
C4 & R1(λ 0.25), Lena Delta within region & 0.06 (5.6 \%) & descriptive \\
C5 & D0(catboost), Lena Delta within region & 0.06 (6.2 \%) & descriptive \\
C6 & R1(cross-validation λ), Canada within region & -0.02 (-2.2 \%) & descriptive \\
C7 & D0(catboost), Canada within region & -0.08 (-8.3 \%) & descriptive \\
C8 & R1(λ 0.25), Lena Delta transfer & 0.03 (3.2 \%) & descriptive \\
C9 & R1(λ 0.25), Alaska transfer & 0.02 (2.0 \%) & descriptive \\
C10 & R1(λ 0.25), Canada transfer & -0.00 (-0.4 \%) & descriptive \\
C11 & R1(λ 0.25), E Russia transfer & -0.08 (-8.5 \%) & descriptive \\
C12 & within region ML over all keys of maximum(air definition) & 0.08 (7.8 \%) & descriptive \\
C13 & within region ML over all keys of maximum(soil definition) & 0.04 (4.0 \%) & descriptive \\
C14 & transfer ML over all keys of maximum(air definition, soil definition same value) & 0.07 (7.2 \%) & descriptive \\
\end{longtable}}
\par\medskip\noindent{\small\textit{d. WF9-a within regions, air $\sqrt{\mathrm{TDD}}$ grouping (C8 README 2.4)}}\par\nopagebreak
{\fontsize{7.5}{9}\selectfont\setlength{\tabcolsep}{4pt}\setlength{\LTpre}{2pt}\setlength{\LTpost}{4pt}
\begin{longtable}{@{}>{\raggedright\arraybackslash}p{\dimexpr 0.061\linewidth-8pt\relax}>{\raggedright\arraybackslash}p{\dimexpr 0.340\linewidth-8pt\relax}>{\raggedright\arraybackslash}p{\dimexpr 0.461\linewidth-8pt\relax}>{\raggedright\arraybackslash}p{\dimexpr 0.138\linewidth-8pt\relax}@{}}
\toprule
\textbf{ID} & \textbf{Contrast or quantity} & \textbf{Value (cm unless stated): cell-weighted [CI] / block-equal [CI]} & \textbf{Verdict} \\
\midrule\endfirsthead
\toprule
\textbf{ID} & \textbf{Contrast or quantity} & \textbf{Value (cm unless stated): cell-weighted [CI] / block-equal [CI]} & \textbf{Verdict} \\
\midrule\endhead
\bottomrule\endlastfoot
D1 & stratified mean n 500(Alaska·Lena Delta) & -0.08 [-0.16, 0.15] /\allowbreak{} 0.22 [0.09, 0.35] & equivalent \\
D2 & stratified mean n 1,000(Alaska·Lena Delta) & -0.02 [-0.12, 0.35] /\allowbreak{} 0.36 [0.19, 0.54] & undetermined \\
D3 & stratified mean all(3 regions) & 0.09 [0.02, 0.40] /\allowbreak{} 0.37 [0.25, 0.50] & higher error \\
D4 & Lena Delta all & 0.05 [-0.18, 0.92] /\allowbreak{} 0.87 [0.51, 1.24] & undetermined \\
D5 & Canada all & 0.23 [0.16, 0.34] /\allowbreak{} 0.21 [0.12, 0.31] & higher error \\
D6 & verdict wording & partial(regions 2/\allowbreak{}3, splits 23/\allowbreak{}24): rejected: higher error contrast exists(n=all) [distinguishable, size below 0.5 cm: n=all] &  \\
\end{longtable}}
\par\medskip\noindent{\small\textit{e. WF9-a, soil $\sqrt{\mathrm{TDD}}$ grouping (sensitivity) (C8 README 2.5)}}\par\nopagebreak
{\fontsize{7.5}{9}\selectfont\setlength{\tabcolsep}{4pt}\setlength{\LTpre}{2pt}\setlength{\LTpost}{4pt}
\begin{longtable}{@{}>{\raggedright\arraybackslash}p{\dimexpr 0.061\linewidth-8pt\relax}>{\raggedright\arraybackslash}p{\dimexpr 0.340\linewidth-8pt\relax}>{\raggedright\arraybackslash}p{\dimexpr 0.461\linewidth-8pt\relax}>{\raggedright\arraybackslash}p{\dimexpr 0.138\linewidth-8pt\relax}@{}}
\toprule
\textbf{ID} & \textbf{Contrast or quantity} & \textbf{Value (cm unless stated): cell-weighted [CI] / block-equal [CI]} & \textbf{Verdict} \\
\midrule\endfirsthead
\toprule
\textbf{ID} & \textbf{Contrast or quantity} & \textbf{Value (cm unless stated): cell-weighted [CI] / block-equal [CI]} & \textbf{Verdict} \\
\midrule\endhead
\bottomrule\endlastfoot
E1 & stratified mean n 500(Alaska·Lena Delta) & 0.27 [0.07, 0.41] /\allowbreak{} 0.28 [0.15, 0.41] & higher error \\
E2 & stratified mean n 1,000(Alaska·Lena Delta) & 0.42 [0.20, 0.64] /\allowbreak{} 0.44 [0.26, 0.61] & higher error \\
E3 & stratified mean all(3 regions) & 0.40 [0.24, 0.59] /\allowbreak{} 0.42 [0.30, 0.55] & higher error \\
E4 & Lena Delta all & 0.99 [0.49, 1.52] /\allowbreak{} 1.02 [0.67, 1.39] & higher error \\
E5 & Canada all & 0.23 [0.16, 0.33] /\allowbreak{} 0.21 [0.11, 0.31] & higher error \\
E6 & verdict wording & partial(regions 2/\allowbreak{}3, splits 23/\allowbreak{}24): rejected: higher error contrast exists(n=500, n=1000, n=all) [distinguishable, size below 0.5 cm: n=500, n=1000, n=all] &  \\
\end{longtable}}
\par\medskip\noindent{\small\textit{f. WF9-c in transfer, air grouping (C8 README 2.6)}}\par\nopagebreak
{\fontsize{7.5}{9}\selectfont\setlength{\tabcolsep}{4pt}\setlength{\LTpre}{2pt}\setlength{\LTpost}{4pt}
\begin{longtable}{@{}>{\raggedright\arraybackslash}p{\dimexpr 0.061\linewidth-8pt\relax}>{\raggedright\arraybackslash}p{\dimexpr 0.343\linewidth-8pt\relax}>{\raggedright\arraybackslash}p{\dimexpr 0.465\linewidth-8pt\relax}>{\raggedright\arraybackslash}p{\dimexpr 0.130\linewidth-8pt\relax}@{}}
\toprule
\textbf{ID} & \textbf{Contrast or quantity} & \textbf{Value (cm unless stated): cell-weighted [CI] / block-equal [CI]} & \textbf{Verdict} \\
\midrule\endfirsthead
\toprule
\textbf{ID} & \textbf{Contrast or quantity} & \textbf{Value (cm unless stated): cell-weighted [CI] / block-equal [CI]} & \textbf{Verdict} \\
\midrule\endhead
\bottomrule\endlastfoot
F1 & stratified mean n 0(Lena Delta and Canada) & -0.01 [-0.04, 0.01] /\allowbreak{} -0.02 [-0.06, 0.00] & equivalent \\
F2 & stratified mean n 10 & -0.01 [-0.04, 0.01] /\allowbreak{} -0.04 [-0.07, -0.00] & equivalent \\
F3 & stratified mean all & -0.11 [-0.25, -0.01] /\allowbreak{} -0.12 [-0.22, -0.03] & lower error \\
F4 & Lena Delta all & -0.27 [-0.52, -0.05] /\allowbreak{} -0.22 [-0.42, -0.04] & lower error \\
F5 & E Russia all(scoring blocks few CI none) & 1.08 [NA, NA] /\allowbreak{} 1.08 [NA, NA] & not decidable \\
F6 & verdict wording & partial(regions 2/\allowbreak{}4): partially supported: lower error contrast = n=all [distinguishable, size below 0.5 cm: n=all] &  \\
\end{longtable}}
\par\medskip\noindent{\small\textit{g. WF9-c, soil grouping (sensitivity) (C8 README 2.7)}}\par\nopagebreak
{\fontsize{7.5}{9}\selectfont\setlength{\tabcolsep}{4pt}\setlength{\LTpre}{2pt}\setlength{\LTpost}{4pt}
\begin{longtable}{@{}>{\raggedright\arraybackslash}p{\dimexpr 0.071\linewidth-8pt\relax}>{\raggedright\arraybackslash}p{\dimexpr 0.246\linewidth-8pt\relax}>{\raggedright\arraybackslash}p{\dimexpr 0.541\linewidth-8pt\relax}>{\raggedright\arraybackslash}p{\dimexpr 0.141\linewidth-8pt\relax}@{}}
\toprule
\textbf{ID} & \textbf{Contrast or quantity} & \textbf{Value (cm unless stated): cell-weighted [CI] / block-equal [CI]} & \textbf{Verdict} \\
\midrule\endfirsthead
\toprule
\textbf{ID} & \textbf{Contrast or quantity} & \textbf{Value (cm unless stated): cell-weighted [CI] / block-equal [CI]} & \textbf{Verdict} \\
\midrule\endhead
\bottomrule\endlastfoot
G1 & stratified mean n 0(Lena Delta and Canada) & 0.00 [-0.04, 0.02] /\allowbreak{} -0.02 [-0.05, 0.01] & equivalent \\
G2 & stratified mean n 10 & 0.01 [-0.03, 0.02] /\allowbreak{} -0.03 [-0.07, 0.00] & equivalent \\
G3 & stratified mean all & -0.04 [-0.20, 0.04] /\allowbreak{} -0.10 [-0.21, -0.01] & equivalent \\
G4 & Lena Delta all & -0.13 [-0.42, 0.05] /\allowbreak{} -0.19 [-0.39, -0.01] & equivalent \\
G5 & E Russia all(CI none) & 1.08 [NA, NA] /\allowbreak{} 1.08 [NA, NA] & not decidable \\
G6 & verdict wording & partial(regions 2/\allowbreak{}4): rejected: lower error contrast none &  \\
\end{longtable}}
\par\medskip\noindent{\small\textit{h. Descriptive contrasts outside the registered hypotheses (C8 README 2.8)}}\par\nopagebreak
{\fontsize{7.5}{9}\selectfont\setlength{\tabcolsep}{4pt}\setlength{\LTpre}{2pt}\setlength{\LTpost}{4pt}
\begin{longtable}{@{}>{\raggedright\arraybackslash}p{\dimexpr 0.055\linewidth-8pt\relax}>{\raggedright\arraybackslash}p{\dimexpr 0.417\linewidth-8pt\relax}>{\raggedright\arraybackslash}p{\dimexpr 0.401\linewidth-8pt\relax}>{\raggedright\arraybackslash}p{\dimexpr 0.126\linewidth-8pt\relax}@{}}
\toprule
\textbf{ID} & \textbf{Contrast or quantity} & \textbf{Value (cm unless stated): cell-weighted [CI] / block-equal [CI]} & \textbf{Verdict} \\
\midrule\endfirsthead
\toprule
\textbf{ID} & \textbf{Contrast or quantity} & \textbf{Value (cm unless stated): cell-weighted [CI] / block-equal [CI]} & \textbf{Verdict} \\
\midrule\endhead
\bottomrule\endlastfoot
H1 & R1(λ 0.25) − P1 within grid, 3 regions stratified mean(air) & -0.16 [-0.17, -0.07] /\allowbreak{} -0.09 [-0.13, -0.04] & lower error \\
H2 & same, Lena Delta(air) & -0.44 [-0.50, -0.19] /\allowbreak{} -0.26 [-0.40, -0.14] & lower error \\
H3 & same, 3 regions stratified mean(soil) & -0.08 [-0.11, -0.02] /\allowbreak{} -0.07 [-0.12, -0.03] & lower error \\
H4 & same, Lena Delta(soil) & -0.19 [-0.30, -0.04] /\allowbreak{} -0.22 [-0.35, -0.10] & lower error \\
H5 & R1(λ 0.25) − P1 between grids, Alaska & -0.99 [-1.11, -0.69] /\allowbreak{} -0.58 [-0.71, -0.45] & lower error \\
H6 & same, Lena Delta & -0.51 [-1.14, -0.35] /\allowbreak{} -1.27 [-1.65, -0.90] & lower error \\
H7 & same, Canada & -0.93 [-1.39, -0.48] /\allowbreak{} -1.39 [-1.69, -1.08] & lower error \\
H8 & D0(catboost) − P1 within grid, 3 regions stratified mean & 0.16 [0.10, 0.42] /\allowbreak{} 0.46 [0.33, 0.59] & higher error \\
H9 & Re(Stefan·CCI anchor, cross-validation λ) − P1 within grid, 3 regions stratified mean & 0.31 [0.21, 0.59] /\allowbreak{} 0.42 [0.28, 0.56] & higher error \\
\end{longtable}}
\par\medskip\noindent{\small\textit{i. WF10 climate extrapolation (warm variant) (C8 README 2.9)}}\par\nopagebreak
{\fontsize{7.5}{9}\selectfont\setlength{\tabcolsep}{4pt}\setlength{\LTpre}{2pt}\setlength{\LTpost}{4pt}
\begin{longtable}{@{}>{\raggedright\arraybackslash}p{\dimexpr 0.067\linewidth-8pt\relax}>{\raggedright\arraybackslash}p{\dimexpr 0.310\linewidth-8pt\relax}>{\raggedright\arraybackslash}p{\dimexpr 0.471\linewidth-8pt\relax}>{\raggedright\arraybackslash}p{\dimexpr 0.152\linewidth-8pt\relax}@{}}
\toprule
\textbf{ID} & \textbf{Contrast or quantity} & \textbf{Value (cm unless stated): cell-weighted [CI] / block-equal [CI]} & \textbf{Verdict} \\
\midrule\endfirsthead
\toprule
\textbf{ID} & \textbf{Contrast or quantity} & \textbf{Value (cm unless stated): cell-weighted [CI] / block-equal [CI]} & \textbf{Verdict} \\
\midrule\endhead
\bottomrule\endlastfoot
I1 & WF10-a verdict wording & not decidable\textsuperscript{\dag} &  \\
I2 & WF10-b verdict wording & not decidable\textsuperscript{\dag} &  \\
I3 & WF10-c verdict wording & not decidable\textsuperscript{\dag} &  \\
I4 & EP(D0 − P1), Alaska warm, n 100 & 7.02 [3.91, 9.51] /\allowbreak{} 0.75 [-1.55, 3.04] & undetermined \\
I5 & same, n 500(descriptive rows) & 4.55 [1.58, 7.09] /\allowbreak{} -1.31 [-3.54, 0.92] & undetermined \\
I6 & same, all & 3.39 [0.61, 5.95] /\allowbreak{} -1.33 [-3.41, 0.77] & undetermined \\
I7 & EP(R1 − D0), Alaska warm, n 100 & -5.56 [-7.87, -3.02] /\allowbreak{} -1.38 [-3.29, 0.52] & undetermined \\
I8 & same, all & -2.45 [-4.73, 0.09] /\allowbreak{} 1.59 [-0.25, 3.38] & undetermined \\
I9 & W R1 − P1, Alaska warm, n 100 & 0.87 [-0.12, 1.16] /\allowbreak{} -1.32 [-1.83, -0.78] & undetermined \\
I10 & same, all & 0.13 [-0.54, 0.41] /\allowbreak{} -0.51 [-0.93, -0.06] & undetermined \\
I11 & Canada warm extrapolation scoring(W) number of blocks & 3–3 &  \\
I12 & Alaska warm W number of blocks & 17–17 &  \\
I13 & incomplete(partial) work units & Alaska warm splits 10 partial &  \\
\end{longtable}}

\clearpage
\phantomsection\label{tab:S11}
\noindent\textbf{Supplementary Table S11.} Additional independent regions (L1e, L4e, L8e, L38, L39).\par\smallskip
\noindent{\footnotesize Licence-verified edition is the main edition; the North Atlantic rows are from the full-data edition and carry a few-blocks flag. All Tibetan scored cells lie above the source elevation range, so Tibetan contrasts are predicted outcomes and are not used as evidence for independent regions. New-region shards ran on the local server and P4 shards on the cloud (cross-environment flag). Method codes (internal identifiers allowed in the Supplementary Information): P0 source-coefficient Stefan model; P* year-matched Stefan model; P1 recalibrated Stefan model ($\kappa = 10$); P2 local least-squares Stefan model; P1* or P1@ed soil-property recalibrated Stefan model; R0 source anchor plus residual ML; R1 recalibrated anchor plus residual ML; R2 augmentation plus anchor plus residual; D0 direct ML; D1 physics pseudo-label augmentation; F1a, F1k, F1n physics-input ML; RM multiplicative residual; Re Stefan-CCI mean anchor plus residual; W selection rule (cross-validation within the target labels); S1 random cells, S2 block stratification, S3 covariate-cluster stratification, S4 covariate max-min distance, S5 source-extrapolation first, S6 sequential selection by prediction variance, S7 $\sqrt{\mathrm{TDD}}$ first. Target suffixes: $|x$ parent region excluded from the source, $|i$ rest of the parent kept, $|r$ within-region design. $n$ = number of target labels; all = all labels of half A. Verdicts: lower error, higher error, equivalent ($\pm$0.5~cm), undetermined; both 95\% block-bootstrap CIs (cell-weighted and block-equal) must agree. \dag~a note of the source row in Korean was omitted; see the cited README section.}\par

\par\medskip\noindent{\small\textit{a. Pool verdicts (L1e, L4e, L8e) (SI README 2.5)}}\par\nopagebreak
{\fontsize{7.5}{9}\selectfont\setlength{\tabcolsep}{4pt}\setlength{\LTpre}{2pt}\setlength{\LTpost}{4pt}
\begin{longtable}{@{}>{\raggedright\arraybackslash}p{\dimexpr 0.180\linewidth-8pt\relax}>{\raggedright\arraybackslash}p{\dimexpr 0.332\linewidth-8pt\relax}>{\raggedright\arraybackslash}p{\dimexpr 0.332\linewidth-8pt\relax}>{\raggedright\arraybackslash}p{\dimexpr 0.157\linewidth-8pt\relax}@{}}
\toprule
\textbf{ID} & \textbf{Filter} & \textbf{Value} & \textbf{Verdict} \\
\midrule\endfirsthead
\toprule
\textbf{ID} & \textbf{Filter} & \textbf{Value} & \textbf{Verdict} \\
\midrule\endhead
\bottomrule\endlastfoot
D-L1e-V-P4, -PE1, -PE2 & test\_\allowbreak{}id=L1e, pool ∈ \{P4, PE1, PE2\}, scope=verdict, column verdict, stat & significant improvement regions 0/\allowbreak{}4, 0/\allowbreak{}5, 1/\allowbreak{}6(Tibet\_\allowbreak{}LGD: n 0, 3, 10, 40) & supported ×3 \\
D-L1e-compare & test\_\allowbreak{}id=L1e, pool=P4·PE1·PE2, scope=compare & '4 regions verdict extended pool(regions 6) is maintained' &  \\
D-L4e-V-P4 & test\_\allowbreak{}id=L4e, pool=P4, scope=verdict & 'net value of ML with sparse labels'(minimum n 10) &  \\
D-L4e-V-PE1 & pool=PE1 & 'partial(regions 2/\allowbreak{}5): net value with sparse labels'(whole pool minimum n none, n ≤ 10 not met) &  \\
D-L4e-V-PE2 & pool=PE2 & 'net value of ML with sparse labels'(minimum n 3), flag 'shrinkage depends on prior misspecification' &  \\
D-L4e-PE1-n10 & test\_\allowbreak{}id=L4e, pool=PE1, item=R1-P1, scope=MEAN, n=10 & −0.13 [−0.43, 0.01]; block-equal −0.19 [−0.42, 0.00] & ns \\
D-L4e-PE2-n10 & pool=PE2, item=R1-P1, scope=MEAN, n=10 & −3.59 [−4.12, −3.13]; block-equal −2.80 [−3.43, −2.19] & improve(flag 'shrinkage depends on prior misspecification') \\
D-L4e-PE2-R1P2-n10 & pool=PE2, item=R1-P2, scope=MEAN, n=10 & 2.61 [0.33, 4.44]; block-equal −0.74 [−3.14, 1.73] & ns \\
D-L4e-pool-P4-10 & test\_\allowbreak{}id=L4e, pool=P4, item='(a)1 pool contrast R1-P1|\allowbreak{}n10', scope=compare\_\allowbreak{}pool & −0.18 [−0.52, −0.02]; block-equal −0.29 [−0.51, −0.08] & lower error \\
D-L4e-pool-PE1-10 & pool=PE1, same item & −0.13 [−0.42, 0.01]; block-equal −0.19 [−0.41, 0.01] & equivalent \\
D-L4e-pool-PE1-new-10 & pool=PE1-new(Russia\_\allowbreak{}C), same item & 0.06 [−0.40, 0.51]; block-equal 0.20 [−0.39, 0.81] & undetermined \\
D-L4e-pool-PE1-new--1 & pool=PE1-new, item='(a)1 pool contrast R1-P1|\allowbreak{}nall' & −1.40 [−2.95, −0.19]; block-equal −1.54 [−3.12, −0.10] & lower error \\
D-L4e-compare & test\_\allowbreak{}id=L4e, pool=P4·PE1·PE2, scope=compare & '4 regions verdict extended pool is not maintained. PE1: L28 weakened' &  \\
D-L8e-P4 & test\_\allowbreak{}id=L8e, pool=P4, item=R1-P0, scope=MEAN & −2.64 [−3.47, −1.90]; block-equal −2.64 [−3.53, −1.67] & improve \\
D-L8e-PE1 & pool=PE1 & −3.19 [−4.19, −2.22]; block-equal −2.62 [−3.71, −1.53] & improve \\
D-L8e-PE2 & pool=PE2 & −26.78 [−29.09, −24.21]; block-equal −23.13 [−26.03, −20.20] & improve(flag 'unblinded(label statistics viewed): Tibetan Plateau included') \\
D-L8e-PE2-rel & pool=PE2, scope=MEAN\_\allowbreak{}rel & −0.16 [−0.19, −0.13](relative units) & improve \\
D-L8e-compare & test\_\allowbreak{}id=L8e, pool=P4·PE1·PE2, scope=compare & '4 regions verdict extended pool(regions 6) is maintained' &  \\
\end{longtable}}
\par\medskip\noindent{\small\textit{b. Region-level contrasts (L38) (SI README 2.5)}}\par\nopagebreak
{\fontsize{7.5}{9}\selectfont\setlength{\tabcolsep}{4pt}\setlength{\LTpre}{2pt}\setlength{\LTpost}{4pt}
\begin{longtable}{@{}>{\raggedright\arraybackslash}p{\dimexpr 0.190\linewidth-8pt\relax}>{\raggedright\arraybackslash}p{\dimexpr 0.228\linewidth-8pt\relax}>{\raggedright\arraybackslash}p{\dimexpr 0.234\linewidth-8pt\relax}>{\raggedright\arraybackslash}p{\dimexpr 0.115\linewidth-8pt\relax}>{\raggedright\arraybackslash}p{\dimexpr 0.234\linewidth-8pt\relax}@{}}
\toprule
\textbf{ID} & \textbf{Item, target, $n$} & \textbf{Value (cm)} & \textbf{Verdict} & \textbf{Note} \\
\midrule\endfirsthead
\toprule
\textbf{ID} & \textbf{Item, target, $n$} & \textbf{Value (cm)} & \textbf{Verdict} & \textbf{Note} \\
\midrule\endhead
\bottomrule\endlastfoot
D-L38-Tibet-D0-P0-n0 & (a) D0-P0, Tibet\_\allowbreak{}LGD|\allowbreak{}x, 0 & −59.98 [−62.71, −57.10]; block-equal −62.91 [−65.29, −60.32] & lower error & P0 RMSE 244.51(rmse\_\allowbreak{}B). unblinded flag \\
D-L38-Tibet-D0-P0-n3, -n10, -n40 & (a) D0-P0, Tibet\_\allowbreak{}LGD|\allowbreak{}x, 3·10·40 & −130.18, −150.92, −157.59 & lower error ×3 & L1 form significant improvement(deep regime) \\
D-L38-Tibet-R1-P0-n-1 & (b) R1-P0, Tibet\_\allowbreak{}LGD|\allowbreak{}x, all & −144.73 [−157.42, −129.90]; block-equal −125.70 [−142.62, −108.64] & lower error &  \\
D-L38-Tibet-P1-P0-n10 & (c) P1-P0, Tibet\_\allowbreak{}LGD|\allowbreak{}x, 10 & −96.37 [−100.70, −91.92]; block-equal −97.75 [−102.91, −92.22] & lower error & P1 RMSE 148.14 \\
D-L38-Tibet-R1-P1-n10 & (c) R1-P1, Tibet\_\allowbreak{}LGD|\allowbreak{}x, 10 & −20.90 [−23.40, −17.89]; block-equal −15.82 [−19.40, −12.33] & lower error & flag 'shrinkage depends on prior misspecification'. R1 127.23, P1 148.14 \\
D-L38-Tibet-R1-P2-n10 & (c) R1-P2, Tibet\_\allowbreak{}LGD|\allowbreak{}x, 10 & 19.11 [5.74, 30.14]; block-equal 2.37 [−11.81, 16.77] & undetermined & (a)10 companion. P2 RMSE 108.12. cell-weighted CI all 0 above block-equal CI 0 'indistinguishable' not used\textsuperscript{\dag} \\
D-L38-Tibet-R1-P3-n10 & (c) R1-P3, Tibet\_\allowbreak{}LGD|\allowbreak{}x, 10 & 16.65 [9.53, 22.48]; block-equal 5.20 [−3.29, 13.58] & undetermined & (a)10 companion. R1 − P2 same form(cell-weighted CI 0 above, block-equal CI 0 included) \\
D-L38-Tibet-P2-P1-n10 & (c) P2-P1, Tibet\_\allowbreak{}LGD|\allowbreak{}x, 10 & −40.02 [−53.21, −24.07]; block-equal −18.19 [−35.84, −0.69] & lower error & flag 'depends on split independence' \\
D-L38-RussiaC-D0-P0-n0, -n3, -n10, -n-1 & (a) D0-P0, Russia\_\allowbreak{}C|\allowbreak{}x, 0·3·10·all & 2.75, 2.44, 1.81, −4.58 & higher error, higher error, undetermined, undetermined & P0 RMSE 32.12 \\
D-L38-RussiaC-R1-P0-n-1 & (b) R1-P0, Russia\_\allowbreak{}C|\allowbreak{}x, all & −5.42 [−8.96, −1.64]; block-equal −2.55 [−6.83, 1.48] & undetermined & replication failed regions(L38 (b)). cell-weighted CI all 0 below block-equal CI 0 weighting conflicting undetermined). cell-weighted point estimate block-equal CI\textsuperscript{\dag} \\
D-L38-RussiaC-P1-P0-n10 & (c) P1-P0, Russia\_\allowbreak{}C|\allowbreak{}x, 10 & −2.11 [−3.37, −0.61]; block-equal −0.70 [−2.24, 0.76] & undetermined &  \\
D-L38-RussiaC-R1-P1-n10 & (c) R1-P1, Russia\_\allowbreak{}C|\allowbreak{}x, 10 & 0.06 [−0.40, 0.51]; block-equal 0.20 [−0.39, 0.81] & undetermined &  \\
D-L38-RussiaW-ref-* & pool=P4 reference, Russia\_\allowbreak{}W|\allowbreak{}x: (b) R1-P0 all; (c) P1-P0 10; (c) R1-P1 10 & −13.98; −11.96; −0.56 [−1.21, −0.02], block-equal −0.24 [−0.84, 0.31] & lower error, lower error, undetermined & LG main run CatBoost shards. C2 reduced basis \\
D-L38-V & test\_\allowbreak{}id=L38, pool=new regions, scope=verdict & 'diverging regions: Russia\_\allowbreak{}C(shallow; (a) met, (b) undetermined); Tibet\_\allowbreak{}LGD(deep; (a) not met, (b) lower error)' &  & stat '(a) met 1/\allowbreak{}2, (b) lower error 1/\allowbreak{}2. symmetric reference(P4): (a) 4/\allowbreak{}4, (b) 1/\allowbreak{}4' \\
D-L38-V-note & same rows, column note & 'Tibetan Plateau scoring cells outside source support share: elevation 1.000, √TDD 0.333 …; Tibetan Plateau result NOVELTY-N7 sentence basis not used' &  &  \\
\end{longtable}}
\par\medskip\noindent{\small\textit{c. Label definition in the Tibetan Plateau (L39) (SI README 2.5)}}\par\nopagebreak
{\fontsize{7.5}{9}\selectfont\setlength{\tabcolsep}{4pt}\setlength{\LTpre}{2pt}\setlength{\LTpost}{4pt}
\begin{longtable}{@{}>{\raggedright\arraybackslash}p{\dimexpr 0.207\linewidth-8pt\relax}>{\raggedright\arraybackslash}p{\dimexpr 0.300\linewidth-8pt\relax}>{\raggedright\arraybackslash}p{\dimexpr 0.300\linewidth-8pt\relax}>{\raggedright\arraybackslash}p{\dimexpr 0.192\linewidth-8pt\relax}@{}}
\toprule
\textbf{ID} & \textbf{Filter} & \textbf{Value} & \textbf{Verdict} \\
\midrule\endfirsthead
\toprule
\textbf{ID} & \textbf{Filter} & \textbf{Value} & \textbf{Verdict} \\
\midrule\endhead
\bottomrule\endlastfoot
D-L39-V & test\_\allowbreak{}id=L39, scope=verdict & 'robust'(Tibetan Plateau labels definition: GPR vs temperature-derived) & robust \\
D-L39-tempderived-R1P1-n10 & test\_\allowbreak{}id=L39, item=R1-P1|\allowbreak{}n10, target=Tibet\_\allowbreak{}LGD\textasciitilde{}L39|\allowbreak{}x & −20.92 [−22.59, −18.23]; block-equal −21.64 [−23.90, −19.20] & lower error(label set at other locations) \\
D-L40-V-* & test\_\allowbreak{}id=L40, scope=verdict, target ∈ \{Russia\_\allowbreak{}W\textasciitilde{}exp\textasciitilde{}lic, Russia\_\allowbreak{}E\textasciitilde{}exp, Canada\textasciitilde{}exp\textasciitilde{}lic, Russia\_\allowbreak{}E\textasciitilde{}expnokyt\}, column verdict, stat & four expanded editions all 'weakened'. four edition all R1-P0|\allowbreak{}nall P1-P0|\allowbreak{}n10 'robust', weakened D0-P0 contrast came from & weakened ×4 \\
D-L41-V-* & test\_\allowbreak{}id=L41, scope=verdict, column verdict, stat, note, same\_\allowbreak{}as; targets number of cells tables/\allowbreak{}lgd\_\allowbreak{}eligibility\_\allowbreak{}v1.csv n\_\allowbreak{}cells, eligible\textsuperscript{\dag} & Tibetan Plateau: (a)(c)(h) targets cells 0 'variant not possible(eligible in the table ineligible)robustmain setting TMx rows main setting sameweakened'(D0-P0|\allowbreak{}n3 higher error → undetermined), (e) 'weakened'(D0-P0|\allowbreak{}n10 undetermined → higher error), (d)(f) 'robust', (c)(g) 'robust' same\_\allowbreak{}as sensitivity\textsuperscript{\dag} & (b)(e)(d)(f) substantive variants \\
D-L42-V-* & test\_\allowbreak{}id=L42, scope=verdict, column verdict, stat & Tibet 'robust', Russia\_\allowbreak{}C 'weakened'(D0-P0|\allowbreak{}n0·n3 higher error → undetermined, R1-P1|\allowbreak{}n10 undetermined → lower error), NAtlantic 'rows none(ineligible)' &  \\
D-elig-* & Delig, spec ∈ \{NAtlantic\_\allowbreak{}lic, Tibet\_\allowbreak{}L42\_\allowbreak{}lic, Russia\_\allowbreak{}C\_\allowbreak{}L42\_\allowbreak{}lic\}, column n\_\allowbreak{}cells, nb\_\allowbreak{}union, min\_\allowbreak{}nb\_\allowbreak{}eval\_\allowbreak{}used, eligible, regime\textsuperscript{\dag} & NAtlantic 22cells, blocks union 6, ineligible; Tibet 132cells, 34, eligible, deep; Russia\_\allowbreak{}C 57cells, 13, eligible, shallow &  \\
\end{longtable}}
\par\medskip\noindent{\small\textit{d. North Atlantic, full-data edition (L38) (SI README 2.5)}}\par\nopagebreak
{\fontsize{7.5}{9}\selectfont\setlength{\tabcolsep}{4pt}\setlength{\LTpre}{2pt}\setlength{\LTpost}{4pt}
\begin{longtable}{@{}>{\raggedright\arraybackslash}p{\dimexpr 0.210\linewidth-8pt\relax}>{\raggedright\arraybackslash}p{\dimexpr 0.281\linewidth-8pt\relax}>{\raggedright\arraybackslash}p{\dimexpr 0.322\linewidth-8pt\relax}>{\raggedright\arraybackslash}p{\dimexpr 0.109\linewidth-8pt\relax}>{\raggedright\arraybackslash}p{\dimexpr 0.079\linewidth-8pt\relax}@{}}
\toprule
\textbf{ID} & \textbf{Item, $n$} & \textbf{Value (cm)} & \textbf{Verdict} & \textbf{Note} \\
\midrule\endfirsthead
\toprule
\textbf{ID} & \textbf{Item, $n$} & \textbf{Value (cm)} & \textbf{Verdict} & \textbf{Note} \\
\midrule\endhead
\bottomrule\endlastfoot
Dfull-L38-NAtl-D0-P0-n0 & (a) D0-P0, 0 & −9.91 [−15.95, −2.85]; block-equal −6.15 [−14.53, 2.40] & undetermined & flag 'few blocks'. P0 RMSE 29.80 \\
Dfull-L38-NAtl-D0-P0-n10 & (a) D0-P0, 10 & −9.15 [−14.87, −1.62]; block-equal −4.37 [−12.73, 4.15] & undetermined &  \\
Dfull-L38-NAtl-R1-P0-n-1 & (b) R1-P0, all & −4.65 [−8.50, 1.42]; block-equal −0.42 [−6.40, 5.86] & undetermined &  \\
Dfull-L38-NAtl-P1-P0-n10, R1-P1-n10 & (c) P1-P0, R1-P1, 10 & −4.22; −0.89 & undetermined ×2 &  \\
Dfull-L1e-V-PE1 & test\_\allowbreak{}id=L1e, pool=PE1, scope=verdict & significant improvement regions 0/\allowbreak{}6(NAtlantic included) & supported &  \\
Dfull-L8e-V-PE1 & test\_\allowbreak{}id=L8e, pool=PE1, scope=verdict & Δ −3.44 [−4.45, −2.14], block-equal [−3.54, −0.87], regions 6/\allowbreak{}6 & supported &  \\
Dfull-L4e-compare & test\_\allowbreak{}id=L4e, pool=P4·PE1·PE2, scope=compare & '4 regions verdict extended pool is not maintained'(R1-P1|\allowbreak{}n10: P4 lower error → PE1 undetermined) &  &  \\
\end{longtable}}

\clearpage
\phantomsection\label{tab:S12}
\noindent\textbf{Supplementary Table S12.} Coverage of 90\% prediction intervals by step: within-region CQR, transfer CQR, pooled generative intervals and hierarchical conformal intervals.\par\smallskip
\noindent{\footnotesize Nominal coverage 0.90. Coverage [95\% CI] and mean width (cm); cell-weighted values. The hierarchical conformal coverage of 0.86 uses a different protocol from the zero-label interval comparison of Fig.~6d and is therefore reported here. Values are copied from paper/claims/SI\_uncertainty/README.md section 2 (no new calculation; section 4 recommendation). Detailed rows follow.}\par

\par\medskip\noindent{\small\textit{a. Steps}}\par\nopagebreak
{\fontsize{7.5}{9}\selectfont\setlength{\tabcolsep}{4pt}\setlength{\LTpre}{2pt}\setlength{\LTpost}{4pt}
\begin{longtable}{@{}>{\raggedright\arraybackslash}p{\dimexpr 0.400\linewidth-8pt\relax}>{\raggedright\arraybackslash}p{\dimexpr 0.140\linewidth-8pt\relax}>{\raggedright\arraybackslash}p{\dimexpr 0.090\linewidth-8pt\relax}>{\raggedright\arraybackslash}p{\dimexpr 0.220\linewidth-8pt\relax}>{\raggedright\arraybackslash}p{\dimexpr 0.150\linewidth-8pt\relax}@{}}
\toprule
\textbf{Step} & \textbf{Coverage [CI]} & \textbf{Width (cm)} & \textbf{Verdict} & \textbf{Source} \\
\midrule\endfirsthead
\toprule
\textbf{Step} & \textbf{Coverage [CI]} & \textbf{Width (cm)} & \textbf{Verdict} & \textbf{Source} \\
\midrule\endhead
\bottomrule\endlastfoot
Within-region CQR (S11, contest; Alaska, 6-fold block out-of-fold) & 0.93 [0.89, 0.97] & 53.64 & descriptive & U1, SI\_uncertainty README 2.2 \\
Transfer CQR trained in Alaska (H17; four main regions, two conditions) & 0.32--0.73 &  & H17 confirmed (upper CI $<$ 0.90 in 4/4) & U2, README 2.3 \\
Pooled generative quantile intervals, nflow (LGU-B1, $\lambda$ 1.0) & 0.76 [0.67, 0.83] &  & supported (replication, unblinded) & U4, README 2.5 \\
Pooled generative quantile intervals, cfm (LGU-B1, $\lambda$ 1.0) & 0.65 [0.56, 0.76] &  & supported (replication, unblinded) & U4, README 2.5 \\
Hierarchical conformal, hier2\_cdf (h37, four main regions) & 0.86 [0.80, 0.92] & 80.54 & partially supported (F10) & U3, README 2.4 \\
Hierarchical conformal, constant normalizer (LGU-B, zero labels) & 0.86 [0.83, 0.90] & 81.29 &  & U4, README 2.5 \\
\end{longtable}}
\par\medskip\noindent{\small\textit{b. U1 within-region intervals in Alaska (SI README 2.2)}}\par\nopagebreak
{\fontsize{7.5}{9}\selectfont\setlength{\tabcolsep}{4pt}\setlength{\LTpre}{2pt}\setlength{\LTpost}{4pt}
\begin{longtable}{@{}>{\raggedright\arraybackslash}p{\dimexpr 0.246\linewidth-8pt\relax}>{\raggedright\arraybackslash}p{\dimexpr 0.278\linewidth-8pt\relax}>{\raggedright\arraybackslash}p{\dimexpr 0.278\linewidth-8pt\relax}>{\raggedright\arraybackslash}p{\dimexpr 0.199\linewidth-8pt\relax}@{}}
\toprule
\textbf{Item} & \textbf{Columns} & \textbf{Value} & \textbf{Verdict} \\
\midrule\endfirsthead
\toprule
\textbf{Item} & \textbf{Columns} & \textbf{Value} & \textbf{Verdict} \\
\midrule\endhead
\bottomrule\endlastfoot
uncalibrated quantile interval(3 seed ensemble interval, quantile bounds mean) & coverage /\allowbreak{} cov\_\allowbreak{}ci\_\allowbreak{}lo /\allowbreak{} cov\_\allowbreak{}ci\_\allowbreak{}hi /\allowbreak{} width\_\allowbreak{}cm & 0.45 /\allowbreak{} 0.32 /\allowbreak{} 0.62 /\allowbreak{} 14.65 & descriptive(contest 'adopted') \\
CQR(3 seed ensemble interval, quantile bounds mean) & coverage /\allowbreak{} cov\_\allowbreak{}ci\_\allowbreak{}lo /\allowbreak{} cov\_\allowbreak{}ci\_\allowbreak{}hi /\allowbreak{} width\_\allowbreak{}cm & 0.93 /\allowbreak{} 0.89 /\allowbreak{} 0.97 /\allowbreak{} 53.64 & descriptive(contest 'adopted') \\
before correction interval score & interval\_\allowbreak{}score /\allowbreak{} is\_\allowbreak{}ci\_\allowbreak{}lo /\allowbreak{} is\_\allowbreak{}ci\_\allowbreak{}hi & 123.91 /\allowbreak{} 88.04 /\allowbreak{} 153.78 & descriptive \\
CQR interval score & interval\_\allowbreak{}score /\allowbreak{} is\_\allowbreak{}ci\_\allowbreak{}lo /\allowbreak{} is\_\allowbreak{}ci\_\allowbreak{}hi & 68.85 /\allowbreak{} 58.64 /\allowbreak{} 94.92 & descriptive(manuscript action: descriptive relaxed) \\
CQR decomposition & is\_\allowbreak{}width\_\allowbreak{}part /\allowbreak{} is\_\allowbreak{}under\_\allowbreak{}part /\allowbreak{} is\_\allowbreak{}over\_\allowbreak{}part & 53.64 /\allowbreak{} 2.66 /\allowbreak{} 12.55 & descriptive \\
CQR coverage(E4 re-scored, S11 same ensemble interval) & coverage /\allowbreak{} cov\_\allowbreak{}ci\_\allowbreak{}lo /\allowbreak{} cov\_\allowbreak{}ci\_\allowbreak{}hi & 0.93 /\allowbreak{} 0.89 /\allowbreak{} 0.96 & descriptive \\
constant width reference interval score & interval\_\allowbreak{}score /\allowbreak{} is\_\allowbreak{}ci\_\allowbreak{}lo /\allowbreak{} is\_\allowbreak{}ci\_\allowbreak{}hi & 65.39 /\allowbreak{} 55.61 /\allowbreak{} 92.19 & descriptive \\
constant width reference decomposition & is\_\allowbreak{}width\_\allowbreak{}part /\allowbreak{} is\_\allowbreak{}under\_\allowbreak{}part /\allowbreak{} is\_\allowbreak{}over\_\allowbreak{}part & 43.46 /\allowbreak{} 3.84 /\allowbreak{} 18.09 & descriptive \\
constant width reference coverage & coverage /\allowbreak{} cov\_\allowbreak{}ci\_\allowbreak{}lo /\allowbreak{} cov\_\allowbreak{}ci\_\allowbreak{}hi & 0.90 /\allowbreak{} 0.85 /\allowbreak{} 0.92 & descriptive \\
constant width half-width(post hoc reference) & half\_\allowbreak{}width\_\allowbreak{}cm & 21.73 & not applicable \\
CQR − constant width(arithmetic difference) & interval\_\allowbreak{}score & +3.46 & descriptive. paired CI none. document [0.6, 6.5] exactly not reproduced(draws dependence, below reading notes)\textsuperscript{\dag} \\
CQR conditional coverage(observed ALT 5quantile) & coverage & 0.85 /\allowbreak{} 1.00 /\allowbreak{} 1.00 /\allowbreak{} 1.00 /\allowbreak{} 0.83 & descriptive \\
CQR width 5quantile both ends & coverage /\allowbreak{} interval\_\allowbreak{}score & W1 0.86 /\allowbreak{} 75.52, W5 0.97 /\allowbreak{} 76.59 & descriptive \\
CQR blocksby coverage(cells 20 or more 33blocks) & coverage(median) /\allowbreak{} lo(10 percentile) /\allowbreak{} hi(90 percentile) /\allowbreak{} interval\_\allowbreak{}score(coverage 0.8 below blocks share) & 0.98 /\allowbreak{} 0.70 /\allowbreak{} 1.00 /\allowbreak{} 0.15 & descriptive \\
reference: M1 edition Alaska within region CQR(x25) & coverage /\allowbreak{} cov\_\allowbreak{}ci\_\allowbreak{}lo /\allowbreak{} cov\_\allowbreak{}ci\_\allowbreak{}hi /\allowbreak{} width\_\allowbreak{}mean /\allowbreak{} interval\_\allowbreak{}score & 0.91 /\allowbreak{} 0.83 /\allowbreak{} 0.94 /\allowbreak{} 48.94 /\allowbreak{} 70.35 & descriptive \\
\end{longtable}}
\par\medskip\noindent{\small\textit{c. U2 transfer CQR (coverage) (SI README 2.3)}}\par\nopagebreak
{\fontsize{7.5}{9}\selectfont\setlength{\tabcolsep}{4pt}\setlength{\LTpre}{2pt}\setlength{\LTpost}{4pt}
\begin{longtable}{@{}>{\raggedright\arraybackslash}p{\dimexpr 0.135\linewidth-8pt\relax}>{\raggedright\arraybackslash}p{\dimexpr 0.091\linewidth-8pt\relax}>{\raggedright\arraybackslash}p{\dimexpr 0.148\linewidth-8pt\relax}>{\raggedright\arraybackslash}p{\dimexpr 0.148\linewidth-8pt\relax}>{\raggedright\arraybackslash}p{\dimexpr 0.148\linewidth-8pt\relax}>{\raggedright\arraybackslash}p{\dimexpr 0.148\linewidth-8pt\relax}>{\raggedright\arraybackslash}p{\dimexpr 0.180\linewidth-8pt\relax}@{}}
\toprule
\textbf{Method} & \textbf{Condition} & \textbf{Lena Delta} & \textbf{Canada} & \textbf{W Russia} & \textbf{E Russia} & \textbf{Verdict} \\
\midrule\endfirsthead
\toprule
\textbf{Method} & \textbf{Condition} & \textbf{Lena Delta} & \textbf{Canada} & \textbf{W Russia} & \textbf{E Russia} & \textbf{Verdict} \\
\midrule\endhead
\bottomrule\endlastfoot
cqr\_\allowbreak{}ak & noinfo & 0.57 [0.48, 0.66] (0.5700) & 0.60 [0.43, 0.74] (0.5953) & 0.32 [0.20, 0.43] (0.3214) & 0.73 [0.56, 0.89] (0.7284) & H17 confirmed(4/\allowbreak{}4 CI upper bound < 0.90) \\
cqr\_\allowbreak{}ak & covonly & 0.57 [0.46, 0.64] (0.5748) & 0.53 [0.35, 0.74] (0.5292) & 0.34 [0.17, 0.45] (0.3415) & 0.66 [0.48, 0.90] (0.6597, upper bound 0.8964) & H17 confirmed(4/\allowbreak{}4) \\
cqr\_\allowbreak{}iw & noinfo & 0.35 [0.26, 0.41] & 0.52 [0.38, 0.68] & 0.30 [0.17, 0.43] & 0.69 [0.53, 0.87] & descriptive(recovery none) \\
cqr\_\allowbreak{}iw & covonly & 0.33 [0.21, 0.40] & 0.47 [0.30, 0.66] & 0.34 [0.15, 0.46] & 0.62 [0.44, 0.87] & descriptive \\
anchor\_\allowbreak{}ak & noinfo & 0.92 [0.79, 0.95] & 0.77 [0.68, 0.87] & 0.57 [0.44, 0.71] & 0.84 [0.71, 1.00] & descriptive \\
anchor\_\allowbreak{}iw & noinfo & 0.93 [0.79, 0.96] & 0.88 [0.79, 0.94] & 0.56 [0.41, 0.70] & 0.84 [0.71, 1.00] & descriptive \\
\end{longtable}}
\par\medskip\noindent{\small\textit{d. U3 zero-label hierarchical conformal intervals (SI README 2.4)}}\par\nopagebreak
{\fontsize{7.5}{9}\selectfont\setlength{\tabcolsep}{4pt}\setlength{\LTpre}{2pt}\setlength{\LTpost}{4pt}
\begin{longtable}{@{}>{\raggedright\arraybackslash}p{\dimexpr 0.181\linewidth-8pt\relax}>{\raggedright\arraybackslash}p{\dimexpr 0.240\linewidth-8pt\relax}>{\raggedright\arraybackslash}p{\dimexpr 0.240\linewidth-8pt\relax}>{\raggedright\arraybackslash}p{\dimexpr 0.338\linewidth-8pt\relax}@{}}
\toprule
\textbf{Item} & \textbf{Columns} & \textbf{Value} & \textbf{Verdict} \\
\midrule\endfirsthead
\toprule
\textbf{Item} & \textbf{Columns} & \textbf{Value} & \textbf{Verdict} \\
\midrule\endhead
\bottomrule\endlastfoot
hier2\_\allowbreak{}cdf(main method) main 4 regions & coverage /\allowbreak{} coverage\_\allowbreak{}lo /\allowbreak{} coverage\_\allowbreak{}hi /\allowbreak{} coverage\_\allowbreak{}beq /\allowbreak{} width\_\allowbreak{}cm /\allowbreak{} interval\_\allowbreak{}score & 0.86 /\allowbreak{} 0.80 /\allowbreak{} 0.92 /\allowbreak{} 0.85 /\allowbreak{} 80.54 /\allowbreak{} 131.98 & partially supported(F10) \\
cell pooling(pooled) & coverage /\allowbreak{} coverage\_\allowbreak{}lo /\allowbreak{} coverage\_\allowbreak{}hi /\allowbreak{} coverage\_\allowbreak{}beq /\allowbreak{} width\_\allowbreak{}cm /\allowbreak{} interval\_\allowbreak{}score & 0.73 /\allowbreak{} 0.65 /\allowbreak{} 0.81 /\allowbreak{} 0.74 /\allowbreak{} 56.34 /\allowbreak{} 152.86 & descriptive \\
hier2\_\allowbreak{}sub & coverage /\allowbreak{} coverage\_\allowbreak{}lo /\allowbreak{} coverage\_\allowbreak{}hi /\allowbreak{} width\_\allowbreak{}cm & 0.83 /\allowbreak{} 0.77 /\allowbreak{} 0.89 /\allowbreak{} 72.86 & descriptive \\
hier2\_\allowbreak{}blk & coverage /\allowbreak{} coverage\_\allowbreak{}lo /\allowbreak{} coverage\_\allowbreak{}hi /\allowbreak{} width\_\allowbreak{}cm & 0.87 /\allowbreak{} 0.81 /\allowbreak{} 0.93 /\allowbreak{} 83.30 & descriptive \\
hier2\_\allowbreak{}cdf Lena Delta & coverage /\allowbreak{} coverage\_\allowbreak{}lo /\allowbreak{} coverage\_\allowbreak{}hi /\allowbreak{} coverage\_\allowbreak{}beq /\allowbreak{} width\_\allowbreak{}cm /\allowbreak{} K\_\allowbreak{}groups & 0.89 /\allowbreak{} 0.80 /\allowbreak{} 0.96 /\allowbreak{} 0.81 /\allowbreak{} 73.03 /\allowbreak{} 4 & descriptive \\
hier2\_\allowbreak{}cdf Canada & same & 0.89 /\allowbreak{} 0.80 /\allowbreak{} 0.97 /\allowbreak{} 0.88 /\allowbreak{} 95.35 /\allowbreak{} 4 & descriptive \\
hier2\_\allowbreak{}cdf W Russia & same & 0.75 /\allowbreak{} 0.59 /\allowbreak{} 0.88 /\allowbreak{} 0.74 /\allowbreak{} 71.13 /\allowbreak{} 4 & descriptive(outside the target band) \\
hier2\_\allowbreak{}cdf E Russia & same & 0.93 /\allowbreak{} 0.80 /\allowbreak{} 1.00 /\allowbreak{} 0.97 /\allowbreak{} 82.64 /\allowbreak{} 4 & descriptive \\
hier2\_\allowbreak{}cdf Alaska(reference) & same & 0.98 /\allowbreak{} 0.96 /\allowbreak{} 0.99 /\allowbreak{} 0.97 /\allowbreak{} 78.54 /\allowbreak{} 4 & descriptive \\
F10 verdict rows & in\_\allowbreak{}band /\allowbreak{} coverage(MAIN6) & 4 /\allowbreak{} 0.88 & partially supported. ci\_\allowbreak{}flag: 'hier2\_\allowbreak{}cdf in [0.85,0.95]: 4/\allowbreak{}6; ref\_\allowbreak{}cqr\_\allowbreak{}ak<0.85: 5/\allowbreak{}6; regions\_\allowbreak{}in\_\allowbreak{}band=Russia\_\allowbreak{}C,Russia\_\allowbreak{}E,Canada,Lena' \\
\end{longtable}}
\par\medskip\noindent{\small\textit{e. U4 normalizers and generative intervals (LGU-B) (SI README 2.5)}}\par\nopagebreak
{\fontsize{7.5}{9}\selectfont\setlength{\tabcolsep}{4pt}\setlength{\LTpre}{2pt}\setlength{\LTpost}{4pt}
\begin{longtable}{@{}>{\raggedright\arraybackslash}p{\dimexpr 0.325\linewidth-8pt\relax}>{\raggedright\arraybackslash}p{\dimexpr 0.259\linewidth-8pt\relax}>{\raggedright\arraybackslash}p{\dimexpr 0.178\linewidth-8pt\relax}>{\raggedright\arraybackslash}p{\dimexpr 0.238\linewidth-8pt\relax}@{}}
\toprule
\textbf{Hypothesis} & \textbf{Cell-weighted} & \textbf{Block-equal} & \textbf{Verdict (blinding)} \\
\midrule\endfirsthead
\toprule
\textbf{Hypothesis} & \textbf{Cell-weighted} & \textbf{Block-equal} & \textbf{Verdict (blinding)} \\
\midrule\endhead
\bottomrule\endlastfoot
LGU-B1 nflow, cqr\_\allowbreak{}pool λ 1.0 & 0.76 [0.67, 0.83] (0.7568 [0.6715, 0.8262]) & 0.79 [0.70, 0.84] (0.7851) & supported(replication (unblinded)) \\
LGU-B1 cfm, cqr\_\allowbreak{}pool λ 1.0 & 0.65 [0.56, 0.76] (0.6529 [0.5636, 0.7645]) & 0.68 [0.60, 0.80] (0.6844) & supported(replication (unblinded)) \\
LGU-B2 interval score nflow − placebo(nflow), confirmatory & +3.29 [−4.17, 8.16] & +0.03 [−6.49, 5.52] & undetermined(blind). verdict\_\allowbreak{}3region undetermined, n\_\allowbreak{}neg\_\allowbreak{}regions 1, coverage\_\allowbreak{}pool 0.89(0.8919) \\
LGU-B3 nflow − cbq & −7.41 [−52.55, 10.00] & −1.46 [−61.11, 23.49] & undetermined(equivalent not testable) \\
LGU-B4 nflow − const(as calibrated) & +3.75 [−10.73, 12.91] & +4.06 [−7.02, 15.18] & undetermined \\
LGU-B4 nflow − const(mean log width matched) & +4.96 [−0.81, 9.00] & +0.72 [−3.80, 6.65] & undetermined \\
LGU-B4 cbq − const(as calibrated) & +11.16 [−11.20, 53.46] & +5.52 [−24.01, 65.10] & undetermined \\
LGU-B4 cbq − const(width matched) & +5.24 [−5.34, 18.66] & +0.18 [−9.62, 12.90] & undetermined \\
LGU-B4 cfm − const(as calibrated) & −1.70 [−17.91, 10.03] & −0.58 [−11.86, 10.63] & undetermined \\
LGU-B4 cfm − const(width matched) & +2.72 [−1.38, 6.88] & −1.00 [−4.62, 3.66] & undetermined \\
LGU-B4 phys − const(as calibrated) & −0.83 [−8.74, 2.54] & −4.64 [−12.13, 0.82] & undetermined \\
LGU-B4 phys − const(width matched) & +0.78 [−4.02, 3.46] & −2.20 [−7.42, 1.80] & undetermined \\
LGU-B6 nflow − const, n 10 /\allowbreak{} 40 /\allowbreak{} 160(range B, Lena Delta and Canada·Alaska) & +10.71 [4.45, 15.57] /\allowbreak{} +10.81 [2.97, 18.57] /\allowbreak{} +10.93 [2.83, 18.93] & +3.50 [−3.52, 10.42] /\allowbreak{} +3.25 [−3.01, 13.03] /\allowbreak{} +2.55 [−3.04, 12.06] & undetermined(cell-weighted higher error side) \\
LGU-B6 cbq − const, n 10 /\allowbreak{} 40 /\allowbreak{} 160 & +8.37 [−8.95, 135.07] /\allowbreak{} +38.34 [1.30, 272.33] /\allowbreak{} +43.00 [3.33, 279.22] & −8.86 [−34.22, 135.12] /\allowbreak{} +18.05 [−24.13, 298.48] /\allowbreak{} +23.93 [−22.35, 300.43] & undetermined \\
\end{longtable}}
\par\medskip\noindent{\small\textit{f. U4 normalizers at zero labels (SI README 2.5)}}\par\nopagebreak
{\fontsize{7.5}{9}\selectfont\setlength{\tabcolsep}{4pt}\setlength{\LTpre}{2pt}\setlength{\LTpost}{4pt}
\begin{longtable}{@{}>{\raggedright\arraybackslash}p{\dimexpr 0.188\linewidth-8pt\relax}>{\raggedright\arraybackslash}p{\dimexpr 0.229\linewidth-8pt\relax}>{\raggedright\arraybackslash}p{\dimexpr 0.216\linewidth-8pt\relax}>{\raggedright\arraybackslash}p{\dimexpr 0.110\linewidth-8pt\relax}>{\raggedright\arraybackslash}p{\dimexpr 0.128\linewidth-8pt\relax}>{\raggedright\arraybackslash}p{\dimexpr 0.128\linewidth-8pt\relax}@{}}
\toprule
\textbf{Normalizer} & \textbf{cov10} & \textbf{Two-stage CI} & \textbf{cov10 block-equal} & \textbf{Width (cm)} & \textbf{Interval score (cm)} \\
\midrule\endfirsthead
\toprule
\textbf{Normalizer} & \textbf{cov10} & \textbf{Two-stage CI} & \textbf{cov10 block-equal} & \textbf{Width (cm)} & \textbf{Interval score (cm)} \\
\midrule\endhead
\bottomrule\endlastfoot
const & 0.86 (0.8640) & [0.83, 0.90] & 0.85 & 81.29 & 131.21 \\
phys & 0.87 (0.8724) & [0.84, 0.90] & 0.87 & 82.98 & 130.38 \\
nflow & 0.89 (0.8919) & [0.86, 0.90] & 0.89 & 96.01 & 134.96 \\
nflow placebo & 0.89 (0.8872) & [0.86, 0.90] & 0.88 & 94.95 & 131.67 \\
cfm & 0.88 (0.8850) & [0.86, 0.90] & 0.88 & 94.34 & 129.50 \\
cbq & 0.95 (0.9533) & [0.91, 0.97] & 0.95 & 132.44 & 142.36 \\
\end{longtable}}
\par\medskip\noindent{\small\textit{g. U5 hierarchical predictive distributions with labels and AB10 (LGU-A) (SI README 2.6)}}\par\nopagebreak
{\fontsize{7.5}{9}\selectfont\setlength{\tabcolsep}{4pt}\setlength{\LTpre}{2pt}\setlength{\LTpost}{4pt}
\begin{longtable}{@{}>{\raggedright\arraybackslash}p{\dimexpr 0.231\linewidth-8pt\relax}>{\raggedright\arraybackslash}p{\dimexpr 0.296\linewidth-8pt\relax}>{\raggedright\arraybackslash}p{\dimexpr 0.296\linewidth-8pt\relax}>{\raggedright\arraybackslash}p{\dimexpr 0.178\linewidth-8pt\relax}@{}}
\toprule
\textbf{Hypothesis} & \textbf{Columns} & \textbf{Value} & \textbf{Verdict} \\
\midrule\endfirsthead
\toprule
\textbf{Hypothesis} & \textbf{Columns} & \textbf{Value} & \textbf{Verdict} \\
\midrule\endhead
\bottomrule\endlastfoot
LGU-A1 (iii) − B4, n 10(confirmatory, AB10) & delta /\allowbreak{} ci\_\allowbreak{}lo /\allowbreak{} ci\_\allowbreak{}hi & +9.51 /\allowbreak{} −0.75 /\allowbreak{} 22.11 & undetermined(blind) \\
same & delta\_\allowbreak{}beq /\allowbreak{} ci\_\allowbreak{}lo\_\allowbreak{}beq /\allowbreak{} ci\_\allowbreak{}hi\_\allowbreak{}beq & +8.64 /\allowbreak{} −2.55 /\allowbreak{} 22.04 & same \\
same & verdict\_\allowbreak{}1stage /\allowbreak{} ci\_\allowbreak{}lo\_\allowbreak{}1stage /\allowbreak{} ci\_\allowbreak{}hi\_\allowbreak{}1stage & higher error /\allowbreak{} 4.55 /\allowbreak{} 13.02 & depends on calibration uncertainty \\
same & eq\_\allowbreak{}state /\allowbreak{} half\_\allowbreak{}width /\allowbreak{} ref\_\allowbreak{}score /\allowbreak{} coverage\_\allowbreak{}pool & not testable(precision insufficient) /\allowbreak{} 12.29 /\allowbreak{} 99.03 /\allowbreak{} 0.90(0.9011) & same \\
LGU-A1 n 40 & delta /\allowbreak{} ci\_\allowbreak{}lo /\allowbreak{} ci\_\allowbreak{}hi & −2.80 /\allowbreak{} −12.50 /\allowbreak{} 7.89 & undetermined \\
same & delta\_\allowbreak{}beq /\allowbreak{} ci\_\allowbreak{}lo\_\allowbreak{}beq /\allowbreak{} ci\_\allowbreak{}hi\_\allowbreak{}beq & −2.48 /\allowbreak{} −13.71 /\allowbreak{} 8.58 & same \\
same & verdict\_\allowbreak{}1stage /\allowbreak{} ci\_\allowbreak{}lo\_\allowbreak{}1stage /\allowbreak{} ci\_\allowbreak{}hi\_\allowbreak{}1stage & lower error /\allowbreak{} −4.95 /\allowbreak{} −0.56 & depends on calibration uncertainty \\
same & eq\_\allowbreak{}state /\allowbreak{} half\_\allowbreak{}width /\allowbreak{} ref\_\allowbreak{}score /\allowbreak{} coverage\_\allowbreak{}pool & not testable(precision insufficient) /\allowbreak{} 11.15 /\allowbreak{} 105.46 /\allowbreak{} 0.90(0.9017) & same \\
LGU-A1 overall & verdict & undetermined & undetermined \\
same contrast \% converted(figure table) & delta\_\allowbreak{}pct /\allowbreak{} ci\_\allowbreak{}lo\_\allowbreak{}pct /\allowbreak{} ci\_\allowbreak{}hi\_\allowbreak{}pct & n 10: 9.60 /\allowbreak{} −0.76 /\allowbreak{} 22.33; n 40: −2.66 /\allowbreak{} −11.85 /\allowbreak{} 7.48; B2: 2.50 /\allowbreak{} −3.17 /\allowbreak{} 6.19 & undetermined \\
AB10 two weighting p & delta /\allowbreak{} p\_\allowbreak{}cell /\allowbreak{} p\_\allowbreak{}beq /\allowbreak{} p\_\allowbreak{}eq /\allowbreak{} verdict4 & 9.51 /\allowbreak{} 0.07 /\allowbreak{} 0.09 /\allowbreak{} 0.80 /\allowbreak{} undetermined & undetermined \\
AB10 bundle rows & delta /\allowbreak{} holm\_\allowbreak{}p /\allowbreak{} holm\_\allowbreak{}p\_\allowbreak{}eq /\allowbreak{} holm\_\allowbreak{}m /\allowbreak{} verdict4 /\allowbreak{} abstract\_\allowbreak{}rule & 9.51 /\allowbreak{} 0.28 /\allowbreak{} 1.00 /\allowbreak{} 10 /\allowbreak{} undetermined /\allowbreak{} '(d) no difference established' & undetermined(J7) \\
LGU-A3 (iii) median − P1, RMSE, main 4 regions n 3 & delta /\allowbreak{} ci\_\allowbreak{}lo /\allowbreak{} ci\_\allowbreak{}hi; delta\_\allowbreak{}beq /\allowbreak{} ci\_\allowbreak{}lo\_\allowbreak{}beq /\allowbreak{} ci\_\allowbreak{}hi\_\allowbreak{}beq & +1.43 /\allowbreak{} 0.68 /\allowbreak{} 1.81; +1.73 /\allowbreak{} 0.96 /\allowbreak{} 2.16 & higher error \\
same n 10 & same & +1.60 /\allowbreak{} 0.57 /\allowbreak{} 1.88; +1.65 /\allowbreak{} 0.76 /\allowbreak{} 2.13 & higher error \\
same W Russia excluded 3 regions n 3 /\allowbreak{} n 10 & delta /\allowbreak{} ci\_\allowbreak{}lo /\allowbreak{} ci\_\allowbreak{}hi & +0.90 /\allowbreak{} 0.41 /\allowbreak{} 1.68; +1.08 /\allowbreak{} 0.32 /\allowbreak{} 1.80 & higher error \\
LGU-A7 (iv) − (iii), interval score, n 10 & delta /\allowbreak{} ci\_\allowbreak{}lo /\allowbreak{} ci\_\allowbreak{}hi & +25.23 /\allowbreak{} 17.46 /\allowbreak{} 34.02 & higher error \\
LGU-A2 infinite interval share, n 3(5rungs) & delta(inf10) maximum & 0.00 & passed(infinite interval 0) \\
\end{longtable}}
\par\medskip\noindent{\small\textit{h. U6 R1 intervals with labels (LGX L25) (SI README 2.7)}}\par\nopagebreak
{\fontsize{7.5}{9}\selectfont\setlength{\tabcolsep}{4pt}\setlength{\LTpre}{2pt}\setlength{\LTpost}{4pt}
\begin{longtable}{@{}>{\raggedright\arraybackslash}p{\dimexpr 0.307\linewidth-8pt\relax}>{\raggedright\arraybackslash}p{\dimexpr 0.089\linewidth-8pt\relax}>{\raggedright\arraybackslash}p{\dimexpr 0.202\linewidth-8pt\relax}>{\raggedright\arraybackslash}p{\dimexpr 0.105\linewidth-8pt\relax}>{\raggedright\arraybackslash}p{\dimexpr 0.120\linewidth-8pt\relax}>{\raggedright\arraybackslash}p{\dimexpr 0.089\linewidth-8pt\relax}>{\raggedright\arraybackslash}p{\dimexpr 0.089\linewidth-8pt\relax}@{}}
\toprule
\textbf{Target} & \textbf{$n$} & \textbf{Coverage} & \textbf{Width (cm)} & \textbf{Width R0, $n$ 0} & \textbf{In band} & \textbf{Narrower} \\
\midrule\endfirsthead
\toprule
\textbf{Target} & \textbf{$n$} & \textbf{Coverage} & \textbf{Width (cm)} & \textbf{Width R0, $n$ 0} & \textbf{In band} & \textbf{Narrower} \\
\midrule\endhead
\bottomrule\endlastfoot
Lena|\allowbreak{}x & 40 & 0.89(0.8940) & 67.88 & 74.62 & 1 & 1 \\
Lena|\allowbreak{}x & 160 & 0.90(0.8984) & 72.64 & 74.62 & 1 & 1 \\
Canada|\allowbreak{}x & 40 & 0.85(0.8475) & 92.39 & 104.84 & 0 & 0 \\
Canada|\allowbreak{}x & 160 & 0.84(0.8351) & 86.00 & 104.84 & 0 & 0 \\
Alaska|\allowbreak{}x & 40 & 0.89(0.8914) & 56.60 & 93.20 & 1 & 1 \\
Alaska|\allowbreak{}x & 160 & 0.90(0.8982) & 51.82 & 93.20 & 1 & 1 \\
\end{longtable}}

\clearpage
\phantomsection\label{tab:S13}
\noindent\textbf{Supplementary Table S13.} Additional registered analyses XA to XL.\par\smallskip
\noindent{\footnotesize XA was registered in commit 2678100 (4 October 2026, 14:22:39) and run locally without any new fit (sealed tables first opened 21:07:28). Both reproduction gates passed. Categories: positive or negative only if the main CI (family-cluster and block joint resampling) excludes zero in both weightings and the $|A| \geq 100$ subset (23 targets) does too; otherwise not established. Where the main and auxiliary categories differ, both are given and the weaker is reported. The registered family list had seven families; the North Atlantic target (point estimate only) was added as an eighth single-target family (implementation deviation, sealed/xa\_hyp.csv). All five hypotheses were robust to the scale variants and to the oracle definition of coefficient error (absolute log ratio of the all-label target coefficient to the source coefficient). Holm family: XA-1 to XA-3. Source: paper/claims/C2\_bias\_diagnosis/XA\_result\_2026-10-04.md sections 4 and 6 (data/processed/xbatch/XA\_c2\_gain\_decomposition/sealed/). Parts c to p: XB to XJ were registered in the same commit (2678100) and run on 5 October 2026; sealed tables were first opened between 03:18 and 05:21 KST (XC 05:20:39) after their reproduction gates passed (docs/EXPERIMENT\_PLAN\_FINAL\_BATCH\_2026-10-04.md section 8). XK and XL were added on 5 October 2026 (commit d721d24) before computation (docs/EXPERIMENT\_PLAN\_FINAL\_BATCH\_ADDENDUM\_XK\_XL\_2026-10-05.md, results). All were designed after earlier results were viewed; further flags are given with each part. $\Delta$ is the RMSE of the first method minus that of the second (cm; negative, the first has lower error), cell-weighted and block-equal with block-bootstrap 95\% CIs (10,000 resamples unless stated). Registered interpretation sentences are rendered in English in the notes. Method codes (internal identifiers allowed in the Supplementary Information): P0 source-coefficient Stefan model; P* year-matched Stefan model; P1 recalibrated Stefan model ($\kappa = 10$); P2 local least-squares Stefan model; P1* or P1@ed soil-property recalibrated Stefan model; R0 source anchor plus residual ML; R1 recalibrated anchor plus residual ML; R2 augmentation plus anchor plus residual; D0 direct ML; D1 physics pseudo-label augmentation; F1a, F1k, F1n physics-input ML; RM multiplicative residual; Re Stefan-CCI mean anchor plus residual; W selection rule (cross-validation within the target labels); S1 random cells, S2 block stratification, S3 covariate-cluster stratification, S4 covariate max-min distance, S5 source-extrapolation first, S6 sequential selection by prediction variance, S7 $\sqrt{\mathrm{TDD}}$ first. Target suffixes: $|x$ parent region excluded from the source, $|i$ rest of the parent kept, $|r$ within-region design. $n$ = number of target labels; all = all labels of half A. Verdicts: lower error, higher error, equivalent ($\pm$0.5~cm), undetermined; both 95\% block-bootstrap CIs (cell-weighted and block-equal) must agree. \dag~a note of the source row in Korean was omitted; see the cited README section.}\par
\par\medskip\noindent{\small\textit{a. XA gain decomposition (post hoc; all labels; coefficient error = absolute ten-label bias; 28 targets, 8 families; 10,000 resamples)}}\par\nopagebreak
{\fontsize{7.5}{9}\selectfont\setlength{\tabcolsep}{4pt}\setlength{\LTpre}{2pt}\setlength{\LTpost}{4pt}
\begin{longtable}{@{}>{\raggedright\arraybackslash}p{\dimexpr 0.108\linewidth-8pt\relax}>{\raggedright\arraybackslash}p{\dimexpr 0.132\linewidth-8pt\relax}>{\raggedright\arraybackslash}p{\dimexpr 0.077\linewidth-8pt\relax}>{\raggedright\arraybackslash}p{\dimexpr 0.135\linewidth-8pt\relax}>{\raggedright\arraybackslash}p{\dimexpr 0.135\linewidth-8pt\relax}>{\raggedright\arraybackslash}p{\dimexpr 0.125\linewidth-8pt\relax}>{\raggedright\arraybackslash}p{\dimexpr 0.077\linewidth-8pt\relax}>{\raggedright\arraybackslash}p{\dimexpr 0.058\linewidth-8pt\relax}>{\raggedright\arraybackslash}p{\dimexpr 0.154\linewidth-8pt\relax}@{}}
\toprule
\textbf{Hypothesis} & \textbf{Gain} & \textbf{$\rho$ cell / block} & \textbf{Main CI (family clusters) cell / block} & \textbf{Auxiliary CI (targets fixed) cell / block} & \textbf{$|A| \geq 100$ CI cell / block} & \textbf{$P_{\mathrm{one}}$ cell, block} & \textbf{Holm $P$} & \textbf{Category main / auxiliary / final} \\
\midrule\endfirsthead
\toprule
\textbf{Hypothesis} & \textbf{Gain} & \textbf{$\rho$ cell / block} & \textbf{Main CI (family clusters) cell / block} & \textbf{Auxiliary CI (targets fixed) cell / block} & \textbf{$|A| \geq 100$ CI cell / block} & \textbf{$P_{\mathrm{one}}$ cell, block} & \textbf{Holm $P$} & \textbf{Category main / auxiliary / final} \\
\midrule\endhead
\bottomrule\endlastfoot
XA-1 & recalibration share P0 $-$ P1 & 0.37 /\allowbreak{} −0.04 & [−0.11, 0.93] /\allowbreak{} [−0.41, 0.86] & [0.14, 0.48] /\allowbreak{} [−0.08, 0.18] & [−0.19, 0.30] /\allowbreak{} [−0.54, −0.17] & 0.06, 0.41 & 0.82 & not established / not established / not established \\
XA-2 & selection-rule share P1 $-$ W & 0.18 /\allowbreak{} 0.01 & [−0.40, 0.66] /\allowbreak{} [−0.60, 0.46] & [−0.08, 0.32] /\allowbreak{} [−0.20, 0.19] & [−0.07, 0.37] /\allowbreak{} [−0.18, 0.30] & 0.34, 0.57 & 0.82 & not established / not established / not established \\
XA-3 & ML share min(P1, P2) $-$ R1(0.25) & 0.28 /\allowbreak{} 0.27 & [−0.19, 0.80] /\allowbreak{} [−0.51, 0.59] & [0.18, 0.54] /\allowbreak{} [0.07, 0.40] & [0.08, 0.47] /\allowbreak{} [0.09, 0.50] & 0.09, 0.22 & 0.66 & not established / positive / not established (two categories, weaker reported) \\
XA-3a (aux.) & ML share P1 $-$ R1(0.25) & 0.39 /\allowbreak{} 0.32 & [−0.08, 0.82] /\allowbreak{} [−0.44, 0.63] & [0.27, 0.58] /\allowbreak{} [0.12, 0.44] & [0.08, 0.47] /\allowbreak{} [0.09, 0.50] & 0.04, 0.17 & 가족 밖 & not established / positive / not established (two categories, weaker reported) \\
XA-4 (aux.) & shrinkage remainder P1 $-$ P2 & −0.14 /\allowbreak{} −0.48 & [−0.64, 0.64] /\allowbreak{} [−0.81, 0.47] & [−0.46, −0.01] /\allowbreak{} [−0.70, −0.18] & [−0.63, −0.25] /\allowbreak{} [−0.82, −0.50] & 0.69, 0.88(음 방향 P(ρ\textasciicircum{}b ≥ 0) 최대 0.31) & 가족 밖 & not established / negative / not established (two categories, weaker reported) \\
descriptive & total P0 $-$ W & 0.38 /\allowbreak{} 0.08 & [−0.07, 0.91] /\allowbreak{} [−0.31, 0.80] & [0.17, 0.54] /\allowbreak{} [−0.02, 0.28] & [−0.13, 0.42] /\allowbreak{} [−0.40, 0.03] & 0.07, 0.26 &  & not established $\times$3 \\
\end{longtable}}
\par\medskip\noindent{\small\textit{b. XA-6 decomposition, four main regions, stratified mean (cm, positive = improvement; post hoc description)}}\par\nopagebreak
{\fontsize{7.5}{9}\selectfont\setlength{\tabcolsep}{4pt}\setlength{\LTpre}{2pt}\setlength{\LTpost}{4pt}
\begin{longtable}{@{}>{\raggedright\arraybackslash}p{\dimexpr 0.060\linewidth-8pt\relax}>{\raggedright\arraybackslash}p{\dimexpr 0.160\linewidth-8pt\relax}>{\raggedright\arraybackslash}p{\dimexpr 0.260\linewidth-8pt\relax}>{\raggedright\arraybackslash}p{\dimexpr 0.260\linewidth-8pt\relax}>{\raggedright\arraybackslash}p{\dimexpr 0.260\linewidth-8pt\relax}@{}}
\toprule
\textbf{$n$} & \textbf{Pool} & \textbf{Recalibration share (cell / block)} & \textbf{Shrinkage remainder} & \textbf{Selection-rule share} \\
\midrule\endfirsthead
\toprule
\textbf{$n$} & \textbf{Pool} & \textbf{Recalibration share (cell / block)} & \textbf{Shrinkage remainder} & \textbf{Selection-rule share} \\
\midrule\endhead
\bottomrule\endlastfoot
All & 4/4 & 2.23 [1.33, 3.05] / 2.13 [1.19, 3.01] & −0.02 [−0.87, 0.57] / −0.43 [−1.23, 0.33] & 1.14 [0.22, 2.22] / 1.43 [0.52, 2.32] \\
10 & 4/4 & 2.45 [1.86, 3.04] / 2.62 [1.88, 3.29] & −0.84 [−1.91, −0.03] / −1.36 [−2.33, −0.42] & 0.43 [−0.49, 1.36] / 0.51 [−0.45, 1.45] \\
40 & 2/4 (Lena Delta, Canada) & −1.99 [−2.71, −0.89] / −2.01 [−2.63, −1.37] & −0.76 [−0.90, −0.57] / −0.84 [−0.98, −0.70] & 1.43 [0.73, 1.90] / 1.55 [1.03, 2.12] \\
160 & 2/4 & −2.14 [−2.93, −0.95] / −2.12 [−2.79, −1.44] & −0.19 [−0.23, −0.14] / −0.21 [−0.24, −0.17] & 1.96 [1.18, 2.51] / 2.06 [1.47, 2.69] \\
\end{longtable}}
\par\medskip\noindent{\small\textit{c. XB label-weighted stacking of physics models and products (transfer: XB-1, XB-2; within region: XB-3, XB-4)}}\par\nopagebreak
{\fontsize{7.5}{9}\selectfont\setlength{\tabcolsep}{4pt}\setlength{\LTpre}{2pt}\setlength{\LTpost}{4pt}
\begin{longtable}{@{}>{\raggedright\arraybackslash}p{\dimexpr 0.165\linewidth-8pt\relax}>{\raggedright\arraybackslash}p{\dimexpr 0.060\linewidth-8pt\relax}>{\raggedright\arraybackslash}p{\dimexpr 0.165\linewidth-8pt\relax}>{\raggedright\arraybackslash}p{\dimexpr 0.136\linewidth-8pt\relax}>{\raggedright\arraybackslash}p{\dimexpr 0.136\linewidth-8pt\relax}>{\raggedright\arraybackslash}p{\dimexpr 0.084\linewidth-8pt\relax}>{\raggedright\arraybackslash}p{\dimexpr 0.078\linewidth-8pt\relax}>{\raggedright\arraybackslash}p{\dimexpr 0.175\linewidth-8pt\relax}@{}}
\toprule
\textbf{Hypothesis and contrast} & \textbf{$n$} & \textbf{Pool} & \textbf{$\Delta$ cell [95\% CI]} & \textbf{$\Delta$ block-equal [95\% CI]} & \textbf{Fallback} & \textbf{Holm $P$} & \textbf{Verdict} \\
\midrule\endfirsthead
\toprule
\textbf{Hypothesis and contrast} & \textbf{$n$} & \textbf{Pool} & \textbf{$\Delta$ cell [95\% CI]} & \textbf{$\Delta$ block-equal [95\% CI]} & \textbf{Fallback} & \textbf{Holm $P$} & \textbf{Verdict} \\
\midrule\endhead
\bottomrule\endlastfoot
XB-1 Stack $-$ P1 (transfer) & 10 & four main regions, partial (3/4; no stacking in E Russia, fallback 0.52) & −0.28 [−1.11, +1.16] & +1.42 [+0.46, +2.39] & 0.36 & 1.00 & undetermined \\
 & 40 & Lena Delta, Canada (2/2) & −1.86 [−2.44, −0.55] & +0.10 [−0.44, +0.64] & 0.14 & 1.00 & undetermined \\
 & 160 & Lena Delta, Canada (2/2) & −2.73 [−3.38, −1.35] & −0.46 [−1.14, +0.22] & 0.10 & 1.00 & undetermined \\
 & all & four main regions (4/4) & −1.54 [−2.28, −0.47] & −0.39 [−1.13, +0.35] & 0.20 & 1.00 & undetermined; hypothesis rejected \\
XB-2 StackR(0.25) $-$ R1(0.25) & 10 & partial (3/4) & −0.33 [−1.10, +0.99] & +1.16 [+0.33, +2.00] & 0.36 & 1.00 & undetermined \\
 & 40 & Lena Delta, Canada (2/2) & −1.65 [−2.15, −0.49] & −0.01 [−0.47, +0.44] & 0.14 & 1.00 & undetermined \\
 & 160 & Lena Delta, Canada (2/2) & −2.27 [−2.82, −1.05] & −0.35 [−0.94, +0.23] & 0.10 & 1.00 & undetermined \\
 & all & four main regions (4/4) & −1.38 [−2.00, −0.41] & −0.43 [−1.05, +0.20] & 0.20 & 1.00 & undetermined; hypothesis rejected \\
XB-3 StackR $-$ R1 (CV $\lambda$), within region & 200 & three targets (3/3) & −0.57 [−0.71, −0.20] & +0.42 [+0.24, +0.60] & 0.08 & 0.004 & undetermined (the two CIs exclude zero on opposite sides) \\
 & 500 & partial (2/3; Canada $|A| < 500$) & −0.55 [−0.69, −0.35] & −0.33 [−0.47, −0.19] & 0.05 & 0.0012 & lower error \\
 & 1,000 & partial (2/3) & −0.498 [−0.62, −0.34] & −0.29 [−0.44, −0.15] & 0.01 & 0.0012 & lower error (distinguishable, size below 0.5 cm) \\
 & all & three targets (3/3) & −0.99 [−1.19, −0.43] & +0.19 [−0.05, +0.43] & 0.01 & 1.00 & undetermined; hypothesis partially supported \\
XB-4 Stack $-$ P1, non-inferiority (margin 0.5 cm), Alaska within region & 200 & Alaska & −0.47 [−0.61, −0.28] & −0.73 [−0.92, −0.55] & 0.16 & $<$0.0001 & non-inferior (four-way lower error, size below 0.5 cm) \\
 & 500 & Alaska & −0.63 [−0.77, −0.44] & −0.86 [−1.06, −0.67] & 0.09 & $<$0.0001 & non-inferior (four-way lower error) \\
 & 1,000 & Alaska & −0.72 [−0.84, −0.51] & −0.85 [−1.05, −0.65] & 0.03 & $<$0.0001 & non-inferior (four-way lower error) \\
 & all & Alaska & −0.75 [−0.89, −0.54] & −0.90 [−1.12, −0.69] & 0.04 & $<$0.0001 & non-inferior (four-way lower error); hypothesis supported \\
\end{longtable}}
\par\noindent{\footnotesize Flags: designed after earlier results were viewed; registered before running; partly unblinded as in the registered hypothesis table. Holm families: XB-1 to XB-3 ($m$ = 12), XB-4 non-inferiority ($m$ = 4); verdicts use the uncorrected CIs. Fallback: share of draws in which stacking was replaced by P1. Registered sentences (rendered): XB-1, `No difference between stacking and the recalibrated Stefan model was established' (40 and 160 labels: Lena Delta and Canada, two regions); the statement that the gain from combining physics models and products is region-conditional (WF7, L29) is retained; at ten labels the stacking fallback rate was 0.36, so XB-1 at ten labels may in effect compare P1 with P1. XB-2, `No difference between stacked residuals and recalibrated-anchor residuals was established' (same qualifiers). XB-3, `Within label-rich regions, stacked residuals had 0.55 cm (cell-weighted) lower error than recalibrated-anchor residuals with 500 labels (partial, 2 of 3 regions)' and `... 0.498 cm lower error with 1,000 labels (partial, 2 of 3 regions); the difference is statistically distinguishable but below 0.5 cm in size'; `no difference was established with 200 and all labels (3 of 3 regions)'. XB-4, `Within Alaska, where the physics model is close to the error floor, the error increase of stacking relative to the recalibrated Stefan model was within 0.5 cm (non-inferior).' XB-3 was not retained when CCI v5 was added as a candidate (sensitivity edition, part c continued); the registered verdict remains the main-run one. Regional rows: no higher-error row. XB-3 Canada rows had opposite signs in the two weightings ($n$ 200 −0.95 [−1.29, +0.14] / +1.68 [+1.24, +2.13]; all −2.05 [−2.60, −0.47] / +1.05 [+0.41, +1.69]; undetermined). Weights are reported descriptively only (no causal sentence): mean weight with all labels, Alaska soil-thaw-index Stefan 0.82, Canada Kudryavtsev 0.57, Lena Delta year-matched Stefan 0.77, Tibetan Plateau affine-calibrated CCI v4 0.92; Spearman $\rho$ of the CCI weight with the P1 RMSE of the target 0.31 ($P$ 0.094, 31 stores). Auxiliary XB-5 (not counted): Stack0 $-$ P0 without labels, four main regions +0.51 [+0.09, +0.69] / +0.09 [−0.14, +0.31], undetermined. Source: data/processed/xbatch/XB\_multisource\_stacking/sealed/ (xb\_tests.csv, xb\_weights\_summary.csv, xb\_fallback.csv); reproduction gate passed (xb\_gate.csv).}\par\smallskip
\par\medskip\noindent{\small\textit{c (continued). XB sensitivity edition with CCI v5 added as an eighth candidate (registered verdicts remain those of the main run)}}\par\nopagebreak
{\fontsize{7.5}{9}\selectfont\setlength{\tabcolsep}{4pt}\setlength{\LTpre}{2pt}\setlength{\LTpost}{4pt}
\begin{longtable}{@{}>{\raggedright\arraybackslash}p{\dimexpr 0.200\linewidth-8pt\relax}>{\raggedright\arraybackslash}p{\dimexpr 0.070\linewidth-8pt\relax}>{\raggedright\arraybackslash}p{\dimexpr 0.140\linewidth-8pt\relax}>{\raggedright\arraybackslash}p{\dimexpr 0.270\linewidth-8pt\relax}>{\raggedright\arraybackslash}p{\dimexpr 0.200\linewidth-8pt\relax}>{\raggedright\arraybackslash}p{\dimexpr 0.120\linewidth-8pt\relax}@{}}
\toprule
\textbf{Hypothesis} & \textbf{$n$} & \textbf{Main-run verdict} & \textbf{CCI v5 edition: cell / block-equal [95\% CI]} & \textbf{Verdict, CCI v5 edition} & \textbf{Holm $P$ (CCI v5 edition)} \\
\midrule\endfirsthead
\toprule
\textbf{Hypothesis} & \textbf{$n$} & \textbf{Main-run verdict} & \textbf{CCI v5 edition: cell / block-equal [95\% CI]} & \textbf{Verdict, CCI v5 edition} & \textbf{Holm $P$ (CCI v5 edition)} \\
\midrule\endhead
\bottomrule\endlastfoot
XB-1 Stack $-$ P1 & 10 & undetermined & +0.27 [−0.78, +1.84] / +1.64 [+0.68, +2.67] & undetermined (fallback 0.37) & 1.00 \\
 & 40 & undetermined & −2.17 [−2.79, −0.41] / +0.70 [+0.06, +1.33] & undetermined (CIs exclude zero on opposite sides) & 0.33 \\
 & 160 & undetermined & −2.79 [−3.43, −1.06] / +0.38 [−0.43, +1.19] & undetermined & 1.00 \\
 & all & undetermined & −1.57 [−2.24, −0.44] / −0.04 [−0.82, +0.75] & undetermined & 1.00 \\
XB-1 hypothesis &  & rejected &  & rejected &  \\
XB-2 StackR(0.25) $-$ R1(0.25) & 10 & undetermined & +0.16 [−0.84, +1.64] / +1.39 [+0.55, +2.28] & undetermined & 1.00 \\
 & 40 & undetermined & −2.02 [−2.60, −0.36] / +0.65 [+0.07, +1.22] & undetermined (opposite sides) & 0.33 \\
 & 160 & undetermined & −2.35 [−2.92, −0.74] / +0.53 [−0.21, +1.28] & undetermined & 1.00 \\
 & all & undetermined & −1.42 [−1.98, −0.38] / −0.04 [−0.71, +0.64] & undetermined & 1.00 \\
XB-2 hypothesis &  & rejected &  & rejected &  \\
XB-3 StackR $-$ R1 (CV $\lambda$), within region & 200 & undetermined & −1.24 [−1.34, −0.67] / +0.32 [+0.11, +0.54] & undetermined (opposite sides) & 0.031 \\
 & 500 & lower error & −0.50 [−0.62, −0.19] / −0.06 [−0.24, +0.11] & undetermined (margin-dependent) & 1.00 \\
 & 1,000 & lower error & −0.36 [−0.54, −0.03] / +0.08 [−0.14, +0.30] & undetermined (margin-dependent) & 1.00 \\
 & all & undetermined & −1.44 [−1.59, −0.74] / +0.17 [−0.10, +0.43] & undetermined & 1.00 \\
XB-3 hypothesis &  & partially supported &  & rejected &  \\
XB-4 Stack $-$ P1, non-inferiority, Alaska & 200 to all & non-inferior at all four $n$ & all −0.70 [−0.82, −0.49] / −0.92 [−1.13, −0.73]; $n$ 200 −0.46 / −0.75 & non-inferior at all four $n$ (four-way lower error) & $<$0.0001 \\
XB-4 hypothesis &  & supported &  & supported &  \\
\end{longtable}}
\par\noindent{\footnotesize Sensitivity edition registered in plan 2.2 as a second-stage candidate: the year-matched CCI v5 map (affine-calibrated; missing cells take the P1 value) added as an eighth candidate, not as a replacement ($K$ = 8 in all 219 store-by-$n$ rows). Local run on 5 October 2026, 05:20 to 05:38 KST (211 units, no failure); the keys shared with the main run agreed within 6.0$\times$10$^{-14}$ cm; sealed tables first opened 05:43:31 KST. Flags as in the main run; designed after earlier results were viewed. The registered verdicts remain those of the main run (XB-1 and XB-2 rejected, XB-3 partially supported, XB-4 supported). In this edition XB-1, XB-2 and XB-4 kept their verdicts, and XB-3 was not retained: the lower error at 500 and 1,000 labels became undetermined (block-equal point estimates −0.06 and +0.08 cm). Regional rows: XB-1 Lena Delta at $n$ 10 changed from undetermined to higher error (+0.31 / +1.98 cm), XB-1 Canada at $n$ 160 and all labels from undetermined to lower error (−5.54 / −1.59 and −6.03 / −1.71 cm), XB-3 Alaska and Lena Delta within region at $n$ 1,000 from lower error to undetermined and Alaska with all labels from equivalent to undetermined; XB-2 and XB-4 regional rows did not change. Weights (descriptive, no causal statement): mean CCI v5 weight with all labels, Canada 0.21, Lena Delta 0.23, Tibetan Plateau 0.32 and Alaska 0.03 in transfer, Canada 0.46, Lena Delta 0.10 and Alaska 0.10 within regions; Spearman $\rho$ of the total CCI weight with the P1 RMSE 0.35 ($P$ 0.051, 31 transfer stores) and 0.50 ($P$ 0.67, three within-region targets). Source: data/processed/xbatch/XB\_multisource\_stacking/sealed/ (xb\_tests\_cci5y.csv, xb\_weights\_summary\_cci5y.csv, xb\_cci\_relation\_cci5y.csv); plan 8.12.}\par\smallskip
\par\medskip\noindent{\small\textit{d. XG existing ALT maps scored on the same blocks (registration deviation)}}\par\nopagebreak
{\fontsize{7.5}{9}\selectfont\setlength{\tabcolsep}{4pt}\setlength{\LTpre}{2pt}\setlength{\LTpost}{4pt}
\begin{longtable}{@{}>{\raggedright\arraybackslash}p{\dimexpr 0.139\linewidth-8pt\relax}>{\raggedright\arraybackslash}p{\dimexpr 0.044\linewidth-8pt\relax}>{\raggedright\arraybackslash}p{\dimexpr 0.110\linewidth-8pt\relax}>{\raggedright\arraybackslash}p{\dimexpr 0.189\linewidth-8pt\relax}>{\raggedright\arraybackslash}p{\dimexpr 0.189\linewidth-8pt\relax}>{\raggedright\arraybackslash}p{\dimexpr 0.090\linewidth-8pt\relax}>{\raggedright\arraybackslash}p{\dimexpr 0.129\linewidth-8pt\relax}>{\raggedright\arraybackslash}p{\dimexpr 0.110\linewidth-8pt\relax}@{}}
\toprule
\textbf{Hypothesis and contrast} & \textbf{$n$} & \textbf{Pool} & \textbf{5 km block mask: cell / block-equal} & \textbf{25 km block mask: cell / block-equal} & \textbf{Holm $P$ (5 / 25 km)} & \textbf{Verdict} & \textbf{Product values imputed (5 / 25 km)} \\
\midrule\endfirsthead
\toprule
\textbf{Hypothesis and contrast} & \textbf{$n$} & \textbf{Pool} & \textbf{5 km block mask: cell / block-equal} & \textbf{25 km block mask: cell / block-equal} & \textbf{Holm $P$ (5 / 25 km)} & \textbf{Verdict} & \textbf{Product values imputed (5 / 25 km)} \\
\midrule\endhead
\bottomrule\endlastfoot
XG-1w B:wei $-$ P0 & 0 & Lena Delta, Canada (2/2) & +1.44 [−0.02, +4.30] / +4.09 [+2.10, +6.12] & +0.97 [−0.41, +3.87] / +4.01 [+1.99, +6.12] & 0.156 / 0.383 & undetermined & 47.4\% / 26.5\% \\
XG-2w R1(0.25) $-$ P1@wei & 10 & Lena Delta, Canada & +1.34 [+0.00, +1.83] / −0.49 [−1.36, +0.37] & +1.25 [−0.07, +1.75] / −0.52 [−1.43, +0.35] & 0.542 / 0.504 & undetermined & 47.4\% / 26.5\% \\
XG-3w R1(0.25) $-$ P1@wei & all & Lena Delta, Canada & +0.87 [−0.41, +1.37] / −0.39 [−1.11, +0.33] & +0.61 [−0.61, +1.19] / −0.43 [−1.17, +0.30] & 0.542 / 0.504 & undetermined & 47.4\% / 26.5\% \\
XG-1c B:cci5y $-$ P0 & 0 & four main regions (4/4) & +23.79 [+15.34, +30.16] / +19.17 [+14.83, +23.72] & same & 0.0006 / 0.0006 & higher error & below 0.1\% \\
XG-2c R1(0.25) $-$ P1@cci5y & 10 & four main regions & −8.37 [−10.97, −7.39] / −12.87 [−14.69, −11.10] & same & 0.0006 / 0.0006 & lower error (HK region-level CI [−16.02, −0.50]) & below 0.1\% \\
XG-3c R1(0.25) $-$ P1@cci5y & all & four main regions & −7.51 [−10.31, −6.27] / −11.75 [−13.63, −9.88] & same & 0.0006 / 0.0006 & lower error (HK CI [−17.05, +2.43]) & below 0.1\% \\
\end{longtable}}
\par\noindent{\footnotesize Flags: registration deviation (WRAPUP 10); designed after earlier results were viewed; XG-1w to XG-3w blind; XG-1c partly unblinded (raw CCI v4 minus P0 viewed in LGX L29); XG-2c and XG-3c partly unblinded (CCI v4 scale recalibration viewed); all confirmatory rows carry a few-blocks flag after masking (fewer than eight scoring blocks per split). The two masks are co-primary and gave the same branch, so no mask-dependence flag applies. CCI v5 has no training sites, so its values are the same under both masks. B:product, the product map without labels; P1@product, the product recalibrated with the same target labels. ALT definitions differ: CCI is the annual maximum thaw depth of a model pixel, Wei v2 is a machine-learning map trained on CALM labels (including CALM sites inside the target regions), and our labels mix GPR and probe measurements. Source-coefficient RMSE before and after masking (cell-weighted): Wei pool 24.93 cm, 25.28 cm (5 km) and 28.34 cm (25 km); Lena Delta 21.69, 21.74 and 27.86 cm; Canada 28.16 and 28.82 cm (both masks); CCI pool 30.64 cm (unchanged). Regional XG-1c rows (5 km): Lena Delta +1.19 [+0.13, +5.38] / +9.34 [+6.57, +12.22] (depends on split independence), Canada +82.38 [+49.01, +99.54] / +53.00 [+41.74, +64.55], E Russia +5.84 [+0.08, +18.20] / +10.44 [+4.25, +17.24] (depends on split independence), all higher error; W Russia +5.78 [−5.86, +17.69] / +3.89 [−6.71, +15.83], undetermined. XG-2c lower error in all four regions; XG-3c undetermined in Canada (−1.47 [−7.55, +0.96] / −10.94 [−13.71, −8.20]) and lower error elsewhere. Wei rows: XG-3w lower error in the Lena Delta only (5 km −1.08 [−1.76, −0.59] / −1.49 [−2.24, −0.77]); the Lena comparison depends strongly on imputed product values (excluding blocks with imputed values removes all Lena blocks, so the Wei pool row cannot be decided; CCI v5 verdicts are unchanged). Cell-level 5 km sensitivity (refitted): XG-1w +2.09 [+0.41, +5.23] / +4.45 [+2.50, +6.46], higher error ($P$ 0.016); XG-2w and XG-3w undetermined; sensitivities do not change the confirmatory verdicts. Registered sentences (rendered): XG-1c, `Without labels, the existing CCI v5 map had 23.79 cm higher error than the source-coefficient Stefan model (4 of 4 regions, after 5 and 25 km block masks of training sites); error was higher in the Lena Delta ($\Delta$ +1.19 cm, CI [0.13, 5.38]), Canada (+82.38 cm, [49.01, 99.54]) and E Russia (+5.84 cm, [0.08, 18.20]).' XG-2c, `With the same ten target labels, the recalibrated-anchor residual had 8.37 cm lower error than CCI v5 recalibrated with those labels' (the HK region-level CI excludes zero, so a region-general sentence is allowed). XG-3c, `... with all target labels, 7.51 cm lower error' (the HK CI includes zero; no region-general sentence). XG-1w, `Without labels, no difference was established between the existing Wei v2 map and the source-coefficient Stefan model (2 of 2 regions, after 5 and 25 km block masks; product values imputed for 47.4\%). Wei v2 also used CALM sites inside the target regions for training, whereas the source-coefficient Stefan model used only labels outside the target region and its 100 km buffer.' XG-2w and XG-3w, `No difference was established between the recalibrated-anchor residual with the same 10 (all) target labels and Wei v2 recalibrated with those labels (product values imputed for 47.4\%).' No sentence states that physics-anchored ML is worse than a product that shares training data (WRAPUP 10). Source: data/processed/xbatch/XG\_product\_comparison/sealed/ (xg\_tests.csv, xg\_confirm\_rows.csv, xg\_tests\_cell.csv, xg\_fill\_sensitivity.csv); gate passed 5 October 2026, 03:42:02.}\par\smallskip
\par\medskip\noindent{\small\textit{e. XH validation ladder: RMSE by scoring scheme (cm, cell-weighted [95\% CI] / block-equal; descriptive)}}\par\nopagebreak
{\fontsize{7.5}{9}\selectfont\setlength{\tabcolsep}{4pt}\setlength{\LTpre}{2pt}\setlength{\LTpost}{4pt}
\begin{longtable}{@{}>{\raggedright\arraybackslash}p{\dimexpr 0.090\linewidth-8pt\relax}>{\raggedright\arraybackslash}p{\dimexpr 0.170\linewidth-8pt\relax}>{\raggedright\arraybackslash}p{\dimexpr 0.170\linewidth-8pt\relax}>{\raggedright\arraybackslash}p{\dimexpr 0.120\linewidth-8pt\relax}>{\raggedright\arraybackslash}p{\dimexpr 0.120\linewidth-8pt\relax}>{\raggedright\arraybackslash}p{\dimexpr 0.170\linewidth-8pt\relax}>{\raggedright\arraybackslash}p{\dimexpr 0.160\linewidth-8pt\relax}@{}}
\toprule
\textbf{Region} & \textbf{Item} & \textbf{W1R random cell} & \textbf{W1S 0.05° site} & \textbf{W1B 0.5° block} & \textbf{W1K kNNDM} & \textbf{V-G region holdout} \\
\midrule\endfirsthead
\toprule
\textbf{Region} & \textbf{Item} & \textbf{W1R random cell} & \textbf{W1S 0.05° site} & \textbf{W1B 0.5° block} & \textbf{W1K kNNDM} & \textbf{V-G region holdout} \\
\midrule\endhead
\bottomrule\endlastfoot
Alaska & median distance to nearest training cell (km) & 0.00 & 0.88 & 26.06 & 58.73 &  \\
 & D0 & 11.54 [9.61, 14.99] / 11.94 & 13.28 / 13.65 & 16.21 / 14.02 & 17.30 [14.15, 19.68] / 14.12 & 35.56 / 31.91 \\
 & P1 (V-G: P0) & 14.24 [12.08, 18.60] / 15.37 & 14.30 / 15.38 & 14.41 / 15.43 & 14.53 [12.43, 18.84] / 15.70 & 14.68 / 14.31 \\
 & R1 (V-G: R0) & 13.27 / 14.39 & 13.82 / 14.88 & 13.97 / 14.97 & 14.07 / 15.41 & 15.00 / 16.76 \\
Lena Delta & median distance (km) & 0.01 & 1.02 & 9.55 & 26.42 &  \\
 & D0 & 13.94 [10.86, 21.12] / 21.44 & 18.28 / 23.25 & 19.77 / 24.51 & 22.10 [17.26, 37.01] / 26.42 & 25.31 / 24.60 \\
 & P1 (V-G: P0) & 20.89 [13.42, 37.22] / 25.92 & 21.11 / 26.07 & 21.18 / 26.23 & 21.40 [13.67, 38.01] / 26.37 & 21.64 / 24.95 \\
 & R1 (V-G: R0) & 18.35 / 24.08 & 19.54 / 24.55 & 20.17 / 24.94 & 21.03 / 25.48 & 21.93 / 24.67 \\
Canada & median distance (km) & 0.01 & 3.86 & 25.33 & 236.53 &  \\
 & D0 & 18.96 [14.01, 24.14] / 20.90 & 23.90 / 20.08 & 24.48 / 21.46 & 33.16 [22.52, 39.09] / 28.48 & 25.24 / 20.22 \\
 & P1 (V-G: P0) & 26.50 [18.43, 33.13] / 20.99 & 27.98 / 21.31 & 29.53 / 21.22 & 30.78 [20.19, 39.06] / 23.40 & 26.48 / 20.96 \\
 & R1 (V-G: R0) & 23.59 / 19.80 & 26.27 / 19.83 & 27.83 / 20.07 & 30.94 / 22.57 & 25.60 / 20.05 \\
\end{longtable}}
\par\medskip\noindent{\small\textit{f. XH optimism difference $\Delta\Delta$ from random-cell to kNNDM splits (cm; descriptive)}}\par\nopagebreak
{\fontsize{7.5}{9}\selectfont\setlength{\tabcolsep}{4pt}\setlength{\LTpre}{2pt}\setlength{\LTpost}{4pt}
\begin{longtable}{@{}>{\raggedright\arraybackslash}p{\dimexpr 0.120\linewidth-8pt\relax}>{\raggedright\arraybackslash}p{\dimexpr 0.130\linewidth-8pt\relax}>{\raggedright\arraybackslash}p{\dimexpr 0.190\linewidth-8pt\relax}>{\raggedright\arraybackslash}p{\dimexpr 0.190\linewidth-8pt\relax}>{\raggedright\arraybackslash}p{\dimexpr 0.190\linewidth-8pt\relax}>{\raggedright\arraybackslash}p{\dimexpr 0.180\linewidth-8pt\relax}@{}}
\toprule
\textbf{Item} & \textbf{Learner} & \textbf{Alaska} & \textbf{Lena Delta} & \textbf{Canada} & \textbf{Three-region stratified mean} \\
\midrule\endfirsthead
\toprule
\textbf{Item} & \textbf{Learner} & \textbf{Alaska} & \textbf{Lena Delta} & \textbf{Canada} & \textbf{Three-region stratified mean} \\
\midrule\endhead
\bottomrule\endlastfoot
XH-1 ($M$ = P1) & CatBoost (main) & +5.46 [+2.13, +8.53] / +1.85 [+0.75, +2.97] & +7.66 [+4.20, +15.86] / +4.52 [+0.59, +8.68] & +9.92 [+3.35, +17.28] / +5.17 [+0.94, +9.26] & +7.68 [+4.51, +11.25] / +3.85 [+1.84, +5.83] \\
 & CatBoost (second setting) & +5.27 / +2.70 & +7.61 / +4.76 & +8.77 / +5.09 & +7.22 [+4.84, +10.42] / +4.18 [+2.28, +6.11] \\
 & random forest & +6.95 / +3.26 & +9.99 / +5.78 & +6.41 / +1.30 & +7.78 [+5.37, +10.26] / +3.44 [+1.97, +5.06] \\
 & CatBoost (main), Canada ±2° &  &  & +7.79 [+1.82, +14.86] / +5.03 [+1.08, +8.92] & +6.97 [+3.99, +10.49] / +3.80 [+1.90, +5.76] \\
XH-2 ($M$ = R1, $\lambda$ 0.25) & CatBoost (main) & +4.96 [+1.51, +8.44] / +1.15 [+0.29, +2.03] & +5.49 [+2.99, +11.04] / +3.58 [+0.62, +6.58] & +6.85 [+2.42, +11.97] / +4.81 [+1.04, +8.58] & +5.77 [+3.25, +8.43] / +3.18 [+1.54, +4.79] \\
 & CatBoost (second setting) & +4.15 / +1.70 & +5.20 / +3.32 & +5.83 / +4.78 & +5.06 [+3.35, +7.32] / +3.27 [+1.74, +4.80] \\
 & random forest & +5.59 / +2.04 & +7.22 / +4.21 & +3.99 / +0.82 & +5.60 [+3.67, +7.45] / +2.36 [+1.24, +3.55] \\
\end{longtable}}
\par\noindent{\footnotesize Flags: designed after earlier results were viewed; descriptive (no verdict words, no multiplicity); Alaska rows are an unblinded replication, Lena Delta and Canada rows blind. Values: RMSE in cm, cell-weighted [95\% CI] / block-equal. D0 direct CatBoost; P1 Stefan coefficient fitted on the training fold; R1 residual ML ($\lambda$ 0.25). V-G rows reuse the region-holdout results of an earlier run (source-coefficient Stefan P0 and residual R0, no target labels); the sealed reuse table was empty, so these values were read from data/processed/lgx/ladder/lgv\_metrics.csv without new fits. $\Delta\Delta$ = [RMSE$_{K}$(D0) $-$ RMSE$_{R}$(D0)] $-$ [RMSE$_{K}$($M$) $-$ RMSE$_{R}$($M$)], in cm, cell-weighted [95\% CI] / block-equal [95\% CI] from 10,000 block resamples of cell-level squared errors. Canada prediction-domain sensitivity ($\pm$2°, W1K): D0 29.99 / 26.97, P1 29.74 / 22.03, R1 29.31 / 21.06 cm. Alaska contest-continuity rows (ridge, W1R to W1K, cell-weighted): direct ridge 12.65 to 14.18 cm, residual ($\lambda$ 0.75) 12.77 to 13.72 cm. Registered descriptive sentence (rendered): `From random-cell to kNNDM splits, the RMSE of direct ML increased by 5.76 cm (Alaska), 8.17 cm (Lena Delta) and 14.20 cm (Canada), and that of the recalibrated Stefan model by 0.29, 0.51 and 4.28 cm (cell-weighted; block-equal, direct ML 2.18, 4.97 and 7.58 cm, recalibrated Stefan 0.32, 0.45 and 2.41 cm; the difference of the two increases, $\Delta\Delta$, was +5.46 [+2.13, +8.53], +7.66 [+4.20, +15.86] and +9.92 [+3.35, +17.28] cm, three-region stratified mean +7.68 [+4.51, +11.25] cm, block-equal three-region +3.85 [+1.84, +5.83] cm).' With random-cell splits, RMSE was 11.54, 13.94 and 18.96 cm for direct ML and 14.24, 20.89 and 26.50 cm for the recalibrated Stefan model (cell-weighted). Accompanying sentence: random validation has been argued to estimate map accuracy validly for probability samples (Wadoux et al., 2021); the labels of this study are not a probability sample. Source: data/processed/xbatch/XH\_validation\_ladder/sealed/ (xh\_ladder.csv, xh\_dd.csv, xh\_stage\_contrasts.csv); gates passed (xh\_gate.csv).}\par\smallskip
\par\medskip\noindent{\small\textit{g. XI warm-block spatial proxy for climate extrapolation, Canada licence-verified expanded edition (not a true extrapolation test)}}\par\nopagebreak
{\fontsize{7.5}{9}\selectfont\setlength{\tabcolsep}{4pt}\setlength{\LTpre}{2pt}\setlength{\LTpost}{4pt}
\begin{longtable}{@{}>{\raggedright\arraybackslash}p{\dimexpr 0.145\linewidth-8pt\relax}>{\raggedright\arraybackslash}p{\dimexpr 0.043\linewidth-8pt\relax}>{\raggedright\arraybackslash}p{\dimexpr 0.164\linewidth-8pt\relax}>{\raggedright\arraybackslash}p{\dimexpr 0.117\linewidth-8pt\relax}>{\raggedright\arraybackslash}p{\dimexpr 0.164\linewidth-8pt\relax}>{\raggedright\arraybackslash}p{\dimexpr 0.125\linewidth-8pt\relax}>{\raggedright\arraybackslash}p{\dimexpr 0.145\linewidth-8pt\relax}>{\raggedright\arraybackslash}p{\dimexpr 0.097\linewidth-8pt\relax}@{}}
\toprule
\textbf{Hypothesis and contrast} & \textbf{$n$} & \textbf{warm\_trim (main): cell / block-equal} & \textbf{Verdict, main} & \textbf{warm (basic): cell / block-equal} & \textbf{Verdict, basic} & \textbf{Reported verdict (weaker)} & \textbf{Holm $P$} \\
\midrule\endfirsthead
\toprule
\textbf{Hypothesis and contrast} & \textbf{$n$} & \textbf{warm\_trim (main): cell / block-equal} & \textbf{Verdict, main} & \textbf{warm (basic): cell / block-equal} & \textbf{Verdict, basic} & \textbf{Reported verdict (weaker)} & \textbf{Holm $P$} \\
\midrule\endhead
\bottomrule\endlastfoot
XI-a EP(D0 $-$ P1), higher-error criterion & 100 & −3.12 [−3.95, −1.02] / +0.33 [−1.52, +2.18] & undetermined & −4.68 [−5.90, −2.50] / −2.29 [−4.55, +0.02] & undetermined & undetermined & 1.00 \\
 & all & −6.17 [−7.00, −3.02] / −1.12 [−3.36, +1.22] & undetermined & −6.49 [−7.66, −3.77] / −3.07 [−5.59, −0.51] & lower error (depends on split independence) & undetermined (both reported); hypothesis rejected & 1.00 \\
XI-b EP(R1 $-$ D0), lower-error criterion & 100 & +0.79 [−0.23, +1.41] / −0.86 [−2.27, +0.54] & undetermined & +2.37 [+1.20, +3.34] / +1.30 [−0.43, +3.04] & undetermined & undetermined & 0.90 \\
 & all & +3.99 [+1.90, +4.61] / +0.49 [−1.47, +2.38] & undetermined & +3.58 [+1.95, +4.57] / +2.30 [+0.31, +4.31] & higher error (depends on split independence) & undetermined (both reported); hypothesis rejected & 1.00 \\
XI-c R1 (CV $\lambda$) $-$ P1 in W, non-inferiority & 100 & −2.44 [−2.83, −1.78] / −2.35 [−2.92, −1.79] & lower error, non-inferior & −2.91 [−3.31, −2.34] / −2.94 [−3.59, −2.23] & lower error, non-inferior & lower error & $<$0.0001 \\
 & all & −2.41 [−3.09, −2.03] / −2.88 [−3.40, −2.36] & lower error, non-inferior & −2.88 [−3.47, −2.27] / −3.20 [−3.98, −2.38] & lower error, non-inferior & lower error; hypothesis supported & $<$0.0001 \\
\end{longtable}}
\par\noindent{\footnotesize Flags: designed after earlier results were viewed; Canada (licence-verified expanded edition) rows blind; Alaska rows an unblinded replication; Holm $m$ = 6 (auxiliary column). This is a spatial proxy with warm blocks of the same period, not a true extrapolation test. EP(A $-$ B) = $\Delta_W$(A $-$ B) $-$ $\Delta_I$(A $-$ B), the contrast in the extrapolation blocks W minus that in the interpolation blocks I (cm). Extrapolation width (median $\sqrt{\mathrm{TDD}}$ of W scoring cells minus the maximum of A, split mean): warm\_trim 0.97 ($^{\circ}$C~d)$^{1/2}$ (all W cells outside A); basic warm $-$0.002 ($^{\circ}$C~d)$^{1/2}$ (58.4\% outside A); cold 1.74; Alaska warm 0.63. RMSE levels (cell-weighted, $n$ 100 and all): in W, P1 44.38 and 45.10 cm, R1 41.95 and 42.66 cm, D0 41.99 and 40.32 cm; in I, P1 21.93 and 21.98 cm, R1 21.94 and 21.81 cm, D0 23.00 and 23.59 cm. Alaska warm and Canada cold rows were undetermined for all three contrasts. Registered sentences (rendered): XI-a, `In a spatial proxy using warm blocks of the same period (extrapolation width 0.97 ($^{\circ}$C~d)$^{1/2}$ in $\sqrt{\mathrm{TDD}}$), no difference was established in the extrapolation loss of direct ML relative to the recalibrated Stefan model (100 labels, all labels)'; XI-b, `In the same spatial proxy, no difference was established between the extrapolation losses of the recalibrated-anchor residual and of direct ML (100 labels, all labels)'; with all labels the two editions disagree (basic edition lower error for XI-a and higher error for XI-b), both are reported and the weaker verdict is used; XI-c, `In the extrapolation blocks of the same proxy, the recalibrated-anchor residual (cross-validated $\lambda$) had 2.44 cm (100 labels) and 2.41 cm (all labels) lower error than the recalibrated Stefan model (cell-weighted) and met the non-inferiority criterion (upper bounds of both CIs below 0.5 cm).' Source: data/processed/xbatch/XI\_climate\_extrapolation\_retest/sealed/ (xi\_hyp.csv, xi\_tests.csv, xi\_holm.csv, xi\_curve.csv); gate passed (xi\_gate.csv).}\par\smallskip
\par\medskip\noindent{\small\textit{h. XE public inputs at 10 to 500 m within regions, stage 1 (xh0; auxiliary contrast XE-c)}}\par\nopagebreak
{\fontsize{7.5}{9}\selectfont\setlength{\tabcolsep}{4pt}\setlength{\LTpre}{2pt}\setlength{\LTpost}{4pt}
\begin{longtable}{@{}>{\raggedright\arraybackslash}p{\dimexpr 0.158\linewidth-8pt\relax}>{\raggedright\arraybackslash}p{\dimexpr 0.060\linewidth-8pt\relax}>{\raggedright\arraybackslash}p{\dimexpr 0.079\linewidth-8pt\relax}>{\raggedright\arraybackslash}p{\dimexpr 0.217\linewidth-8pt\relax}>{\raggedright\arraybackslash}p{\dimexpr 0.170\linewidth-8pt\relax}>{\raggedright\arraybackslash}p{\dimexpr 0.316\linewidth-8pt\relax}@{}}
\toprule
\textbf{Contrast (within-grid RMSE, R1 with CV $\lambda$)} & \textbf{$n$} & \textbf{Pool} & \textbf{Three-target mean: cell / block-equal} & \textbf{Four-way verdict} & \textbf{Target rows} \\
\midrule\endfirsthead
\toprule
\textbf{Contrast (within-grid RMSE, R1 with CV $\lambda$)} & \textbf{$n$} & \textbf{Pool} & \textbf{Three-target mean: cell / block-equal} & \textbf{Four-way verdict} & \textbf{Target rows} \\
\midrule\endhead
\bottomrule\endlastfoot
XE-c R1(xh0) $-$ R1(x25) & 200 & 3/3 & −0.09 [−0.10, +0.06] / +0.06 [+0.01, +0.12] & equivalent & Alaska, Canada equivalent; Lena Delta equivalent (margin-dependent, depends on split independence, few blocks) \\
 & 500 & partial (2/3) & −0.16 [−0.18, +0.05] / +0.05 [−0.05, +0.15] & equivalent & Alaska equivalent; Lena Delta equivalent (same flags) \\
 & 1,000 & partial (2/3) & −0.17 [−0.22, +0.05] / −0.01 [−0.11, +0.11] & equivalent & Alaska equivalent; Lena Delta equivalent (same flags) \\
 & all & 3/3 & −0.12 [−0.16, +0.02] / −0.08 [−0.18, +0.02] & equivalent & Alaska, Canada equivalent; Lena Delta undetermined (−0.34 [−0.45, +0.05] / −0.22 [−0.50, +0.06]) \\
XE-c R1(xh0) $-$ P1 & 200 & 3/3 & −0.02 [−0.06, +0.21] / +0.23 [+0.15, +0.30] & equivalent (margin-dependent) & Alaska equivalent; Lena Delta undetermined; Canada higher error (+0.33 [+0.25, +0.45] / +0.30 [+0.22, +0.39], size below 0.5 cm) \\
 & 500 & partial (2/3) & −0.24 [−0.30, +0.15] / +0.27 [+0.12, +0.41] & equivalent (margin-dependent; depends on split independence) & Alaska equivalent; Lena Delta undetermined \\
 & 1,000 & partial (2/3) & −0.19 [−0.27, +0.33] / +0.36 [+0.18, +0.54] & undetermined (margin-dependent) & Alaska equivalent; Lena Delta undetermined \\
 & all & 3/3 & −0.03 [−0.08, +0.35] / +0.29 [+0.16, +0.43] & equivalent (margin-dependent; depends on split independence) & Alaska equivalent; Lena Delta undetermined; Canada higher error (+0.25 [+0.16, +0.38] / +0.21 [+0.12, +0.31], size below 0.5 cm) \\
XE-a, XE-b, XE-a0 &  &  & \textcolor{blue}{[PENDING: XE-a, XE-b, XE-a0, satellite inputs (stage 2) and the xt2 table; registered deadline 8 October 2026, 14:22 KST]} &  &  \\
\end{longtable}}
\par\noindent{\footnotesize Flags: designed after earlier results were viewed; xh0 partly unblinded (the ten H0 inputs were inputs of H19, and the within-grid R1 $-$ P1 contrast of WF9-a had been viewed). All rows are the auxiliary contrast XE-c (four-way verdict, not counted as confirmatory, outside the XE Holm family). xh0 = the 25 covariates plus ten H0 inputs (topographic wetness index and its standard deviation, flow accumulation, two curvatures, relief, surface-water occurrence and distance, tree cover). Within-grid share explained, 1 $-$ SSE$_w$(M)/SSE$_w$(P1), all labels (descriptive, no CI): Alaska 0.07\% (x25) and 0.69\% (xh0), Lena Delta $-$0.06\% and 3.88\%, Canada $-$2.18\% and $-$2.15\%. Total RMSE, x25 to xh0: Alaska 14.06 to 14.00 cm, Lena Delta 20.89 to 20.31 cm, Canada 29.72 to 26.75 cm (between-grid component 23.62 to 19.75 cm), P1 14.53, 20.62 and 29.10 cm; these total and between-grid contrasts are not registered and not judged. No pre-fixed sentence applies to XE-c; only facts are reported. The reproduction gate was rescored at the registered local-to-cloud tolerance (gate level corrected, not a design change). Source: data/processed/xbatch/XE\_hires\_covariates/sealed/ (xe\_r1b\_xh0\_tests.csv, xe\_r1b\_xh0\_decomp.csv).}\par\smallskip
\par\medskip\noindent{\small\textit{h (continued). XE-e diagnostic: share of the within-grid residual of P1 explained by on-site soil moisture (block cross-validation)}}\par\nopagebreak
{\fontsize{7.5}{9}\selectfont\setlength{\tabcolsep}{4pt}\setlength{\LTpre}{2pt}\setlength{\LTpost}{4pt}
\begin{longtable}{@{}>{\raggedright\arraybackslash}p{\dimexpr 0.100\linewidth-8pt\relax}>{\raggedright\arraybackslash}p{\dimexpr 0.220\linewidth-8pt\relax}>{\raggedright\arraybackslash}p{\dimexpr 0.190\linewidth-8pt\relax}>{\raggedright\arraybackslash}p{\dimexpr 0.190\linewidth-8pt\relax}>{\raggedright\arraybackslash}p{\dimexpr 0.150\linewidth-8pt\relax}>{\raggedright\arraybackslash}p{\dimexpr 0.150\linewidth-8pt\relax}@{}}
\toprule
\textbf{Region} & \textbf{Model} & \textbf{Non-GPR VWC, registered subset (main)} & \textbf{Non-GPR VWC, finite-predictor cells} & \textbf{GPR VWC, registered subset} & \textbf{GPR VWC, finite cells} \\
\midrule\endfirsthead
\toprule
\textbf{Region} & \textbf{Model} & \textbf{Non-GPR VWC, registered subset (main)} & \textbf{Non-GPR VWC, finite-predictor cells} & \textbf{GPR VWC, registered subset} & \textbf{GPR VWC, finite cells} \\
\midrule\endhead
\bottomrule\endlastfoot
Alaska & linear regression & −3.54\% & −0.79\% & −0.54\% & −5.56\% \\
 & CatBoost (mean of seeds 0 and 1) & −4.89\% & −2.40\% & +0.82\% & −4.16\% \\
Canada & linear regression & −4.58\% & +2.55\% & n.a. & n.a. \\
 & CatBoost & −4.48\% & −6.59\% & n.a. & n.a. \\
\end{longtable}}
\par\noindent{\footnotesize XE-e, diagnostic (exploratory): share of the within-grid residual of the recalibrated Stefan model P1 explained by on-site volumetric water content (shallow and deep layers) under five-fold block cross-validation. Flags: designed after earlier results were viewed; blind (the relation of VWC with P1 residuals had not been computed); the registered rule says XE-e changes no verdict. Site data: ABoVE soil thaw depth and moisture validation, version 2 (15 m radius). Registered subsets: Alaska 2,468 cells in 31 blocks, Canada 602 cells in 15 blocks; shares computed on cells in grid groups of at least two cells (2,466 and 599); non-GPR VWC finite in 1,329 and 595 cells. GPR VWC shares error with ALT and is a sensitivity only; Canada has no GPR VWC (n.a.). Reference: a 0.5 cm reduction of within-grid RMSE corresponds to explained shares of 7.9\% (Alaska) and 5.0\% (Canada). Description: under block cross-validation on-site VWC did not explain the within-grid residual of P1 (main shares −4.89\% to −3.54\%, sensitivities −6.59\% to +2.55\%). Run locally on 5 October 2026 (05:41, 8 s); sealed table first opened 05:42:35 KST. Source: data/processed/xbatch/XE\_hires\_covariates/sealed/xe\_e\_xe\_e.csv and xe\_e\_xe\_e\_meta.json; plan 8.6.}\par\smallskip
\par\medskip\noindent{\small\textit{i. XD placement algorithm (XD-1 to XD-3; R1 $\lambda$ 0.25, cm; exploratory)}}\par\nopagebreak
{\fontsize{7.5}{9}\selectfont\setlength{\tabcolsep}{4pt}\setlength{\LTpre}{2pt}\setlength{\LTpost}{4pt}
\begin{longtable}{@{}>{\raggedright\arraybackslash}p{\dimexpr 0.149\linewidth-8pt\relax}>{\raggedright\arraybackslash}p{\dimexpr 0.040\linewidth-8pt\relax}>{\raggedright\arraybackslash}p{\dimexpr 0.109\linewidth-8pt\relax}>{\raggedright\arraybackslash}p{\dimexpr 0.129\linewidth-8pt\relax}>{\raggedright\arraybackslash}p{\dimexpr 0.129\linewidth-8pt\relax}>{\raggedright\arraybackslash}p{\dimexpr 0.219\linewidth-8pt\relax}>{\raggedright\arraybackslash}p{\dimexpr 0.145\linewidth-8pt\relax}>{\raggedright\arraybackslash}p{\dimexpr 0.080\linewidth-8pt\relax}@{}}
\toprule
\textbf{Hypothesis and contrast} & \textbf{$n$} & \textbf{Pool or target} & \textbf{$\Delta$ cell [95\% CI]} & \textbf{$\Delta$ block-equal [95\% CI]} & \textbf{Four-way verdict} & \textbf{Non-inferiority ($P\times2$)} & \textbf{Holm $P$} \\
\midrule\endfirsthead
\toprule
\textbf{Hypothesis and contrast} & \textbf{$n$} & \textbf{Pool or target} & \textbf{$\Delta$ cell [95\% CI]} & \textbf{$\Delta$ block-equal [95\% CI]} & \textbf{Four-way verdict} & \textbf{Non-inferiority ($P\times2$)} & \textbf{Holm $P$} \\
\midrule\endhead
\bottomrule\endlastfoot
XD-1 S8a $-$ S1 & 10 & PE1 (5/5) & −0.33 [−1.20, +0.66] & −0.16 [−0.91, +0.57] & undetermined; hypothesis rejected &  & 1.00 \\
 & 40 & PE1, partial (2/5) & −0.70 [−1.61, +0.47] & −0.96 [−1.76, −0.22] & undetermined &  & 1.00 \\
XD-2 S8w $-$ S8a (post hoc design, Lena target) & 10 & Lena Delta (transfer) & −1.86 [−3.30, +0.10] & +0.49 [−0.96, +1.97] & undetermined; hypothesis rejected (higher error in eight Alaska and Canada rows) &  & 1.00 \\
 & 40 & Lena Delta (transfer) & −0.18 [−1.32, +0.67] & −0.51 [−1.49, +0.41] & undetermined &  & 1.00 \\
XD-3 S8a $-$ S2 & 10 & PE1 (5/5) & +0.32 [+0.01, +0.61] & +0.28 [+0.00, +0.56] & higher error (depends on split independence; size below 0.5 cm) & not met (0.20) & 1.00 \\
 & 40 & PE1, partial (2/5) & −0.63 [−1.06, −0.21] & −0.62 [−0.98, −0.27] & lower error & met ($<$0.0001) & $<$0.0001 \\
XD-3 S8a $-$ S4 & 10 & PE1 (5/5) & +0.09 [−0.36, +0.43] & −0.07 [−0.40, +0.27] & equivalent (depends on split independence) & met (0.020; before correction only) & 0.27 \\
 & 40 & PE1, partial (2/5) & +0.22 [−0.20, +0.62] & +0.06 [−0.27, +0.39] & undetermined (margin-dependent; depends on split independence) & not met (0.16) & 1.00 \\
XD-3 S8b $-$ S2 & 10 & PE1 (5/5) & +0.45 [+0.12, +0.79] & +0.42 [+0.13, +0.72] & higher error (size below 0.5 cm) & not met (0.75) & 1.00 \\
 & 40 & PE1, partial (2/5) & −0.48 [−0.85, −0.08] & −0.47 [−0.80, −0.17] & lower error & met ($<$0.0001) & $<$0.0001 \\
XD-3 S8b $-$ S4 & 10 & PE1 (5/5) & +0.22 [−0.15, +0.53] & +0.07 [−0.22, +0.38] & undetermined (margin-dependent; draw-dependent) & not met (0.074) & 0.89 \\
 & 40 & PE1, partial (2/5) & +0.38 [−0.01, +0.79] & +0.21 [−0.13, +0.57] & undetermined (margin-dependent) & not met (0.61) & 1.00 \\
XD-3 S8c $-$ S2 & 10 & PE1 (5/5) & −0.36 [−0.87, +0.01] & −0.34 [−0.71, +0.03] & undetermined (margin-dependent; draw-dependent) & met ($<$0.0001) & $<$0.0001 \\
 & 40 & PE1, partial (2/5) & −1.13 [−1.67, −0.60] & −0.90 [−1.37, −0.42] & lower error & met ($<$0.0001) & $<$0.0001 \\
XD-3 S8c $-$ S4 & 10 & PE1 (5/5) & −0.59 [−1.21, −0.19] & −0.68 [−1.09, −0.29] & lower error (depends on split independence) & met ($<$0.0001) & $<$0.0001 \\
 & 40 & PE1, partial (2/5) & −0.28 [−0.57, −0.00] & −0.22 [−0.47, +0.04] & undetermined (margin-dependent; draw-dependent) & met ($<$0.0001) & $<$0.0001 \\
XD-3 S8p $-$ S2 & 10 & PE1 (5/5) & +0.52 [−0.37, +1.28] & +0.31 [−0.44, +1.06] & undetermined & not met (0.93) & 1.00 \\
 & 40 & PE1, partial (2/5) & −0.75 [−1.92, +0.12] & −0.50 [−1.33, +0.34] & undetermined & met (0.019; before correction only) & 0.27 \\
XD-3 S8p $-$ S4 & 10 & PE1 (5/5) & +0.29 [−0.49, +0.88] & −0.04 [−0.63, +0.57] & undetermined (margin-dependent) & not met (0.39) & 1.00 \\
 & 40 & PE1, partial (2/5) & +0.10 [−0.92, +0.77] & +0.18 [−0.42, +0.78] & undetermined (margin-dependent) & not met (0.31) & 1.00 \\
\end{longtable}}
\par\medskip\noindent{\small\textit{j. XD-4 learned placement policy S9* (S9-DS) against Algorithm P (= S1) in test tasks outside Alaska (descriptive, no test)}}\par\nopagebreak
{\fontsize{7.5}{9}\selectfont\setlength{\tabcolsep}{4pt}\setlength{\LTpre}{2pt}\setlength{\LTpost}{4pt}
\begin{longtable}{@{}>{\raggedright\arraybackslash}p{\dimexpr 0.400\linewidth-8pt\relax}>{\raggedright\arraybackslash}p{\dimexpr 0.160\linewidth-8pt\relax}>{\raggedright\arraybackslash}p{\dimexpr 0.160\linewidth-8pt\relax}>{\raggedright\arraybackslash}p{\dimexpr 0.280\linewidth-8pt\relax}@{}}
\toprule
\textbf{Test task or family} & \textbf{$v$ cell-weighted} & \textbf{$v$ block-equal} & \textbf{Direction} \\
\midrule\endfirsthead
\toprule
\textbf{Test task or family} & \textbf{$v$ cell-weighted} & \textbf{$v$ block-equal} & \textbf{Direction} \\
\midrule\endhead
\bottomrule\endlastfoot
LE-1 & +0.408 & +0.149 & higher error in both weightings \\
LE-2 & −0.027 & −0.006 & lower error in both weightings \\
CA-2 & −0.008 & +0.005 & weightings disagree \\
CA-3 & −0.036 & −0.004 & lower error in both weightings \\
Tibetan Plateau & +0.132 & +0.086 & higher error in both weightings \\
Central Russia (not counted; no $n$ 40) & +0.031 & +0.018 & descriptive \\
Family Lena Delta (LE-1, LE-2) & +0.190 & +0.071 & higher error in both weightings \\
Family Canada (CA-2, CA-3) & −0.022 & +0.0002 & not consistent \\
Family Tibetan Plateau & +0.132 & +0.086 & higher error in both weightings \\
\end{longtable}}
\par\medskip\noindent{\small\textit{k. XD-6 region traits and placement effect (descriptive)}}\par\nopagebreak
{\fontsize{7.5}{9}\selectfont\setlength{\tabcolsep}{4pt}\setlength{\LTpre}{2pt}\setlength{\LTpost}{4pt}
\begin{longtable}{@{}>{\raggedright\arraybackslash}p{\dimexpr 0.125\linewidth-8pt\relax}>{\raggedright\arraybackslash}p{\dimexpr 0.067\linewidth-8pt\relax}>{\raggedright\arraybackslash}p{\dimexpr 0.092\linewidth-8pt\relax}>{\raggedright\arraybackslash}p{\dimexpr 0.124\linewidth-8pt\relax}>{\raggedright\arraybackslash}p{\dimexpr 0.149\linewidth-8pt\relax}>{\raggedright\arraybackslash}p{\dimexpr 0.221\linewidth-8pt\relax}>{\raggedright\arraybackslash}p{\dimexpr 0.221\linewidth-8pt\relax}@{}}
\toprule
\textbf{Target (transfer)} & \textbf{Blocks} & \textbf{Effective blocks} & \textbf{Covariate between-block share} & \textbf{ln(cell-weighted $E$ / block-mean $E$)} & \textbf{S8a $-$ S1, $n$ 10} & \textbf{S8a $-$ S1, $n$ 40} \\
\midrule\endfirsthead
\toprule
\textbf{Target (transfer)} & \textbf{Blocks} & \textbf{Effective blocks} & \textbf{Covariate between-block share} & \textbf{ln(cell-weighted $E$ / block-mean $E$)} & \textbf{S8a $-$ S1, $n$ 10} & \textbf{S8a $-$ S1, $n$ 40} \\
\midrule\endhead
\bottomrule\endlastfoot
Alaska & 36.2 & 3.49 & 0.78 & +0.095 & −0.96 / −1.97 (lower error) & −1.14 / −2.23 (lower error) \\
Canada & 17.8 & 3.89 & 0.83 & −0.044 & −0.81 / +0.52 (undetermined) & −2.39 / −1.71 (lower error) \\
Lena Delta & 7.6 & 1.72 & 0.39 & −0.069 & +1.70 / +1.02 (undetermined) & +0.98 / −0.22 (undetermined) \\
W Russia & 10.6 & 7.73 & 0.92 & +0.114 & −0.06 / −0.05 (undetermined) &  \\
E Russia & 10.8 & 7.81 & 0.90 & −0.301 & +0.56 / −0.72 (undetermined) &  \\
Central Russia & 5.4 & 2.87 & 0.76 & +0.075 & −3.03 / −1.55 (undetermined) &  \\
Tibetan Plateau & 16.8 & 10.25 & 0.78 & +0.009 & +2.30 / +2.37 (undetermined) & +3.01 / +1.15 (undetermined) \\
AL-1 & 3.75 & 1.26 & 0.35 & +0.006 & +0.35 / −1.65 (undetermined) & −1.11 / −1.82 (lower error) \\
AL-2 & 8.6 & 1.17 & 0.32 & −0.102 & −0.31 / −0.71 (lower error) & −0.34 / −0.81 (lower error) \\
AL-3 & 10.2 & 2.89 & 0.71 & −0.009 & +1.43 / +1.30 (higher error) & +0.35 / +0.15 (undetermined) \\
AL-4 & 1.33 & 1.00 & 0.18 & −0.001 & not decidable & not decidable \\
AL-5 & 6.8 & 1.95 & 0.72 & +0.204 & −2.10 / −2.90 (lower error) & −2.39 / −4.30 (undetermined) \\
AL-6 & 3.25 & 1.00 & 0.14 & −0.000 & not decidable & not decidable \\
\end{longtable}}
\par\noindent{\footnotesize Flags: exploratory (Supplementary Information); designed after earlier results were viewed; XD-1 and XD-3 partly unblinded (S2 and S4 of WF8 viewed); XD-2 partly unblinded and a post hoc design targeting the Lena Delta loss; XD-4 blind (placements of the test tasks never computed before), descriptive, no verdict words and no $P$ values (three families; minimum family-level sign-test $P$ 0.125). Contrasts use R1 ($\lambda$ 0.25); contrasts between different label sets use two-stage CIs, S8w $-$ S8a uses block resampling with the same labels; Holm $m$ = 20 (auxiliary column). S8a, $\gamma$ 0 with farthest-first block order and the cell nearest the block centre first; S8b, farthest-first within blocks; S8c, $\gamma$ 0.5; S8p, $\gamma$ 1; S8w, S8a labels with design-weighted coefficients. Algorithm P (appendix XC-0): rule 2 excluded S2, S4, S8a, S8b, S8c and S8p on Alaskan targets, so Algorithm P = S1 (cell-random sampling); XC-5b, XC-5c and S-XC3 were therefore not tested and remain in the XC Holm family with $P$ = 1. XD-2 higher-error rows: Alaska transfer $n$ 10 +2.02 [+1.36, +2.80] / +1.94 [+1.02, +2.81], $n$ 40 +1.44 / +3.06; Canada transfer $n$ 10 +4.89 [+1.07, +6.65] / +2.82, $n$ 40 +1.99 / +1.21 (both depend on split independence); Alaska within region $n$ 20 +1.21 / +2.35; Canada within region $n$ 50 +1.78 / +1.07, $n$ 100 +1.30 / +0.87 (both depend on split independence), $n$ 200 +1.02 / +0.86. Post hoc rows outside the registered hypotheses (not counted): S8c $-$ S1 lower error in the four-way verdict at $n$ 10 and 40 but undetermined with the two-stage common CI; these rows do not change Algorithm P. XD-4: $v$ = mean over splits, draws, $n \in \{20, 40\}$ and policy seeds of (RMSE(S9*) $-$ RMSE(Algorithm P)) / RMSE(Algorithm P); task RMSE differences (cm, cell-weighted [two-stage CI] (block-equal)): LE-1 $n$ 20 +7.60 [+0.61, +9.12] (+2.43), $n$ 40 +8.82 [+1.30, +10.43] (+3.49); LE-2 $n$ 20 −0.62, $n$ 40 −0.29; CA-2 $n$ 20 +0.19 [−2.01, +2.14], $n$ 40 −0.34 [−1.31, +1.23]; CA-3 $n$ 20 −1.11 [−1.89, +0.30], $n$ 40 −1.18 [−2.06, +0.20]; Tibetan Plateau $n$ 20 +21.72 [+16.47, +26.48] (+13.65), $n$ 40 +7.45 [+4.03, +10.32] (+2.65). Registered sentences (rendered): XD-1, `No difference between placement combining block allocation and covariate coverage and random placement was established with 10 labels' and `... with 40 labels' (transfer, PE1); XD-2, `Design-weighted coefficients increased error in Alaska (transfer $n$ 10 and 40, within region $n$ 20) and Canada (transfer $n$ 10 and 40, within region $n$ 50, 100 and 200)' (post hoc design, Lena target); XD-3, `The error increase of S8 variant V relative to S2 (or S4) was within 0.5 cm' for S8a $-$ S2 ($n$ 40), S8a $-$ S4 ($n$ 10; significant before correction only), S8b $-$ S2 ($n$ 40), S8c $-$ S2 ($n$ 10, 40), S8c $-$ S4 ($n$ 10, 40) and S8p $-$ S2 ($n$ 40; significant before correction only), `... was larger' for S8a $-$ S2 and S8b $-$ S2 at $n$ 10 (statistically distinguishable but below 0.5 cm in size), and `non-inferiority was not established' for the six other rows; XD-4, `A placement policy learned and frozen in Alaskan sub-regions had lower error than Algorithm P in 2 of 5 test tasks outside Alaska (0 of 3 families; both weightings, no test)' and `... higher error in 2 of 5 test tasks (2 of 3 families; both weightings, no test).' Common limitation: the candidate pool consists of already measured cells, so gains are not guaranteed to transfer to field candidates (all land grid cells). XD-6 is descriptive (trait means over splits; effects as cell-weighted / block-equal cm); correlations of traits with effects were not computed. XD-5: part k (continued). Source: data/processed/xbatch/XD\_placement\_policy/sealed/ (xd\_alg\_tests.csv, xd\_alg\_contrasts.csv, xd\_alg\_holm.csv, xd6\_traits.csv; xd4\_summary.json, xd4\_tasks.csv, xd4\_families.csv, xd4\_ci.csv); gate passed (gate/xd\_gate.csv).}\par\smallskip
\par\medskip\noindent{\small\textit{k (continued). XD-5 regret against the best of 200 candidate sets (descriptive)}}\par\nopagebreak
{\fontsize{7.5}{9}\selectfont\setlength{\tabcolsep}{4pt}\setlength{\LTpre}{2pt}\setlength{\LTpost}{4pt}
\begin{longtable}{@{}>{\raggedright\arraybackslash}p{\dimexpr 0.220\linewidth-8pt\relax}>{\raggedright\arraybackslash}p{\dimexpr 0.110\linewidth-8pt\relax}>{\raggedright\arraybackslash}p{\dimexpr 0.140\linewidth-8pt\relax}>{\raggedright\arraybackslash}p{\dimexpr 0.110\linewidth-8pt\relax}>{\raggedright\arraybackslash}p{\dimexpr 0.160\linewidth-8pt\relax}>{\raggedright\arraybackslash}p{\dimexpr 0.130\linewidth-8pt\relax}>{\raggedright\arraybackslash}p{\dimexpr 0.130\linewidth-8pt\relax}@{}}
\toprule
\textbf{Test task} & \textbf{U(S9*)} & \textbf{U(Algorithm P = S1)} & \textbf{U(oracle)} & \textbf{Median of 200 candidate sets} & \textbf{Regret, S9*} & \textbf{Regret, Algorithm P} \\
\midrule\endfirsthead
\toprule
\textbf{Test task} & \textbf{U(S9*)} & \textbf{U(Algorithm P = S1)} & \textbf{U(oracle)} & \textbf{Median of 200 candidate sets} & \textbf{Regret, S9*} & \textbf{Regret, Algorithm P} \\
\midrule\endhead
\bottomrule\endlastfoot
CA-2 & −0.005 & 0 & −0.089 & +0.002 & 0.084 & 0.089 \\
CA-3 & −0.024 & 0 & −0.085 & −0.028 & 0.061 & 0.085 \\
LE-1 & +0.275 & 0 & −0.069 & +0.001 & 0.344 & 0.069 \\
LE-2 & −0.017 & 0 & −0.051 & −0.024 & 0.034 & 0.051 \\
Tibetan Plateau & +0.106 & 0 & −0.090 & −0.009 & 0.195 & 0.090 \\
Central Russia (not counted) & +0.023 & 0 & n.a. & n.a. & n.a. & n.a. \\
\end{longtable}}
\par\medskip\noindent{\small\textit{k (continued). XD-5 cross-validation utility of the learned policies (post hoc)}}\par\nopagebreak
{\fontsize{7.5}{9}\selectfont\setlength{\tabcolsep}{4pt}\setlength{\LTpre}{2pt}\setlength{\LTpost}{4pt}
\begin{longtable}{@{}>{\raggedright\arraybackslash}p{\dimexpr 0.160\linewidth-8pt\relax}>{\raggedright\arraybackslash}p{\dimexpr 0.160\linewidth-8pt\relax}>{\raggedright\arraybackslash}p{\dimexpr 0.170\linewidth-8pt\relax}>{\raggedright\arraybackslash}p{\dimexpr 0.170\linewidth-8pt\relax}>{\raggedright\arraybackslash}p{\dimexpr 0.170\linewidth-8pt\relax}>{\raggedright\arraybackslash}p{\dimexpr 0.170\linewidth-8pt\relax}@{}}
\toprule
\textbf{Task} & \textbf{Family} & \textbf{S9-DS, leave family out} & \textbf{S9-DS, leave task out} & \textbf{S9-GBM, leave family out} & \textbf{S9-GBM, leave task out} \\
\midrule\endfirsthead
\toprule
\textbf{Task} & \textbf{Family} & \textbf{S9-DS, leave family out} & \textbf{S9-DS, leave task out} & \textbf{S9-GBM, leave family out} & \textbf{S9-GBM, leave task out} \\
\midrule\endhead
\bottomrule\endlastfoot
LE-1 & Lena Delta & +0.074 & +0.034 & −0.024 & −0.013 \\
LE-2 & Lena Delta & −0.015 & −0.008 & −0.029 & −0.020 \\
CA-2 & Canada & +0.066 & +0.081 & −0.004 & +0.082 \\
CA-3 & Canada & −0.031 & −0.037 & −0.044 & −0.025 \\
Tibetan Plateau & Tibetan Plateau & +0.053 & +0.072 & −0.009 & −0.009 \\
\end{longtable}}
\par\noindent{\footnotesize XD-5, descriptive (Supplementary Information; no verdict words and no $P$ values), built after the XD-4 tables were opened; local run 5 October 2026, 03:30 to 05:37 KST; sealed tables first opened 05:38:08 KST. U is the relative change of RMSE against the mean of S1 draws, averaged over the two weightings (negative, lower error); regret = U $-$ U(oracle), with the oracle the best of 200 candidate label sets (mean over splits and $n$). Because Algorithm P = S1, the S9* $-$ S1 description equals XD-4. S9* had lower regret than Algorithm P in 3 of 5 test tasks (CA-2, CA-3, LE-2) and higher regret in 2 (LE-1, Tibetan Plateau); the LE-1 mean is set by four rows of splits 2 and 4 (U(S9*) +0.38, +0.73, +0.71, +0.64; the other four rows −0.11 to −0.03). Cross-validation utility (post hoc comparison after S9* = S9-DS had been frozen): leaving a family out, S9-DS was negative in 2 of 5 tasks (LE-2, CA-3) and S9-GBM in 5 of 5 (−0.044 to −0.004). For the Tibetan Plateau both schemes use the same training set; S9-GBM gave the same value (−0.009) and S9-DS gave +0.053 and +0.072, because the minibatch-order seed includes the fold name. The S9-DS utility is therefore sensitive to the minibatch-order seed; this is not a determinism failure (the frozen model gave identical selections on two GPUs, Jaccard 1.000). Alaska-family values (leave the Alaska family out) are in plan 8.11. Common limitation: the candidate pool consists of already measured cells, so gains are not guaranteed to transfer to field candidates. Source: data/processed/xbatch/XD\_placement\_policy/sealed/ (xd5\_regret.csv, xd5\_cv\_ds.csv, xd5\_cv\_gbm.csv, xd5\_meta.json); plan 8.11.}\par\smallskip
\par\medskip\noindent{\small\textit{l. XJ temperature-derived auxiliary labels}}\par\nopagebreak
{\fontsize{7.5}{9}\selectfont\setlength{\tabcolsep}{4pt}\setlength{\LTpre}{2pt}\setlength{\LTpost}{4pt}
\begin{longtable}{@{}>{\raggedright\arraybackslash}p{\dimexpr 0.300\linewidth-8pt\relax}>{\raggedright\arraybackslash}p{\dimexpr 0.050\linewidth-8pt\relax}>{\raggedright\arraybackslash}p{\dimexpr 0.050\linewidth-8pt\relax}>{\raggedright\arraybackslash}p{\dimexpr 0.170\linewidth-8pt\relax}>{\raggedright\arraybackslash}p{\dimexpr 0.170\linewidth-8pt\relax}>{\raggedright\arraybackslash}p{\dimexpr 0.150\linewidth-8pt\relax}>{\raggedright\arraybackslash}p{\dimexpr 0.110\linewidth-8pt\relax}@{}}
\toprule
\textbf{Hypothesis} & \textbf{$w$} & \textbf{$n$} & \textbf{$\Delta$ cell [95\% CI]} & \textbf{$\Delta$ block-equal [95\% CI]} & \textbf{Four-way verdict} & \textbf{Holm $P$ ($m$ 15)} \\
\midrule\endfirsthead
\toprule
\textbf{Hypothesis} & \textbf{$w$} & \textbf{$n$} & \textbf{$\Delta$ cell [95\% CI]} & \textbf{$\Delta$ block-equal [95\% CI]} & \textbf{Four-way verdict} & \textbf{Holm $P$ ($m$ 15)} \\
\midrule\endhead
\bottomrule\endlastfoot
XJ-1 R1@$w$ $-$ R1, four main regions, $n$ 0 and 10, $w$ 0.1, 0.3, 1 (six rows) &  &  &  &  & not tested (licence basis not recorded) & 1.00 \\
XJ-2 D1@$w$ $-$ D1, $n$ 0, $w$ 0.1, 0.3, 1 (three rows) &  &  &  &  & not tested (licence basis not recorded) & 1.00 \\
XJ-3 P1\_aux($w$) $-$ P1, Tibetan Plateau & 0.1 & 3 & −20.85 [−21.80, −19.89] & −21.81 [−23.20, −20.11] & lower error & 0.0015 \\
 & 0.1 & 10 & −8.14 [−8.88, −7.10] & −5.94 [−7.67, −4.02] & lower error & 0.0015 \\
 & 0.3 & 3 & −47.54 [−50.12, −44.50] & −45.18 [−50.11, −39.44] & lower error & 0.0015 \\
 & 0.3 & 10 & −19.63 [−21.74, −16.53] & −12.66 [−17.30, −7.67] & lower error & 0.0015 \\
 & 1 & 3 & −83.17 [−90.21, −73.41] & −66.86 [−79.47, −53.35] & lower error & 0.0015 \\
 & 1 & 10 & −37.50 [−43.30, −28.88] & −20.57 [−31.04, −9.55] & lower error & 0.004 \\
\end{longtable}}
\par\noindent{\footnotesize Flags: designed after earlier results were viewed; XJ-3 partly unblinded (L39). Target-side auxiliary rows: 39 temperature-derived rows (F3, CC BY 4.0), of which rows outside the B blocks of a split entered the A pool (15, 16, 8, 24 and 28 rows in splits 1 to 5); only direct-label cells are scored. Baseline RMSE (cell-weighted): P0 244.51 cm; P1 201.05 cm ($n$ 3) and 148.14 cm ($n$ 10); P1\_aux($w$ 1) 117.89 and 110.63 cm. Descriptive $n$ 0 rows: $-$34.85 ($w$ 0.1), $-$75.81 (0.3) and $-$124.60 cm (1). Registered sentences (rendered): XJ-3, `Temperature-derived auxiliary rows (weight $w$) reduced the transfer error with $k$ labels by $a$ cm. Temperature-derived labels differ in definition from direct measurements (L39), so this result is not generalized as evidence for label augmentation', with ($w$, $k$, $a$) = (0.1, 3, 20.85), (0.1, 10, 8.14), (0.3, 3, 47.54), (0.3, 10, 19.63), (1, 3, 83.17), (1, 10, 37.50); XJ-1 and XJ-2, `The source-side design of temperature-derived auxiliary rows was not tested (licence basis not recorded).' All rows concern one target (Tibetan Plateau). Source: data/processed/xbatch/XJ\_tempderived\_aux\_labels/sealed/ (xj\_hyp.csv, xj\_tests.csv, xj\_holm.csv); gate passed (xj\_gate.csv).}\par\smallskip
\par\medskip\noindent{\small\textit{m. XK error by grid size within regions (added registration; descriptive)}}\par\nopagebreak
{\fontsize{7.5}{9}\selectfont\setlength{\tabcolsep}{4pt}\setlength{\LTpre}{2pt}\setlength{\LTpost}{4pt}
\begin{longtable}{@{}>{\raggedright\arraybackslash}p{\dimexpr 0.089\linewidth-8pt\relax}>{\raggedright\arraybackslash}p{\dimexpr 0.061\linewidth-8pt\relax}>{\raggedright\arraybackslash}p{\dimexpr 0.079\linewidth-8pt\relax}>{\raggedright\arraybackslash}p{\dimexpr 0.069\linewidth-8pt\relax}>{\raggedright\arraybackslash}p{\dimexpr 0.129\linewidth-8pt\relax}>{\raggedright\arraybackslash}p{\dimexpr 0.129\linewidth-8pt\relax}>{\raggedright\arraybackslash}p{\dimexpr 0.297\linewidth-8pt\relax}>{\raggedright\arraybackslash}p{\dimexpr 0.148\linewidth-8pt\relax}@{}}
\toprule
\textbf{Region} & \textbf{Grid} & \textbf{Grid cells} & \textbf{Blocks} & \textbf{P1 cell / block} & \textbf{R1 cell / block} & \textbf{P1 $-$ R1 cell [95\% CI] / block-equal [95\% CI]} & \textbf{CCI v4 cell / block} \\
\midrule\endfirsthead
\toprule
\textbf{Region} & \textbf{Grid} & \textbf{Grid cells} & \textbf{Blocks} & \textbf{P1 cell / block} & \textbf{R1 cell / block} & \textbf{P1 $-$ R1 cell [95\% CI] / block-equal [95\% CI]} & \textbf{CCI v4 cell / block} \\
\midrule\endhead
\bottomrule\endlastfoot
Alaska & 1 km & 13,606 & 74 & 14.30 / 18.16 & 13.80 / 17.31 & 0.50 [0.18, 1.10] / 0.86 [0.29, 1.54] & 19.90 / 28.22 \\
 & 0.05° & 105 & 36 & 12.94 / 15.29 & 11.27 / 13.60 & 1.67 [0.60, 2.69] / 1.70 [0.67, 2.68] & 20.96 / 26.78 \\
 & 0.1° & 83 & 35 & 12.76 / 15.13 & 11.22 / 13.40 & 1.54 [0.50, 2.54] / 1.72 [0.68, 2.69] & 21.30 / 26.66 \\
 & 0.25° & 56 & 35 & 13.16 / 14.78 & 11.55 / 13.05 & 1.61 [0.54, 2.72] / 1.72 [0.64, 2.81] & 23.13 / 26.35 \\
Lena Delta & 1 km & 2,958 & 20 & 21.02 / 34.20 & 21.01 / 32.70 & 0.01 [−1.39, 1.41] / 1.50 [−0.38, 3.20] & 22.45 / 40.23 \\
 & 0.05° & 56 & 13 & 24.66 / 26.74 & 24.97 / 25.90 & −0.30 [−3.32, 1.32] / 0.84 [−2.19, 3.48] & 28.65 / 31.17 \\
 & 0.1° & 37 & 12 & 26.38 / 31.93 & 26.75 / 31.20 & −0.38 [−3.77, 1.72] / 0.72 [−2.23, 3.70] & 30.62 / 36.25 \\
 & 0.25° & 22 & 12 & 15.85 / 18.64 & 15.07 / 17.39 & 0.78 [−1.25, 3.04] / 1.25 [−1.60, 4.26] & 20.14 / 23.33 \\
Canada & 1 km & 747 & 36 & 28.84 / 27.00 & 26.22 / 26.64 & 2.62 [0.00, 4.42] / 0.36 [−3.38, 3.34] & 30.05 / 40.77 \\
 & 0.05° & 37 & 16 & 20.29 / 18.75 & 19.18 / 16.40 & 1.11 [−1.15, 4.49] / 2.34 [0.34, 4.64] & 32.50 / 29.61 \\
 & 0.1° & 26 & 15 & 20.16 / 19.27 & 18.80 / 16.84 & 1.35 [−0.50, 4.22] / 2.43 [0.53, 4.70] & 31.78 / 29.51 \\
 & 0.25° & 19 & 14 & 15.85 / 16.27 & 14.89 / 13.53 & 0.96 [−1.93, 4.45] / 2.74 [−0.04, 4.87] & 33.87 / 28.93 \\
\end{longtable}}
\par\medskip\noindent{\small\textit{m (continued). Products against R1 at the same grid size, RMSE(product) $-$ RMSE(R1)}}\par\nopagebreak
{\fontsize{7.5}{9}\selectfont\setlength{\tabcolsep}{4pt}\setlength{\LTpre}{2pt}\setlength{\LTpost}{4pt}
\begin{longtable}{@{}>{\raggedright\arraybackslash}p{\dimexpr 0.120\linewidth-8pt\relax}>{\raggedright\arraybackslash}p{\dimexpr 0.053\linewidth-8pt\relax}>{\raggedright\arraybackslash}p{\dimexpr 0.160\linewidth-8pt\relax}>{\raggedright\arraybackslash}p{\dimexpr 0.160\linewidth-8pt\relax}>{\raggedright\arraybackslash}p{\dimexpr 0.160\linewidth-8pt\relax}>{\raggedright\arraybackslash}p{\dimexpr 0.180\linewidth-8pt\relax}>{\raggedright\arraybackslash}p{\dimexpr 0.170\linewidth-8pt\relax}@{}}
\toprule
\textbf{Region (cells)} & \textbf{Grid} & \textbf{CCI v4} & \textbf{CCI v5} & \textbf{Wei 2026} & \textbf{Wei 2026 (5 km mask)} & \textbf{Yi-Kimball} \\
\midrule\endfirsthead
\toprule
\textbf{Region (cells)} & \textbf{Grid} & \textbf{CCI v4} & \textbf{CCI v5} & \textbf{Wei 2026} & \textbf{Wei 2026 (5 km mask)} & \textbf{Yi-Kimball} \\
\midrule\endhead
\bottomrule\endlastfoot
Alaska (13,605) & 1 km & 6.10 [2.54, 9.11] / 10.87 & 54.22 [12.08, 90.76] / 53.21 & 2.23 [0.69, 5.01] / −0.94 & 3.02 [0.83, 5.73] / 3.32 (9,131 cells) & 46.09 [10.97, 99.91] / 51.90 (10,605 cells) \\
 & 0.1° & 10.08 [6.42, 14.12] / 13.26 & 53.20 [34.47, 71.94] / 63.02 & 3.88 [0.77, 6.78] / 3.14 & 4.03 [0.74, 7.71] / 3.08 & 39.10 [26.43, 52.86] / 48.49 \\
Lena Delta (1,342) & 1 km & 0.78 [−0.73, 5.49] / 7.69 & −0.31 [−1.87, 1.44] / 1.79 & 9.91 [−4.04, 14.05] / −2.08 & 10.00 [−4.81, 14.13] / −2.03 (1,330 cells) &  \\
 & 0.1° & 3.14 [−0.04, 6.40] / 4.08 & −0.61 [−4.72, 2.64] / 0.09 & −1.14 [−6.72, 8.57] / −1.33 & −1.26 [−7.10, 9.24] / −1.32 &  \\
Canada (745) & 1 km & 3.76 [−2.75, 16.12] / 13.75 & 54.25 [23.70, 78.39] / 61.55 & 2.99 [−3.66, 10.78] / −2.61 & 3.32 [−4.23, 11.33] / 7.36 (726 cells) &  \\
 & 0.1° & 12.98 [1.58, 23.37] / 12.67 & 60.09 [31.09, 87.72] / 67.23 & 5.42 [−0.14, 13.07] / 7.27 & 5.42 [0.36, 13.70] / 7.27 &  \\
\end{longtable}}
\par\noindent{\footnotesize Flags: added registration of 5 October 2026 (commit d721d24, before computation); designed after earlier results were viewed; descriptive, no verdict words. Out-of-block predictions from 0.5° block five-fold cross-validation; grid cells with fewer than three label cells dropped; 1,000 block resamples. First table: models set (cells where P1, R1 and CCI v4 are finite), no mask; RMSE in cm, cell-weighted / block-equal. Second table: RMSE(product) $-$ RMSE(R1) in cm, cell-weighted [95\% CI] / block-equal, common set (all product values finite; Alaska Yi-Kimball on its own set), no mask, plus the 5 km mask edition for Wei 2026. Bias at label cells (product $-$ label, 1 km, cell-weighted, common set), reported as fact without a cause: CCI v5 +35.63 cm (Alaska), +58.60 cm (Canada) and +0.27 cm (Lena Delta); CCI v4 $-$5.09, +10.35 and $-$6.04 cm. Pre-fixed interpretation rules: rule 1 (all RMSE decrease with grid size) does not hold, so its sentence is not used; rule 2 does not fix the weighting and the two weightings pick different grid sizes in all three regions, so only the fact is reported: `In Alaska the RMSE difference between P1 and R1 was 0.50 cm at 1 km and 1.54 to 1.67 cm at 0.05°, 0.1° and 0.25° (cell-weighted).' Products are compared only as RMSE differences at the same grid size, for example `At the 0.1° grid in Alaska, the RMSE of CCI v4, CCI v5, Wei 2026 and Yi-Kimball was larger than that of R1 by 10.08, 53.20, 3.88 and 39.10 cm' and `At the 0.1° grid in the Lena Delta, the RMSE of CCI v5 and Wei 2026 was smaller than that of R1 by 0.61 and 1.14 cm (CIs include zero)' (cell-weighted); product rankings are not generalized, and Wei 2026 used CALM sites for training (WRAPUP 10). Supplementary Fig.~S15. Source: data/processed/xbatch/XK\_support\_scale/xk\_support\_scale\_v1.csv and xk\_meta.json.}\par\smallskip
\par\medskip\noindent{\small\textit{n. XL 1 km maps compared with ALT products (added registration; descriptive)}}\par\nopagebreak
{\fontsize{7.5}{9}\selectfont\setlength{\tabcolsep}{4pt}\setlength{\LTpre}{2pt}\setlength{\LTpost}{4pt}
\begin{longtable}{@{}>{\raggedright\arraybackslash}p{\dimexpr 0.078\linewidth-8pt\relax}>{\raggedright\arraybackslash}p{\dimexpr 0.233\linewidth-8pt\relax}>{\raggedright\arraybackslash}p{\dimexpr 0.068\linewidth-8pt\relax}>{\raggedright\arraybackslash}p{\dimexpr 0.117\linewidth-8pt\relax}>{\raggedright\arraybackslash}p{\dimexpr 0.101\linewidth-8pt\relax}>{\raggedright\arraybackslash}p{\dimexpr 0.101\linewidth-8pt\relax}>{\raggedright\arraybackslash}p{\dimexpr 0.078\linewidth-8pt\relax}>{\raggedright\arraybackslash}p{\dimexpr 0.088\linewidth-8pt\relax}>{\raggedright\arraybackslash}p{\dimexpr 0.136\linewidth-8pt\relax}@{}}
\toprule
\textbf{Region} & \textbf{Map} & \textbf{Mean} & \textbf{Quartiles (25--75\%)} & \textbf{Mean difference} & \textbf{Mean absolute difference} & \textbf{Pearson $r$} & \textbf{Spearman $\rho$} & \textbf{Share of mapped cells with values} \\
\midrule\endfirsthead
\toprule
\textbf{Region} & \textbf{Map} & \textbf{Mean} & \textbf{Quartiles (25--75\%)} & \textbf{Mean difference} & \textbf{Mean absolute difference} & \textbf{Pearson $r$} & \textbf{Spearman $\rho$} & \textbf{Share of mapped cells with values} \\
\midrule\endhead
\bottomrule\endlastfoot
Alaska & this study (recalibrated anchor plus residual ML) & 56.09 & 48.41–63.61 &  &  &  &  &  \\
 & recalibrated Stefan & 57.46 & 48.92–66.06 & −1.38 & 2.57 & 0.956 & 0.961 & 1.000 \\
 & this map, 0.1° cell means & 56.09 & 48.49–63.53 & 0.00 & 1.37 & 0.983 & 0.983 & 1.000 \\
 & CCI v5 (1997–2023 mean) & 105.46 & 61.30–133.89 & −49.37 & 50.45 & 0.464 & 0.642 & 0.999 \\
 & Wei 2026 (2000–2024 mean) & 70.79 & 57.64–71.53 & −14.66 & 18.14 & −0.154 & 0.090 & 0.994 \\
 & Aalto 2018 (2000–2014) & 78.82 & 60.40–102.61 & −22.63 & 37.16 & −0.104 & 0.106 & 0.991 \\
 & Yi-Kimball (2001–2015 mean) & 105.15 & 54.33–137.00 & −49.33 & 53.37 & 0.338 & 0.461 & 0.862 \\
Lena Delta & this study & 42.18 & 40.90–43.19 &  &  &  &  &  \\
 & recalibrated Stefan & 41.09 & 40.12–42.10 & 1.10 & 1.61 & 0.361 & 0.387 & 1.000 \\
 & this map, 0.1° cell means & 42.18 & 41.14–43.16 & 0.00 & 0.81 & 0.783 & 0.816 & 1.000 \\
 & CCI v5 & 41.24 & 32.96–48.15 & 0.94 & 8.99 & −0.226 & −0.262 & 1.000 \\
 & Wei 2026 & 58.64 & 53.32–63.07 & −16.50 & 16.51 & 0.045 & 0.068 & 0.940 \\
 & Aalto 2018 & 55.16 & 48.81–60.46 & −12.98 & 13.09 & −0.015 & −0.023 & 0.986 \\
\end{longtable}}
\par\noindent{\footnotesize Flags: added registration of 5 October 2026 (commit d721d24, before computation); designed after earlier results were viewed; descriptive, no verdict. Mapped cells (cm; difference = this map minus the compared map): Alaska 858,517 of 899,633 domain cells (permafrost fraction at least 50\%), Lena Delta 38,086 of 53,011. The 0.1° cell means explain 97\% (Alaska, $r$ 0.983) and 61\% (Lena Delta, $r$ 0.783) of the variance of the 1 km map; the standard deviation of 1 km values within a 0.1° cell is 1.8 cm in Alaska (total 9.5 cm) and 1.1 cm in the Lena Delta (total 1.8 cm). Registered map sentence (rendered): `The 1 km map is a display resolution; its information resolution is bounded by the climate (about 10 km) and soil (about 5 km) inputs.' The maps carry no accuracy claim; accuracy is compared only at label cells (XG and XK). Areas of large difference are named without a cause: Interior Alaska (CCI v5 and Yi-Kimball, more than 50 cm), about 61–63° N, 141–150° W (Wei 2026), and west of about 148° W at 62–64° N (this map deeper than Aalto 2018); in the Lena Delta, the northern and eastern delta (this map deeper than CCI v5) and most of the grid (this map shallower than Wei 2026 and Aalto 2018). In Alaska, 73.2\% of mapped cells lie outside the training range of at least one of the 25 covariates (mainly terrain roughness, slope and topographic position); the constant-width split-conformal interval is $\pm$20.15 cm (Alaska) and $\pm$25.39 cm (Lena Delta), and its coverage is not guaranteed in extrapolated cells. Supplementary Figs~S10 to S14. Source: data/processed/xbatch/XL\_map\_products/xl\_summary\_v1.csv and xl\_meta.json.}\par\smallskip
\par\medskip\noindent{\small\textit{o. XC workflow applied in sequence, in new random splits of the same regions (PE1 pool; Algorithm P = S1)}}\par\nopagebreak
{\fontsize{7.5}{9}\selectfont\setlength{\tabcolsep}{4pt}\setlength{\LTpre}{2pt}\setlength{\LTpost}{4pt}
\begin{longtable}{@{}>{\raggedright\arraybackslash}p{\dimexpr 0.193\linewidth-8pt\relax}>{\raggedright\arraybackslash}p{\dimexpr 0.043\linewidth-8pt\relax}>{\raggedright\arraybackslash}p{\dimexpr 0.116\linewidth-8pt\relax}>{\raggedright\arraybackslash}p{\dimexpr 0.125\linewidth-8pt\relax}>{\raggedright\arraybackslash}p{\dimexpr 0.125\linewidth-8pt\relax}>{\raggedright\arraybackslash}p{\dimexpr 0.116\linewidth-8pt\relax}>{\raggedright\arraybackslash}p{\dimexpr 0.108\linewidth-8pt\relax}>{\raggedright\arraybackslash}p{\dimexpr 0.174\linewidth-8pt\relax}@{}}
\toprule
\textbf{Hypothesis and contrast} & \textbf{$n$} & \textbf{Pool} & \textbf{$\Delta$ cell [95\% CI]} & \textbf{$\Delta$ block-equal [95\% CI]} & \textbf{Four-way verdict} & \textbf{Holm $P$ (rank, multiplier)} & \textbf{Branch} \\
\midrule\endfirsthead
\toprule
\textbf{Hypothesis and contrast} & \textbf{$n$} & \textbf{Pool} & \textbf{$\Delta$ cell [95\% CI]} & \textbf{$\Delta$ block-equal [95\% CI]} & \textbf{Four-way verdict} & \textbf{Holm $P$ (rank, multiplier)} & \textbf{Branch} \\
\midrule\endhead
\bottomrule\endlastfoot
XC-1b A1 $-$ A2 (workflow $-$ random placement with fixed recipe) & 40 & partial (2/5: Lena Delta, Canada) & −0.86 [−1.03, −0.39] & −0.36 [−0.72, +0.003] & undetermined & 0.21 (5, 4) & undetermined (minimum detectable effect about 0.5 cm; two-stage 0.46 cm) \\
XC-1c A1 $-$ A2 & 160 & partial (2/5) & −0.93 [−1.25, −0.40] & −0.54 [−0.94, −0.13] & lower error & 0.040 (4, 5) & lower error \\
XC-2a A1 $-$ A3 (workflow $-$ recalibrated Stefan), non-inferiority 0.5 cm & 10 & 5/5 & +0.08 [−0.07, +0.20] & −0.005 [−0.15, +0.14] & equivalent & $<$0.0001 (1, 8) & non-inferior; XC-2s not superior \\
XC-2b A1 $-$ A3 & 40 & partial (2/5) & −1.35 [−1.64, −0.80] & −1.00 [−1.42, −0.59] & lower error & $<$0.0001 (2, 7) & non-inferior and superior (XC-2s adjusted $P$ 0.0008) \\
XC-2c A1 $-$ A3 & 160 & partial (2/5) & −1.51 [−1.99, −0.88] & −1.29 [−1.78, −0.80] & lower error & $<$0.0001 (3, 6) & non-inferior and superior (XC-2s adjusted $P$ 0.0008) \\
XC-5b A1 $-$ A5 (placement under the selection rule) & 40 &  &  &  & not tested & 1.00 (6, 3) & not decidable (Algorithm P = S1, rule 5) \\
XC-5c A1 $-$ A5 & 160 &  &  &  & not tested & 1.00 (7, 2) & not decidable (Algorithm P = S1, rule 5) \\
XC-F3 A1 $-$ A2, external family F3 & 10 & F3 &  &  & not tested & 1.00 (8, 1) & not decidable ($A1 = A2$ at $n$ 10); \textcolor{blue}{[PENDING: appendix XC-F3 after the XF data deadline, 11 October 2026, 14:22 KST]} \\
\end{longtable}}
\par\medskip\noindent{\small\textit{o (continued). Regional rows}}\par\nopagebreak
{\fontsize{7.5}{9}\selectfont\setlength{\tabcolsep}{4pt}\setlength{\LTpre}{2pt}\setlength{\LTpost}{4pt}
\begin{longtable}{@{}>{\raggedright\arraybackslash}p{\dimexpr 0.240\linewidth-8pt\relax}>{\raggedright\arraybackslash}p{\dimexpr 0.050\linewidth-8pt\relax}>{\raggedright\arraybackslash}p{\dimexpr 0.170\linewidth-8pt\relax}>{\raggedright\arraybackslash}p{\dimexpr 0.170\linewidth-8pt\relax}>{\raggedright\arraybackslash}p{\dimexpr 0.170\linewidth-8pt\relax}>{\raggedright\arraybackslash}p{\dimexpr 0.200\linewidth-8pt\relax}@{}}
\toprule
\textbf{Contrast} & \textbf{$n$} & \textbf{Region} & \textbf{Cell-weighted [95\% CI]} & \textbf{Block-equal [95\% CI]} & \textbf{Verdict} \\
\midrule\endfirsthead
\toprule
\textbf{Contrast} & \textbf{$n$} & \textbf{Region} & \textbf{Cell-weighted [95\% CI]} & \textbf{Block-equal [95\% CI]} & \textbf{Verdict} \\
\midrule\endhead
\bottomrule\endlastfoot
A1 $-$ A2 & 40 & Lena Delta & +0.21 [+0.02, +0.50] & +0.11 [−0.10, +0.32] & undetermined (margin-dependent) \\
 & 40 & Canada & −1.94 [−2.27, −1.07] & −0.83 [−1.52, −0.15] & lower error \\
 & 160 & Lena Delta & +0.54 [−0.05, +0.82] & −0.10 [−0.49, +0.27] & undetermined (margin-dependent) \\
 & 160 & Canada & −2.41 [−2.77, −1.31] & −0.98 [−1.69, −0.28] & lower error \\
A1 $-$ A3 & 10 & Lena Delta & +0.08 [−0.26, +0.14] & −0.35 [−0.54, −0.18] & undetermined \\
 & 10 & Canada & −0.84 & −0.85 & lower error \\
 & 10 & W Russia & −0.28 & −0.27 & undetermined \\
 & 10 & E Russia & +1.34 [+0.95, +1.84] & +1.30 [+0.82, +1.78] & higher error \\
 & 10 & Central Russia & +0.10 & +0.14 & equivalent (depends on split independence) \\
 & 40 & Lena Delta & +0.25 [−0.13, +0.41] & −0.24 [−0.49, −0.01] & equivalent (depends on split independence) \\
 & 40 & Canada & −2.96 & −1.75 & lower error \\
 & 160 & Lena Delta & +0.53 [−0.35, +0.75] & −0.61 [−1.15, −0.10] & undetermined \\
 & 160 & Canada & −3.55 & −1.96 & lower error \\
A1 $-$ A3, selection family (S-XC4 b, S-XC8) & 40 & Alaska (parent excluded) & +1.51 [+0.73, +1.79] & +0.48 [+0.28, +0.67] & higher error \\
 & 160 & Alaska (parent excluded) & +2.34 [+0.73, +2.86] & +0.68 [+0.36, +1.04] & higher error \\
\end{longtable}}
\par\noindent{\footnotesize Flags: designed after earlier results were viewed; re-test in reused regions (new random splits of the same cells, seeds 201 to 210); XC-1, XC-2 and XC-5 partly unblinded (WF2, WF4, WF8, LGX-N4), XC-F3 blind (registered before running). Holm $m$ = 8; the Holm-adjusted $P$ sets the wording; no row needed the 'significant before correction only' flag (XC-1c 0.040). Run on the local server, 5 October 2026, 03:16 to 05:19 KST (208 units, no failure); sealed tables first opened 05:20:39 KST; reproduction gate passed (128 keys, local-to-cloud tolerance). Algorithm P = S1 (appendix XC-0, rule 5): the placement step of the workflow A1 was random nested labels, so A1 = A2 at $n$ 10 (both R1 with $\lambda$ 0.25) and A1 = A5 at $n$ 40 and 160 (both the selection rule W); XC-1b and XC-1c therefore equal the auxiliary contrast S-XC1 at $n$ 40 and 160 (rule W against the fixed recipe on the same labels), and the placement step contributed nothing. XC-5b, XC-5c, XC-F3, S-XC0 and S-XC3 were not tested and stay in the Holm family with $P$ = 1. Because the two arms of the main contrasts share label sets, the main CI is the same-label block-resampling CI; common-resampling and two-stage upper bounds gave the same four-way and non-inferiority verdicts. HK region-level intervals include zero for all five tested rows (XC-2a [−0.92, +1.04] cm), so no region-general sentence is written. Leave one region out (auxiliary): XC-1b and XC-1c are undetermined with the Lena Delta only (+0.21, +0.54) and lower error with Canada only (−1.94, −2.41); without Canada, XC-2a has higher error (+0.31 [+0.12, +0.44] / +0.21 [+0.04, +0.37], both upper bounds below 0.5 cm); without E Russia, XC-2a has lower error (−0.24 / −0.33). S-XC4 (a), targets more than 2 cm worse than P0 (A1 / A2 / A3), independent-region family: $n$ 10 1 (E Russia) / 1 (E Russia) / 0; $n$ 40 0 / 0 / 1 (Canada); $n$ 160 1 (Alaska) / 1 (Canada) / 1 (Canada); all labels 1 (Alaska) / 2 (Canada, E Russia) / 1 (Canada). S-XC4 (b), regions more than 0.5 cm worse than the recalibrated Stefan model (cell-weighted point estimate): $n$ 10, 1 of 7 (E Russia +1.34); $n$ 40, 1 of 4 (Alaska, selection family, +1.51); $n$ 160, 2 of 3 (Lena Delta +0.53, Alaska +2.34); Alaska at $n$ 10 was +0.05 cell-weighted but +1.75 [+1.34, +2.16] block-equal; Tibetan Plateau $n$ 10 −21.00 / −17.54, $n$ 40 −29.68 / −23.46. Auxiliary contrasts (verdict words, not counted): S-XC1 (rule W $-$ fixed R1, same labels) $n$ 10 −0.87 [−1.38, −0.38] / −0.89 [−1.43, −0.35], lower error (depends on split independence); S-XC2 (stacking candidates added to W) $n$ 40 −0.12 / +0.30 equivalent, $n$ 160 −0.41 / +0.03 undetermined; S-XC5 regret A1 $-$ A6 $n$ 10 +1.73 / +1.41 (higher error), $n$ 40 +0.24 / +0.21, $n$ 160 +0.42 / +0.20 (undetermined); S-XC6 (post hoc threshold) median RMSE(P0) $-$ RMSE(A1 at the largest $n$) 0.37 / 0.21 cm for 16 targets with $|b_{10}| \geq 8$ cm and −1.81 / −0.57 cm for 5 targets below; S-XC8, PE1 all labels A1 $-$ A2 −1.29 / −1.33 and A1 $-$ A3 −1.43 / −1.56 (lower error); S-XC9, equal and $|b_{10}|$-proportional budgets gave the same allocation; S-XC10 (A1 $-$ P0, margin 0.5 cm) $n$ 10 −1.87 / −1.58 (lower error, non-inferior), $n$ 40 −0.22 / −0.14 and $n$ 160 −0.06 / −0.20 (undetermined, non-inferior, margin-dependent), E Russia $n$ 10 +2.30 [+1.35, +4.54] / +4.00 (higher error); S-XC11 (descriptive) share of reducible error removed, Canada 0.20 to 0.25, Lena Delta −0.04 to −0.09, Alaska −0.42 to −1.64; S-XC12 $\rho$ 0.17 [−0.30, 0.86] / −0.06 [−0.49, 0.65] (21 targets, 7 families). Within-region auxiliary arm XC-r (descriptive verdicts, outside the Holm family): rule W $-$ R1 (cross-validated $\lambda$) three targets $n$ 200 −0.35 [−0.49, +0.06] / +0.15 [−0.14, +0.46] equivalent (depends on split independence, margin-dependent), $n$ 500 +0.03 / +0.14 and $n$ 1,000 +0.36 / +0.05 undetermined (Alaska rows equivalent at all three). Registered sentences (rendered; all with 'in new random splits of the same regions'): common sentence, `Algorithm P was chosen in the Alaska family, and the cells of the evaluation regions were used to design the components (placement, rule W); the evaluation splits are new partitions of the same cells', with the fact that the placement step was random placement; XC-1b, `the workflow that applies the placement algorithm, the ten-label bias diagnostic and the label-count method rule in sequence established no difference from random placement with the fixed residual recipe at 40 labels (minimum detectable effect about 0.5 cm)', Lena Delta row +0.21 cm [0.02, 0.50], undetermined; XC-1c, `... had 0.93 cm (cell-weighted) lower error than random placement with the fixed residual recipe at 160 labels (Lena Delta and Canada, two regions)', Lena Delta row +0.54 cm [−0.05, 0.82], undetermined; XC-2a, `the error increase of the workflow relative to the recalibrated Stefan model was within 0.5 cm (non-inferior, 10 labels)', `in E Russia the error was larger ($\Delta$ +1.34 cm, CI [0.95, 1.84])', `of seven independent regions (five PE1 regions, Alaska, Tibetan Plateau) the workflow was more than 0.5 cm worse than the recalibrated Stefan model in one (E Russia, +1.34 cm)' and `non-inferiority refers to the pooled mean and did not hold in E Russia', Lena Delta row +0.08, undetermined; XC-2b and XC-2c, non-inferior at 40 and 160 labels and `1.35 cm' and `1.51 cm lower error than the recalibrated Stefan model', with one of four evaluated regions (Alaska, selection family, +1.51 cm) and two of three (Lena Delta +0.53 cm, Alaska +2.34 cm) more than 0.5 cm worse, and `non-inferiority refers to the pooled mean and did not hold' in those regions; Lena Delta rows +0.25 (equivalent) and +0.53 (undetermined); XC-5b, XC-5c, not tested (Algorithm P = S1); XC-F3, not tested ($A1 = A2$ at $n$ 10); the registered 'no eligible region' sentence is written after the XF data deadline. Source: data/processed/xbatch/XC\_workflow\_end\_to\_end/sealed/ (xc\_hyp.csv, xc\_contrasts\_main.csv, xc\_holm.csv, xc\_sentences.json, xc\_harm\_a.csv, xc\_harm\_b.csv, xc\_loo.csv, xc\_contrasts\_aux.csv, xc\_xcr\_contrasts.csv); docs/EXPERIMENT\_PLAN\_FINAL\_BATCH\_2026-10-04.md 8.10.}\par\smallskip
\par\medskip\noindent{\small\textit{p. XF}}\par\nopagebreak
{\fontsize{7.5}{9}\selectfont\setlength{\tabcolsep}{4pt}\setlength{\LTpre}{2pt}\setlength{\LTpost}{4pt}
\begin{longtable}{@{}>{\raggedright\arraybackslash}p{\dimexpr 0.500\linewidth-8pt\relax}>{\raggedright\arraybackslash}p{\dimexpr 0.500\linewidth-8pt\relax}@{}}
\toprule
\textbf{Analysis} & \textbf{Status} \\
\midrule\endfirsthead
\toprule
\textbf{Analysis} & \textbf{Status} \\
\midrule\endhead
\bottomrule\endlastfoot
XF new independent regions & \textcolor{blue}{[PENDING: XF, new public regions; registered data deadline 11 October 2026, 14:22 KST; no eligible new macro region found so far]} \\
\end{longtable}}

\clearpage
\phantomsection\label{tab:S14}
\noindent\textbf{Supplementary Table S14.} Deployment scenarios SC1w to SC3w and related wrap-up tests.\par\smallskip
\noindent{\footnotesize Registered in WRAPUP (316714c, 30 September 2026, 03:02:42), partly unblinded (earlier shrinkage curves were known). The verdict column translates the registered verdict string; values are copied from the stat column. Source: paper/claims/C1\_label0\_safety/tables/lgw\_tests.csv (copy of data/processed/lgw/lgw\_tests.csv) and docs/MANUSCRIPT\_DRAFT\_RESULTS\_ABSTRACT\_2026-09-30.md (Deployment scenarios). Method codes (internal identifiers allowed in the Supplementary Information): P0 source-coefficient Stefan model; P* year-matched Stefan model; P1 recalibrated Stefan model ($\kappa = 10$); P2 local least-squares Stefan model; P1* or P1@ed soil-property recalibrated Stefan model; R0 source anchor plus residual ML; R1 recalibrated anchor plus residual ML; R2 augmentation plus anchor plus residual; D0 direct ML; D1 physics pseudo-label augmentation; F1a, F1k, F1n physics-input ML; RM multiplicative residual; Re Stefan-CCI mean anchor plus residual; W selection rule (cross-validation within the target labels); S1 random cells, S2 block stratification, S3 covariate-cluster stratification, S4 covariate max-min distance, S5 source-extrapolation first, S6 sequential selection by prediction variance, S7 $\sqrt{\mathrm{TDD}}$ first. Target suffixes: $|x$ parent region excluded from the source, $|i$ rest of the parent kept, $|r$ within-region design. $n$ = number of target labels; all = all labels of half A. Verdicts: lower error, higher error, equivalent ($\pm$0.5~cm), undetermined; both 95\% block-bootstrap CIs (cell-weighted and block-equal) must agree. \dag~a note of the source row in Korean was omitted; see the cited README section.}\par
{\fontsize{7.5}{9}\selectfont\setlength{\tabcolsep}{4pt}\setlength{\LTpre}{2pt}\setlength{\LTpost}{4pt}
\begin{longtable}{@{}>{\raggedright\arraybackslash}p{\dimexpr 0.091\linewidth-8pt\relax}>{\raggedright\arraybackslash}p{\dimexpr 0.139\linewidth-8pt\relax}>{\raggedright\arraybackslash}p{\dimexpr 0.238\linewidth-8pt\relax}>{\raggedright\arraybackslash}p{\dimexpr 0.228\linewidth-8pt\relax}>{\raggedright\arraybackslash}p{\dimexpr 0.304\linewidth-8pt\relax}@{}}
\toprule
\textbf{ID} & \textbf{Role} & \textbf{Question} & \textbf{Registered verdict} & \textbf{Values} \\
\midrule\endfirsthead
\toprule
\textbf{ID} & \textbf{Role} & \textbf{Question} & \textbf{Registered verdict} & \textbf{Values} \\
\midrule\endhead
\bottomrule\endlastfoot
SC1w & main & Few-label recipe R1 non-inferior to P0 (0.5 cm) at 3 and 10 labels in five regions & safety not confirmed (non-inferior cells 4/10); no higher-error cell & non-inferior: Lena Delta $n$ 3, Canada $n$ 3, W Russia $n$ 3 and 10; not confirmed: Lena Delta $n$ 10, Canada $n$ 10, E Russia $n$ 3 and 10, Alaska $n$ 3 and 10 \\
SC1w-d10 & auxiliary ($\delta$ 1.0 cm) & Same with a 1.0 cm margin & safety not confirmed (non-inferior cells 5/10) & not confirmed: Canada $n$ 10, E Russia $n$ 3 and 10, Alaska $n$ 3 and 10 \\
SC1w-P & main & Recalibration alone (P1) non-inferior to P0 & rejected & Canada $n$ 3: $\Delta$ +0.97 [+0.26, +1.30], Holm $P$ 0.021; Canada $n$ 10: +2.26 [+0.71, +2.93], Holm $P$ 0.008 \\
SC2w & main & Residual learning versus recalibrated Stefan at 40 labels (Lena Delta, Canada, Alaska) & sentence A: lower error in 2 of 3 regions & Lena Delta lower; Canada lower; Alaska undetermined \\
SC2w-n160 & auxiliary & Same at 160 labels & sentence A: lower error in 2 of 3 regions & Lena Delta lower; Canada lower; Alaska undetermined \\
SC3w & main & Sequential stopping rule ($\tau$ 0.05) halves labels against always 40 with error increase $\leq$ 0.3 cm and no higher-error region & rule not operating (label ratio 0.91); $\Delta$ described only & label ratio 0.909; three-region mean $\Delta$ $-$0.048 [$-$0.079, $-$0.002], block-equal $-$0.077 [$-$0.105, $-$0.049]; condition (i) not met, (ii) to (iv) met \\
SC3w, $\tau$ 0.025 & descriptive & $\tau$ sensitivity & rule not operating (label ratio 0.99) & label ratio 0.988; $\Delta$ $-$0.001 [$-$0.006, +0.004], block-equal $-$0.009 [$-$0.014, $-$0.003] \\
SC3w, $\tau$ 0.1 & descriptive & $\tau$ sensitivity & rejected (condition (i) not met) & label ratio 0.542; $\Delta$ $-$0.177 [$-$0.291, $-$0.055], block-equal $-$0.310 [$-$0.394, $-$0.225] \\
SC3w, $\tau$ 0.2 & descriptive & $\tau$ sensitivity & all four conditions met (not the main verdict) & label ratio 0.185; $\Delta$ $-$0.677 [$-$0.909, $-$0.353], block-equal $-$0.845 [$-$1.041, $-$0.640] \\
S-a & sensitivity (interpretation rule fixed before results) & Label unit: point labels versus 1 km location means, P1 $-$ P0 & with point labels the few-label recalibration gain was smaller (Canada, $n$ 10, $-$0.86 cm); deployment label unit is the 1 km location mean & median $\Delta_{\mathrm{unit}}$: Lena Delta $n$ 3 +0.31, $n$ 10 +0.34; Canada $n$ 3 $-$0.47, $n$ 10 $-$0.86; Alaska $n$ 3 +0.07, $n$ 10 +0.09 \\
S-b & sensitivity & Label year: single-year versus multi-year mean labels (W and E Russia) & recalibration gain with 3 to 10 labels maintained & median $\Delta_{\mathrm{year}}$: W Russia $n$ 3 +0.02, $n$ 10 +0.05; E Russia $n$ 3 +0.01, $n$ 10 +0.04 \\
L43 & main & Block-dispersed versus cell-random placement for P1 & placement effect not established & P1 $n$ 10 undetermined (regions 4/4, Holm $P$ 1); $n$ 40 undetermined (regions 2/4, Holm $P$ 1) \\
AK1w & main & Combinations that beat the physics model in Alaskan sub-regions with an Alaskan source also do so with an outside-only source & majority ($r$ 0.57, 68 pairs, 5 targets) & $r$ 0.574; $r'$ 0.779; target-level majority 2/5; sensitivity $r$ 0.619 \\
\end{longtable}}

\clearpage
\phantomsection\label{tab:S15}
\noindent\textbf{Supplementary Table S15.} Selection and withdrawal history and earlier negative results.\par\smallskip
\noindent{\footnotesize Part a lists claims, hypotheses and rules that were dropped, withdrawn, replaced or corrected, with the commit that recorded the action and its timing relative to the first retrieval of LG-family results (T\_res, 30 September 2026, 07:20:22). Part b lists earlier negative results reclassified under the four-way verdict rule; its result columns are not yet filled. Source: docs/MANUSCRIPT\_DRAFT\_SUPPORT\_2026-09-30.md M4.5 and M4.4 (draft of 30 September 2026, not re-verified at assembly).}\par
\par\medskip\noindent{\small\textit{a. Selection and withdrawal history}}\par\nopagebreak
{\fontsize{7.5}{9}\selectfont\setlength{\tabcolsep}{4pt}\setlength{\LTpre}{2pt}\setlength{\LTpost}{4pt}
\begin{longtable}{@{}>{\raggedright\arraybackslash}p{\dimexpr 0.033\linewidth-8pt\relax}>{\raggedright\arraybackslash}p{\dimexpr 0.176\linewidth-8pt\relax}>{\raggedright\arraybackslash}p{\dimexpr 0.096\linewidth-8pt\relax}>{\raggedright\arraybackslash}p{\dimexpr 0.128\linewidth-8pt\relax}>{\raggedright\arraybackslash}p{\dimexpr 0.279\linewidth-8pt\relax}>{\raggedright\arraybackslash}p{\dimexpr 0.167\linewidth-8pt\relax}>{\raggedright\arraybackslash}p{\dimexpr 0.120\linewidth-8pt\relax}@{}}
\toprule
\textbf{\#} & \textbf{Item} & \textbf{Type} & \textbf{Origin (commit, KST)} & \textbf{Action and reason} & \textbf{Recorded in (commit, KST)} & \textbf{Timing relative to LG-family results} \\
\midrule\endfirsthead
\toprule
\textbf{\#} & \textbf{Item} & \textbf{Type} & \textbf{Origin (commit, KST)} & \textbf{Action and reason} & \textbf{Recorded in (commit, KST)} & \textbf{Timing relative to LG-family results} \\
\midrule\endhead
\bottomrule\endlastfoot
1 & "Component decomposition: 72–95 \% of the few-label gain in level regions is coefficient recalibration" (NOVELTY N3; FINAL plan §10.3) & dropped claim & FINAL plan and session records, 2026-09-26 & Rejected as a contribution claim (2 reviewers reject, 1 conditional). Reasons: the residual CI includes 0 only in Russia W; CA-2 degrades under recalibration (+7.7 cm); the share keeps only targets where recalibration worked; the classification uses scored labels; at α = 1 the residual does not respond to target labels. Remaining form: case values in cm, no percentage & NOVELTY §1, §5, §10 (`d8355d2` 2026-09-29 18:18:36) & before (no LG result) \\
2 & "The ranking of ML and physics reverses with the validation scheme" (NOVELTY N8) & dropped claim & earlier internal answers & Rejected. The paired Alaskan block contrast failed the two-weighting rule (cell-weighted +1.66 [−1.32, 3.93], block-equal −1.87 [−3.15, −0.72] cm); direct ridge was below Stefan in every scheme; no CI; pattern already in Gautam 2025. Remaining form: methods justification. The registered test is L19, whose rule forbids the word "reversal" unless direct ML wins under random splits & NOVELTY §5 (`d8355d2` 2026-09-29 18:18:36); L19 (`8513cc2` 18:17:29) & before \\
3 & NOVELTY N5 placement clause (dispersed versus concentrated labels) & dropped clause & NOVELTY draft & Deleted: no block-equal CI on mean rows; sign of the n = 160 value depends on the target set. A placement sentence is allowed only through L43 & NOVELTY §5 (`d8355d2`); L43 (`316714c` 2026-09-30 03:02:42) & before \\
4 & NOVELTY N2 (two baselines) as an independent contribution & demoted & NOVELTY draft & Moved to Methods; depends on N1 and has precedents & NOVELTY §5 (`d8355d2`) & before \\
5 & "Break-even n predicted by abs(log E ratio), ρ −0.84"; "H24 k-center adopted"; "Lena and Canada cannot exceed at n ≤ 320"; contest result "CCI combination cancels bias"; control ladder "all rungs significant" & dropped claims & H24, H25, H27, contest 2026-07 & Discarded or withdrawn, no reuse & RESEARCH\_\allowbreak{}FRAME A.4 (`6830277` 2026-09-29 13:18:02); AUDIT (`f098031` 2026-09-26 21:42:16) & before \\
6 & SC1–SC3 (NEXT P3; RESEARCH\_\allowbreak{}FRAME §5.1 rank 5, §5.2) & withdrawn hypotheses & NEXT `f227fc0` 2026-09-28 17:40:56; RESEARCH\_\allowbreak{}FRAME `6830277` & Withdrawn and replaced by SC1w–SC3w. Reasons: the level/\allowbreak{}structure classification became a post hoc variable; the 3-split h25b protocol was replaced by the LG grid. Never computed & WRAPUP 3.1 (`316714c` 2026-09-30 03:02:42); NEXT revision line (`ff1be02` 11:44:30) & before (withdrawal); the NEXT record line after T\_\allowbreak{}res \\
7 & AK1 (NEXT §8.3) & withdrawn hypothesis & NEXT `6830277` 2026-09-29 13:18:02 & Withdrawn and replaced by AK1w. Reason: stage 2 re-scores the same Alaskan sub-regions with another source, so "transfers to new regions" was broader than the design. Never computed & WRAPUP 5 (`316714c`); NEXT revision line (`ff1be02`) & before; NEXT line after T\_\allowbreak{}res \\
8 & SC1w decision rule "no inferior cell and ≥ 8 decided cells" & replaced rule & WRAPUP initial draft (uncommitted, file hash 5976615c…) & Replaced by cell-wise non-inferiority before commit; the draft rule counted undetermined cells toward support & WRAPUP revision 1 (2) (`316714c`) & before \\
9 & Three-way classification in WRAPUP 7.2 (a)3 & withdrawn rule & WRAPUP initial draft & Withdrawn before commit; L38 follows the LG 6B.5 wording & WRAPUP revision 1 (3) (`316714c`) & before \\
10 & Abstract bundle selection: the review example (L1, L2, L4 n = 10, L8, P1 − P0 n = 10, L10, L15, L29, L30) became AB1–AB10 & selection & WRAPUP\_\allowbreak{}REVIEW item 4 (`422c34d` 2026-09-30 01:45:19) & L30 removed (learner robustness is written as a qualifier); AB9 (second baseline) and AB10 (uncertainty) added. Selection made with knowledge of earlier experiments (M1, P2, W3, E1–E3, h25b, a2, b1, b2) & WRAPUP 1.1 (`316714c`) & before \\
11 & LGX X7: existing product comparison and CALM trend reproduction & excluded analysis & LG 6A.10 (`8513cc2` 2026-09-29 18:17:29) & Excluded; product comparison later registered only as a conditional descriptive SI table (WRAPUP 10, executed only if the map is in the main text) & WRAPUP 10 (`316714c`) & before \\
12 & LG 6A.10 exclusions of the TabPFN label grid and of independent-region expansion & exclusions replaced & LG 6A.10 (`8513cc2`) & Replaced by LGT (LG revision 9) and LGD (LG revision 10); the Mongolia and Central Asia exclusion is kept & `6799b74` 2026-09-29 22:43:12 & before \\
13 & WRAPUP 11.1: local re-run of running Rescale jobs; LG refit with site-level draws; time-axis validation, campaign hold-out, CALM trends; site-scale GIPL2 or CryoGrid runs; selection rules at n ≥ 80 and a measurement-priority map; maps for several regions; Ran 2022 ESSD and 15 pending data sets; ERA5-Land 1990–2009 year matching; residual kriging baseline (conditional on LGU-C1); n = 80 grid point (conditional on L4 minimum n = 160); LGF lower-priority tiers; a new ML model of E; Mongolia in the confirmatory pool & excluded analyses & WRAPUP 11.1 (`316714c` 2026-09-30 03:02:42) & Not performed; reasons as in WRAPUP 11.1. The two conditional items (execution plan C12, C13) follow the registration format of LG revision 15 (n) & WRAPUP 11.1 (`316714c`); LG revision 15 (n) (`ff1be02`) & before \\
14 & Exception to row 13, first item: remaining FT-Transformer shards of the stopped Qjpbeb job & scoped exception & user directive 2026-09-30 08:45 & Run locally as a cross-environment auxiliary only; the primary verdict closes on Rescale shards; auxiliary verdicts are not used in confirmatory sentences or the abstract & WRAPUP revision 2 (d); LG revision 15 (c)(d)(f) (`ff1be02` 11:44:30) & after T\_\allowbreak{}res, bundles not unpacked \\
15 & Covariate-space uniform selection rule (RESEARCH\_\allowbreak{}FRAME 5.2, FAILURE\_\allowbreak{}ANALYSIS 6), Gautam 2025 set-up reproduction, R0 covariate-support check rule & excluded analyses & RESEARCH\_\allowbreak{}FRAME 5.2 (`6830277`) & Not performed in this paper; LGX X3a covers random versus region hold-out; map panel (e) marks extrapolation & WRAPUP revision 2 (e) (`ff1be02` 11:44:30) & after T\_\allowbreak{}res, not unpacked \\
16 & NOVELTY 7.4 option 3 (exclusion of ALT values below 20 cm; 138 Alaskan and 89 Lena rows) and option 4 (source-target TDD range overlap table) & excluded analyses & NOVELTY 7.4 (`d8355d2` 2026-09-29 18:18:36) & Not registered by the user; sensitivity analyses unrelated to confirmatory verdicts & WRAPUP revision 2 (e); LG revision 15 (o) (`ff1be02`) & after T\_\allowbreak{}res, not unpacked \\
17 & Region-equal-weight E0 sensitivity & not registered & WRAPUP\_\allowbreak{}REVIEW item 11 (`422c34d` 2026-09-30 01:45:19) & Not registered and not performed; only the E interpretation paragraph is written (M3) & execution plan decision 9; WRAPUP revision 2 (e); LG revision 15 (o) (`ff1be02`) & after T\_\allowbreak{}res; would have required registration before LG P0 retrieval \\
18 & LGX-8.6 options: (a) Rescale continuation and (e) Rescale summary job & removed options & WRAPUP 8.6 (`316714c`) & Removed by the user directive of 2026-09-30 08:45 (no further Rescale cost); option (b) local merge into the primary verdict was already withdrawn by LG revision 14 (4) after cross-environment check (i) & LG revision 15 (a) (`ff1be02`); revision 14 (`028900a` 04:19:28) & (a), (e) after T\_\allowbreak{}res; (b) before \\
19 & AB10 resampling count "2,000" & corrected & WRAPUP 1.1 (`316714c`) & Corrected to 1,000 (LGU revision 2, p resolution 1e-3); the WRAPUP 1.1 body keeps the old text & LGU revision 2 (`582be9a` 03:36:00); WRAPUP revision 2 (b) (`ff1be02`) & LGU change before; WRAPUP record after T\_\allowbreak{}res \\
20 & LGU candidate sentence "labels correct the interval location but do not reduce its width" & dropped sentence & LGU review document & Not used; width statements belong to LGX L25 & LGU initial (`a112cef` 2026-09-29 23:30:02) & before \\
21 & Label "revision after retrieval, bundles not unpacked and not viewed" & status label & WRAPUP 0.3 fourth item & Four amendments (LG revision 15, LGF revision 6, LGU revision 4, WRAPUP revision 2) were committed after T\_\allowbreak{}res (07:20:22) in `ff1be02` (11:44:30). WRAPUP 0.3 would label result-dependent rules "registered after viewing results". The amendments change only technical corrections, a map-mask interpretation, the status of auxiliary local runs and the licence-verified data rule; at commit time no bundle was unpacked and no result table was created or opened. They therefore carry the label above instead. The user has not yet confirmed this judgement (execution plan decision 11); if the user decides otherwise, the affected rows carry "registered after viewing results" & LG revision 15 (l); WRAPUP revision 2 (g) (`ff1be02`) & after T\_\allowbreak{}res; LGX bundles unpacked at 11:44:39 \\
22 & LG revision 15 (b) statement of the remaining nn FT-Transformer shards (16 main plus 35 others; Greenland has no unit) & corrected record & LG revision 15 (`ff1be02` 11:44:30) & After unpacking, the 51 shards were 16 main, AL-1 4, six targets × 5 = 30, and Greenland split 2 = 1; the Greenland sentence was wrong & LG revision 15 post-unpack note (`1a86df0` 2026-09-30 11:46:13) & after unpacking (file names and status fields only) \\
\end{longtable}}
\par\medskip\noindent{\small\textit{b. Past negative results, reclassified (frame)}}\par\nopagebreak
{\fontsize{7.5}{9}\selectfont\setlength{\tabcolsep}{4pt}\setlength{\LTpre}{2pt}\setlength{\LTpost}{4pt}
\begin{longtable}{@{}>{\raggedright\arraybackslash}p{\dimexpr 0.071\linewidth-8pt\relax}>{\raggedright\arraybackslash}p{\dimexpr 0.210\linewidth-8pt\relax}>{\raggedright\arraybackslash}p{\dimexpr 0.085\linewidth-8pt\relax}>{\raggedright\arraybackslash}p{\dimexpr 0.078\linewidth-8pt\relax}>{\raggedright\arraybackslash}p{\dimexpr 0.127\linewidth-8pt\relax}>{\raggedright\arraybackslash}p{\dimexpr 0.078\linewidth-8pt\relax}>{\raggedright\arraybackslash}p{\dimexpr 0.123\linewidth-8pt\relax}>{\raggedright\arraybackslash}p{\dimexpr 0.229\linewidth-8pt\relax}@{}}
\toprule
\textbf{Past ID} & \textbf{Source table} & \textbf{CI type} & \textbf{Original resamples} & \textbf{Reclassification} & \textbf{Value} & \textbf{Related registered test} & \textbf{Audit and withdrawal history} \\
\midrule\endfirsthead
\toprule
\textbf{Past ID} & \textbf{Source table} & \textbf{CI type} & \textbf{Original resamples} & \textbf{Reclassification} & \textbf{Value} & \textbf{Related registered test} & \textbf{Audit and withdrawal history} \\
\midrule\endhead
\bottomrule\endlastfoot
M1 H-series & `m1/\allowbreak{}m1\_\allowbreak{}sc\_\allowbreak{}tests.csv` & \textcolor{blue}{[VERIFY]} & not recorded & \textcolor{blue}{[RESULT]} & \textcolor{blue}{[RESULT]} & H13 to L8 and AB8; H3 to L15 and AB3; H15, H17 to LGU-B1 & \textcolor{blue}{[VERIFY]} \\
FP-F4 & `h4/\allowbreak{}b1\_\allowbreak{}summary.csv` (d\_\allowbreak{}phys columns), `b1\_\allowbreak{}protocol.csv` & \textcolor{blue}{[VERIFY]} & 1,000 & \textcolor{blue}{[RESULT]} & \textcolor{blue}{[RESULT]} & SC1w, SC3w (b1 "always shrink plus residual") & FINAL verdict rejected: the 3-label two-stage protocol was not better than always shrink plus residual (`f098031` 2026-09-26 21:42:16) \\
FP-F5, FP-F6 & `h4/\allowbreak{}b2\_\allowbreak{}curve.csv` (d\_\allowbreak{}phys, d\_\allowbreak{}shrink columns), `b2\_\allowbreak{}regional.csv` (cell-weighted only), `b2\_\allowbreak{}f5.csv`, `b2\_\allowbreak{}theory\_\allowbreak{}corr.csv` & mixed & 1,000 & \textcolor{blue}{[RESULT]} & \textcolor{blue}{[RESULT]} & AB4, L28 (κ), P3 & FINAL verdicts rejected: adaptive pooling did not beat κ = 10 (F5); theoretical label count uncorrelated (F6) \\
FP-F7 & `h4/\allowbreak{}b3\_\allowbreak{}summary.csv` (blk\_\allowbreak{}d\_\allowbreak{}phys, blk\_\allowbreak{}d\_\allowbreak{}rand columns), `b3\_\allowbreak{}region.csv`, `b3\_\allowbreak{}f7.csv` & \textcolor{blue}{[VERIFY]} & 1,000 & \textcolor{blue}{[RESULT]} & \textcolor{blue}{[RESULT]} & none (selection rules at n ≥ 80 excluded, LG 6A.10) & FINAL verdict rejected: no gain from D-optimal selection \\
FP-F8 & `h4/\allowbreak{}b4\_\allowbreak{}tests.csv` & \textcolor{blue}{[VERIFY]} & not recorded & \textcolor{blue}{[RESULT]} & \textcolor{blue}{[RESULT]} & L32, L33 (TabPFN) & FINAL verdict supported; "TabPFN on par" wording to be corrected (J10) \\
FP-F9 & `h4/\allowbreak{}c1\_\allowbreak{}tests.csv` & \textcolor{blue}{[VERIFY]} & 1,000 & \textcolor{blue}{[RESULT]} & \textcolor{blue}{[RESULT]} & L29, AB2 (physics and product baselines) & FINAL verdict rejected: no gain from multi-layer CCI combination \\
H18–H22 & `h2/\allowbreak{}h\_\allowbreak{}tests\_\allowbreak{}all.csv` (block-equal CIs present, 1,055 rows); predictions `m1/\allowbreak{}h1819\_\allowbreak{}shard*\_\allowbreak{}preds.npz`, `h2/\allowbreak{}h21\_\allowbreak{}preds.npz` (includes H20), `h2/\allowbreak{}h22\_\allowbreak{}preds.npz` & two weightings & not recorded & \textcolor{blue}{[RESULT]} & re-scored \textcolor{blue}{[RESULT]} & none & WRAPUP 1.4 (b) body text corrected by revision 2 (c) \\
H23, H24 & `h2/\allowbreak{}h23\_\allowbreak{}tests.csv`, `h2/\allowbreak{}h24\_\allowbreak{}tests.csv` & row resampling & not recorded & protocol differs, descriptive & \textcolor{blue}{[RESULT]} & L43 (placement) & H24 verdict "A (adopted)" withdrawn (AUDIT action 4, `f098031` 2026-09-26 21:42:16); "H24 k-center adopted" discarded (RESEARCH\_\allowbreak{}FRAME A.4, `6830277` 2026-09-29 13:18:02) \\
H25 & `h3/\allowbreak{}h25b\_\allowbreak{}curve.csv`, `h3/\allowbreak{}h25x\_\allowbreak{}curve.csv` (two weightings); `h3/\allowbreak{}h25\_\allowbreak{}curve.csv` (row CI); break-even tables (descriptive) & mixed & 1,000 (h25b, h25x); not recorded (h25) & \textcolor{blue}{[RESULT]} & \textcolor{blue}{[RESULT]} & L4, AK1w & "break-even n set by abs(log E ratio), ρ −0.84" redefined (AUDIT action 1, `f098031`) and discarded (RESEARCH\_\allowbreak{}FRAME A.4, `6830277`) \\
H26 & registered paired test not produced & none & none & source missing & none & none & rejected without the registered paired CI; criterion unattainable at registration (AUDIT workflow:F5, `f098031`) \\
H27 & `h3/\allowbreak{}h27*\_\allowbreak{}rule.csv` (4 files) & rule tables & not recorded & protocol differs, descriptive & \textcolor{blue}{[RESULT]} & none & redefined with H25 (AUDIT action 1) \\
H28 & `h3/\allowbreak{}h28\_\allowbreak{}tests.csv` & \textcolor{blue}{[VERIFY]} & not recorded & \textcolor{blue}{[RESULT]} & \textcolor{blue}{[RESULT]} & L3 (nested α selection) & H28 = rejected (AUDIT; RESEARCH\_\allowbreak{}FRAME A.4) \\
H29 & `h3/\allowbreak{}h29\_\allowbreak{}summary.csv`, `h3/\allowbreak{}h29b\_\allowbreak{}summary.csv` & summary & not recorded & protocol differs, descriptive & \textcolor{blue}{[RESULT]} & L43 (h29b harmful single-block share) & block-CI re-run h29b (`f098031` 2026-09-26 21:42:16) \\
H30 & `h3/\allowbreak{}h30\_\allowbreak{}deploy\_\allowbreak{}tests.csv` & \textcolor{blue}{[VERIFY]} & not recorded & \textcolor{blue}{[RESULT]} & \textcolor{blue}{[RESULT]} & \textcolor{blue}{[VERIFY]} & \textcolor{blue}{[VERIFY]} \\
\end{longtable}}

\clearpage
\phantomsection\label{tab:S16}
\noindent\textbf{Supplementary Table S16.} Software and execution environments.\par\smallskip
\noindent{\footnotesize LG main run, LGX-A and LGX-B ran on cloud nodes (Rescale iolite-4: 64 CPU cores, four NVIDIA T4 GPUs). LGT, LGF, LGU, LGD, the local continuation of LGX and the aggregations ran on a shared local server with NVIDIA RTX 3090 GPUs (24 GB) at low priority (nice 10, at most 32 CPU threads, one process per GPU). Package versions of the WF runs are those of the Code availability statement (results/rescale\_wf*/logs\_wf/deps.log). Each shard records the code hash, the configuration hash, the package versions and the device. Source: docs/MANUSCRIPT\_DRAFT\_METHODS\_INTRO\_2026-09-30.md Table 4 and methods.tex (Code availability).}\par
{\fontsize{7.5}{9}\selectfont\setlength{\tabcolsep}{4pt}\setlength{\LTpre}{2pt}\setlength{\LTpost}{4pt}
\begin{longtable}{@{}>{\raggedright\arraybackslash}p{\dimexpr 0.220\linewidth-8pt\relax}>{\raggedright\arraybackslash}p{\dimexpr 0.260\linewidth-8pt\relax}>{\raggedright\arraybackslash}p{\dimexpr 0.260\linewidth-8pt\relax}>{\raggedright\arraybackslash}p{\dimexpr 0.260\linewidth-8pt\relax}@{}}
\toprule
\textbf{Item} & \textbf{Local server} & \textbf{Cloud (LG, LGX)} & \textbf{Cloud (WF), local (XA)} \\
\midrule\endfirsthead
\toprule
\textbf{Item} & \textbf{Local server} & \textbf{Cloud (LG, LGX)} & \textbf{Cloud (WF), local (XA)} \\
\midrule\endhead
\bottomrule\endlastfoot
Python & 3.9.18 & 3.11.10 & \textcolor{blue}{[S16: Python version of the WF and X runs]} \\
NumPy /\allowbreak{} pandas /\allowbreak{} scikit-learn /\allowbreak{} SciPy & 1.26.4 /\allowbreak{} 2.1.4 /\allowbreak{} 1.3.0 /\allowbreak{} 1.11.4 & 2.4.6 /\allowbreak{} 2.3.3 /\allowbreak{} 1.9.1 /\allowbreak{} 1.17.1 & NumPy 2.1.1 (third WF run 2.4.6) / pandas 2.3.3 / scikit-learn 1.9.1 / SciPy 1.17.1 \\
CatBoost /\allowbreak{} PyTorch /\allowbreak{} pytabkit & 1.2.10 /\allowbreak{} 2.6.0 (CUDA 12.4) /\allowbreak{} 1.7.3 & 1.2.10 /\allowbreak{} 2.4.1 (CUDA 12.1) /\allowbreak{} 1.7.3 & CatBoost 1.2.10 / \textcolor{blue}{[S16]} / \textcolor{blue}{[S16]} \\
Foundation models & tabpfn 8.0.7, tabicl 2.0.2 & none & none \\
GPU & NVIDIA RTX 3090, 24 GB & NVIDIA T4 & CPU nodes (WF); XA local, no fit \\
\end{longtable}}

\clearpage
\phantomsection\label{tab:S17}
\noindent\textbf{Supplementary Table S17.} Data sources, licences and excluded data.\par\smallskip
\noindent{\footnotesize Licences as recorded at access; rows from sources without a confirmed licence are excluded from the deposited tables and from the licence-verified edition used for the main verdicts. Values from these sources can be obtained from the original providers. Source: docs/MANUSCRIPT\_DRAFT\_METHODS\_INTRO\_2026-09-30.md (Data availability, 30 September 2026) and front.tex (Data availability); licence status per data/processed/fidelity\_base\_v4\_meta.json (by\_source).}\par
{\fontsize{7.5}{9}\selectfont\setlength{\tabcolsep}{4pt}\setlength{\LTpre}{2pt}\setlength{\LTpost}{4pt}
\begin{longtable}{@{}>{\raggedright\arraybackslash}p{\dimexpr 0.118\linewidth-8pt\relax}>{\raggedright\arraybackslash}p{\dimexpr 0.272\linewidth-8pt\relax}>{\raggedright\arraybackslash}p{\dimexpr 0.356\linewidth-8pt\relax}>{\raggedright\arraybackslash}p{\dimexpr 0.145\linewidth-8pt\relax}>{\raggedright\arraybackslash}p{\dimexpr 0.109\linewidth-8pt\relax}@{}}
\toprule
\textbf{Use} & \textbf{Data set} & \textbf{Repository and identifier} & \textbf{Licence} & \textbf{Status} \\
\midrule\endfirsthead
\toprule
\textbf{Use} & \textbf{Data set} & \textbf{Repository and identifier} & \textbf{Licence} & \textbf{Status} \\
\midrule\endhead
\bottomrule\endlastfoot
Labels, main regions & GTN-P/CALM compilation (Streletskiy et al., 2025) & PANGAEA, doi:10.1594/PANGAEA.972777 & CC BY 4.0 & used \\
Labels, main regions & ALLena thaw depth, Lena River Delta (Veremeeva et al., 2025) & PANGAEA, doi:10.1594/PANGAEA.973813; 10.1594/PANGAEA.974408 & CC BY 4.0 & used \\
Labels, main regions & ABoVE soil moisture and ALT, version 2 (Moore et al., 2025) & ORNL DAAC, doi:10.3334/ORNLDAAC/2369 & \textcolor{blue}{[S17: licence]} & used \\
Labels, new regions & GPR and pit survey points, Tibetan Plateau (Du et al., 2026) & Zenodo, doi:10.5281/zenodo.21999366; field data doi:10.12072/ncdc.permafrost.db7703.2026 & MIT (Zenodo record) \textcolor{blue}{[미확인: field data licence]} & used \\
Labels, new regions & Ilulissat probing (Scheer et al., 2024) & PANGAEA, doi:10.1594/PANGAEA.964306 & CC BY 4.0 & used \\
Labels, new regions & T-MOSAiC myThaw 2021--2023 (Martin et al., 2023; Boike et al., 2024; Hammar et al., 2025) & PANGAEA, doi:10.1594/PANGAEA.956039; .971586; .974461 & CC BY 4.0 & used \\
Labels, new regions & Additional CALM events of the PANGAEA compilation & PANGAEA, doi:10.1594/PANGAEA.972777 & CC BY 4.0 & used \\
Labels, new regions & Disko Island (Zastruzny et al., 2024) & PANGAEA, doi:10.1594/PANGAEA.967139 & CC BY 4.0 & used \\
Labels, new regions & Tavvavuoma (Sannel, 2020) & Bolin Centre Database, doi:10.17043/sannel-2020-temperature-1 & CC BY 4.0 & used \\
Labels, new regions & Cape Mamontov Klyk (Grosse, 2007) & PANGAEA, doi:10.1594/PANGAEA.611409 & CC BY 3.0 & used \\
Labels, new regions & Pedon data (Palmtag et al., 2022) & Bolin Centre Database, doi:10.17043/palmtag-2022-pedon-1 & ODC-By & used \\
Labels, new regions & FireALT, measured depths only (Talucci et al., 2024) & NSF Arctic Data Center, doi:10.18739/A2RN3092P & CC BY 4.0 & used \\
Labels, new regions & Syrdakh observatory (Pohl et al., 2026) & Zenodo, doi:10.5281/zenodo.19890671 & CC BY 4.0 & used \\
Labels, new regions & Kolyma Water-Balance Station (Makarieva et al., 2017) & PANGAEA, doi:10.1594/PANGAEA.881754 & CC BY 3.0 & used \\
Labels, new regions & Indigirka tree mounds (Liang et al., 2023) & PANGAEA, doi:10.1594/PANGAEA.961876 & CC BY 4.0 & used \\
Labels, new regions & Yamal transect (Walker et al., 2009) & PANGAEA, doi:10.1594/PANGAEA.842711 & CC BY 3.0 & used (PANGAEA rows) \\
Label definition test & Temperature-derived labels (Fu, 2025) & figshare, doi:10.6084/m9.figshare.29206613.v1 & CC BY 4.0 & L39 only \\
Labels, licence not confirmed & CALM web files, Abisko and Kapp Linn\'e (\AA kerman, 1998; \AA kerman and Johansson, 2008) & https://www2.gwu.edu/\textasciitilde calm/data/north.htm & not confirmed & full-data edition only \\
Labels, licence not confirmed & CUSP v1.1 synthesis (Jorgenson and Kanevskiy, 2025; Petrone et al., 2016) & github.com/jonschwenk/cusp v1.1; doi:10.18739/A27P8TG0G (CC0 1.0); doi:10.1594/PANGAEA.845258 (CC BY-NC-SA 3.0) & not confirmed & full-data edition only \\
Labels, licence not confirmed & NSIDC GGD353 thaw tubes (Nixon, 2003) & doi:10.7265/7m84-k262 & provider asks to be consulted & full-data edition only \\
Labels, licence not confirmed & Yamal transect, ten report-derived rows & report & not confirmed & full-data edition only \\
Descriptive only & NSIDC GGD402; Yang and Qiu (2026) &  &  & parsed, no analysis \\
Covariates & ERA5-Land monthly averaged data (Mu\~noz-Sabater et al., 2021) & Copernicus Climate Data Store, doi:10.24381/cds.68d2bb30 & \textcolor{blue}{[S17: licence]} & used \\
Covariates & ESA CCI Permafrost ALT, version 4.0 (Westermann et al., 2024) & CEDA, doi:10.5285/d34330ce3f604e368c06d76de1987ce5 & \textcolor{blue}{[S17: licence]} & used \\
Covariates & SoilGrids 2.0 (Poggio et al., 2021) & ISRIC WCS service & CC BY 4.0 \textcolor{blue}{[미확인: release tag at access]} & used \\
Covariates & Copernicus DEM GLO-30 & doi:10.5270/ESA-c5d3d65 & \copyright{} DLR e.V. 2010--2014 and Airbus 2014--2018, Copernicus programme & used \\
Map masks & ESA CCI Permafrost extent, 1997--2021 mean & \textcolor{blue}{[미확인: data set DOI]} &  & maps only \\
Map masks & JRC Global Surface Water occurrence (Pekel et al., 2016) & \textcolor{blue}{[미확인: citation]} &  & maps only \\
Optional & KPDC Council case study (KOPRI-KPDC-00002125, -00002707, -00002955) & KPDC & KPDC access conditions & \textcolor{blue}{[DECISION: KPDC 자료, SI 포함 여부와 식별자, 이용 조건]} \\
\end{longtable}}

\clearpage
\phantomsection\label{tab:S18}
\noindent\textbf{Supplementary Table S18.} Survey of public data for independent regions (XF).\par\smallskip
\noindent{\footnotesize Public field ALT data sets surveyed for independent regions (survey of 4 October 2026, docs/research/2026-10-04/open\_regions.md) and the outcome of the eligibility rules of Supplementary Methods 2 applied blind in XF. Rows are filled after the XF run (placement fixed before results: Supplementary Information). \textcolor{blue}{[PENDING: XF, new public regions; registered data deadline 11 October 2026, 14:22 KST; no eligible new macro region found so far]}}\par
{\fontsize{7.5}{9}\selectfont\setlength{\tabcolsep}{4pt}\setlength{\LTpre}{2pt}\setlength{\LTpost}{4pt}
\begin{longtable}{@{}>{\raggedright\arraybackslash}p{\dimexpr 0.315\linewidth-8pt\relax}>{\raggedright\arraybackslash}p{\dimexpr 0.120\linewidth-8pt\relax}>{\raggedright\arraybackslash}p{\dimexpr 0.086\linewidth-8pt\relax}>{\raggedright\arraybackslash}p{\dimexpr 0.120\linewidth-8pt\relax}>{\raggedright\arraybackslash}p{\dimexpr 0.154\linewidth-8pt\relax}>{\raggedright\arraybackslash}p{\dimexpr 0.103\linewidth-8pt\relax}>{\raggedright\arraybackslash}p{\dimexpr 0.103\linewidth-8pt\relax}@{}}
\toprule
\textbf{Data set} & \textbf{Identifier} & \textbf{Licence} & \textbf{Method and period} & \textbf{Size and relation to existing regions} & \textbf{Survey decision} & \textbf{XF eligibility} \\
\midrule\endfirsthead
\toprule
\textbf{Data set} & \textbf{Identifier} & \textbf{Licence} & \textbf{Method and period} & \textbf{Size and relation to existing regions} & \textbf{Survey decision} & \textbf{XF eligibility} \\
\midrule\endhead
\bottomrule\endlastfoot
\textcolor{blue}{[PENDING: XF survey rows (docs/research/2026-10-04/open\_regions.md 4.1--4.4) and eligibility after the XF run; registered data deadline 11 October 2026, 14:22 KST]} &  &  &  &  &  &  \\
\end{longtable}}


\clearpage
\section*{Supplementary Figures}

Supplementary figures are planned from the result maps and panels moved out of the main figures (MANUSCRIPT\_SPEC 8). Each figure below is a placeholder with its planned content.

\begin{figure}[htbp]\centering
\fbox{\parbox[c][45mm][c]{0.85\linewidth}{\centering Supplementary Fig.~S1 in preparation}}
\caption{Validation ladder (Results, evaluation design). Split maps of the validation ladder; pairs of distance distributions in Alaska for each scoring scheme (test-to-training and prediction-grid-to-label distances, as in Meyer and Pebesma 2022). \textcolor{blue}{[SI figure: to be drawn]}}\label{fig:S1}
\end{figure}

\begin{figure}[htbp]\centering
\fbox{\parbox[c][45mm][c]{0.85\linewidth}{\centering Supplementary Fig.~S2 in preparation}}
\caption{Transfer without target labels. Target-wise maps of the error change of the low-weight residual and of the anchor plus residual relative to the source-coefficient model; learner-wise detail. \textcolor{blue}{[SI figure: to be drawn]}}\label{fig:S2}
\end{figure}

\begin{figure}[htbp]\centering
\fbox{\parbox[c][45mm][c]{0.85\linewidth}{\centering Supplementary Fig.~S3 in preparation}}
\caption{Coefficient recalibration. Target-wise maps of the source-coefficient error and of the ten-label recalibration gain; curves of the soil-property recalibrated Stefan model; regional contrast forest plots. \textcolor{blue}{[SI figure: to be drawn]}}\label{fig:S3}
\end{figure}

\begin{figure}[htbp]\centering
\fbox{\parbox[c][45mm][c]{0.85\linewidth}{\centering Supplementary Fig.~S4 in preparation}}
\caption{Method selection. Target-wise maps of the method selected by the rule and of the rule minus the fixed recipe; matrix and rank scatter of the source cross-validation tuning. \textcolor{blue}{[SI figure: to be drawn]}}\label{fig:S4}
\end{figure}

\begin{figure}[htbp]\centering
\fbox{\parbox[c][45mm][c]{0.85\linewidth}{\centering Supplementary Fig.~S5 in preparation}}
\caption{Gains within label-rich regions. Maps of the error change by scoring block in the Lena Delta and Canada; figures of XE, XI and XB. \textcolor{blue}{[SI figure: to be drawn]}}\label{fig:S5}
\end{figure}

\begin{figure}[htbp]\centering
\fbox{\parbox[c][45mm][c]{0.85\linewidth}{\centering Supplementary Fig.~S6 in preparation}}
\caption{Placement of new labels. Target-wise maps of strategy minus random cells; proximity and distance contrasts (L23, L24); learned placement policy (XD). \textcolor{blue}{[SI figure: to be drawn]}}\label{fig:S6}
\end{figure}

\begin{figure}[htbp]\centering
\fbox{\parbox[c][45mm][c]{0.85\linewidth}{\centering Supplementary Fig.~S7 in preparation}}
\caption{Independent regions. Target-wise maps of the error change and prediction maps for the independent regions. \textcolor{blue}{[SI figure: to be drawn]}}\label{fig:S7}
\end{figure}

\begin{figure}[htbp]\centering
\fbox{\parbox[c][45mm][c]{0.85\linewidth}{\centering Supplementary Fig.~S8 in preparation}}
\caption{Prediction intervals. Maps of predicted ALT, 90\% interval width and the extrapolation area for the Lena Delta (outputs/maps/transfer\_lena/). \textcolor{blue}{[SI figure: to be drawn]}}\label{fig:S8}
\end{figure}

\begin{figure}[htbp]\centering
\fbox{\parbox[c][45mm][c]{0.85\linewidth}{\centering Supplementary Fig.~S9 in preparation}}
\caption{Deployment scenarios. Scenario matrix, stopping-rule curves and forest plot of the abstract contrast bundle (panels of the earlier Fig.~7). \textcolor{blue}{[SI figure: to be drawn]}}\label{fig:S9}
\end{figure}

\clearpage
% si_maps.tex : generated by tools/make_si_maps.py (do not edit by hand).
% Supplementary Figs S1-S6 (former S10-S15): 1 km ALT maps (display resolution, no accuracy claim) and the XK/XL figures.
% S1 (former S10): Alaska_ALT_map_v3.pdf (file time 2026-10-05 13:03); caption from Alaska_ALT_map_v3_legend.md
\begin{figure}[p]\centering
\IfFileExists{../../../outputs/figures/paper/v3_restructure/maps/Alaska_ALT_map_v3.pdf}{\includegraphics[width=\linewidth,height=0.72\textheight,keepaspectratio]{../../../outputs/figures/paper/v3_restructure/maps/Alaska_ALT_map_v3.pdf}}{\fbox{\parbox[c][45mm][c]{0.85\linewidth}{\centering Supplementary Fig.~S1 in preparation}}}
\caption{\textbf{Active-layer thickness (ALT) of the Alaskan permafrost zone at 1 km display resolution: recalibrated Stefan model, residual machine-learning correction and final map.} a, Recalibrated Stefan ALT, E·√TDD with E = 1.617 (least squares on all 13,606 label cells, shrunk towards the source coefficient). b, Residual ML correction λ·g (final minus recalibrated Stefan; positive, deeper), where g is a CatBoost residual model on 25 covariates (mean of two seeds) and λ = 0.5 was selected from {0.25, 0.5, 1.0} by 0.5° block five-fold cross-validation. c, Final ALT, a plus b. a and c share one colour scale (35–75 cm, 1st–99th percentile of both maps). d, Number of covariates outside the 0.5–99.5 percentile range of the training labels (hatching density), with label cells as dots; 73\% of mapped cells have at least one, mostly terrain roughness, slope and topographic position. The 90\% prediction interval of c is ±20.2 cm everywhere (constant-width split conformal on block cross-validation residuals; cross-validated RMSE 13.8 cm, recalibrated Stefan 14.3 cm; lower bound truncated at 0 cm in 5,384 cells). Its nominal coverage is an average over cells similar to the labels and is not guaranteed in hatched cells, so the width is not mapped. The 1 km grid is a display resolution: air thawing index and climate covariates come from ERA5-Land (0.1°, about 4–6 × 11 km here), soil covariates from SoilGrids windows at about 5 km, and the residual correction acts mainly between climate-grid cells (main text, Fig.~6c). Cells are drawn as their native 0.009° rectangles. Domain: land (Natural Earth 10 m) with ESA CCI permafrost fraction of at least 50\% (1997–2021 mean), 59–71.5° N, 168–141° W; 858,517 of 899,633 cells mapped. White, sea or water (Global Surface Water occurrence of at least 50\%); grey, missing CCI ALT or soil thaw index. NAD83 Alaska Albers (EPSG:3338); inset, location. Data: ABoVE ALT v2 (Moore et al. 2025, ORNL DAAC) and GTN-P/CALM (Streletskiy et al. 2025) labels; ERA5-Land (Copernicus C3S); SoilGrids 2.0 (ISRIC); Copernicus DEM GLO-30 (ESA, AWS Open Data); ESA CCI Permafrost v4.0; JRC Global Surface Water; Natural Earth; Cartopy 0.23.}\label{fig:S10}
\end{figure}
\clearpage
% S2 (former S11): Lena_ALT_map_v3.pdf (file time 2026-10-05 13:03); caption from Lena_ALT_map_v3_legend.md
\begin{figure}[p]\centering
\IfFileExists{../../../outputs/figures/paper/v3_restructure/maps/Lena_ALT_map_v3.pdf}{\includegraphics[width=\linewidth,height=0.72\textheight,keepaspectratio]{../../../outputs/figures/paper/v3_restructure/maps/Lena_ALT_map_v3.pdf}}{\fbox{\parbox[c][45mm][c]{0.85\linewidth}{\centering Supplementary Fig.~S2 in preparation}}}
\caption{\textbf{Active-layer thickness (ALT) of the Lena River Delta at 1 km display resolution: recalibrated Stefan model, residual machine-learning correction and final map.} a, Recalibrated Stefan ALT, E·√TDD with E = 1.435 (least squares on all 3,037 Lena label cells, shrunk towards the source coefficient 1.619). b, Residual ML correction 0.25·g (final minus recalibrated Stefan; positive, deeper), where g is a CatBoost residual model (mean of two seeds) fitted to residuals of source-region cells beyond a 100 km buffer and of the Lena labels. c, Final ALT, a plus b. a and c share one colour scale (35–50 cm, 1st–99th percentile of both maps). d, Number of covariates outside the 0.5–99.5 percentile range of the training rows of c (source cells and Lena labels), with label cells as dots; 84\% of mapped cells have at least one, mostly ERA5-Land coldest-month temperature, mean annual air temperature and freezing index. Relative to source cells alone, all have three or more. The 90\% prediction interval of c is ±25.4 cm everywhere (constant-width split conformal on 0.5° block five-fold cross-validation residuals within the Lena labels; cross-validated RMSE 20.8 cm, recalibrated Stefan 21.5 cm). Its nominal coverage is an average over cells similar to the labels and is not guaranteed in hatched cells, so the width is not mapped. The 1 km grid is a display resolution: air thawing index and climate covariates come from ERA5-Land (0.1°, about 3 × 11 km here), soil covariates from SoilGrids windows at about 5 km, and the residual correction acts mainly at the climate-grid scale (main text, Fig.~6c); the horizontal bands in a follow ERA5-Land rows. Cells are drawn as their native 0.009° rectangles. Grid: 71.5–73.6° N, 123.3–130.1° E; 38,086 mapped cells. White, sea or water (Global Surface Water occurrence of at least 50\%); grey, missing CCI ALT. Polar stereographic projection (central meridian 126.7° E, true scale at 70° N); inset, location. Data: ALLena thaw-depth compilation (Veremeeva et al. 2025, PANGAEA); ERA5-Land (Copernicus C3S); SoilGrids 2.0 (ISRIC); Copernicus DEM GLO-30 (ESA, AWS Open Data); ESA CCI Permafrost v4.0; JRC Global Surface Water; Natural Earth 10 m; Cartopy 0.23.}\label{fig:S11}
\end{figure}
\clearpage
% S3 (former S12): XL_display_resolution.pdf (file time 2026-10-05 13:04); caption from XL_display_resolution_legend.md
\begin{figure}[p]\centering
\IfFileExists{../../../outputs/figures/paper/v3_restructure/maps/XL_display_resolution.pdf}{\includegraphics[width=\linewidth,height=0.72\textheight,keepaspectratio]{../../../outputs/figures/paper/v3_restructure/maps/XL_display_resolution.pdf}}{\fbox{\parbox[c][45mm][c]{0.85\linewidth}{\centering Supplementary Fig.~S3 in preparation}}}
\caption{\textbf{The same residual ML ALT map at its 1 km cells and as ERA5-Land 0.1° cell means.} a, b, Alaska, 50 × 50 km around the 0.5° block with the most label cells (69.1° N, 148.8° W; 2,892 cells); white lines, ERA5-Land 0.1° cell boundaries. c, d, Lena Delta, full grid. a and c show the 1 km cells; b and d colour each cell with the mean of the mapped cells in its ERA5-Land 0.1° cell (boundaries at grid point ± 0.05°). All cells are drawn as their native 0.009° rectangles without resampling. One colour scale per region from the 1st–99th percentile of both panels (Alaska 42–54 cm, Lena Delta 38–48 cm), narrower than in Supplementary Figs~S1 and S2 so that within-cell detail is visible. The 0.1° cell means explain 97\% (Alaska, whole domain) and 61\% (Lena Delta) of the variance of the 1 km map; the standard deviation of the 1 km values around their 0.1° cell mean is 1.8 cm in Alaska (total 9.5 cm) and 1.1 cm in the Lena Delta (total 1.8 cm). The 1 km map is a display resolution; its information resolution is bound to the climate input (ERA5-Land, about 10 km) and the soil input (SoilGrids, about 5 km). Within-cell detail comes from the terrain (Copernicus DEM 30 m windows) and CCI ALT inputs. Projections as in Supplementary Figs~S1 and S2 (NAD83 Alaska Albers; polar stereographic, central meridian 126.7° E). Data as in Supplementary Figs~S1 and S2; Cartopy 0.23.}\label{fig:S12}
\end{figure}
\clearpage
% S4 (former S13): XL_Alaska_product_differences.pdf (file time 2026-10-05 13:03); caption from XL_Alaska_product_differences_legend.md
\begin{figure}[p]\centering
\IfFileExists{../../../outputs/figures/paper/v3_restructure/maps/XL_Alaska_product_differences.pdf}{\includegraphics[width=\linewidth,height=0.72\textheight,keepaspectratio]{../../../outputs/figures/paper/v3_restructure/maps/XL_Alaska_product_differences.pdf}}{\fbox{\parbox[c][45mm][c]{0.85\linewidth}{\centering Supplementary Fig.~S4 in preparation}}}
\caption{\textbf{The 1 km Alaska ALT map compared with four ALT products and with the recalibrated Stefan model.} a–e, ALT of the residual ML map (a, as in Supplementary Fig.~S1) and of CCI v5 (1997–2023 mean, b), Wei 2026 (2000–2024 mean, c), Aalto 2018 (2000–2014 baseline, d) and Yi-Kimball (2001–2015 mean, e) on one colour scale (0–240 cm, 2nd–98th percentile of all five maps over mapped cells). Grey text, area mean over mapped cells with a product value and Pearson correlation r with the residual ML map. f, Recalibrated Stefan − residual ML on its own scale (±10 cm, the scale of panel b of Supplementary Figs~S1 and S2 with the sign reversed). g–j, Product − residual ML (positive, product deeper) on one scale (±190 cm, 98th percentile of the absolute differences of g–j). Products are read at the 1 km cell centres (nearest pixel, period mean). Mean differences are +49.4 cm (CCI v5), +14.7 (Wei 2026), +22.6 (Aalto 2018), +49.3 (Yi-Kimball) and +1.4 cm (Stefan). Large differences with CCI v5 and Yi-Kimball are widespread in Interior Alaska; the largest differences with Wei 2026 lie at about 61–63° N, 141–150° W; Aalto 2018 is shallower than the map at about 62–64° N west of 148° W. These maps carry no accuracy statement; accuracy is compared only at label cells (Supplementary Table S2 and Supplementary Fig.~S6). The 1 km map is a display resolution; its information resolution is bound to the climate input (ERA5-Land, about 10 km) and the soil input (SoilGrids, about 5 km). Grey cells, no product value (Yi-Kimball covers 86\% of mapped cells). NAD83 Alaska Albers (EPSG:3338). Data: CCI v5 (ESA CCI Permafrost, CEDA, doi:10.5285/a6fbedd8ee5b472c8e84e55f746c1704); Wei et al. 2026 (Zenodo 21667583, CC BY 4.0); Aalto et al. 2018 (Zenodo 5003804, CC0); Yi and Kimball (ORNL DAAC 1760); covariates and labels as in Supplementary Fig.~S1; Natural Earth 10 m; Cartopy 0.23.}\label{fig:S13}
\end{figure}
\clearpage
% S5 (former S14): XL_Lena_product_differences.pdf (file time 2026-10-05 13:04); caption from XL_Lena_product_differences_legend.md
\begin{figure}[p]\centering
\IfFileExists{../../../outputs/figures/paper/v3_restructure/maps/XL_Lena_product_differences.pdf}{\includegraphics[width=\linewidth,height=0.72\textheight,keepaspectratio]{../../../outputs/figures/paper/v3_restructure/maps/XL_Lena_product_differences.pdf}}{\fbox{\parbox[c][45mm][c]{0.85\linewidth}{\centering Supplementary Fig.~S5 in preparation}}}
\caption{\textbf{The 1 km Lena Delta ALT map compared with three ALT products and with the recalibrated Stefan model.} a–d, ALT of the residual ML map (a, as in Supplementary Fig.~S2) and of CCI v5 (1997–2023 mean, b), Wei 2026 (2000–2024 mean, c) and Aalto 2018 (2000–2014 baseline, d) on one colour scale (25–75 cm, 2nd–98th percentile of all four maps over mapped cells). Grey text, area mean over mapped cells with a product value and Pearson correlation r with the residual ML map. e, Recalibrated Stefan − residual ML on its own scale (±10 cm, the scale of panel b of Supplementary Figs~S1 and S2 with the sign reversed). f–h, Product − residual ML (positive, product deeper) on one scale (±35 cm, 98th percentile of the absolute differences of f–h). Products are read at the 1 km cell centres (nearest pixel, period mean). Mean differences are −0.9 cm (CCI v5), +16.5 (Wei 2026, 94\% of cells with values), +13.0 (Aalto 2018) and −1.1 cm (Stefan). CCI v5 is shallower than the map in the northern and eastern delta and deeper in the south-western part of the grid; Wei 2026 and Aalto 2018 are deeper over most of the grid. These maps carry no accuracy statement; accuracy is compared only at label cells. The 1 km map is a display resolution; its information resolution is bound to the climate input (ERA5-Land, about 10 km) and the soil input (SoilGrids, about 5 km). Grey cells, no product value. Polar stereographic projection (central meridian 126.7° E, true scale at 70° N). Data: CCI v5 (ESA CCI Permafrost, CEDA); Wei et al. 2026 (Zenodo 21667583, CC BY 4.0); Aalto et al. 2018 (Zenodo 5003804, CC0); covariates and labels as in Supplementary Fig.~S2; Natural Earth 10 m; Cartopy 0.23.}\label{fig:S14}
\end{figure}
\clearpage
% S6 (former S15): XK_support_scale.pdf (file time 2026-10-05 13:04); caption from XK_support_scale_legend.md
\begin{figure}[p]\centering
\IfFileExists{../../../outputs/figures/paper/v3_restructure/maps/XK_support_scale.pdf}{\includegraphics[width=\linewidth,height=0.72\textheight,keepaspectratio]{../../../outputs/figures/paper/v3_restructure/maps/XK_support_scale.pdf}}{\fbox{\parbox[c][45mm][c]{0.85\linewidth}{\centering Supplementary Fig.~S6 in preparation}}}
\caption{\textbf{Prediction error of the recalibrated Stefan model, the residual ML model and four ALT products after aggregation to larger grid cells.} Model values are within-region out-of-block predictions from 0.5° block five-fold cross-validation; the Stefan coefficient is refitted on the training folds, and λ of the residual model is selected by an inner block cross-validation within the training folds. Product values are period means at the label cells: CCI v4 (1997–2021), CCI v5 (1997–2023), Wei 2026 (2000–2024) and, for Alaska, Yi-Kimball (2001–2015). Labels and predictions are averaged in grid cells of 0.05°, 0.1° and 0.25° whose boundaries follow the ERA5-Land cells (0.1° = ERA5-Land cell); grid cells with fewer than three label cells are dropped; 1 km denotes single label cells. Grey numbers, grid cells per size. a–c, Cell-weighted RMSE (logarithmic axis) on the cells where all shown values exist (Alaska 10,605 cells in 60 blocks, Lena Delta 1,342 in 15, Canada 745 in 34); bands, 95\% block bootstrap intervals for the two models (1,000 resamples of 0.5° blocks). d–f, RMSE of the recalibrated Stefan model minus RMSE of the residual ML model on all label cells with model values (Alaska 13,606, Lena Delta 2,958, Canada 747), cell-weighted (filled) and block-equal (open) with 95\% block bootstrap intervals. In Alaska the cell-weighted difference is 0.50 cm (0.18 to 1.10) for single cells and 1.54 to 1.67 cm (lower bounds 0.50 to 0.60) at 0.05°, 0.1° and 0.25°. In Canada it is 2.62 cm (0.00 to 4.42) for single cells and 0.96 to 1.35 cm at larger sizes, with intervals that include zero. In the Lena Delta all intervals include zero. Removing cells within 5 km of Wei 2026 training sites changes the Alaska differences by at most 0.17 cm (data/processed/xbatch/XK\_support\_scale/). Products are compared only as RMSE differences at the same grid size.}\label{fig:S15}
\end{figure}
\clearpage

